% !TEX program = xelatex
% -----------------------------------------------------------------------------
% Lecture note skeleton
% Solvable Quantum Many-Body Models in Condensed Matter
% -----------------------------------------------------------------------------
% Recommended compiler: XeLaTeX or LuaLaTeX.
% This file is an English lecture-note scaffold. It is intended to be expanded
% chapter by chapter into a pedagogical set of notes.
% -----------------------------------------------------------------------------

\documentclass[12pt,openany]{book}

% -----------------------------------------------------------------------------
% Packages
% -----------------------------------------------------------------------------
\usepackage[a4paper,margin=0.9in]{geometry}
\usepackage{amsmath,amssymb,amsthm,mathtools}
\usepackage{bm}
\usepackage{booktabs}
\usepackage{array}
\usepackage{makecell}
\usepackage{ragged2e}
\usepackage{adjustbox}
\usepackage{microtype}
\usepackage{longtable}
\usepackage{tabularx}
\usepackage{enumitem}
\usepackage{xcolor}
\usepackage{hyperref}
\usepackage[nameinlink,capitalize]{cleveref}

\hypersetup{
    colorlinks=true,
    linkcolor=blue!60!black,
    citecolor=blue!60!black,
    urlcolor=blue!60!black,
    pdftitle={Solvable Quantum Many-Body Models in Condensed Matter},
    pdfauthor={Qian--Rui Lee}
}

% -----------------------------------------------------------------------------
% Theorem-like environments
% -----------------------------------------------------------------------------
\theoremstyle{definition}
\newtheorem{definition}{Definition}[chapter]
\newtheorem{example}{Example}[chapter]
\newtheorem{remark}{Remark}[chapter]
\newtheorem{principle}{Principle}[chapter]

\theoremstyle{plain}
\newtheorem{proposition}{Proposition}[chapter]
\newtheorem{theorem}{Theorem}[chapter]
\newtheorem{lemma}{Lemma}[chapter]

% -----------------------------------------------------------------------------
% Macros
% -----------------------------------------------------------------------------
\newcommand{\ii}{\mathrm{i}}
\newcommand{\dd}{\mathrm{d}}
\newcommand{\ee}{\mathrm{e}}
\newcommand{\Tr}{\operatorname{Tr}}
\newcommand{\Pf}{\operatorname{Pf}}
\newcommand{\sgn}{\operatorname{sgn}}
\newcommand{\diag}{\operatorname{diag}}
\newcommand{\id}{\mathbf{1}}
\newcommand{\ZZ}{\mathbb{Z}}
\newcommand{\RR}{\mathbb{R}}
\newcommand{\CC}{\mathbb{C}}
\newcommand{\calH}{\mathcal{H}}
\newcommand{\calZ}{\mathcal{Z}}
\newcommand{\calT}{\mathcal{T}}
\newcommand{\ket}[1]{\left|#1\right\rangle}
\newcommand{\bra}[1]{\left\langle #1\right|}
\newcommand{\braket}[2]{\left\langle #1\middle|#2\right\rangle}
\newcommand{\abs}[1]{\left|#1\right|}
\newcommand{\norm}[1]{\left\|#1\right\|}
\newcommand{\todo}[1]{\textcolor{red!70!black}{\textbf{[TODO:} #1\textbf{]}}}

\newcolumntype{Y}{>{\RaggedRight\arraybackslash}X}
\newcolumntype{L}[1]{>{\RaggedRight\arraybackslash}p{#1}}
\setlength{\tabcolsep}{4pt}
\renewcommand{\arraystretch}{1.18}
\setlength{\emergencystretch}{3em}
\newcommand{\boxedpar}[1]{%
\begin{center}
\fbox{\begin{minipage}{0.93\linewidth}\centering #1\end{minipage}}
\end{center}}

% -----------------------------------------------------------------------------
% Title
% -----------------------------------------------------------------------------
\title{\bfseries Solvable Quantum Many-Body Models in Condensed Matter\\[4pt]
\large Mother Models, Equivalence Maps, and Unified Solution Methods}
\author{Qian--Rui Lee}
\date{\today}

\begin{document}
\pagestyle{plain}

\maketitle

\frontmatter

\chapter*{Preface}
\addcontentsline{toc}{chapter}{Preface}

The purpose of these lecture notes is not to present a long catalogue of isolated solvable models. Instead, the central aim is to construct a \emph{mother-model and equivalence-map approach} to exactly solvable quantum many-body systems in condensed matter physics. The guiding strategy is to identify a small number of mathematically general and physically representative structures, solve them in detail, and then recover many familiar models as parameter limits, dual descriptions, transfer-matrix limits, Jordan--Wigner images, Bogoliubov--de Gennes representations, gauge-sector reductions, or special cases of Bethe-ansatz/Yang--Baxter integrability.

The defining principle of the notes is therefore
\boxedpar{%
\emph{Solve a small number of mother models}\par
\vspace{2pt}
$\Longrightarrow$\par
\vspace{2pt}
\emph{cover whole families by parameter choices or equivalence maps.}}

Thus the intended structure is not
\[
\text{Ising}\to\text{XY}\to\text{Kitaev}\to\text{SSH}\to\cdots,
\]
but rather
\boxedpar{%
\begin{tabular}{@{}c@{}}
free fermions/BdG
$\to$ transfer matrices
$\to$ Bethe ansatz/Yang--Baxter\\
$\to$ commuting projectors
$\to$ Majorana--gauge solvability.
\end{tabular}}

\section*{How to read these notes}

A productive reading route is the following.
\begin{enumerate}[label=(\roman*)]
    \item First master the unified tools in Part~I: spin algebra, the Jordan--Wigner transformation, Fourier transformation on a lattice, BdG formalism, Bogoliubov diagonalization, transfer matrices, and duality.
    \item Then study the transverse-field XY chain in Part~II. This is the first genuine mother model in the notes; it contains the TFIM, the XX chain, and the Kitaev chain after fermionization.
    \item Next read Part~III to understand the transfer-matrix correspondence between one-dimensional quantum systems and two-dimensional classical statistical models.
    \item Then move to Part~IV, which treats interacting but integrable systems such as the XXZ spin chain and the six-vertex model.
    \item Finally, read Part~V for two-dimensional quantum exact solvability, including the toric code, the Kitaev honeycomb model, Rokhsar--Kivelson points, and exact ground-state constructions.
\end{enumerate}

\tableofcontents

% -----------------------------------------------------------------------------
\mainmatter
% -----------------------------------------------------------------------------

\chapter{Introduction: Solvability by Mother Models and Equivalence Maps}

\section{Motivation}

Exactly solvable models in condensed matter physics have two complementary roles. First, they provide rare nonperturbative laboratories in which the spectrum, quasiparticles, correlation functions, order parameters, topological invariants, or partition functions can be computed with controlled mathematics. Second, they serve as conceptual fixed points around which more general many-body phenomena can be organized: quantum criticality, topological phases, fractionalization, duality, integrability, and entanglement structure.

The distinctive feature of these notes is that solvable models are organized by \emph{mother models}. A mother model is a Hamiltonian, transfer matrix, algebraic structure, or exactly solvable construction general enough that several familiar models can be obtained from it through parameter choices, limiting procedures, duality transformations, fermionization, or gauge-sector reduction.

\begin{principle}[Mother-model principle]
If a family of models shares the same diagonalization mechanism, transfer matrix, quadratic fermion/BdG structure, Yang--Baxter algebra, stabilizer algebra, or Majorana-gauge reduction, one should first solve the mother structure rather than repeatedly solve each special case from scratch.
\end{principle}

\section{Four notions of exact solvability}

The phrase ``exactly solvable'' is not unique. These notes distinguish four useful levels of solvability.

{\small\setlength{\tabcolsep}{3pt}\begin{longtable}{L{0.15\textwidth}L{0.39\textwidth}L{0.31\textwidth}}
\caption{Different strengths of exact solvability.}\label{tab:solvability-levels}\\
\toprule
Level & Meaning & Representative models \\
\midrule
\endfirsthead
\toprule
Level & Meaning & Representative models \\
\midrule
\endhead
A. Full spectral diagonalization & One can obtain the full spectrum, quasiparticles, partition function, or correlation-function structure. & XY chain, TFIM, Kitaev chain, 2D Ising model \\
B. Free-fermion/\newline Pfaffian solvability & An interacting spin model or classical statistical model can be mapped by Jordan--Wigner, Majorana, Kasteleyn, or Pfaffian methods to a quadratic fermion problem. & TFIM, XY chain, Kitaev chain, dimer models, 2D Ising model, Kitaev honeycomb model in a fixed flux sector \\
C. Bethe/Yang--Baxter integrability & The model is not necessarily free, but it is solvable because of factorized scattering, commuting transfer matrices, or an underlying Yang--Baxter structure. & XXX, XXZ, XYZ, six-vertex, eight-vertex, 1D Hubbard model \\
D. Exact ground state or commuting-projector solvability & The ground state or the whole spectrum is controlled by projector, stabilizer, or frustration-free structure, although generic dynamics need not be integrable. & Toric code, cluster model, AKLT chain, RK dimer model \\
\bottomrule
\end{longtable}}

\section{Core equivalence map}

The central map of these notes is shown in \cref{fig:core-equivalence-map}.

\begin{figure}[htbp]
\centering
\small
\fbox{%
\begin{minipage}{0.96\linewidth}
\begin{align*}
&\text{TFIM}
\subset
\text{XY chain}
\xrightarrow{\text{JW}}
\text{Kitaev chain}
\xrightarrow{\text{continuum}}
\text{Majorana CFT}
\\[4pt]
&\text{SSH chain}
\subset
\text{two-band chiral fermion}
\sim
\substack{\text{Dirac Hamiltonian}\\\text{with winding number}}
\\[4pt]
&\text{1D TFIM}
\longleftrightarrow
\text{2D classical Ising}
\longleftrightarrow
\text{free Majorana/Pfaffian formulation}
\\[4pt]
&\text{XXZ chain}
\longleftrightarrow
\text{six-vertex model}
\longleftrightarrow
\text{Yang--Baxter equation}
\longleftrightarrow
\text{Bethe ansatz}
\\[4pt]
&\text{XYZ chain}
\longleftrightarrow
\text{eight-vertex model}
\\[4pt]
&\text{toric code}
\subset
\text{commuting-projector/quantum-double models}
\\[4pt]
&\text{Kitaev honeycomb model}
\longleftrightarrow
\substack{\text{free Majorana fermions}\\\text{in static }\ZZ_2\text{ gauge fields}}
\\[4pt]
&\text{AKLT/RK/cluster models}
\subset
\substack{\text{frustration-free}\\\text{exact-ground-state models}}
\end{align*}
\end{minipage}}
\caption{Core equivalence map among the main solvable structures.}
\label{fig:core-equivalence-map}
\end{figure}

\section{Primary mother structures}

The full notes are organized around five primary mother structures:
\begin{equation}
\boxed{
\begin{array}{c}
\text{free-fermion/BdG solvability},\\
\text{transfer matrices and quantum--classical correspondence},\\
\text{Bethe ansatz and Yang--Baxter integrability},\\
\text{commuting-projector/stabilizer solvability},\\
\text{Majorana fermions coupled to static }\ZZ_2\text{ gauge fields}.
\end{array}
}
\end{equation}

% -----------------------------------------------------------------------------
\chapter{Model Taxonomy by Equivalence Class}
% -----------------------------------------------------------------------------

\section{One-dimensional quantum models}

{\footnotesize\setlength{\tabcolsep}{3pt}\begin{longtable}{L{0.20\textwidth}L{0.24\textwidth}L{0.27\textwidth}L{0.16\textwidth}}
\caption{Mother classes, representative models, and solution methods for one-dimensional quantum many-body systems.}\label{tab:1d-quantum-taxonomy}\\
\toprule
Mother class & Representative models & Main solution method & Role in the notes \\
\midrule
\endfirsthead
\toprule
Mother class & Representative models & Main solution method & Role in the notes \\
\midrule
\endhead
Quadratic fermion / BdG & TFIM, transverse-field XY chain, XX chain, Kitaev chain & Jordan--Wigner transformation, Fourier transform, Bogoliubov rotation & Essential main line \\
Two-band chiral free fermion & SSH chain, dimerized hopping chain, dimerized XX chain & Bloch Hamiltonian and winding number & Essential main line \\
Bethe-ansatz integrability & XXX chain, XXZ chain & Coordinate Bethe ansatz and algebraic Bethe ansatz & Essential main line \\
Yang--Baxter / vertex-model integrability & XXZ $\leftrightarrow$ six-vertex, XYZ $\leftrightarrow$ eight-vertex & Commuting transfer matrices, $R$-matrix, Yang--Baxter equation & Middle-to-advanced main line \\
Nested Bethe ansatz & 1D Hubbard model & Lieb--Wu equations, nested Bethe ansatz & Advanced chapter \\
Long-range integrability & Haldane--Shastry model, Calogero--Sutherland model & Inverse-square interaction, fractional statistics, special conserved charges & Advanced chapter \\
Frustration-free / MPS solvability & AKLT chain, cluster model, Majumdar--Ghosh chain & Projector decomposition, exact MPS, stabilizer algebra & Supplementary but recommended \\
Continuum low-energy field theory & Luttinger liquid, Ising CFT, Majorana CFT, compact boson & Bosonization, CFT, scaling limit & Unifying low-energy chapter \\
\bottomrule
\end{longtable}}

\section{Two-dimensional classical models relevant to quantum many-body physics}

{\footnotesize\setlength{\tabcolsep}{3pt}\begin{longtable}{L{0.20\textwidth}L{0.23\textwidth}L{0.24\textwidth}L{0.18\textwidth}}
\caption{Two-dimensional classical exactly solved models and their relations to quantum many-body systems.}\label{tab:2d-classical-taxonomy}\\
\toprule
Mother class & Representative models & Main solution method & Relation to quantum systems \\
\midrule
\endfirsthead
\toprule
Mother class & Representative models & Main solution method & Relation to quantum systems \\
\midrule
\endhead
Transfer matrix / free fermion & 2D classical Ising model at zero field & Onsager solution, transfer matrix, fermionization & Anisotropic limit gives the 1D TFIM \\
Vertex integrability & Six-vertex model & Bethe ansatz, Yang--Baxter equation & Quantum Hamiltonian limit gives the XXZ chain \\
Vertex integrability & Eight-vertex model & Baxter solution, elliptic parametrization & Quantum Hamiltonian limit gives the XYZ chain \\
Pfaffian free-fermion structure & Classical dimer model & Kasteleyn orientation, Pfaffian determinant & Related to RK dimers and Majorana/Pfaffian methods \\
Loop/height/\newline Coulomb-gas representation & $O(n)$ loop model, SOS/RSOS models & Height mapping, Coulomb gas, integrable points & Related to CFT, minimal models, and critical lattice models \\
\bottomrule
\end{longtable}}

\section{Two-dimensional quantum models}

{\footnotesize\setlength{\tabcolsep}{3pt}\begin{longtable}{L{0.20\textwidth}L{0.23\textwidth}L{0.25\textwidth}L{0.17\textwidth}}
\caption{Classification of two-dimensional quantum exactly solvable models.}\label{tab:2d-quantum-taxonomy}\\
\toprule
Mother class & Representative models & Main solution method & Remarks \\
\midrule
\endfirsthead
\toprule
Mother class & Representative models & Main solution method & Remarks \\
\midrule
\endhead
Commuting projector & Toric code & Stabilizer diagonalization, local commuting projectors & Exactly solvable fixed point of topological order \\
Quantum double & $D(G)$ models, generalized toric codes & Group algebra, quantum double, anyon fusion & Non-Abelian generalization of the toric code structure \\
Majorana + static gauge field & Kitaev honeycomb model & Majorana representation, conserved $\ZZ_2$ gauge fields, free Majorana Hamiltonian in each flux sector & Central exactly solvable 2D spin-liquid model \\
Free Chern/BdG bands & Chern insulator, $p+\ii p$ superconductor & Band/BdG diagonalization, Chern number & Noninteracting many-body exact models \\
Rokhsar--Kivelson structure & Quantum dimer models at the RK point & Equal-weight ground state, stochastic/classical correspondence & Exact ground state; not necessarily full-spectrum solvable \\
Stabilizer SPT & 2D cluster states and related graph states & Commuting stabilizers, symmetry-protected order & Useful examples for quantum information and SPT phases \\
\bottomrule
\end{longtable}}

% -----------------------------------------------------------------------------
\part{Unified Mathematical Tools}
% -----------------------------------------------------------------------------

\chapter{\texorpdfstring{Spin-$\frac12$ Algebra and Lattice Conventions}{Spin-1/2 Algebra and Lattice Conventions}}

\section{Purpose of this chapter}

Before discussing Jordan--Wigner transformations, Bogoliubov rotations, transfer
matrices, or Bethe ansatz, we must fix a small number of elementary conventions.
For spin chains, many sign errors do not come from difficult physics; they come
from inconsistent choices of
\begin{equation}
    \sigma^\pm,
    \qquad
    |\uparrow\rangle, |\downarrow\rangle,
    \qquad
    \sigma^z \text{ versus } n,
    \qquad
    \text{open versus periodic boundary conditions}.
\end{equation}
This chapter fixes these conventions once and for all. The main convention used
throughout the free-fermion part of the notes is
\begin{equation}
    |0\rangle_j\equiv |\uparrow\rangle_j,
    \qquad
    |1\rangle_j\equiv |\downarrow\rangle_j,
    \qquad
    n_j=\frac{1-\sigma_j^z}{2}.
\end{equation}
Thus the fermionic occupied state corresponds to a down spin. With this choice,
the Jordan--Wigner convention in the next chapter is
\begin{equation}
    c_j=
    \left(\prod_{\ell=1}^{j-1}\sigma_\ell^z\right)\sigma_j^+,
    \qquad
    c_j^\dagger=
    \left(\prod_{\ell=1}^{j-1}\sigma_\ell^z\right)\sigma_j^- .
\end{equation}
The present chapter should therefore be read as the local algebraic dictionary on
which all later fermionization formulae are based.

\section{Local Hilbert space}

A single spin-$1/2$ degree of freedom has local Hilbert space
\begin{equation}
    \calH_j\cong \CC^2.
\end{equation}
We use the ordered basis
\begin{equation}
    |\uparrow\rangle
    =
    \begin{pmatrix}1\\0\end{pmatrix},
    \qquad
    |\downarrow\rangle
    =
    \begin{pmatrix}0\\1\end{pmatrix}.
\end{equation}
The Pauli matrices are
\begin{equation}
    \sigma^x=
    \begin{pmatrix}
        0&1\\
        1&0
    \end{pmatrix},
    \qquad
    \sigma^y=
    \begin{pmatrix}
        0&-\ii\\
        \ii&0
    \end{pmatrix},
    \qquad
    \sigma^z=
    \begin{pmatrix}
        1&0\\
        0&-1
    \end{pmatrix}.
\end{equation}
They act on the basis states as
\begin{equation}
    \sigma^z|\uparrow\rangle=|\uparrow\rangle,
    \qquad
    \sigma^z|\downarrow\rangle=-|\downarrow\rangle.
\end{equation}
For a chain of length $L$, the many-body Hilbert space is the ordered tensor
product
\begin{equation}
    \calH_L
    =
    \bigotimes_{j=1}^{L}\calH_j
    \cong
    (\CC^2)^{\otimes L}.
\end{equation}
A basis vector is written as
\begin{equation}
    |s_1,s_2,\ldots,s_L\rangle,
    \qquad
    s_j\in\{\uparrow,\downarrow\}.
\end{equation}
Equivalently, with the occupation convention introduced above,
\begin{equation}
    |s_1,s_2,\ldots,s_L\rangle
    \equiv
    |n_1,n_2,\ldots,n_L\rangle,
    \qquad
    n_j=\begin{cases}
        0, & s_j=\uparrow,\\
        1, & s_j=\downarrow.
    \end{cases}
\end{equation}

The operator acting as $\sigma^\alpha$ on site $j$ and as the identity elsewhere is
\begin{equation}
    \sigma_j^\alpha
    =
    \id_1\otimes\cdots\otimes\id_{j-1}
    \otimes\sigma^\alpha
    \otimes\id_{j+1}\otimes\cdots\otimes\id_L,
    \qquad
    \alpha=x,y,z.
\end{equation}
Operators supported on distinct sites commute:
\begin{equation}
    [\sigma_i^\alpha,\sigma_j^\beta]=0,
    \qquad
    i\neq j.
\end{equation}

\section{Pauli algebra and spin operators}

On one site, the Pauli matrices obey
\begin{equation}
\label{eq:pauli-product-algebra}
    \sigma^\alpha\sigma^\beta
    =
    \delta_{\alpha\beta}\id
    +
    \ii\sum_{\gamma}\epsilon_{\alpha\beta\gamma}\sigma^\gamma,
    \qquad
    \alpha,\beta\in\{x,y,z\}.
\end{equation}
Equivalently,
\begin{equation}
    [\sigma^\alpha,\sigma^\beta]
    =
    2\ii\sum_\gamma\epsilon_{\alpha\beta\gamma}\sigma^\gamma,
    \qquad
    \{\sigma^\alpha,\sigma^\beta\}
    =
    2\delta_{\alpha\beta}\id .
\end{equation}
The physical spin operators are
\begin{equation}
    S^\alpha_j=\frac12\sigma^\alpha_j.
\end{equation}
Therefore
\begin{equation}
    [S_i^\alpha,S_j^\beta]
    =
    \ii\delta_{ij}
    \sum_\gamma\epsilon_{\alpha\beta\gamma}S_j^\gamma.
\end{equation}
The total spin components are
\begin{equation}
    S_{\rm tot}^\alpha
    =
    \sum_{j=1}^{L}S_j^\alpha
    =
    \frac12\sum_{j=1}^{L}\sigma_j^\alpha.
\end{equation}

\begin{remark}[Pauli normalization versus spin normalization]
Many condensed-matter texts write Heisenberg-type Hamiltonians using $S^\alpha$,
whereas the free-fermion spin-chain literature often writes them using Pauli
matrices. Since $S^\alpha=\sigma^\alpha/2$, one has
\begin{equation}
    \bm{S}_j\cdot\bm{S}_{j+1}
    =
    \frac14
    \left(
        \sigma_j^x\sigma_{j+1}^x
        +
        \sigma_j^y\sigma_{j+1}^y
        +
        \sigma_j^z\sigma_{j+1}^z
    \right).
\end{equation}
Thus coupling constants differ by factors of $4$ when translating between the two
normalizations. Unless explicitly stated otherwise, the spin-chain Hamiltonians in
these notes are written in Pauli normalization.
\end{remark}

Useful trace identities are
\begin{equation}
    \Tr\,\sigma^\alpha=0,
    \qquad
    \Tr(\sigma^\alpha\sigma^\beta)=2\delta_{\alpha\beta},
\end{equation}
and
\begin{equation}
    \Tr(\sigma^\alpha\sigma^\beta\sigma^\gamma)
    =
    2\ii\epsilon_{\alpha\beta\gamma}.
\end{equation}
These identities are especially useful when deriving transfer matrices, density
matrices, and two-band Bloch Hamiltonians.

\section{Raising and lowering operators}

We define
\begin{equation}
\label{eq:sigma-pm-definition}
    \sigma^+
    =
    \frac{\sigma^x+\ii\sigma^y}{2}
    =
    \begin{pmatrix}
        0&1\\
        0&0
    \end{pmatrix},
    \qquad
    \sigma^-
    =
    \frac{\sigma^x-\ii\sigma^y}{2}
    =
    \begin{pmatrix}
        0&0\\
        1&0
    \end{pmatrix}.
\end{equation}
Hence
\begin{equation}
    \sigma^+|\downarrow\rangle=|\uparrow\rangle,
    \qquad
    \sigma^-|\uparrow\rangle=|\downarrow\rangle,
\end{equation}
while
\begin{equation}
    \sigma^+|\uparrow\rangle=0,
    \qquad
    \sigma^-|\downarrow\rangle=0.
\end{equation}
The inverse relations are
\begin{equation}
\label{eq:sigma-x-y-pm}
    \sigma^x=\sigma^++\sigma^-,
    \qquad
    \sigma^y=\ii(\sigma^- -\sigma^+).
\end{equation}
The local algebra is
\begin{equation}
    (\sigma^+)^2=(\sigma^-)^2=0,
    \qquad
    \{\sigma^+,\sigma^-\}=\id,
    \qquad
    [\sigma^+,\sigma^-]=\sigma^z,
\end{equation}
and
\begin{equation}
    [\sigma^z,\sigma^+]=2\sigma^+,
    \qquad
    [\sigma^z,\sigma^-]=-2\sigma^-.
\end{equation}
The two local projectors are
\begin{equation}
\label{eq:spin-projectors}
    \sigma^+\sigma^-=\frac{1+\sigma^z}{2}=|\uparrow\rangle\langle\uparrow|,
    \qquad
    \sigma^-\sigma^+=\frac{1-\sigma^z}{2}=|\downarrow\rangle\langle\downarrow|.
\end{equation}

\section{Occupation convention and hard-core bosons}

The Jordan--Wigner transformation used later identifies the spin-down state with
fermion occupation. We therefore define
\begin{equation}
\label{eq:spin-occupation-convention}
    |0\rangle_j\equiv |\uparrow\rangle_j,
    \qquad
    |1\rangle_j\equiv |\downarrow\rangle_j.
\end{equation}
The local number operator is
\begin{equation}
\label{eq:number-spin-convention}
    n_j
    =
    |\downarrow\rangle_j\langle\downarrow|
    =
    \sigma_j^-\sigma_j^+
    =
    \frac{1-\sigma_j^z}{2}.
\end{equation}
Thus
\begin{equation}
\label{eq:sigmaz-number-convention}
    \sigma_j^z=1-2n_j.
\end{equation}

It is convenient to introduce local hard-core boson operators
\begin{equation}
\label{eq:hard-core-boson-convention}
    b_j=\sigma_j^+,
    \qquad
    b_j^\dagger=\sigma_j^-.
\end{equation}
With the present convention, $b_j^\dagger$ creates a down spin, hence creates an
occupied site:
\begin{equation}
    b_j^\dagger |0\rangle_j=|1\rangle_j,
    \qquad
    b_j |1\rangle_j=|0\rangle_j.
\end{equation}
The number operator is
\begin{equation}
    n_j=b_j^\dagger b_j.
\end{equation}
On the same site, hard-core bosons obey
\begin{equation}
    b_j^2=(b_j^\dagger)^2=0,
    \qquad
    \{b_j,b_j^\dagger\}=1,
    \qquad
    [b_j,b_j^\dagger]=1-2n_j=\sigma_j^z.
\end{equation}
On different sites, however, they commute rather than anticommute:
\begin{equation}
    [b_i,b_j]=0,
    \qquad
    [b_i,b_j^\dagger]=0,
    \qquad
    [b_i^\dagger,b_j^\dagger]=0,
    \qquad
    i\neq j.
\end{equation}
This mixed behavior is the reason a nonlocal Jordan--Wigner string is needed to
convert hard-core bosons into canonical fermions.

\begin{remark}[Why $b=\sigma^+$ but $b^\dagger=\sigma^-$]
This convention may look reversed if one thinks of $\sigma^+$ as a creation
operator. Here the created object is not spin-up, but an occupied fermionic site.
Since occupation means spin down, the operator that creates occupation must send
$|\uparrow\rangle$ to $|\downarrow\rangle$. Therefore $b^\dagger=\sigma^-$.
\end{remark}

\section{Useful two-site identities}

The following identities are used repeatedly in the XY, XX, XXZ, and Kitaev-chain
chapters. For nearest-neighbour sites $j$ and $j+1$,
\begin{align}
\label{eq:xy-raising-lowering-identity}
    \sigma_j^x\sigma_{j+1}^x+
    \sigma_j^y\sigma_{j+1}^y
    &=
    2\left(
        \sigma_j^+\sigma_{j+1}^-
        +
        \sigma_j^-\sigma_{j+1}^+
    \right),
    \\
\label{eq:pairing-raising-lowering-identity}
    \sigma_j^x\sigma_{j+1}^x-
    \sigma_j^y\sigma_{j+1}^y
    &=
    2\left(
        \sigma_j^+\sigma_{j+1}^+
        +
        \sigma_j^-\sigma_{j+1}^-
    \right).
\end{align}
Equivalently,
\begin{align}
    \sigma_j^x\sigma_{j+1}^x
    &=
    \sigma_j^+\sigma_{j+1}^+
    +
    \sigma_j^+\sigma_{j+1}^-
    +
    \sigma_j^-\sigma_{j+1}^+
    +
    \sigma_j^-\sigma_{j+1}^-,
    \\
    \sigma_j^y\sigma_{j+1}^y
    &=
    -\sigma_j^+\sigma_{j+1}^+
    +
    \sigma_j^+\sigma_{j+1}^-
    +
    \sigma_j^-\sigma_{j+1}^+
    -
    \sigma_j^-\sigma_{j+1}^- .
\end{align}
In terms of the occupation convention,
\begin{equation}
\label{eq:zz-number-identity}
    \sigma_j^z\sigma_{j+1}^z
    =
    (1-2n_j)(1-2n_{j+1})
    =
    1-2n_j-2n_{j+1}+4n_jn_{j+1}.
\end{equation}
This last identity explains why an Ising interaction in the $z$ direction becomes a
density--density interaction after fermionization, whereas the $xx$ and $yy$ terms
of the XY chain become hopping and pairing terms.

\section{Lattice conventions}

Unless stated otherwise, a one-dimensional chain has sites
\begin{equation}
    j=1,2,\ldots,L,
\end{equation}
with lattice spacing set to
\begin{equation}
    a=1.
\end{equation}
A nearest-neighbour Hamiltonian is written using one of two bond sets:
\begin{align}
    \mathcal{B}_{\rm OBC}
    &=
    \{(j,j+1)\,|\,j=1,2,\ldots,L-1\},
    \\
    \mathcal{B}_{\rm PBC}
    &=
    \{(j,j+1)\,|\,j=1,2,\ldots,L\},
    \qquad
    L+1\equiv 1.
\end{align}
Thus
\begin{equation}
    \sum_{\langle j,j+1\rangle\in\mathcal{B}_{\rm OBC}}
    \quad\text{means}\quad
    \sum_{j=1}^{L-1},
\end{equation}
and
\begin{equation}
    \sum_{\langle j,j+1\rangle\in\mathcal{B}_{\rm PBC}}
    \quad\text{means}\quad
    \sum_{j=1}^{L}\quad\text{with}\quad L+1\equiv1.
\end{equation}
When no boundary condition is explicitly stated, formulae for local bulk
Hamiltonian densities should be read as open-chain formulae first. Periodic
boundary conditions must be treated with additional care after the Jordan--Wigner
transformation because the boundary spin term carries a fermion-parity-dependent
sign.

For translation-invariant periodic chains, the one-site translation operator $T$
is defined by
\begin{equation}
    T\sigma_j^\alpha T^{-1}=\sigma_{j+1}^\alpha,
    \qquad
    T^L=1.
\end{equation}
Momentum eigenstates satisfy
\begin{equation}
    T|\psi_k\rangle=\ee^{\ii k}|\psi_k\rangle,
\end{equation}
where the allowed values of $k$ depend on the boundary condition and, after the
Jordan--Wigner transformation, on the fermion-parity sector. The detailed Fourier
convention is fixed in the Fourier-transform chapter.

\begin{remark}[Ordering convention]
The Jordan--Wigner string depends on a one-dimensional ordering of sites. In these
notes the ordering is always
\begin{equation}
    1<2<\cdots<L.
\end{equation}
This is natural for one-dimensional chains. In higher-dimensional spin models, a
linear ordering can still be imposed, but it usually destroys locality. This is one
reason why two-dimensional exact solvability often requires other structures, such
as commuting projectors, transfer matrices, or Majorana fermions coupled to static
$\ZZ_2$ gauge fields.
\end{remark}

\section{Common spin Hamiltonians}

We now list the spin-chain Hamiltonians that will appear repeatedly in these
notes. The purpose here is not to solve them, but to fix signs and normalizations.
The default Pauli-normalized nearest-neighbour XYZ Hamiltonian is
\begin{equation}
\label{eq:xyz-general-convention}
    H_{\rm XYZ}
    =
    -\sum_{\langle j,j+1\rangle}
    \left(
        J_x\sigma_j^x\sigma_{j+1}^x
        +
        J_y\sigma_j^y\sigma_{j+1}^y
        +
        J_z\sigma_j^z\sigma_{j+1}^z
    \right)
    -
    h\sum_{j=1}^{L}\sigma_j^z .
\end{equation}
With this convention, positive $J_\alpha$ favors ferromagnetic alignment in the
$\alpha$ direction.

\subsection{Ising and transverse-field Ising chains}

The Ising interaction used in the free-fermion part of the notes is chosen along
$x$, while the transverse field is along $z$:
\begin{equation}
\label{eq:tfim-convention-chapter}
    H_{\rm TFIM}
    =
    -J\sum_{\langle j,j+1\rangle}\sigma_j^x\sigma_{j+1}^x
    -
    h\sum_{j=1}^{L}\sigma_j^z .
\end{equation}
The choice of $x$ as the ordering direction and $z$ as the transverse-field
direction is adapted to the occupation convention $n_j=(1-\sigma_j^z)/2$. It is
unitarily equivalent to the alternative textbook convention
$-J\sum_j\sigma_j^z\sigma_{j+1}^z-h\sum_j\sigma_j^x$ by a global spin rotation.

The purely classical Ising interaction is recovered by setting $h=0$:
\begin{equation}
    H_{\rm Ising}^{x}
    =
    -J\sum_{\langle j,j+1\rangle}\sigma_j^x\sigma_{j+1}^x .
\end{equation}
A longitudinal field along the ordering direction,
\begin{equation}
    -h_x\sum_j\sigma_j^x,
\end{equation}
can be added, but it destroys the standard free-fermion solvability of the
one-dimensional TFIM.

\subsection{XY and XX chains}

The transverse-field XY chain is the central free-fermion mother model of these
notes:
\begin{equation}
\label{eq:xy-convention-chapter}
    H_{\rm XY}
    =
    -\sum_{\langle j,j+1\rangle}
    \left[
        \frac{J(1+\gamma)}{2}\sigma_j^x\sigma_{j+1}^x
        +
        \frac{J(1-\gamma)}{2}\sigma_j^y\sigma_{j+1}^y
    \right]
    -
    h\sum_{j=1}^{L}\sigma_j^z .
\end{equation}
The parameter $\gamma$ is the XY anisotropy. Important limits are
\begin{align}
    \gamma=1
    &\quad\Longrightarrow\quad
    H_{\rm XY}=H_{\rm TFIM},
    \\
    \gamma=0
    &\quad\Longrightarrow\quad
    H_{\rm XX}
    =
    -\frac{J}{2}
    \sum_{\langle j,j+1\rangle}
    \left(
        \sigma_j^x\sigma_{j+1}^x
        +
        \sigma_j^y\sigma_{j+1}^y
    \right)
    -h\sum_j\sigma_j^z .
\end{align}
Using \cref{eq:xy-raising-lowering-identity}, the XX interaction can also be
written as
\begin{equation}
    H_{\rm XX}
    =
    -J\sum_{\langle j,j+1\rangle}
    \left(
        \sigma_j^+\sigma_{j+1}^-
        +
        \sigma_j^-\sigma_{j+1}^+
    \right)
    -h\sum_j\sigma_j^z .
\end{equation}
This form makes the conservation of total $z$-magnetization manifest.

\subsection{XXZ and Heisenberg chains}

In the Pauli normalization compatible with the above XX convention, we write
\begin{equation}
\label{eq:xxz-pauli-convention}
    H_{\rm XXZ}
    =
    -\frac{J}{2}
    \sum_{\langle j,j+1\rangle}
    \left(
        \sigma_j^x\sigma_{j+1}^x
        +
        \sigma_j^y\sigma_{j+1}^y
        +
        \Delta\sigma_j^z\sigma_{j+1}^z
    \right)
    -h\sum_j\sigma_j^z .
\end{equation}
The isotropic point $\Delta=1$ gives a Heisenberg chain in Pauli normalization:
\begin{equation}
    H_{\rm Heis}
    =
    -\frac{J}{2}
    \sum_{\langle j,j+1\rangle}
    \bm{\sigma}_j\cdot\bm{\sigma}_{j+1}
    -h\sum_j\sigma_j^z .
\end{equation}
In spin-operator normalization, the same interaction is written as
\begin{equation}
    H_{\rm Heis}^{S}
    =
    \mathcal{J}\sum_{\langle j,j+1\rangle}
    \bm{S}_j\cdot\bm{S}_{j+1}
    -h_S\sum_j S_j^z.
\end{equation}
The conversion is
\begin{equation}
    \bm{S}_j\cdot\bm{S}_{j+1}
    =
    \frac14\bm{\sigma}_j\cdot\bm{\sigma}_{j+1},
    \qquad
    S_j^z=\frac12\sigma_j^z.
\end{equation}
The Bethe-ansatz chapters may use the conventional antiferromagnetic sign
$\mathcal{J}>0$; when this happens, the sign and normalization will be stated
explicitly.

\subsection{Bond-dependent Kitaev-type spin Hamiltonians}

In two dimensions, the Kitaev honeycomb model uses bond-dependent Ising
interactions:
\begin{equation}
\label{eq:kitaev-honeycomb-convention}
    H_{\rm Kitaev}
    =
    -J_x\sum_{\langle ij\rangle_x}\sigma_i^x\sigma_j^x
    -J_y\sum_{\langle ij\rangle_y}\sigma_i^y\sigma_j^y
    -J_z\sum_{\langle ij\rangle_z}\sigma_i^z\sigma_j^z.
\end{equation}
Here $\langle ij\rangle_\alpha$ denotes a nearest-neighbour bond of type $\alpha$.
This model is not solved by the one-dimensional Jordan--Wigner transformation.
Instead, it becomes exactly solvable through a Majorana representation and static
$\ZZ_2$ gauge fields.

\section{Symmetries and conserved quantities}

Symmetry is one of the main tools for block diagonalization. We collect the most
important symmetries here.

\subsection{Translation symmetry}

For periodic boundary conditions, a translation-invariant Hamiltonian satisfies
\begin{equation}
    [H,T]=0.
\end{equation}
The Hilbert space can then be decomposed into momentum sectors labelled by the
eigenvalues $\ee^{\ii k}$ of $T$. Translation symmetry is essential for the Fourier
transform and the reduction of quadratic Hamiltonians into independent momentum
blocks.

\subsection{Fermion parity and Ising spin flip}

Define the global $z$-parity operator
\begin{equation}
\label{eq:global-z-parity}
    P_z
    =
    \prod_{j=1}^{L}\sigma_j^z.
\end{equation}
It satisfies
\begin{equation}
    P_z^2=1,
    \qquad
    P_z\sigma_j^xP_z=-\sigma_j^x,
    \qquad
    P_z\sigma_j^yP_z=-\sigma_j^y,
    \qquad
    P_z\sigma_j^zP_z=\sigma_j^z.
\end{equation}
Therefore the TFIM and the XY chain commute with $P_z$:
\begin{equation}
    [H_{\rm TFIM},P_z]=0,
    \qquad
    [H_{\rm XY},P_z]=0.
\end{equation}
Under the occupation convention \cref{eq:number-spin-convention},
\begin{equation}
    P_z
    =
    \prod_{j=1}^{L}(1-2n_j)
    =
    (-1)^N,
    \qquad
    N=\sum_{j=1}^{L}n_j.
\end{equation}
Thus the global Ising spin-flip symmetry of the spin chain becomes fermion parity
after the Jordan--Wigner transformation.

\subsection{\texorpdfstring{$U(1)$}{U(1)} spin-rotation symmetry}

The total $z$-magnetization is
\begin{equation}
    M_z=\sum_{j=1}^{L}\sigma_j^z=L-2N.
\end{equation}
A Hamiltonian has $U(1)$ spin-rotation symmetry about the $z$ axis if
\begin{equation}
    [H,M_z]=0.
\end{equation}
The XX, XXZ, and Heisenberg chains have this symmetry. Equivalently, they are
invariant under
\begin{equation}
    U_z(\theta)
    =
    \exp\!\left(-\frac{\ii\theta}{2}\sum_j\sigma_j^z\right),
\end{equation}
which acts as
\begin{equation}
    U_z(\theta)\sigma_j^\pm U_z(\theta)^{-1}
    =
    \ee^{\mp\ii\theta}\sigma_j^\pm.
\end{equation}
The anisotropic XY chain with $\gamma\neq0$ does not conserve $M_z$ because
\cref{eq:pairing-raising-lowering-identity} contains pair creation and pair
annihilation terms. Nevertheless, it still conserves $P_z=(-1)^N$.

\subsection{Full spin-rotation symmetry}

The isotropic Heisenberg chain at zero field has full $SU(2)$ symmetry:
\begin{equation}
    [H_{\rm Heis},S_{\rm tot}^\alpha]=0,
    \qquad
    \alpha=x,y,z.
\end{equation}
An external field along $z$ breaks this to the $U(1)$ subgroup generated by
$S_{\rm tot}^z$.

\subsection{Boundary-sector subtlety after Jordan--Wigner}

For an open chain, the Jordan--Wigner transformation maps local nearest-neighbour
bulk terms to local fermion bilinears without any boundary ambiguity. For a
periodic spin chain, the boundary term connecting $L$ and $1$ depends on the parity
sector
\begin{equation}
    P_z=(-1)^N=p,
    \qquad
    p=\pm1.
\end{equation}
With the Jordan--Wigner convention used in these notes, the effective fermionic
boundary condition in a fixed parity sector is
\begin{equation}
\label{eq:chapter-effective-boundary-condition}
    c_{L+1}=-p\,c_1.
\end{equation}
Thus
\begin{equation}
    p=+1
    \quad\Longrightarrow\quad
    c_{L+1}=-c_1,
    \qquad
    p=-1
    \quad\Longrightarrow\quad
    c_{L+1}=c_1.
\end{equation}
That is, even fermion parity corresponds to anti-periodic fermions, while odd
fermion parity corresponds to periodic fermions. The derivation is given in the
Jordan--Wigner chapter.

\section{Summary of conventions}

The conventions fixed in this chapter are summarized as follows:
\begin{equation}
    |0\rangle_j=|\uparrow\rangle_j,
    \qquad
    |1\rangle_j=|\downarrow\rangle_j,
    \qquad
    n_j=\frac{1-\sigma_j^z}{2}.
\end{equation}
The raising and lowering operators are
\begin{equation}
    \sigma^+=\frac{\sigma^x+\ii\sigma^y}{2},
    \qquad
    \sigma^-=\frac{\sigma^x-\ii\sigma^y}{2},
\end{equation}
with
\begin{equation}
    \sigma^+|\downarrow\rangle=|\uparrow\rangle,
    \qquad
    \sigma^-|\uparrow\rangle=|\downarrow\rangle.
\end{equation}
The hard-core boson convention is
\begin{equation}
    b_j=\sigma_j^+,
    \qquad
    b_j^\dagger=\sigma_j^-,
    \qquad
    n_j=b_j^\dagger b_j.
\end{equation}
The principal free-fermion mother model is the transverse-field XY chain
\begin{equation}
    H_{\rm XY}
    =
    -\sum_{\langle j,j+1\rangle}
    \left[
        \frac{J(1+\gamma)}{2}\sigma_j^x\sigma_{j+1}^x
        +
        \frac{J(1-\gamma)}{2}\sigma_j^y\sigma_{j+1}^y
    \right]
    -h\sum_j\sigma_j^z,
\end{equation}
which contains the TFIM at $\gamma=1$ and the XX chain at $\gamma=0$. The global
operator
\begin{equation}
    P_z=\prod_j\sigma_j^z=(-1)^N
\end{equation}
will become fermion parity in the Jordan--Wigner representation. These conventions
are the reference point for all later equivalence maps.


\chapter{Jordan--Wigner Transformation}

\section{Definition}

The Jordan--Wigner transformation maps a one-dimensional spin-$1/2$ chain to a
chain of spinless fermions. The essential difficulty is that spin operators on
different sites commute, whereas fermionic operators on different sites must
anticommute. The nonlocal Jordan--Wigner string supplies precisely the required
fermionic signs.

Throughout this chapter we use the convention
\begin{equation}
    |0\rangle_j \equiv |\uparrow\rangle_j,
    \qquad
    |1\rangle_j \equiv |\downarrow\rangle_j .
\end{equation}
Thus the occupied fermionic state corresponds to spin down. The Pauli matrices are
\begin{equation}
    \sigma^x =
    \begin{pmatrix}
        0 & 1\\
        1 & 0
    \end{pmatrix},
    \qquad
    \sigma^y =
    \begin{pmatrix}
        0 & -\ii\\
        \ii & 0
    \end{pmatrix},
    \qquad
    \sigma^z =
    \begin{pmatrix}
        1 & 0\\
        0 & -1
    \end{pmatrix}.
\end{equation}
The spin raising and lowering operators are
\begin{equation}
    \sigma^+
    =
    \frac{\sigma^x+\ii\sigma^y}{2}
    =
    \begin{pmatrix}
        0 & 1\\
        0 & 0
    \end{pmatrix},
    \qquad
    \sigma^-
    =
    \frac{\sigma^x-\ii\sigma^y}{2}
    =
    \begin{pmatrix}
        0 & 0\\
        1 & 0
    \end{pmatrix}.
\end{equation}
They act as
\begin{equation}
    \sigma^+|\downarrow\rangle = |\uparrow\rangle,
    \qquad
    \sigma^-|\uparrow\rangle = |\downarrow\rangle .
\end{equation}

With the convention that spin down is occupied, it is natural to identify the local
hard-core boson operators as
\begin{equation}
    b_j=\sigma_j^+,
    \qquad
    b_j^\dagger=\sigma_j^- .
\end{equation}
The local number operator is then
\begin{equation}
    n_j
    =
    b_j^\dagger b_j
    =
    \sigma_j^-\sigma_j^+
    =
    \frac{1-\sigma_j^z}{2}.
\end{equation}
Equivalently,
\begin{equation}
    \sigma_j^z = 1-2n_j.
\end{equation}

We define the Jordan--Wigner string by
\begin{equation}
\label{eq:jw-string-definition}
    K_j
    =
    \prod_{\ell=1}^{j-1}\sigma_\ell^z
    =
    \prod_{\ell=1}^{j-1}(1-2n_\ell).
\end{equation}
Since \(n_\ell\) has eigenvalues \(0\) and \(1\), we may also write
\begin{equation}
    K_j
    =
    \exp\!\left(
        \ii\pi\sum_{\ell=1}^{j-1}n_\ell
    \right).
\end{equation}
This exponential form is often useful because
\begin{equation}
    \ee^{\ii\pi n_\ell}=1-2n_\ell=\sigma_\ell^z.
\end{equation}

The Jordan--Wigner transformation used in these notes is
\begin{equation}
\label{eq:jw-definition-product}
    c_j
    =
    \left(\prod_{\ell=1}^{j-1}\sigma_\ell^z\right)\sigma_j^+,
    \qquad
    c_j^\dagger
    =
    \left(\prod_{\ell=1}^{j-1}\sigma_\ell^z\right)\sigma_j^- .
\end{equation}
Equivalently,
\begin{equation}
\label{eq:jw-definition-compact}
    c_j = K_j\sigma_j^+ = K_j b_j,
    \qquad
    c_j^\dagger = K_j\sigma_j^- = K_j b_j^\dagger .
\end{equation}
The inverse transformation is
\begin{equation}
\label{eq:jw-inverse}
    \sigma_j^+ = K_j c_j,
    \qquad
    \sigma_j^- = K_j c_j^\dagger,
\end{equation}
or, equivalently,
\begin{equation}
    b_j=K_jc_j,
    \qquad
    b_j^\dagger=K_jc_j^\dagger .
\end{equation}

The string satisfies
\begin{equation}
    K_j^\dagger=K_j,
    \qquad
    K_j^2=1,
    \qquad
    K_{j+1}=K_j\sigma_j^z=K_j(1-2n_j).
\end{equation}
Because \(K_j\) only acts on sites strictly to the left of \(j\), it commutes with all
operators on site \(j\):
\begin{equation}
    [K_j,\sigma_j^\pm]=0,
    \qquad
    [K_j,\sigma_j^z]=0,
    \qquad
    [K_j,c_j]=0,
    \qquad
    [K_j,c_j^\dagger]=0.
\end{equation}

Using the inverse transformation, the spin operators become
\begin{equation}
\label{eq:spin-operators-jw}
    \sigma_j^x
    =
    K_j(c_j^\dagger+c_j),
    \qquad
    \sigma_j^y
    =
    K_j\,\ii(c_j^\dagger-c_j),
    \qquad
    \sigma_j^z
    =
    1-2c_j^\dagger c_j .
\end{equation}

\begin{remark}
A common source of sign confusion is the alternative string
\(\prod_{\ell<j}(-\sigma_\ell^z)\). This differs from the convention used here by
a site-dependent sign:
\begin{equation}
    \prod_{\ell=1}^{j-1}(-\sigma_\ell^z)
    =
    (-1)^{j-1}
    \prod_{\ell=1}^{j-1}\sigma_\ell^z .
\end{equation}
It therefore corresponds to a staggered fermion redefinition
\begin{equation}
    c_j^{\mathrm{alt}} = (-1)^{j-1}c_j .
\end{equation}
In this chapter and in the rest of these notes, we consistently use
\(K_j=\prod_{\ell<j}\sigma_\ell^z\).
\end{remark}

\section{Canonical anticommutation relations}

We now prove that the operators defined by the Jordan--Wigner transformation obey
the canonical fermionic anticommutation relations
\begin{equation}
\label{eq:canonical-fermion-algebra}
    \{c_i,c_j^\dagger\}=\delta_{ij},
    \qquad
    \{c_i,c_j\}=0,
    \qquad
    \{c_i^\dagger,c_j^\dagger\}=0.
\end{equation}

First consider the same site. Since \(K_j\) commutes with \(\sigma_j^\pm\) and
\(K_j^2=1\),
\begin{align}
    c_jc_j^\dagger
    &=
    K_j\sigma_j^+K_j\sigma_j^-
    =
    \sigma_j^+\sigma_j^-,
    \\
    c_j^\dagger c_j
    &=
    K_j\sigma_j^-K_j\sigma_j^+
    =
    \sigma_j^-\sigma_j^+.
\end{align}
Therefore
\begin{equation}
    \{c_j,c_j^\dagger\}
    =
    \sigma_j^+\sigma_j^-
    +
    \sigma_j^-\sigma_j^+
    =
    1.
\end{equation}
Also,
\begin{equation}
    c_j^2=(\sigma_j^+)^2=0,
    \qquad
    (c_j^\dagger)^2=(\sigma_j^-)^2=0.
\end{equation}

Now take \(i<j\). Define
\begin{equation}
    R_{ij}
    =
    \prod_{\ell=i+1}^{j-1}\sigma_\ell^z,
\end{equation}
so that
\begin{equation}
    K_j = K_i\sigma_i^zR_{ij}.
\end{equation}
All operators at different spin sites commute. The only nontrivial sign comes from
the factor \(\sigma_i^z\), which is present in the string of the operator at site \(j\).

For two annihilation operators,
\begin{align}
    c_i c_j
    &=
    K_i\sigma_i^+K_j\sigma_j^+
    \nonumber\\
    &=
    K_i\sigma_i^+K_i\sigma_i^zR_{ij}\sigma_j^+
    \nonumber\\
    &=
    \sigma_i^+\sigma_i^zR_{ij}\sigma_j^+
    \nonumber\\
    &=
    -\sigma_i^+R_{ij}\sigma_j^+,
\end{align}
where we used
\begin{equation}
    \sigma_i^+\sigma_i^z=-\sigma_i^+.
\end{equation}
On the other hand,
\begin{align}
    c_j c_i
    &=
    K_j\sigma_j^+K_i\sigma_i^+
    \nonumber\\
    &=
    K_i\sigma_i^zR_{ij}\sigma_j^+K_i\sigma_i^+
    \nonumber\\
    &=
    \sigma_i^z\sigma_i^+R_{ij}\sigma_j^+
    \nonumber\\
    &=
    +\sigma_i^+R_{ij}\sigma_j^+,
\end{align}
where
\begin{equation}
    \sigma_i^z\sigma_i^+=+\sigma_i^+.
\end{equation}
Hence
\begin{equation}
    c_ic_j=-c_jc_i,
    \qquad i\neq j.
\end{equation}
Thus
\begin{equation}
    \{c_i,c_j\}=0.
\end{equation}

Similarly,
\begin{align}
    c_i c_j^\dagger
    &=
    K_i\sigma_i^+K_j\sigma_j^-
    \nonumber\\
    &=
    \sigma_i^+\sigma_i^zR_{ij}\sigma_j^-
    \nonumber\\
    &=
    -\sigma_i^+R_{ij}\sigma_j^-,
\end{align}
while
\begin{align}
    c_j^\dagger c_i
    &=
    K_j\sigma_j^-K_i\sigma_i^+
    \nonumber\\
    &=
    \sigma_i^z\sigma_i^+R_{ij}\sigma_j^-
    \nonumber\\
    &=
    +\sigma_i^+R_{ij}\sigma_j^-.
\end{align}
Therefore
\begin{equation}
    c_ic_j^\dagger=-c_j^\dagger c_i,
    \qquad i\neq j.
\end{equation}
Together with the same-site result, this proves
\begin{equation}
    \{c_i,c_j^\dagger\}=\delta_{ij}.
\end{equation}

Finally, for two creation operators,
\begin{align}
    c_i^\dagger c_j^\dagger
    &=
    K_i\sigma_i^-K_j\sigma_j^-
    \nonumber\\
    &=
    \sigma_i^-\sigma_i^zR_{ij}\sigma_j^-
    \nonumber\\
    &=
    +\sigma_i^-R_{ij}\sigma_j^-,
\end{align}
where
\begin{equation}
    \sigma_i^-\sigma_i^z=+\sigma_i^-,
\end{equation}
whereas
\begin{align}
    c_j^\dagger c_i^\dagger
    &=
    K_j\sigma_j^-K_i\sigma_i^-
    \nonumber\\
    &=
    \sigma_i^z\sigma_i^-R_{ij}\sigma_j^-
    \nonumber\\
    &=
    -\sigma_i^-R_{ij}\sigma_j^-.
\end{align}
Hence
\begin{equation}
    c_i^\dagger c_j^\dagger=-c_j^\dagger c_i^\dagger,
    \qquad i\neq j.
\end{equation}
This completes the proof of the canonical fermionic anticommutation relations.

\section{Boundary conditions and fermion parity}

For an open chain, the Jordan--Wigner transformation is unambiguous: the string
always extends from the left boundary to the site immediately before \(j\). No operator
wraps around the chain.

For a periodic spin chain, however, a boundary term connecting site \(L\) to site \(1\)
carries a fermion-parity-dependent sign. The total fermion number and fermion parity
are
\begin{equation}
    N=\sum_{j=1}^L n_j,
    \qquad
    P_F=(-1)^N.
\end{equation}
Since \(1-2n_j=\sigma_j^z\), we have
\begin{equation}
\label{eq:fermion-parity-spin}
    P_F
    =
    \prod_{j=1}^L(1-2n_j)
    =
    \prod_{j=1}^L\sigma_j^z.
\end{equation}
The parity operator satisfies
\begin{equation}
    P_F^2=1,
    \qquad
    P_F c_j P_F=-c_j,
    \qquad
    P_F c_j^\dagger P_F=-c_j^\dagger.
\end{equation}

The formal Jordan--Wigner string at site \(L+1\) is
\begin{equation}
    K_{L+1}
    =
    \prod_{j=1}^L\sigma_j^z
    =
    P_F.
\end{equation}
Thus, if the spin variables are periodic,
\begin{equation}
    \sigma_{L+1}^\pm\equiv\sigma_1^\pm,
\end{equation}
the formal extension of the Jordan--Wigner definition gives
\begin{equation}
\label{eq:formal-boundary-fermion}
    c_{L+1}^{\mathrm{form}}
    =
    K_{L+1}\sigma_{L+1}^+
    =
    P_F\sigma_1^+
    =
    P_Fc_1 .
\end{equation}
This equation is a formal extension of the operator definition. It should not yet
be identified with the effective boundary condition of a quadratic fermion
Hamiltonian. The effective boundary condition is obtained only after the boundary
spin bilinears are rewritten in the same algebraic form as the bulk nearest-neighbour
bilinears.

For example,
\begin{align}
    b_L^\dagger b_1
    &=
    K_Lc_L^\dagger c_1
    \nonumber\\
    &=
    \left(\prod_{\ell=1}^{L-1}\sigma_\ell^z\right)c_L^\dagger c_1
    \nonumber\\
    &=
    P_F\sigma_L^z c_L^\dagger c_1
    \nonumber\\
    &=
    -P_F c_L^\dagger c_1.
\end{align}
Similarly,
\begin{equation}
    b_L^\dagger b_1^\dagger
    =
    -P_F c_L^\dagger c_1^\dagger.
\end{equation}
Therefore, within a fixed fermion-parity sector \(P_F=p=\pm1\), rewriting the boundary
term in the same form as the bulk terms is equivalent to imposing the effective
fermionic boundary condition
\begin{equation}
\label{eq:effective-jw-boundary-condition}
    c_{L+1}=-p\,c_1.
\end{equation}
Consequently,
\begin{equation}
    p=+1
    \quad\Longrightarrow\quad
    c_{L+1}=-c_1,
\end{equation}
so the even-parity sector corresponds to anti-periodic fermions, whereas
\begin{equation}
    p=-1
    \quad\Longrightarrow\quad
    c_{L+1}=c_1,
\end{equation}
so the odd-parity sector corresponds to periodic fermions.

\section{Useful Jordan--Wigner identities}

This section collects the identities most frequently used when translating
one-dimensional spin-chain Hamiltonians into fermionic language. All formulae below
use the convention
\begin{equation}
    K_j=\prod_{\ell=1}^{j-1}\sigma_\ell^z
    =
    \prod_{\ell=1}^{j-1}(1-2n_\ell),
    \qquad
    c_j=K_j\sigma_j^+,
    \qquad
    c_j^\dagger=K_j\sigma_j^-.
\end{equation}

\subsection*{Local spin, hard-core boson, and number operators}

The hard-core boson convention is
\begin{equation}
    \sigma_j^+=b_j,
    \qquad
    \sigma_j^-=b_j^\dagger,
    \qquad
    n_j=b_j^\dagger b_j=c_j^\dagger c_j.
\end{equation}
Therefore
\begin{equation}
    \sigma_j^x=b_j^\dagger+b_j,
    \qquad
    \sigma_j^y=\ii(b_j^\dagger-b_j),
    \qquad
    \sigma_j^z=1-2b_j^\dagger b_j=1-2c_j^\dagger c_j.
\end{equation}
The local hard-core algebra is
\begin{equation}
    \{b_j,b_j^\dagger\}=1,
    \qquad
    \{b_j,b_j\}=0,
    \qquad
    \{b_j^\dagger,b_j^\dagger\}=0,
\end{equation}
while for \(i\neq j\),
\begin{equation}
    [b_i,b_j]=0,
    \qquad
    [b_i,b_j^\dagger]=0,
    \qquad
    [b_i^\dagger,b_j^\dagger]=0.
\end{equation}
The fermions satisfy
\begin{equation}
    \{c_i,c_j^\dagger\}=\delta_{ij},
    \qquad
    \{c_i,c_j\}=0,
    \qquad
    \{c_i^\dagger,c_j^\dagger\}=0.
\end{equation}

\subsection*{String identities}

The string identities most often used are
\begin{equation}
    K_j^2=1,
    \qquad
    K_j^\dagger=K_j,
    \qquad
    K_{j+1}=K_j(1-2n_j).
\end{equation}
Moreover,
\begin{equation}
    c_j^\dagger(1-2n_j)=c_j^\dagger,
    \qquad
    c_j(1-2n_j)=-c_j,
\end{equation}
and
\begin{equation}
    (1-2n_j)c_j^\dagger=-c_j^\dagger,
    \qquad
    (1-2n_j)c_j=c_j.
\end{equation}
These four identities are the main source of the signs in nearest-neighbour JW
formulae.

\subsection*{Hard-core boson bilinears}

Using \(b_j=K_jc_j\), \(b_j^\dagger=K_jc_j^\dagger\), and
\(K_{j+1}=K_j(1-2n_j)\), one obtains
\begin{align}
    b_j^\dagger b_j
    &=
    c_j^\dagger c_j
    =
    n_j,
    \\
    b_j^\dagger b_{j+1}^\dagger
    &=
    c_j^\dagger(1-2n_j)c_{j+1}^\dagger
    =
    c_j^\dagger c_{j+1}^\dagger,
    \\
    b_j^\dagger b_{j+1}
    &=
    c_j^\dagger(1-2n_j)c_{j+1}
    =
    c_j^\dagger c_{j+1},
    \\
    b_j b_{j+1}
    &=
    c_j(1-2n_j)c_{j+1}
    =
    -c_jc_{j+1},
    \\
    b_j b_{j+1}^\dagger
    &=
    c_j(1-2n_j)c_{j+1}^\dagger
    =
    -c_jc_{j+1}^\dagger.
\end{align}
Equivalently, after reordering fermions on different sites,
\begin{equation}
    -c_jc_{j+1}^\dagger
    =
    c_{j+1}^\dagger c_j,
    \qquad
    -c_jc_{j+1}
    =
    c_{j+1}c_j.
\end{equation}

\subsection*{Spin operators in fermions}

The single-site spin operators are
\begin{equation}
    \sigma_j^x
    =
    K_j(c_j^\dagger+c_j),
    \qquad
    \sigma_j^y
    =
    K_j\,\ii(c_j^\dagger-c_j),
    \qquad
    \sigma_j^z
    =
    1-2n_j.
\end{equation}

The nearest-neighbour \(xx\) bilinear is
\begin{align}
    \sigma_j^x\sigma_{j+1}^x
    &=
    (b_j^\dagger+b_j)(b_{j+1}^\dagger+b_{j+1})
    \nonumber\\
    &=
    c_j^\dagger c_{j+1}^\dagger
    +
    c_j^\dagger c_{j+1}
    +
    c_{j+1}^\dagger c_j
    +
    c_{j+1}c_j
    \nonumber\\
    &=
    c_j^\dagger c_{j+1}^\dagger
    +
    c_j^\dagger c_{j+1}
    +
    \text{H.c.}.
\end{align}
The nearest-neighbour \(yy\) bilinear is
\begin{align}
    \sigma_j^y\sigma_{j+1}^y
    &=
    - (b_j^\dagger-b_j)(b_{j+1}^\dagger-b_{j+1})
    \nonumber\\
    &=
    -c_j^\dagger c_{j+1}^\dagger
    +
    c_j^\dagger c_{j+1}
    +
    c_{j+1}^\dagger c_j
    -
    c_{j+1}c_j
    \nonumber\\
    &=
    -\left(
        c_j^\dagger c_{j+1}^\dagger
        -
        c_j^\dagger c_{j+1}
        +
        \text{H.c.}
    \right).
\end{align}
The nearest-neighbour \(zz\) bilinear is
\begin{align}
    \sigma_j^z\sigma_{j+1}^z
    &=
    (1-2n_j)(1-2n_{j+1})
    \nonumber\\
    &=
    1-2n_j-2n_{j+1}+4n_jn_{j+1}.
\end{align}

Useful linear combinations are
\begin{align}
    \sigma_j^x\sigma_{j+1}^x
    +
    \sigma_j^y\sigma_{j+1}^y
    &=
    2\left(
        c_j^\dagger c_{j+1}
        +
        c_{j+1}^\dagger c_j
    \right),
    \\
    \sigma_j^x\sigma_{j+1}^x
    -
    \sigma_j^y\sigma_{j+1}^y
    &=
    2\left(
        c_j^\dagger c_{j+1}^\dagger
        +
        c_{j+1}c_j
    \right).
\end{align}

\subsection*{Tensor-product representation}

For direct matrix construction on a finite chain of length \(L\), the JW operators are
\begin{equation}
    c_j
    =
    \left(
        \bigotimes_{\ell=1}^{j-1}\sigma_\ell^z
    \right)
    \otimes
    \sigma_j^+
    \otimes
    \left(
        \bigotimes_{\ell=j+1}^{L}\id_\ell
    \right),
\end{equation}
and
\begin{equation}
    c_j^\dagger
    =
    \left(
        \bigotimes_{\ell=1}^{j-1}\sigma_\ell^z
    \right)
    \otimes
    \sigma_j^-
    \otimes
    \left(
        \bigotimes_{\ell=j+1}^{L}\id_\ell
    \right).
\end{equation}
For example, for \(L=4\),
\begin{align}
    c_1
    &=
    \sigma_1^+\otimes\id_2\otimes\id_3\otimes\id_4,
    \\
    c_2
    &=
    \sigma_1^z\otimes\sigma_2^+\otimes\id_3\otimes\id_4,
    \\
    c_3
    &=
    \sigma_1^z\otimes\sigma_2^z\otimes\sigma_3^+\otimes\id_4,
    \\
    c_4
    &=
    \sigma_1^z\otimes\sigma_2^z\otimes\sigma_3^z\otimes\sigma_4^+.
\end{align}
The creation operators are obtained by replacing \(\sigma_j^+\) with \(\sigma_j^-\).

\subsection*{Compact summary}

The essential dictionary is
\begin{equation}
    n_j
    =
    c_j^\dagger c_j
    =
    \frac{1-\sigma_j^z}{2},
    \qquad
    \sigma_j^z=1-2n_j,
\end{equation}
\begin{equation}
    K_j
    =
    \prod_{\ell=1}^{j-1}\sigma_\ell^z
    =
    \prod_{\ell=1}^{j-1}(1-2n_\ell),
\end{equation}
\begin{equation}
    c_j=K_j\sigma_j^+,
    \qquad
    c_j^\dagger=K_j\sigma_j^-,
\end{equation}
\begin{equation}
    \sigma_j^x=K_j(c_j^\dagger+c_j),
    \qquad
    \sigma_j^y=K_j\,\ii(c_j^\dagger-c_j),
    \qquad
    \sigma_j^z=1-2c_j^\dagger c_j,
\end{equation}
and
\begin{equation}
    \{c_i,c_j^\dagger\}=\delta_{ij},
    \qquad
    \{c_i,c_j\}=0,
    \qquad
    \{c_i^\dagger,c_j^\dagger\}=0.
\end{equation}
These identities are the basic toolkit for translating one-dimensional spin chains
into fermionic language.

\chapter{Fourier Transform and Momentum-Space Blocks}
\label{ch:fourier-transform-momentum-blocks}

\section{Purpose of this chapter}

After the Jordan--Wigner transformation, the spin chains of interest become
fermionic lattice Hamiltonians.  For open chains, the resulting quadratic
Hamiltonians can be diagonalized directly in real space.  For periodic and
translation-invariant chains, however, the natural variables are momentum modes.
The purpose of this chapter is to fix the Fourier-transform convention used in
these notes and to explain why quadratic Hamiltonians decompose into independent
momentum-space blocks.

There are three convention-dependent points that must be fixed carefully:
\begin{enumerate}[label=(\roman*)]
    \item the sign convention in the Fourier transform;
    \item the allowed momenta, which depend on the fermionic boundary condition;
    \item the harmless but useful global phase that may be included in the
    Fourier transform.
\end{enumerate}
The third point is especially convenient in superconducting chains, where the
pairing gap may be complex.  A uniform phase of the pairing gap can often be moved
between the Hamiltonian and the definition of the momentum-space fermions by a
constant gauge rotation.

\section{Discrete Fourier transform with a boundary twist}

Consider a chain of length $L$ with fermionic operators $c_j,c_j^\dagger$,
$j=1,\ldots,L$, satisfying
\begin{equation}
    \{c_i,c_j^\dagger\}=\delta_{ij},
    \qquad
    \{c_i,c_j\}=\{c_i^\dagger,c_j^\dagger\}=0.
\end{equation}
We impose the general boundary condition
\begin{equation}
\label{eq:general-fermion-boundary-condition}
    c_{j+L}=\ee^{\ii\vartheta}c_j,
    \qquad
    \vartheta\in[0,2\pi).
\end{equation}
The phase $\vartheta$ is a boundary twist.  Periodic boundary conditions
correspond to $\vartheta=0$, while anti-periodic boundary conditions correspond
to $\vartheta=\pi$.

The allowed momenta are the solutions of
\begin{equation}
    \ee^{\ii k L}=\ee^{\ii\vartheta}.
\end{equation}
Thus
\begin{equation}
\label{eq:allowed-momenta-with-twist}
    k_m
    =
    \frac{2\pi m+\vartheta}{L},
    \qquad
    m=0,1,\ldots,L-1,
\end{equation}
where momenta are understood modulo $2\pi$.  We denote this finite set by
\begin{equation}
    \mathcal{K}_\vartheta
    =
    \left\{
    \frac{2\pi m+\vartheta}{L}\;\middle|\;m=0,1,\ldots,L-1
    \right\}.
\end{equation}

\begin{definition}[Fourier transform used in these notes]
\label{def:fourier-transform-with-phase}
For $k\in\mathcal{K}_\vartheta$, we define
\begin{equation}
\label{eq:fourier-transform-with-phase}
    c_k
    =
    \frac{\ee^{-\ii\phi}}{\sqrt{L}}
    \sum_{j=1}^{L}\ee^{-\ii k j}c_j,
    \qquad
    c_j
    =
    \frac{\ee^{\ii\phi}}{\sqrt{L}}
    \sum_{k\in\mathcal{K}_\vartheta}\ee^{\ii k j}c_k.
\end{equation}
Here $\phi$ is a constant Fourier-gauge phase.  Unless otherwise stated, the
default convention is $\phi=0$.
\end{definition}

The phase $\phi$ in \cref{eq:fourier-transform-with-phase} is the phase shown in
the convention
\begin{equation}
    c_k
    =
    \frac{\ee^{-\ii\phi}}{\sqrt{L}}
    \sum_{j=1}^{L}\ee^{-\ii k j}c_j,
    \qquad
    c_j
    =
    \frac{\ee^{\ii\phi}}{\sqrt{L}}
    \sum_{k}\ee^{\ii k j}c_k.
\end{equation}
It is not a boundary twist.  The boundary twist is $\vartheta$, which changes the
allowed values of $k$.  The constant phase $\phi$ only redefines every momentum
operator by the same $U(1)$ phase.

The inverse transform follows from the finite orthogonality relations
\begin{equation}
\label{eq:finite-fourier-orthogonality}
    \frac{1}{L}\sum_{j=1}^{L}\ee^{\ii(k-k')j}
    =
    \delta_{k,k'},
    \qquad
    \frac{1}{L}\sum_{k\in\mathcal{K}_\vartheta}\ee^{\ii k(j-j')}
    =
    \delta_{j,j'} .
\end{equation}
The second identity is valid for $j,j'=1,\ldots,L$.  The common phase $\phi$ drops
out of the orthogonality relations.

\begin{proposition}[Canonical algebra in momentum space]
The operators defined in \cref{eq:fourier-transform-with-phase} obey
\begin{equation}
    \{c_k,c_{k'}^\dagger\}=\delta_{k,k'},
    \qquad
    \{c_k,c_{k'}\}=\{c_k^\dagger,c_{k'}^\dagger\}=0 .
\end{equation}
\end{proposition}

\begin{proof}
Using the real-space canonical anticommutation relations,
\begin{align}
    \{c_k,c_{k'}^\dagger\}
    &=
    \frac{1}{L}
    \sum_{j,j'=1}^{L}
    \ee^{-\ii k j}\ee^{\ii k'j'}
    \{c_j,c_{j'}^\dagger\}
    \nonumber\\
    &=
    \frac{1}{L}\sum_{j=1}^{L}\ee^{\ii(k'-k)j}
    =
    \delta_{k,k'}.
\end{align}
The phases $\ee^{-\ii\phi}$ and $\ee^{\ii\phi}$ cancel.  The other
anticommutators vanish in the same way.
\end{proof}

\begin{remark}[Choice of lattice origin]
The transform above uses $\ee^{\ii k j}$ with $j=1,\ldots,L$.  Some books use
$\ee^{\ii k(j-1)}$ instead.  This is only a $k$-dependent redefinition
$c_k\mapsto \ee^{\ii k}c_k$.  It can shift phase factors between the Fourier
transform and the momentum-space Hamiltonian, but it cannot change spectra,
correlation functions, or topological invariants.  In these notes we keep the
$j=1,\ldots,L$ convention to match the Jordan--Wigner chapters.
\end{remark}

\section{Allowed momenta in Jordan--Wigner parity sectors}

For a periodic spin chain, the Jordan--Wigner transformation does not produce one
single fermionic boundary condition.  Instead, the effective boundary condition
depends on the fermion-parity sector.  With the convention used in these notes,
\cref{eq:effective-jw-boundary-condition} gives
\begin{equation}
    c_{L+1}=-p\,c_1,
    \qquad
    p=(-1)^N=\pm1.
\end{equation}
Therefore
\begin{equation}
\label{eq:jw-parity-momentum-sectors}
\begin{array}{c|c|c|c}
    \text{fermion parity} & p & \text{boundary condition} & \text{allowed momenta} \\
    \hline
    \text{even} & +1 & c_{L+1}=-c_1
    & k_m=\dfrac{(2m+1)\pi}{L} \\
    \text{odd} & -1 & c_{L+1}=c_1
    & k_m=\dfrac{2\pi m}{L}
\end{array}
\end{equation}
with $m=0,1,\ldots,L-1$.

This distinction is negligible for many bulk thermodynamic quantities in the
limit $L\to\infty$, but it is essential for finite-size spectra, exact diagonal
forms, and parity-sensitive ground-state questions.  In particular, periodic
fermions may contain self-conjugate momenta satisfying
\begin{equation}
    k\equiv -k\pmod{2\pi}.
\end{equation}
For even $L$ and periodic boundary conditions, these are $k=0$ and $k=\pi$.  Such
modes must be treated separately in a BdG decomposition because they are not paired
with a distinct momentum $-k$.

In the thermodynamic limit, for a smooth function $f(k)$,
\begin{equation}
\label{eq:thermodynamic-limit-momentum-sum}
    \frac{1}{L}\sum_{k\in\mathcal{K}_\vartheta}f(k)
    \longrightarrow
    \int_{-\pi}^{\pi}\frac{\dd k}{2\pi}\,f(k).
\end{equation}
The boundary angle $\vartheta$ only shifts the mesh by an amount of order $1/L$.

\section{Translation-invariant hopping terms}

The main reason for using Fourier modes is that translation-invariant hopping
terms are diagonal in momentum.  Consider
\begin{equation}
\label{eq:translation-invariant-hopping-real-space}
    H_t
    =
    \sum_{j=1}^{L}\sum_{r}t_r\,c_j^\dagger c_{j+r},
\end{equation}
where the boundary condition is understood and the hopping amplitudes satisfy
\begin{equation}
    t_{-r}=t_r^*
\end{equation}
so that $H_t$ is Hermitian.  Substituting the Fourier transform gives
\begin{align}
    \sum_{j=1}^{L}c_j^\dagger c_{j+r}
    &=
    \sum_{k\in\mathcal{K}_\vartheta}
    \ee^{\ii k r}c_k^\dagger c_k.
\end{align}
Therefore
\begin{equation}
\label{eq:hopping-dispersion-general}
    H_t
    =
    \sum_{k\in\mathcal{K}_\vartheta}\xi(k)c_k^\dagger c_k,
    \qquad
    \xi(k)=\sum_{r}t_r\ee^{\ii k r}.
\end{equation}
The Fourier-gauge phase $\phi$ cancels from all number-conserving bilinears.

\begin{example}[Nearest-neighbour tight-binding chain]
For
\begin{equation}
    H_t=-t\sum_j\left(c_j^\dagger c_{j+1}+c_{j+1}^\dagger c_j\right)
    -\mu\sum_j c_j^\dagger c_j,
\end{equation}
we obtain
\begin{equation}
    H_t=
    \sum_k\xi(k)c_k^\dagger c_k,
    \qquad
    \xi(k)=-\mu-2t\cos k.
\end{equation}
This is the normal-state dispersion that appears in the Kitaev chain and in the
fermionic image of the transverse-field XY chain.
\end{example}

\section{Translation-invariant pairing terms}

Superconducting or Jordan--Wigner-transformed anisotropic spin chains contain
terms that do not conserve fermion number but do conserve fermion parity.  A
convenient translation-invariant pairing Hamiltonian is
\begin{equation}
\label{eq:general-pairing-creation-form}
    H_\Delta
    =
    \frac{1}{2}\sum_{j=1}^{L}\sum_r
    \left(
    \Delta_r c_j^\dagger c_{j+r}^\dagger
    +
    \Delta_r^* c_{j+r}c_j
    \right).
\end{equation}
For spinless fermions, only the antisymmetric part of the pairing is physical,
so one may take
\begin{equation}
    \Delta_{-r}=-\Delta_r.
\end{equation}
The Fourier transform gives
\begin{equation}
\label{eq:momentum-space-pairing-gap-general}
    H_\Delta
    =
    \frac{1}{2}\sum_{k\in\mathcal{K}_\vartheta}
    \left[
    \Delta_\phi(k)c_k^\dagger c_{-k}^\dagger
    +
    \Delta_\phi(k)^*c_{-k}c_k
    \right],
\end{equation}
where
\begin{equation}
\label{eq:gap-function-with-fourier-phase}
    \Delta_\phi(k)
    =
    \ee^{-2\ii\phi}
    \sum_r\Delta_r\ee^{\ii k r}.
\end{equation}
Thus the additional Fourier phase does not affect hopping terms, but it rotates
the superconducting gap by twice the phase.  This factor of two is simply the
charge-two nature of a Cooper-pair operator: a pair creation operator contains two
fermion creation operators.

If the pairing is written instead in the annihilation form
\begin{equation}
    \sum_{j,r}\left(\widetilde{\Delta}_r c_jc_{j+r}
    +\widetilde{\Delta}_r^*c_{j+r}^\dagger c_j^\dagger\right),
\end{equation}
then the coefficient of the annihilation term transforms with $\ee^{+2\ii\phi}$.
The two statements are the same convention expressed from opposite sides of the
Hermitian conjugate pair.

\section{BdG blocks for spinless superconducting chains}

Combining the hopping and pairing terms, define the Nambu spinor
\begin{equation}
\label{eq:nambu-spinor-fourier-chapter}
    \Psi_k
    =
    \begin{pmatrix}
        c_k\\
        c_{-k}^\dagger
    \end{pmatrix}.
\end{equation}
Then a general spinless quadratic Hamiltonian can be written as
\begin{equation}
\label{eq:general-bdg-momentum-form-fourier-chapter}
    H
    =
    \frac{1}{2}\sum_{k\in\mathcal{K}_\vartheta}
    \Psi_k^\dagger\calH_{\rm BdG}(k)\Psi_k
    +E_0,
\end{equation}
with
\begin{equation}
\label{eq:bdg-block-general-fourier-chapter}
    \calH_{\rm BdG}(k)
    =
    \begin{pmatrix}
        \xi(k) & \Delta_\phi(k)\\
        \Delta_\phi(k)^* & -\xi(-k)
    \end{pmatrix}.
\end{equation}
For most nearest-neighbour chains considered in these notes, $\xi(-k)=\xi(k)$ and
$\Delta_\phi(-k)=-\Delta_\phi(k)$.

Equivalently, after dropping any term proportional to the identity matrix, the BdG block can be
written as
\begin{equation}
\label{eq:bdg-d-vector-fourier-chapter}
    \calH_{\rm BdG}(k)
    =
    d_x(k)\tau_x+d_y(k)\tau_y+d_z(k)\tau_z,
\end{equation}
with
\begin{equation}
    d_x(k)=\operatorname{Re}\Delta_\phi(k),
    \qquad
    d_y(k)=-\operatorname{Im}\Delta_\phi(k),
    \qquad
    d_z(k)=\xi(k).
\end{equation}
The quasiparticle energies are
\begin{equation}
\label{eq:bdg-spectrum-fourier-chapter}
    E_\pm(k)=\pm\sqrt{\xi(k)^2+\abs{\Delta_\phi(k)}^2}.
\end{equation}
Hence a constant Fourier-gauge phase changes the orientation of the vector
$(d_x,d_y)$ in the pairing plane but does not change the spectrum.

\section{Independent momentum-space blocks}

A translation-invariant number-conserving quadratic Hamiltonian has the form
\begin{equation}
    H=\sum_k\xi(k)c_k^\dagger c_k.
\end{equation}
It is already decomposed into independent one-mode blocks labelled by $k$.

A pairing Hamiltonian is different: it couples $k$ and $-k$.  Therefore its basic
blocks are not single momenta but pairs $(k,-k)$.  Let $\mathcal{K}_+$ be a choice
of one representative from each pair $\{k,-k\}$, excluding self-conjugate momenta.
Then
\begin{equation}
\label{eq:paired-momentum-block-decomposition}
    H
    =
    \sum_{k\in\mathcal{K}_+}H_{k,-k}
    +H_{\rm self}
    +E_0.
\end{equation}
Here $H_{\rm self}$ contains possible self-conjugate modes such as $k=0$ and
$k=\pi$ under periodic boundary conditions.  For odd-parity spinless pairing,
$\Delta_\phi(k)=0$ at self-conjugate momenta, so these modes remain unpaired.

For a non-self-conjugate pair, the even-parity local Hilbert space is spanned by
\begin{equation}
    \ket{0}_{k,-k},
    \qquad
    c_k^\dagger c_{-k}^\dagger\ket{0}_{k,-k}.
\end{equation}
In this subspace, the BdG block describes a two-level system.  This is why the
transverse-field Ising chain, the transverse-field XY chain, and the Kitaev chain
all reduce to products of independent two-level problems in momentum space.

\begin{principle}[Momentum-space block principle]
A translation-invariant quadratic fermion Hamiltonian decomposes into independent
momentum blocks.  Number-conserving Hamiltonians decompose into $k$ blocks.
Spinless superconducting Hamiltonians decompose into $(k,-k)$ BdG blocks, plus
possible self-conjugate modes.
\end{principle}

\section{Physical meaning of the extra Fourier phase}

The constant phase $\phi$ in \cref{eq:fourier-transform-with-phase} implements the
operator redefinition
\begin{equation}
    c_k^{(\phi)}=\ee^{-\ii\phi}c_k^{(0)}.
\end{equation}
In Nambu space,
\begin{equation}
    \Psi_k^{(\phi)}
    =
    \begin{pmatrix}
        \ee^{-\ii\phi} & 0\\
        0 & \ee^{\ii\phi}
    \end{pmatrix}
    \Psi_k^{(0)}
    =
    \ee^{-\ii\phi\tau_z}\Psi_k^{(0)}.
\end{equation}
Therefore BdG matrices in the two conventions are related by the unitary gauge
transformation
\begin{equation}
\label{eq:bdg-gauge-transformation-fourier-phase}
    \calH_{\rm BdG}^{(\phi)}(k)
    =
    \ee^{-\ii\phi\tau_z}
    \calH_{\rm BdG}^{(0)}(k)
    \ee^{\ii\phi\tau_z}.
\end{equation}
This rotates the $\tau_x$--$\tau_y$ components of the BdG vector by a constant
angle and leaves $d_z(k)$ invariant.

The distinction between $\vartheta$ and $\phi$ is important:
\begin{itemize}
    \item $\vartheta$ changes the boundary condition and therefore changes the
    allowed momenta at finite $L$;
    \item $\phi$ is a uniform gauge choice for momentum-space fermions and does
    not change the allowed momenta.
\end{itemize}
A boundary twist can represent a physical flux through a ring.  A constant
Fourier-gauge phase by itself cannot.

\section{Example: complex pairing phase in the Kitaev chain}

Consider the Kitaev chain written in the creation-pair convention
\begin{equation}
\label{eq:kitaev-chain-complex-gap-fourier-chapter}
    H_{\rm K}
    =
    -\mu\sum_j\left(c_j^\dagger c_j-\frac{1}{2}\right)
    -t\sum_j\left(c_j^\dagger c_{j+1}+c_{j+1}^\dagger c_j\right)
    +\sum_j\left(
        \Delta c_j^\dagger c_{j+1}^\dagger
        +
        \Delta^* c_{j+1}c_j
    \right),
\end{equation}
where
\begin{equation}
    \Delta=\abs{\Delta}\ee^{\ii\theta_\Delta}.
\end{equation}
Using \cref{eq:gap-function-with-fourier-phase}, the off-diagonal BdG gap is
\begin{equation}
\label{eq:kitaev-gap-with-fourier-phase}
    \Delta_\phi(k)
    =
    2\ii\abs{\Delta}\,
    \ee^{\ii(\theta_\Delta-2\phi)}\sin k.
\end{equation}
The normal dispersion is
\begin{equation}
    \xi(k)=-\mu-2t\cos k.
\end{equation}
Thus
\begin{equation}
\label{eq:kitaev-bdg-complex-gap-fourier-chapter}
    \calH_{\rm K}(k)
    =
    \begin{pmatrix}
        \xi(k) & 2\ii\abs{\Delta}\ee^{\ii(\theta_\Delta-2\phi)}\sin k\\
        -2\ii\abs{\Delta}\ee^{-\ii(\theta_\Delta-2\phi)}\sin k & -\xi(k)
    \end{pmatrix}.
\end{equation}
The spectrum is
\begin{equation}
    E_\pm(k)
    =
    \pm\sqrt{\left(\mu+2t\cos k\right)^2
    +4\abs{\Delta}^2\sin^2k},
\end{equation}
which is independent of $\theta_\Delta$ and $\phi$.

Equation \cref{eq:kitaev-gap-with-fourier-phase} shows the role of the extra
Fourier phase.  Choosing
\begin{equation}
    \phi=\frac{\theta_\Delta}{2}
\end{equation}
removes the explicit superconducting phase from the gap and gives the standard
purely imaginary $p$-wave form
\begin{equation}
    \Delta_\phi(k)=2\ii\abs{\Delta}\sin k.
\end{equation}
Alternatively, one may choose a different $\phi$ to make the off-diagonal gap point
along another fixed direction in the $\tau_x$--$\tau_y$ plane.

For an isolated uniform Kitaev chain, the global phase $\theta_\Delta$ is therefore
not an independent bulk observable.  It is a superconducting gauge phase.  The
bulk spectrum and the topological phase boundary
\begin{equation}
    \abs{\mu}=2\abs{t}
\end{equation}
are unchanged by it.  What is physical is a gauge-invariant relative phase: for
example, the phase difference between two superconductors in a Josephson junction,
a spatial gradient of the superconducting phase, or a boundary twist that cannot
be removed globally.  In such situations the phase is tied to current, flux, or
junction physics and cannot be discarded as a mere convention.

From the BdG-vector viewpoint, a uniform complex phase of $\Delta$ rotates the
closed curve
\begin{equation}
    k\mapsto (d_x(k),d_y(k),d_z(k))
\end{equation}
around the $d_z$ axis.  This does not change whether the curve winds around the
origin in the pairing plane.  Hence the topological classification of the uniform
Kitaev chain is unaffected by the constant gap phase, even though the explicit
matrix representation of the BdG block changes.

\section{Summary of conventions}

The Fourier convention used in the rest of these notes is
\begin{equation}
    c_k
    =
    \frac{\ee^{-\ii\phi}}{\sqrt{L}}
    \sum_{j=1}^{L}\ee^{-\ii k j}c_j,
    \qquad
    c_j
    =
    \frac{\ee^{\ii\phi}}{\sqrt{L}}
    \sum_{k\in\mathcal{K}_\vartheta}\ee^{\ii k j}c_k,
\end{equation}
with default $\phi=0$.  The allowed momenta are
\begin{equation}
    k_m=\frac{2\pi m+\vartheta}{L},
    \qquad
    m=0,1,\ldots,L-1.
\end{equation}
For Jordan--Wigner fermions obtained from a periodic spin chain,
\begin{equation}
\begin{array}{c|c|c}
    \text{sector} & \text{boundary condition} & \text{momenta} \\
    \hline
    (-1)^N=+1 & \text{anti-periodic} & k_m=\dfrac{(2m+1)\pi}{L} \\
    (-1)^N=-1 & \text{periodic} & k_m=\dfrac{2\pi m}{L}
\end{array}
\end{equation}
Number-conserving hopping terms give one-mode blocks labelled by $k$.  Pairing
terms give BdG blocks labelled by $(k,-k)$.  The extra phase $\phi$ is a useful
Fourier-gauge convention: it cancels from hopping terms but shifts the phase of a
pairing gap by twice its value.


\chapter{Bogoliubov--de Gennes Formalism}

The Jordan--Wigner transformation maps several spin-$1/2$ chains to quadratic
fermion Hamiltonians.  If the fermion number is conserved, Fourier
transformation already diagonalizes the problem into one-mode momentum blocks. If
pairing terms are present, however, the Hamiltonian conserves only fermion
parity, and the elementary block is no longer a single mode $k$ but a paired
sector $(k,-k)$.  The Bogoliubov--de Gennes (BdG) formalism is the systematic
language for this situation.

The purpose of this chapter is to fix the algebraic conventions used later in
the XY chain, the transverse-field Ising chain, and the Kitaev chain.  We do not
yet perform the full Bogoliubov rotation; that will be done in the next chapter.
Here we explain how a quadratic pairing Hamiltonian is encoded as a matrix in
Nambu space, why the resulting matrix has a built-in particle--hole redundancy,
and how this Nambu-space structure differs from an ordinary two-band Bloch
Hamiltonian such as the SSH model.

\section{Real-space quadratic fermion Hamiltonians}

Let $c_j,c_j^\dagger$ be canonical spinless fermions on a chain of length $L$,
obeying
\begin{equation}
    \{c_i,c_j^\dagger\}=\delta_{ij},
    \qquad
    \{c_i,c_j\}=\{c_i^\dagger,c_j^\dagger\}=0.
\end{equation}
The most general quadratic Hamiltonian preserving fermion parity is
\begin{equation}
\label{eq:bdg-real-space-quadratic-hamiltonian}
    H
    =
    \sum_{i,j=1}^{L} h_{ij}c_i^\dagger c_j
    +\frac12\sum_{i,j=1}^{L}
    \left(
        \Delta_{ij}c_i^\dagger c_j^\dagger
        +\Delta_{ij}^*c_j c_i
    \right)
    +E_{\rm cl}.
\end{equation}
Hermiticity requires
\begin{equation}
    h^\dagger=h.
\end{equation}
Fermi statistics require the pairing matrix to be antisymmetric,
\begin{equation}
\label{eq:bdg-pairing-antisymmetry-real-space}
    \Delta^T=-\Delta,
\end{equation}
because only the antisymmetric part of $\Delta_{ij}c_i^\dagger c_j^\dagger$
survives.

It is useful to introduce the real-space Nambu spinor
\begin{equation}
\label{eq:bdg-real-space-nambu-spinor}
    \Upsilon
    =
    \begin{pmatrix}
        c_1\\
        \vdots\\
        c_L\\
        c_1^\dagger\\
        \vdots\\
        c_L^\dagger
    \end{pmatrix},
    \qquad
    \Upsilon^\dagger
    =
    \begin{pmatrix}
        c_1^\dagger&\cdots&c_L^\dagger&c_1&\cdots&c_L
    \end{pmatrix}.
\end{equation}
Then \cref{eq:bdg-real-space-quadratic-hamiltonian} can be written as
\begin{equation}
\label{eq:bdg-real-space-matrix-form}
    H
    =
    \frac12\Upsilon^\dagger\mathbb{H}_{\rm BdG}\Upsilon
    +E_{\rm const},
\end{equation}
where
\begin{equation}
\label{eq:bdg-real-space-matrix}
    \mathbb{H}_{\rm BdG}
    =
    \begin{pmatrix}
        h & \Delta\\
        -\Delta^* & -h^T
    \end{pmatrix}.
\end{equation}
The constant differs from the normal-ordering constant in
\cref{eq:bdg-real-space-quadratic-hamiltonian}; explicitly,
\begin{equation}
    E_{\rm const}=E_{\rm cl}+\frac12\Tr h.
\end{equation}
This shift has no effect on eigenvectors or quasiparticle energies, but it must
be retained when computing the absolute ground-state energy.

\begin{remark}[Why the factor of $1/2$ appears]
The Nambu spinor contains both particles and holes.  Therefore the degrees of
freedom are doubled.  The prefactor $1/2$ in
\cref{eq:bdg-real-space-matrix-form} compensates this doubling.  Without this
factor the normal hopping term and the pairing term would be counted twice.
\end{remark}

\section{Intrinsic particle--hole redundancy}

The BdG matrix is not an arbitrary Hermitian $2L\times 2L$ matrix.  It satisfies
an intrinsic particle--hole constraint.  Let
\begin{equation}
    \tau_x=
    \begin{pmatrix}
        0&\id_L\\
        \id_L&0
    \end{pmatrix},
    \qquad
    \mathcal{C}=\tau_x K,
\end{equation}
where $K$ denotes complex conjugation.  Then
\begin{equation}
\label{eq:bdg-real-space-phs}
    \tau_x\mathbb{H}_{\rm BdG}^*\tau_x
    =
    -\mathbb{H}_{\rm BdG},
    \qquad
    \text{equivalently}\qquad
    \mathcal{C}\mathbb{H}_{\rm BdG}\mathcal{C}^{-1}
    =
    -\mathbb{H}_{\rm BdG}.
\end{equation}
Thus if
\begin{equation}
    \mathbb{H}_{\rm BdG}\Phi=E\Phi,
\end{equation}
then
\begin{equation}
    \mathbb{H}_{\rm BdG}(\mathcal{C}\Phi)=-E(\mathcal{C}\Phi).
\end{equation}
The spectrum is symmetric about zero.  This symmetry is not an optional physical
symmetry imposed on the model; it is a redundancy of the Nambu description.

\begin{principle}[BdG redundancy principle]
A BdG Hamiltonian doubles the Hilbert-space variables by treating $c$ and
$c^\dagger$ as independent components of a Nambu spinor.  The particle--hole
constraint removes this artificial doubling: the negative-energy BdG states are
particle--hole partners of the positive-energy states, not independent physical
single-particle bands.
\end{principle}

\section{Translation-invariant BdG blocks}

For a translation-invariant chain, the Fourier convention of the previous
chapter is
\begin{equation}
\label{eq:bdg-chapter-fourier-convention}
    c_k
    =
    \frac{\ee^{-\ii\phi}}{\sqrt{L}}
    \sum_{j=1}^{L}\ee^{-\ii kj}c_j,
    \qquad
    c_j
    =
    \frac{\ee^{\ii\phi}}{\sqrt{L}}
    \sum_{k\in\mathcal{K}_\vartheta}\ee^{\ii kj}c_k,
\end{equation}
with
\begin{equation}
    \mathcal{K}_\vartheta
    =
    \left\{
        k_m=\frac{2\pi m+\vartheta}{L}
        \ \middle|\
        m=0,1,\ldots,L-1
    \right\}.
\end{equation}
The boundary twist $\vartheta$ fixes the allowed momenta, while the constant
phase $\phi$ is a Fourier-gauge convention.

Consider the momentum-space Hamiltonian
\begin{equation}
\label{eq:bdg-chapter-general-momentum-hamiltonian}
    H
    =
    \sum_{k\in\mathcal{K}_\vartheta}\xi(k)c_k^\dagger c_k
    +\frac12\sum_{k\in\mathcal{K}_\vartheta}
    \left[
        \Delta_\phi(k)c_k^\dagger c_{-k}^\dagger
        +\Delta_\phi(k)^*c_{-k}c_k
    \right]
    +E_{\rm cl}.
\end{equation}
For spinless fermions,
\begin{equation}
\label{eq:bdg-chapter-odd-gap}
    \Delta_\phi(-k)=-\Delta_\phi(k).
\end{equation}
The natural Nambu spinor for the pair $(k,-k)$ is
\begin{equation}
\label{eq:bdg-chapter-nambu-spinor}
    \Psi_k
    =
    \begin{pmatrix}
        c_k\\
        c_{-k}^\dagger
    \end{pmatrix},
    \qquad
    \Psi_k^\dagger
    =
    \begin{pmatrix}
        c_k^\dagger&c_{-k}
    \end{pmatrix}.
\end{equation}
Then
\begin{equation}
\label{eq:bdg-chapter-block-form}
    H
    =
    \frac12\sum_{k\in\mathcal{K}_\vartheta}
    \Psi_k^\dagger\mathcal{H}_{\rm BdG}(k)\Psi_k
    +E_0,
\end{equation}
where
\begin{equation}
\label{eq:bdg-chapter-block-matrix}
    \mathcal{H}_{\rm BdG}(k)
    =
    \begin{pmatrix}
        \xi(k) & \Delta_\phi(k)\\
        \Delta_\phi(k)^* & -\xi(-k)
    \end{pmatrix}.
\end{equation}
The corresponding particle--hole constraint is
\begin{equation}
\label{eq:bdg-chapter-phs-momentum}
    \tau_x\mathcal{H}_{\rm BdG}(-k)^*\tau_x
    =
    -\mathcal{H}_{\rm BdG}(k).
\end{equation}
For the nearest-neighbour chains studied later, one usually has
\begin{equation}
    \xi(-k)=\xi(k),
    \qquad
    \Delta_\phi(-k)=-\Delta_\phi(k),
\end{equation}
and the matrix simplifies to
\begin{equation}
\label{eq:bdg-chapter-even-dispersion-block}
    \mathcal{H}_{\rm BdG}(k)
    =
    \begin{pmatrix}
        \xi(k) & \Delta_\phi(k)\\
        \Delta_\phi(k)^* & -\xi(k)
    \end{pmatrix}.
\end{equation}

\section{Two-level representation}

Let $\tau_x,\tau_y,\tau_z$ denote Pauli matrices acting in Nambu space, not in the
original spin space.  For the common even-dispersion case
\cref{eq:bdg-chapter-even-dispersion-block}, the BdG block can be written as
\begin{equation}
\label{eq:bdg-chapter-d-vector-form}
    \mathcal{H}_{\rm BdG}(k)
    =
    d_x(k)\tau_x+d_y(k)\tau_y+d_z(k)\tau_z,
\end{equation}
where
\begin{equation}
\label{eq:bdg-chapter-d-vector-components}
    d_x(k)=\operatorname{Re}\Delta_\phi(k),
    \qquad
    d_y(k)=-\operatorname{Im}\Delta_\phi(k),
    \qquad
    d_z(k)=\xi(k).
\end{equation}
The BdG spectrum is therefore
\begin{equation}
\label{eq:bdg-chapter-spectrum}
    E_\pm(k)
    =
    \pm E(k),
    \qquad
    E(k)=\sqrt{\xi(k)^2+\abs{\Delta_\phi(k)}^2}
    =\sqrt{d_x(k)^2+d_y(k)^2+d_z(k)^2}.
\end{equation}
A bulk gap closes only if
\begin{equation}
\label{eq:bdg-chapter-gap-closing-condition}
    d_x(k)=d_y(k)=d_z(k)=0
\end{equation}
for some allowed momentum $k$ in the thermodynamic limit.

\begin{remark}[Meaning of the negative branch]
The negative branch $-E(k)$ is not a second independent physical species.  It is
the particle--hole partner of the positive branch $E(k)$.  The physical
quasiparticle spectrum after Bogoliubov diagonalization consists of the positive
energies $E(k)$, together with a ground-state energy shift.
\end{remark}

\section[Independent (k,-k) sectors]{Independent $(k,-k)$ sectors}

For $k\not\equiv -k$ modulo the reciprocal lattice, the four-dimensional Fock
space generated by $c_k^\dagger$ and $c_{-k}^\dagger$ splits into even and odd
fermion-parity sectors.  The even sector is
\begin{equation}
\label{eq:bdg-chapter-even-sector-basis}
    \mathcal{H}_{k,-k}^{\rm even}
    =
    \operatorname{span}\left\{
        \ket{0}_{k,-k},
        c_k^\dagger c_{-k}^\dagger\ket{0}_{k,-k}
    \right\},
\end{equation}
whereas the odd sector is
\begin{equation}
    \mathcal{H}_{k,-k}^{\rm odd}
    =
    \operatorname{span}\left\{
        c_k^\dagger\ket{0}_{k,-k},
        c_{-k}^\dagger\ket{0}_{k,-k}
    \right\}.
\end{equation}
Pairing terms connect the two states in the even sector but do not connect states
of opposite fermion parity.  Thus the BdG problem is equivalent to a two-level
problem in each even-parity $(k,-k)$ sector.

If the allowed momentum set contains self-conjugate momenta satisfying
\begin{equation}
    k\equiv -k \pmod{2\pi},
\end{equation}
then $k=0$ or $k=\pi$.  For spinless odd-parity pairing,
\begin{equation}
    \Delta_\phi(0)=0,
    \qquad
    \Delta_\phi(\pi)=0,
\end{equation}
so these modes are not paired by the superconducting term.  They contribute
ordinary number-operator pieces of the form
\begin{equation}
    \xi(k)c_k^\dagger c_k.
\end{equation}
Such self-conjugate modes are often irrelevant in the thermodynamic bulk
spectrum, but they matter in finite-size calculations and in the parity-sector
bookkeeping of Jordan--Wigner-transformed spin chains.

\section{Gauge phase in Nambu space}

The previous chapter allowed a constant phase $\phi$ in the Fourier transform.
This phase acts on momentum-space fermions as
\begin{equation}
    c_k^{(\phi)}=\ee^{-\ii\phi}c_k^{(0)}.
\end{equation}
Consequently,
\begin{equation}
\label{eq:bdg-chapter-nambu-gauge-transformation}
    \Psi_k^{(\phi)}
    =
    \ee^{-\ii\phi\tau_z}\Psi_k^{(0)}.
\end{equation}
The BdG matrix transforms as
\begin{equation}
\label{eq:bdg-chapter-bdg-gauge-transformation}
    \mathcal{H}_{\rm BdG}^{(\phi)}(k)
    =
    \ee^{-\ii\phi\tau_z}
    \mathcal{H}_{\rm BdG}^{(0)}(k)
    \ee^{\ii\phi\tau_z}.
\end{equation}
Therefore $\phi$ rotates the $\tau_x$--$\tau_y$ components of the BdG vector while
leaving $d_z(k)=\xi(k)$ and the spectrum invariant.  Equivalently, the pairing
gap transforms as
\begin{equation}
\label{eq:bdg-chapter-gap-gauge-phase}
    \Delta_\phi(k)=\ee^{-2\ii\phi}\Delta_0(k).
\end{equation}
The factor of two reflects the fact that the pairing operator carries charge
$2$ under the fermion $U(1)$ phase rotation.

This distinction is important:
\begin{equation}
    \vartheta \text{ changes the allowed momenta, whereas } \phi
    \text{ only changes the Nambu gauge.}
\end{equation}
A boundary twist $\vartheta$ can represent a physical flux through a ring.  A
constant $\phi$ is a basis convention unless it is tied to a relative phase
between distinct superconducting regions or to a spatially varying order
parameter.

\section{Example: Kitaev-chain BdG matrix}

Consider the spinless Kitaev chain
\begin{equation}
\label{eq:bdg-chapter-kitaev-real-space}
    H_{\rm K}
    =
    -\mu\sum_j\left(c_j^\dagger c_j-\frac12\right)
    -t\sum_j\left(c_j^\dagger c_{j+1}+c_{j+1}^\dagger c_j\right)
    +\sum_j
    \left(
        \Delta c_j^\dagger c_{j+1}^\dagger
        +\Delta^* c_{j+1}c_j
    \right),
\end{equation}
with
\begin{equation}
    \Delta=\abs{\Delta}\ee^{\ii\theta_\Delta}.
\end{equation}
Using the Fourier convention \cref{eq:bdg-chapter-fourier-convention}, one finds
\begin{equation}
\label{eq:bdg-chapter-kitaev-xi-gap}
    \xi(k)=-\mu-2t\cos k,
    \qquad
    \Delta_\phi(k)
    =
    2\ii\abs{\Delta}\ee^{\ii(\theta_\Delta-2\phi)}\sin k.
\end{equation}
The BdG block is
\begin{equation}
\label{eq:bdg-chapter-kitaev-bdg-block}
    \mathcal{H}_{\rm K}(k)
    =
    \begin{pmatrix}
        -\mu-2t\cos k
        &
        2\ii\abs{\Delta}\ee^{\ii(\theta_\Delta-2\phi)}\sin k
        \\
        -2\ii\abs{\Delta}\ee^{-\ii(\theta_\Delta-2\phi)}\sin k
        &
        \mu+2t\cos k
    \end{pmatrix}.
\end{equation}
The spectrum is
\begin{equation}
\label{eq:bdg-chapter-kitaev-spectrum}
    E_\pm(k)
    =
    \pm\sqrt{(\mu+2t\cos k)^2+4\abs{\Delta}^2\sin^2 k}.
\end{equation}
It is independent of the uniform phase $\theta_\Delta$ and of the Fourier-gauge
phase $\phi$.

Choosing
\begin{equation}
    \phi=\frac{\theta_\Delta}{2}
\end{equation}
removes the explicit superconducting phase from the off-diagonal matrix element:
\begin{equation}
    \Delta_\phi(k)=2\ii\abs{\Delta}\sin k.
\end{equation}
For a single uniform Kitaev chain this is only a gauge choice.  In contrast,
relative superconducting phases between different segments, Josephson junctions,
or spatial phase gradients are physical because they cannot be removed by one
constant global choice of $\phi$.

For $\phi=\theta_\Delta/2$, the BdG vector is
\begin{equation}
    \bm{d}(k)
    =
    \bigl(0,-2\abs{\Delta}\sin k,-\mu-2t\cos k\bigr),
\end{equation}
using the convention $d_y=-\operatorname{Im}\Delta_\phi$.  The gap closes at
\begin{equation}
    \sin k=0,
    \qquad
    \mu+2t\cos k=0,
\end{equation}
so the thermodynamic phase boundaries are
\begin{equation}
\label{eq:bdg-chapter-kitaev-phase-boundaries}
    \mu=-2t
    \quad (k=0),
    \qquad
    \mu=2t
    \quad (k=\pi).
\end{equation}
Equivalently, for real nonzero $t$,
\begin{equation}
    \abs{\mu}=2\abs{t}
\end{equation}
marks the bulk gap closing between the trivial and topological phases.

\section{Example: BdG form of the transverse-field XY chain}

The transverse-field XY chain later studied as a mother model has the Pauli-form
Hamiltonian
\begin{equation}
\label{eq:bdg-chapter-xy-spin-hamiltonian}
    H_{\rm XY}
    =
    -\sum_{j=1}^{L}
    \left[
        \frac{1+\gamma}{2}\sigma_j^x\sigma_{j+1}^x
        +
        \frac{1-\gamma}{2}\sigma_j^y\sigma_{j+1}^y
    \right]
    -h\sum_{j=1}^{L}\sigma_j^z,
\end{equation}
up to the boundary-sector subtleties already fixed in the Jordan--Wigner chapter.
With the convention
\begin{equation}
    n_j=c_j^\dagger c_j=\frac{1-\sigma_j^z}{2},
\end{equation}
the bulk Jordan--Wigner image is a quadratic pairing Hamiltonian.  In the common
normalization used in these notes, its BdG vector may be chosen as
\begin{equation}
\label{eq:bdg-chapter-xy-d-vector}
    \bm{d}_{\rm XY}(k)
    =
    \bigl(0,-2\gamma\sin k,2(h-\cos k)\bigr),
\end{equation}
so that
\begin{equation}
\label{eq:bdg-chapter-xy-spectrum}
    E_{\rm XY}(k)
    =
    2\sqrt{(h-\cos k)^2+\gamma^2\sin^2 k}.
\end{equation}
This form makes the common free-fermion structure transparent:
\begin{equation}
    \text{XY chain}
    \quad\longleftrightarrow\quad
    \text{spinless BdG chain}
    \quad\longleftrightarrow\quad
    \text{independent }(k,-k)\text{ two-level systems}.
\end{equation}
The transverse-field Ising chain is the special case $\gamma=1$.  The XX chain is
the special case $\gamma=0$, where the pairing component vanishes and the problem
reduces to a number-conserving free-fermion hopping model.

\begin{remark}[Normalization warning]
Different references place factors of $2$ in different locations in
\cref{eq:bdg-chapter-xy-spin-hamiltonian}.  The invariant statement is the shape
of the spectrum and the gap-closing conditions.  With the above convention, the
critical lines are obtained from
\begin{equation}
    h-\cos k=0,
    \qquad
    \gamma\sin k=0.
\end{equation}
For $\gamma\neq0$, this gives $h=1$ at $k=0$ and $h=-1$ at $k=\pi$.
\end{remark}

\section{Nambu space versus sublattice space}

A two-component Hamiltonian of the form
\begin{equation}
    \mathcal{H}(k)=d_x(k)\tau_x+d_y(k)\tau_y+d_z(k)\tau_z
\end{equation}
may describe very different physics depending on what the two components mean.
For the Kitaev and XY chains, the two components are Nambu components,
\begin{equation}
    \Psi_k=(c_k,c_{-k}^\dagger)^T.
\end{equation}
For the SSH chain, the two components are sublattice components,
\begin{equation}
    \psi_k=(a_k,b_k)^T,
\end{equation}
where $a_k$ and $b_k$ annihilate fermions on two distinct sites within the unit
cell.  The formal $2\times2$ matrix notation is similar, but the physical meaning
is different.

\begin{center}
\begin{tabularx}{0.96\linewidth}{@{}L{0.24\linewidth}YY@{}}
\toprule
 & BdG/Nambu two-level system & SSH/sublattice two-band system \\
\midrule
Two-component spinor
& $(c_k,c_{-k}^\dagger)^T$
& $(a_k,b_k)^T$ \\
Physical components
& particle and hole of related momenta
& two independent orbitals/sublattices in one unit cell \\
Reason for two bands
& artificial Nambu doubling
& genuine Bloch-band structure \\
Built-in symmetry
& particle--hole redundancy
$\tau_x\mathcal{H}(-k)^*\tau_x=-\mathcal{H}(k)$
& no BdG redundancy; SSH has chiral symmetry only for special hopping structure \\
Prefactor in Hamiltonian
& $\frac12\sum_k\Psi_k^\dagger\mathcal{H}_{\rm BdG}(k)\Psi_k$
& $\sum_k\psi_k^\dagger h_{\rm Bloch}(k)\psi_k$ \\
Meaning of negative energy
& hole partner of positive-energy quasiparticle
& physical valence band \\
Ground state
& Bogoliubov/BCS quasiparticle vacuum
& filled-band Slater determinant at fixed particle number \\
\bottomrule
\end{tabularx}
\end{center}

This distinction is central to the organization of these notes.  The Kitaev,
TFIM, and XY chains are solved through BdG pairing blocks.  The SSH model is a
number-conserving free-fermion Bloch problem.  Both can be represented by maps
from momentum space to a vector $\bm{d}(k)$, and both can carry topological
winding numbers in one dimension, but the Hilbert-space interpretation of the
same-looking Pauli matrices is not the same.

\section{Summary}

The BdG formalism rewrites a quadratic parity-preserving fermion Hamiltonian in
Nambu space.  In real space,
\begin{equation}
    \mathbb{H}_{\rm BdG}
    =
    \begin{pmatrix}
        h&\Delta\\
        -\Delta^*&-h^T
    \end{pmatrix},
    \qquad
    h^\dagger=h,
    \qquad
    \Delta^T=-\Delta.
\end{equation}
In a translation-invariant spinless chain, the basic momentum block is
\begin{equation}
    \mathcal{H}_{\rm BdG}(k)
    =
    \begin{pmatrix}
        \xi(k)&\Delta_\phi(k)\\
        \Delta_\phi(k)^*&-\xi(-k)
    \end{pmatrix}.
\end{equation}
For $\xi(-k)=\xi(k)$ this becomes a two-level Hamiltonian
\begin{equation}
    \mathcal{H}_{\rm BdG}(k)=\bm{d}(k)\cdot\bm{\tau},
    \qquad
    \bm{d}(k)=
    \bigl(\operatorname{Re}\Delta_\phi(k),
    -\operatorname{Im}\Delta_\phi(k),
    \xi(k)\bigr),
\end{equation}
with eigenvalues
\begin{equation}
    E_\pm(k)=\pm\sqrt{\xi(k)^2+\abs{\Delta_\phi(k)}^2}.
\end{equation}
The particle--hole relation
\begin{equation}
    \tau_x\mathcal{H}_{\rm BdG}(-k)^*\tau_x=-\mathcal{H}_{\rm BdG}(k)
\end{equation}
is an intrinsic Nambu redundancy.  The next chapter will use this structure to
construct the explicit Bogoliubov transformation, the quasiparticle operators,
and the BCS product ground state.


\chapter{Bogoliubov Transformation and Quasiparticles}
\label{chap:bogoliubov-transformation-quasiparticles}

The previous chapter formulated quadratic pairing Hamiltonians as BdG matrices.
A BdG matrix makes the particle--hole structure transparent, but it is still not
the final solution of the many-body problem.  The final step is the
\emph{Bogoliubov transformation}: a canonical change of fermionic operators that
turns each BdG block into a collection of independent quasiparticle number operators.

This chapter gives the explicit diagonalization used throughout the later
chapters on the transverse-field XY chain, the transverse-field Ising chain, and
the Kitaev chain.  The main idea is simple.  For each independent momentum pair
$(k,-k)$, pairing mixes the empty state and the doubly occupied state.  The
Bogoliubov transformation is the fermionic rotation that chooses the correct
linear combinations of particles and holes so that the Hamiltonian becomes a collection of free quasiparticles.

\section{Independent positive-momentum sectors}
\label{sec:bt-independent-momentum-sectors}

For the spinless BdG chains considered in these notes, the momentum-space
Hamiltonian has the form
\begin{equation}
\label{eq:bt-full-momentum-hamiltonian}
    H
    =
    \sum_{k\in\mathcal{K}_\vartheta}\xi(k)c_k^\dagger c_k
    +\frac12\sum_{k\in\mathcal{K}_\vartheta}
    \left[
        \Delta_\phi(k)c_k^\dagger c_{-k}^\dagger
        +\Delta_\phi(k)^*c_{-k}c_k
    \right]
    +E_{\rm cl},
\end{equation}
where
\begin{equation}
\label{eq:bt-even-xi-odd-gap}
    \xi(-k)=\xi(k),
    \qquad
    \Delta_\phi(-k)=-\Delta_\phi(k).
\end{equation}
The allowed set $\mathcal{K}_\vartheta$ is fixed by the boundary twist
introduced in the Fourier-transform chapter.  Let $\mathcal{K}_+$ denote a set
containing exactly one representative from each non-self-conjugate pair
$\{k,-k\}$.  Self-conjugate momenta, if present, satisfy
\begin{equation}
    k\equiv -k \pmod{2\pi},
\end{equation}
and are therefore $k=0$ or $k=\pi$.

For $k\in\mathcal{K}_+$, the contribution of the pair $(k,-k)$ is
\begin{equation}
\label{eq:bt-pair-hamiltonian}
    H_k
    =
    \xi_k\left(c_k^\dagger c_k+c_{-k}^\dagger c_{-k}-1\right)
    +\Delta_k c_k^\dagger c_{-k}^\dagger
    +\Delta_k^*c_{-k}c_k,
\end{equation}
where, for compactness,
\begin{equation}
    \xi_k\equiv \xi(k),
    \qquad
    \Delta_k\equiv \Delta_\phi(k).
\end{equation}
Equivalently, with the Nambu spinor
\begin{equation}
\label{eq:bt-nambu-spinor-positive-k}
    \Psi_k
    =
    \begin{pmatrix}
        c_k\\
        c_{-k}^\dagger
    \end{pmatrix},
    \qquad
    \Psi_k^\dagger
    =
    \begin{pmatrix}
        c_k^\dagger&c_{-k}
    \end{pmatrix},
\end{equation}
one has
\begin{equation}
\label{eq:bt-pair-hamiltonian-nambu}
    H_k=\Psi_k^\dagger\mathcal{H}_{\rm BdG}(k)\Psi_k,
    \qquad
    \mathcal{H}_{\rm BdG}(k)
    =
    \begin{pmatrix}
        \xi_k&\Delta_k\\
        \Delta_k^*&-\xi_k
    \end{pmatrix}.
\end{equation}
Notice that no additional factor of $1/2$ appears in
\cref{eq:bt-pair-hamiltonian-nambu}, because $\mathcal{K}_+$ already contains
only one momentum from each pair.  This is equivalent to the all-momentum BdG
formula with a factor of $1/2$.

The positive BdG energy is
\begin{equation}
\label{eq:bt-positive-energy}
    E_k
    =
    \sqrt{\xi_k^2+\abs{\Delta_k}^2}.
\end{equation}
The goal is to find new canonical fermions $\gamma_k,\gamma_{-k}$ such that
\begin{equation}
\label{eq:bt-diagonal-goal}
    H_k
    =
    E_k\left(\gamma_k^\dagger\gamma_k+
    \gamma_{-k}^\dagger\gamma_{-k}-1\right).
\end{equation}

\section{Bogoliubov angle and coherence factors}
\label{sec:bt-coherence-factors}

Write the complex pairing amplitude as
\begin{equation}
\label{eq:bt-gap-phase-definition}
    \Delta_k=\abs{\Delta_k}\ee^{\ii\alpha_k},
    \qquad
    \alpha_k=\operatorname{Arg}\Delta_k,
\end{equation}
whenever $\Delta_k\neq0$.  Introduce the polar angle $\theta_k\in[0,\pi]$ by
\begin{equation}
\label{eq:bt-bogoliubov-angle}
    \cos\theta_k=\frac{\xi_k}{E_k},
    \qquad
    \sin\theta_k=\frac{\abs{\Delta_k}}{E_k}.
\end{equation}
The standard coherence factors are
\begin{equation}
\label{eq:bt-coherence-factors}
    u_k=\cos\frac{\theta_k}{2}
    =\sqrt{\frac{E_k+\xi_k}{2E_k}},
    \qquad
    v_k=-\ee^{\ii\alpha_k}\sin\frac{\theta_k}{2}
    =-\ee^{\ii\alpha_k}\sqrt{\frac{E_k-\xi_k}{2E_k}}.
\end{equation}
They obey
\begin{equation}
\label{eq:bt-coherence-normalization}
    \abs{u_k}^2+\abs{v_k}^2=1,
\end{equation}
and, with the above phase convention,
\begin{equation}
\label{eq:bt-coherence-ratio}
    \frac{v_k}{u_k}
    =
    -\frac{\Delta_k}{E_k+\xi_k}
\end{equation}
whenever $E_k+\xi_k\neq0$.

\begin{remark}[Gauge convention at singular points]
The expression \cref{eq:bt-coherence-factors} is a convenient local gauge.  At
points where $E_k+\xi_k=0$, the ratio in \cref{eq:bt-coherence-ratio} is not a
good coordinate even though the physical state is well defined.  One may instead
use an equivalent gauge based on $E_k-\xi_k$.  This is the same type of gauge
singularity encountered when representing a two-level eigenvector on a Bloch
sphere.
\end{remark}

\section{Canonical Bogoliubov transformation}
\label{sec:bt-canonical-transformation}

Define the quasiparticle operators by
\begin{equation}
\label{eq:bt-bogoliubov-transformation}
    \gamma_k
    =
    u_k c_k-v_k c_{-k}^\dagger,
    \qquad
    \gamma_{-k}
    =
    u_k c_{-k}+v_k c_k^\dagger.
\end{equation}
Equivalently,
\begin{equation}
\label{eq:bt-bogoliubov-transformation-matrix}
    \begin{pmatrix}
        \gamma_k\\
        \gamma_{-k}^\dagger
    \end{pmatrix}
    =
    U_k
    \begin{pmatrix}
        c_k\\
        c_{-k}^\dagger
    \end{pmatrix},
    \qquad
    U_k
    =
    \begin{pmatrix}
        u_k&-v_k\\
        v_k^*&u_k
    \end{pmatrix}.
\end{equation}
The matrix $U_k$ is unitary because of
\cref{eq:bt-coherence-normalization}.  Therefore the transformation preserves the
canonical anticommutation relations:
\begin{equation}
\label{eq:bt-quasiparticle-car}
    \{\gamma_q,\gamma_{q'}^\dagger\}=\delta_{q,q'},
    \qquad
    \{\gamma_q,\gamma_{q'}\}=0.
\end{equation}
The inverse transformation is
\begin{equation}
\label{eq:bt-inverse-bogoliubov-transformation}
    c_k=u_k\gamma_k+v_k\gamma_{-k}^\dagger,
    \qquad
    c_{-k}=u_k\gamma_{-k}-v_k\gamma_k^\dagger.
\end{equation}

\begin{proposition}[Diagonalization of one BdG block]
\label{prop:bt-one-block-diagonalization}
With $u_k,v_k$ given by \cref{eq:bt-coherence-factors}, the Bogoliubov
transformation \cref{eq:bt-bogoliubov-transformation} diagonalizes the pair
Hamiltonian:
\begin{equation}
\label{eq:bt-pair-diagonalized}
    H_k
    =
    E_k\left(\gamma_k^\dagger\gamma_k+
    \gamma_{-k}^\dagger\gamma_{-k}-1\right).
\end{equation}
\end{proposition}

\begin{proof}
Using the Nambu notation of \cref{eq:bt-nambu-spinor-positive-k}, the
transformation is
\begin{equation}
    \Gamma_k=U_k\Psi_k,
    \qquad
    \Gamma_k=
    \begin{pmatrix}
        \gamma_k\\
        \gamma_{-k}^\dagger
    \end{pmatrix}.
\end{equation}
A direct multiplication gives
\begin{equation}
\label{eq:bt-unitary-diagonalization}
    U_k\mathcal{H}_{\rm BdG}(k)U_k^\dagger
    =
    \begin{pmatrix}
        E_k&0\\
        0&-E_k
    \end{pmatrix}
    =E_k\tau_z.
\end{equation}
Hence
\begin{align}
    H_k
    &
    =\Psi_k^\dagger\mathcal{H}_{\rm BdG}(k)\Psi_k
      =\Gamma_k^\dagger E_k\tau_z\Gamma_k\notag\\
    &
    =E_k\gamma_k^\dagger\gamma_k
     -E_k\gamma_{-k}\gamma_{-k}^\dagger
      =E_k\left(\gamma_k^\dagger\gamma_k+
      \gamma_{-k}^\dagger\gamma_{-k}-1\right),
\end{align}
where the last equality uses
$\gamma_{-k}\gamma_{-k}^\dagger=1-\gamma_{-k}^\dagger\gamma_{-k}$.
\end{proof}

Thus the BdG eigenvalues $\pm E_k$ are not two independent particle species.
They are the particle and hole branches of one canonical pair of quasiparticle
modes.  The physical excitation energies are nonnegative and are given by
$E_k$.

\section{Even-sector two-level problem}
\label{sec:bt-even-sector-two-level}

The same diagonalization can be seen directly in the many-body Hilbert space.
For $k\not\equiv -k$, the even-parity subspace is spanned by
\begin{equation}
\label{eq:bt-even-basis}
    \ket{0}_{k,-k},
    \qquad
    \ket{11}_{k,-k}=c_k^\dagger c_{-k}^\dagger\ket{0}_{k,-k}.
\end{equation}
In this basis the pair Hamiltonian is represented by
\begin{equation}
\label{eq:bt-even-sector-matrix}
    H_k^{\rm even}
    =
    \begin{pmatrix}
        -\xi_k&\Delta_k^*\\
        \Delta_k&\xi_k
    \end{pmatrix}.
\end{equation}
Its eigenvalues are $-E_k$ and $+E_k$.  The lower eigenstate is
\begin{equation}
\label{eq:bt-pair-ground-state}
    \ket{\mathrm{GS}}_k
    =
    u_k\ket{0}_{k,-k}
    +v_k c_k^\dagger c_{-k}^\dagger\ket{0}_{k,-k},
\end{equation}
with the same $u_k,v_k$ as in \cref{eq:bt-coherence-factors}.  Indeed,
\begin{equation}
    \frac{v_k}{u_k}=-\frac{\Delta_k}{E_k+\xi_k}
\end{equation}
is exactly the lower-eigenvector equation for
\cref{eq:bt-even-sector-matrix}.

The odd-parity states
\begin{equation}
    c_k^\dagger\ket{0}_{k,-k},
    \qquad
    c_{-k}^\dagger\ket{0}_{k,-k}
\end{equation}
are not mixed by the pairing term.  Relative to the ground state of the pair,
they are one-quasiparticle excitations with energy cost $E_k$.  The doubly
excited state $\gamma_k^\dagger\gamma_{-k}^\dagger\ket{\mathrm{GS}}_k$ has energy
cost $2E_k$.

\section{BCS product ground state}
\label{sec:bt-bcs-product-ground-state}

The full quasiparticle vacuum is defined by
\begin{equation}
\label{eq:bt-quasiparticle-vacuum-definition}
    \gamma_q\ket{\mathrm{GS}}=0
    \qquad
    \text{for all quasiparticle modes }q.
\end{equation}
For the non-self-conjugate momentum pairs, it is the product of the pair ground
states:
\begin{equation}
\label{eq:bt-bcs-product-state}
    \ket{\mathrm{GS}}
    =
    \prod_{k\in\mathcal{K}_+}
    \left(
        u_k+v_k c_k^\dagger c_{-k}^\dagger
    \right)
    \ket{0},
\end{equation}
up to the separate treatment of self-conjugate modes.  The ordering of the
product over $k$ must be fixed once and for all, since fermion creation operators
anticommute.  A different ordering changes only an overall sign convention for
the many-body state.

The diagonal Hamiltonian is
\begin{equation}
\label{eq:bt-full-diagonal-hamiltonian}
    H
    =
    E_{\rm vac}
    +
    \sum_{k\in\mathcal{K}_+}E_k
    \left(
        \gamma_k^\dagger\gamma_k+
        \gamma_{-k}^\dagger\gamma_{-k}
    \right)
    +H_{\rm self},
\end{equation}
where
\begin{equation}
\label{eq:bt-paired-sector-vacuum-energy}
    E_{\rm vac}=E_{\rm cl}-\sum_{k\in\mathcal{K}_+}E_k
\end{equation}
modulo the constant convention used in the original Hamiltonian.  The term
$H_{\rm self}$ denotes possible $k=0$ or $k=\pi$ contributions, discussed below.

\begin{principle}[Meaning of the BCS product]
The BCS ground state is not a fixed-particle-number state.  It is a coherent
superposition of states with different fermion numbers but the same fermion
parity.  This is natural for BdG Hamiltonians, which conserve fermion parity
rather than fermion number.
\end{principle}

For each pair $(k,-k)$, the occupation and anomalous expectation values are
\begin{equation}
\label{eq:bt-bcs-basic-contractions}
    \langle c_k^\dagger c_k\rangle=\abs{v_k}^2,
    \qquad
    \langle c_{-k}^\dagger c_{-k}\rangle=\abs{v_k}^2,
    \qquad
    \langle c_{-k}c_k\rangle=u_kv_k.
\end{equation}
Using \cref{eq:bt-coherence-factors}, these become
\begin{equation}
\label{eq:bt-bcs-basic-contractions-explicit}
    \abs{v_k}^2=\frac12\left(1-\frac{\xi_k}{E_k}\right),
    \qquad
    u_kv_k=-\frac{\Delta_k}{2E_k}.
\end{equation}
These contractions are the input for the covariance-matrix and Pfaffian methods
developed in the next chapter.

\section{Self-conjugate momenta and parity bookkeeping}
\label{sec:bt-self-conjugate-momenta}

If the allowed momentum set contains $q=0$ or $q=\pi$, then $q=-q$ modulo
$2\pi$.  For spinless odd-parity pairing,
\begin{equation}
    \Delta_q=0,
\end{equation}
so such a mode cannot be paired with itself.  Its BdG contribution is instead an
ordinary one-mode problem.  In a symmetric BdG convention one writes
\begin{equation}
\label{eq:bt-self-conjugate-bdg-term}
    H_q=\xi_q\left(c_q^\dagger c_q-\frac12\right).
\end{equation}
Equivalently,
\begin{equation}
\label{eq:bt-self-conjugate-diagonal}
    H_q=\abs{\xi_q}\left(\eta_q^\dagger\eta_q-\frac12\right),
\end{equation}
where
\begin{equation}
    \eta_q=
    \begin{cases}
        c_q, & \xi_q>0,\\
        c_q^\dagger, & \xi_q<0.
    \end{cases}
\end{equation}
Thus the ground-state occupation of the self-conjugate mode is
\begin{equation}
\label{eq:bt-self-conjugate-occupation}
    n_q^{\rm GS}
    =
    \begin{cases}
        0, & \xi_q>0,\\
        1, & \xi_q<0.
    \end{cases}
\end{equation}
These modes are negligible for most thermodynamic bulk quantities but essential
for finite-size spectra and fermion-parity selection.

In Jordan--Wigner images of spin chains, the physical Hilbert space splits into
fermion-parity sectors.  The parity sector fixes the boundary condition, and the
boundary condition fixes whether self-conjugate momenta are present.  Therefore a
careful finite-size solution requires the following sequence:
\begin{equation}
\label{eq:bt-finite-size-sequence}
\begin{aligned}
    \text{spin parity sector}
    &\longrightarrow
    \text{fermion boundary condition}
    \longrightarrow
    \text{allowed momentum set}\\
    &\longrightarrow
    \text{Bogoliubov diagonalization}.
\end{aligned}
\end{equation}
Only after this sequence can one compare the vacuum energies of different parity
sectors to determine the true finite-size ground state.

\section{Role of the Fourier-gauge phase}
\label{sec:bt-fourier-gauge-phase}

The Fourier convention used in these notes allows
\begin{equation}
\label{eq:bt-fourier-phase-convention}
    c_k^{(\phi)}=\ee^{-\ii\phi}c_k^{(0)},
    \qquad
    c_j=\frac{\ee^{\ii\phi}}{\sqrt{L}}
    \sum_k \ee^{\ii kj}c_k^{(\phi)}.
\end{equation}
As explained in the previous chapter, this is a Nambu gauge choice.  It changes
\begin{equation}
\label{eq:bt-gap-phase-under-fourier-gauge}
    \Delta_k^{(\phi)}=\ee^{-2\ii\phi}\Delta_k^{(0)},
\end{equation}
but leaves the spectrum invariant:
\begin{equation}
    E_k=\sqrt{\xi_k^2+\abs{\Delta_k^{(\phi)}}^2}.
\end{equation}
Since the coherence factor $v_k$ carries the phase of $\Delta_k$, it transforms
as
\begin{equation}
\label{eq:bt-v-gauge-transformation}
    v_k^{(\phi)}=\ee^{-2\ii\phi}v_k^{(0)},
    \qquad
    u_k^{(\phi)}=u_k^{(0)},
\end{equation}
within the convention \cref{eq:bt-coherence-factors}.  The BCS factor transforms
covariantly:
\begin{equation}
\label{eq:bt-bcs-factor-gauge-invariance}
    u_k^{(\phi)}+v_k^{(\phi)}
    c_k^{(\phi)\dagger}c_{-k}^{(\phi)\dagger}
    =
    u_k^{(0)}+v_k^{(0)}c_k^{(0)\dagger}c_{-k}^{(0)\dagger}.
\end{equation}
Thus the many-body state is unchanged; only its coordinate representation in
Nambu space has been rotated.

This observation is especially useful for a Kitaev chain with complex pairing,
\begin{equation}
\label{eq:bt-complex-kitaev-gap}
    \Delta=\abs{\Delta}\ee^{\ii\theta_\Delta}.
\end{equation}
The gap function derived in the previous chapter is
\begin{equation}
\label{eq:bt-kitaev-gap-with-phase}
    \Delta_\phi(k)
    =
    2\ii\abs{\Delta}\ee^{\ii(\theta_\Delta-2\phi)}\sin k.
\end{equation}
Choosing
\begin{equation}
\label{eq:bt-kitaev-phase-gauge-choice}
    \phi=\frac{\theta_\Delta}{2}
\end{equation}
sets
\begin{equation}
    \Delta_\phi(k)=2\ii\abs{\Delta}\sin k.
\end{equation}
For a single uniform Kitaev chain, this removes the uniform superconducting phase
from the BdG block.  The phase becomes physical only when it is relative to
another phase, spatially dependent, or tied to a boundary twist or flux that
cannot be removed by a single global gauge choice.

\section{Application: quasiparticles of the Kitaev chain}
\label{sec:bt-kitaev-quasiparticles}

For the Kitaev chain in the convention of the previous chapter,
\begin{equation}
\label{eq:bt-kitaev-xi-delta}
    \xi_k=-\mu-2t\cos k,
    \qquad
    \Delta_k=2\ii\abs{\Delta}\ee^{\ii(\theta_\Delta-2\phi)}\sin k.
\end{equation}
The quasiparticle dispersion is
\begin{equation}
\label{eq:bt-kitaev-quasiparticle-spectrum}
    E_k=
    \sqrt{(\mu+2t\cos k)^2+4\abs{\Delta}^2\sin^2 k}.
\end{equation}
The Bogoliubov angle is determined by
\begin{equation}
\label{eq:bt-kitaev-angle}
    \cos\theta_k=\frac{-\mu-2t\cos k}{E_k},
    \qquad
    \sin\theta_k=\frac{2\abs{\Delta}\abs{\sin k}}{E_k}.
\end{equation}
The phase entering $v_k$ is
\begin{equation}
\label{eq:bt-kitaev-alpha}
    \alpha_k
    =
    \operatorname{Arg}\left[
        2\ii\abs{\Delta}\ee^{\ii(\theta_\Delta-2\phi)}\sin k
    \right].
\end{equation}
In the gauge $\phi=\theta_\Delta/2$, this reduces to the phase of
$\ii\sin k$.  Hence the only remaining phase discontinuity is the sign change of
$\sin k$, which reflects the odd-parity nature of spinless $p$-wave pairing.

The gap closes when both terms inside \cref{eq:bt-kitaev-quasiparticle-spectrum}
vanish:
\begin{equation}
    \sin k=0,
    \qquad
    \mu+2t\cos k=0.
\end{equation}
Thus, for real nonzero $t$, the bulk phase boundaries are
\begin{equation}
    \mu=-2t
    \quad(k=0),
    \qquad
    \mu=2t
    \quad(k=\pi).
\end{equation}
The Bogoliubov quasiparticles become gapless at these transition points.

\section{Application: quasiparticles of the XY and TFIM chains}
\label{sec:bt-xy-tfim-quasiparticles}

For the transverse-field XY chain in the normalization fixed earlier, the BdG
vector can be chosen as
\begin{equation}
\label{eq:bt-xy-d-vector}
    \bm{d}_{\rm XY}(k)
    =
    \bigl(0,-2\gamma\sin k,2(h-
    \cos k)\bigr).
\end{equation}
Therefore
\begin{equation}
\label{eq:bt-xy-xi-gap}
    \xi_k=2(h-\cos k),
    \qquad
    \Delta_k=2\ii\gamma\sin k,
\end{equation}
up to the same Nambu-gauge conventions discussed above.  The quasiparticle
energy is
\begin{equation}
\label{eq:bt-xy-quasiparticle-spectrum}
    E_k=2\sqrt{(h-\cos k)^2+\gamma^2\sin^2 k}.
\end{equation}
The transverse-field Ising chain is obtained by setting $\gamma=1$:
\begin{equation}
\label{eq:bt-tfim-spectrum}
    E_k^{\rm TFIM}
    =2\sqrt{(h-\cos k)^2+\sin^2 k}
    =2\sqrt{1+h^2-2h\cos k}.
\end{equation}
The gap closes at $h=1$ and $k=0$ in the thermodynamic limit.  With the
normalization used here, the excitation gap behaves as
\begin{equation}
\label{eq:bt-tfim-gap}
    \Delta_{\rm gap}=2\abs{1-h}
\end{equation}
near the critical point.

The XX chain corresponds to $\gamma=0$.  In that case the pairing vanishes,
\begin{equation}
    \Delta_k=0,
\end{equation}
and the Bogoliubov transformation reduces to an ordinary choice of occupied and
empty fermion modes.  Thus the XX chain is a number-conserving free-fermion
problem, while the TFIM and generic XY chain are genuine BdG pairing problems.

\section{Summary}
\label{sec:bt-summary}

Each non-self-conjugate momentum pair $(k,-k)$ of a spinless BdG chain is solved
by a canonical Bogoliubov transformation.  Starting from
\begin{equation}
    H_k
    =
    \xi_k\left(n_k+n_{-k}-1\right)
    +\Delta_k c_k^\dagger c_{-k}^\dagger
    +\Delta_k^*c_{-k}c_k,
\end{equation}
one defines
\begin{equation}
    E_k=\sqrt{\xi_k^2+\abs{\Delta_k}^2},
    \qquad
    u_k=\sqrt{\frac{E_k+\xi_k}{2E_k}},
    \qquad
    v_k=-\ee^{\ii\operatorname{Arg}\Delta_k}
    \sqrt{\frac{E_k-\xi_k}{2E_k}}.
\end{equation}
The quasiparticles are
\begin{equation}
    \gamma_k=u_kc_k-v_kc_{-k}^\dagger,
    \qquad
    \gamma_{-k}=u_kc_{-k}+v_kc_k^\dagger,
\end{equation}
and the Hamiltonian becomes
\begin{equation}
    H_k=E_k\left(\gamma_k^\dagger\gamma_k+
    \gamma_{-k}^\dagger\gamma_{-k}-1\right).
\end{equation}
The quasiparticle vacuum is the BCS product state
\begin{equation}
    \ket{\mathrm{GS}}
    =
    \prod_{k\in\mathcal{K}_+}
    \left(u_k+v_kc_k^\dagger c_{-k}^\dagger\right)\ket{0},
\end{equation}
with separate finite-size bookkeeping for possible $k=0$ and $k=\pi$ modes.
The phase of the pairing gap appears in $v_k$ and in the anomalous correlator,
but a uniform Fourier-gauge phase changes only the representation, not the
physical spectrum or state.  This diagonal form is the starting point for the
Gaussian correlation-function machinery developed in the next chapter.


\chapter{Correlation Functions, Covariance Matrices, and Pfaffians}
\chaptermark{Correlation Functions and Pfaffians}
\label{chap:correlation-functions-covariance-pfaffians}

The preceding chapters reduced the transverse-field Ising chain, the XY chain,
and the Kitaev chain to quadratic fermion or BdG Hamiltonians.  Diagonalizing the
Hamiltonian gives the spectrum and the quasiparticle vacuum.  For many-body
physics, however, the spectrum is only half of the solution.  One also wants
order parameters, spin-spin correlators, string correlators, entanglement
entropies, and finite-temperature response functions.

The essential reason quadratic fermion models are exactly solvable at the level
of correlation functions is that their ground states and thermal states are
\emph{fermionic Gaussian states}.  In such states all higher-order correlation
functions are fixed by two-point functions.  For local fermionic observables this
statement is Wick's theorem.  For spin observables obtained after the
Jordan--Wigner transformation, the nonlocal strings turn spin correlators into
products of many Majorana operators, and Wick's theorem organizes those products
as Pfaffians.  In translation-invariant one-dimensional problems, these
Pfaffians often further reduce to Toeplitz determinants.

This chapter gives the correlation-function machinery used throughout the rest
of the notes.  The conventions are those of the previous chapters:
\begin{equation}
    n_j=c_j^\dagger c_j=\frac{1-\sigma_j^z}{2},
    \qquad
    \sigma_j^z=1-2n_j,
\end{equation}
and the Jordan--Wigner convention is
\begin{equation}
    c_j=\left(\prod_{\ell<j}\sigma_\ell^z\right)\sigma_j^+,
    \qquad
    c_j^\dagger=\left(\prod_{\ell<j}\sigma_\ell^z\right)\sigma_j^- .
\end{equation}
Thus the occupied fermion state corresponds to spin down.

\section{Gaussian fermionic states and Wick theorem}
\label{sec:gaussian-fermionic-states-wick}

A fermionic state is called Gaussian if its density matrix can be written, up to
normalization, as the exponential of a quadratic form in fermion operators.  In a
complex-fermion basis this means schematically
\begin{equation}
    \rho
    =\frac{1}{Z}
    \exp\left[-\sum_{i,j}c_i^\dagger K_{ij}c_j
    -\frac12\sum_{i,j}\left(
        L_{ij}c_i^\dagger c_j^\dagger+L_{ij}^*c_jc_i
    \right)\right],
\end{equation}
where $K^\dagger=K$ and $L^T=-L$.  Equivalently, in a Majorana basis to be
introduced below,
\begin{equation}
    \rho=\frac{1}{Z}\exp\left(\frac{\ii}{4}\sum_{m,n}W_{mn}a_ma_n\right),
    \qquad
    W^T=-W,
\end{equation}
with $W$ real for a physical Hermitian quadratic exponent.

The ground state of any gapped quadratic Hamiltonian is Gaussian.  A thermal
state of a quadratic Hamiltonian is also Gaussian.  This is the algebraic reason
why free fermion and BdG models are not merely spectrally solvable, but also
correlation-function solvable.

\begin{theorem}[Wick theorem for parity-even Gaussian states]
\label{thm:wick-majorana}
Let $\rho$ be a parity-even fermionic Gaussian state, and let
$O_1,\ldots,O_{2N}$ be fermionic operators linear in $c_j$ and $c_j^\dagger$.
Then
\begin{equation}
\label{eq:wick-general-linear-fermion}
    \left\langle O_1O_2\cdots O_{2N}\right\rangle
    =
    \sum_{\text{pairings }P}
    \operatorname{sgn}(P)
    \prod_{(r,s)\in P}\left\langle O_rO_s\right\rangle .
\end{equation}
All odd products have vanishing expectation value:
\begin{equation}
    \left\langle O_1O_2\cdots O_{2N+1}\right\rangle=0.
\end{equation}
The sign $\operatorname{sgn}(P)$ is the fermionic sign obtained by bringing the
paired operators next to one another.
\end{theorem}

For example,
\begin{align}
\label{eq:wick-four-point-example}
    \langle O_1O_2O_3O_4\rangle
    &=
    \langle O_1O_2\rangle\langle O_3O_4\rangle
    -
    \langle O_1O_3\rangle\langle O_2O_4\rangle
    +
    \langle O_1O_4\rangle\langle O_2O_3\rangle .
\end{align}
The alternating signs are not optional.  They are the main bookkeeping issue in
fermionic correlation functions.

\begin{proof}[Sketch of proof]
For a Gaussian density matrix, introduce Grassmann sources $\eta_m$ coupled to a
basis of Majorana operators $a_m$.  The generating function has the form
\begin{equation}
    Z[\eta]
    =\left\langle
    \exp\left(\sum_m\eta_ma_m\right)
    \right\rangle
    =\exp\left(\frac12\sum_{m,n}\eta_mG_{mn}\eta_n\right),
\end{equation}
where $G_{mn}$ is the two-point contraction matrix.  Differentiating $Z[\eta]$
with respect to the Grassmann sources gives all correlation functions.  Since
$Z[\eta]$ is the exponential of a quadratic form, every nonzero derivative pairs
sources into two-point contractions.  This gives \cref{eq:wick-general-linear-fermion}.
Odd derivatives vanish because the generating function is even in the Grassmann
sources.
\end{proof}

For practical calculations it is usually inefficient to explicitly enumerate all
pairings.  The compact object that automatically includes all pairings and their
fermionic signs is the Pfaffian.

\section{Pfaffians}
\label{sec:pfaffians-basic-definition}

Let $M$ be a $2N\times2N$ antisymmetric matrix.  Its Pfaffian is defined by
\begin{equation}
\label{eq:pfaffian-definition}
    \Pf(M)
    =
    \frac{1}{2^N N!}
    \sum_{\pi\in S_{2N}}
    \sgn(\pi)
    \prod_{r=1}^{N}
    M_{\pi(2r-1),\pi(2r)} .
\end{equation}
Equivalently, it is the signed sum over all pairings of the $2N$ indices.
It obeys the fundamental identity
\begin{equation}
\label{eq:pfaffian-determinant-identity}
    \Pf(M)^2=\det M .
\end{equation}
For the smallest cases,
\begin{equation}
    \Pf
    \begin{pmatrix}
        0&m_{12}\\
        -m_{12}&0
    \end{pmatrix}
    =m_{12},
\end{equation}
and
\begin{equation}
\label{eq:pfaffian-four-by-four}
    \Pf
    \begin{pmatrix}
        0&m_{12}&m_{13}&m_{14}\\
        -m_{12}&0&m_{23}&m_{24}\\
        -m_{13}&-m_{23}&0&m_{34}\\
        -m_{14}&-m_{24}&-m_{34}&0
    \end{pmatrix}
    =m_{12}m_{34}-m_{13}m_{24}+m_{14}m_{23}.
\end{equation}
Thus \cref{eq:pfaffian-four-by-four} is precisely the four-point Wick expansion
\cref{eq:wick-four-point-example}.

If $a_m$ are Majorana operators and the indices
$m_1,\ldots,m_{2N}$ are all distinct, Wick theorem can be written as
\begin{equation}
\label{eq:majorana-product-pfaffian-contraction}
    \left\langle a_{m_1}a_{m_2}\cdots a_{m_{2N}}\right\rangle
    =
    \Pf\left(G_{\mathcal{I}}\right),
    \qquad
    \mathcal{I}=\{m_1,\ldots,m_{2N}\},
\end{equation}
where the antisymmetric contraction matrix is
\begin{equation}
    (G_{\mathcal{I}})_{rs}
    =
    \begin{cases}
    \left\langle a_{m_r}a_{m_s}\right\rangle,& r<s,\\
    -\left\langle a_{m_s}a_{m_r}\right\rangle,& r>s,\\
    0,&r=s.
    \end{cases}
\end{equation}
The next section rewrites this formula in terms of the covariance matrix
$\Gamma$.

\section{Majorana covariance matrix}
\label{sec:majorana-covariance-matrix}

For each lattice site define two Hermitian Majorana operators
\begin{equation}
\label{eq:majorana-definition-this-chapter}
    a_{2j-1}=c_j+c_j^\dagger,
    \qquad
    a_{2j}=\ii(c_j-c_j^\dagger).
\end{equation}
They satisfy
\begin{equation}
\label{eq:majorana-clifford-algebra}
    a_m^\dagger=a_m,
    \qquad
    \{a_m,a_n\}=2\delta_{mn}.
\end{equation}
The inverse relations are
\begin{equation}
\label{eq:complex-fermion-from-majorana}
    c_j=\frac12\left(a_{2j-1}-\ii a_{2j}\right),
    \qquad
    c_j^\dagger=\frac12\left(a_{2j-1}+\ii a_{2j}\right).
\end{equation}

\begin{definition}[Majorana covariance matrix]
\label{def:majorana-covariance-matrix}
For a state $\rho$, the Majorana covariance matrix is the real antisymmetric
matrix
\begin{equation}
\label{eq:majorana-covariance-definition}
    \Gamma_{mn}
    =
    \frac{\ii}{2}\left\langle [a_m,a_n]\right\rangle .
\end{equation}
Equivalently,
\begin{equation}
\label{eq:majorana-two-point-from-gamma}
    \left\langle a_ma_n\right\rangle
    =
    \delta_{mn}-\ii\Gamma_{mn}.
\end{equation}
\end{definition}

The covariance matrix is the natural object for Gaussian-state calculations.
For a parity-even Gaussian state, all correlation functions are determined by
$\Gamma$.  In particular, for distinct Majorana indices
$m_1,\ldots,m_{2N}$,
\begin{equation}
\label{eq:majorana-product-pfaffian-gamma}
    \left\langle a_{m_1}a_{m_2}\cdots a_{m_{2N}}\right\rangle
    =
    \Pf\left(-\ii\Gamma_{\mathcal{I}}\right),
    \qquad
    \mathcal{I}=\{m_1,\ldots,m_{2N}\}.
\end{equation}
Here $\Gamma_{\mathcal{I}}$ denotes the submatrix of $\Gamma$ restricted to the
ordered list $\mathcal{I}=(m_1,\ldots,m_{2N})$.  The order matters: permuting the
indices changes the Pfaffian by the sign of the permutation, exactly as required
for fermionic operators.

\begin{remark}[Pure Gaussian states]
For a pure Gaussian state, the covariance matrix obeys
\begin{equation}
\label{eq:pure-gaussian-gamma-square}
    \Gamma^2=-\id.
\end{equation}
For a mixed Gaussian state, its eigenvalues come in pairs $\pm \ii\nu_\alpha$ with
$0\leq \nu_\alpha\leq1$.  The pure-state limit corresponds to
$\nu_\alpha=1$ for every mode.
\end{remark}

\subsection*{Relation to normal and anomalous correlators}

It is often useful to express $\Gamma$ in terms of the complex-fermion
correlation matrices
\begin{equation}
\label{eq:normal-anomalous-correlators-definition}
    C_{ij}=\langle c_i^\dagger c_j\rangle,
    \qquad
    F_{ij}=\langle c_i c_j\rangle,
    \qquad
    F_{ij}=-F_{ji}.
\end{equation}
Using \cref{eq:majorana-definition-this-chapter}, one obtains
\begin{align}
\label{eq:gamma-block-odd-odd}
    \Gamma_{2i-1,2j-1}
    &=
    \ii\left(C_{ij}-C_{ji}+F_{ij}-F_{ij}^*\right),\\
\label{eq:gamma-block-even-even}
    \Gamma_{2i,2j}
    &=
    \ii\left(C_{ij}-C_{ji}-F_{ij}+F_{ij}^*\right),\\
\label{eq:gamma-block-odd-even}
    \Gamma_{2i-1,2j}
    &=
    \delta_{ij}-C_{ij}-C_{ji}-F_{ij}+F_{ij}^*.
\end{align}
The remaining block follows from antisymmetry:
\begin{equation}
    \Gamma_{2i,2j-1}=-\Gamma_{2j-1,2i}.
\end{equation}

For a translation-invariant BCS ground state written as
\begin{equation}
    \ket{\mathrm{GS}}
    =
    \prod_{k\in\mathcal{K}_+}
    \left(u_k+v_kc_k^\dagger c_{-k}^\dagger\right)\ket{0},
\end{equation}
with the Fourier convention of the previous chapter,
\begin{equation}
    c_j=\frac{\ee^{\ii\phi}}{\sqrt{L}}\sum_k\ee^{\ii kj}c_k,
\end{equation}
the normal and anomalous kernels may be written as
\begin{align}
\label{eq:translation-invariant-C-r}
    C_r
    &=
    \langle c_j^\dagger c_{j+r}\rangle
    =
    \frac{1}{L}\sum_k\ee^{\ii kr}\abs{v_k}^2,\\
\label{eq:translation-invariant-F-r}
    F_r
    &=
    \langle c_j c_{j+r}\rangle
    =
    \frac{\ee^{2\ii\phi}}{L}\sum_k\ee^{-\ii kr}\kappa_k,
    \qquad
    \kappa_k=\langle c_kc_{-k}\rangle=-u_kv_k.
\end{align}
The phase $\phi$ is the Fourier-gauge phase discussed earlier.  It rotates the
representation of anomalous correlators but leaves gauge-invariant spin and
fermion-parity observables unchanged.

\subsection*{Covariance matrix from a quadratic Majorana Hamiltonian}

A general quadratic Majorana Hamiltonian can be written as
\begin{equation}
\label{eq:majorana-quadratic-hamiltonian}
    H=\frac{\ii}{4}\sum_{m,n=1}^{2L}A_{mn}a_ma_n,
    \qquad
    A^T=-A,
\end{equation}
where $A$ is real antisymmetric.  For its thermal state at inverse temperature
$\beta$, the covariance matrix is
\begin{equation}
\label{eq:thermal-gamma-majorana-hamiltonian}
    \Gamma_\beta
    =
    \ii\tanh\left(\frac{\beta\,\ii A}{2}\right).
\end{equation}
In the zero-temperature limit for a gapped Hamiltonian,
\begin{equation}
\label{eq:ground-state-gamma-sign-function}
    \Gamma_0=\ii\,\sgn(\ii A).
\end{equation}
This formula is useful for numerical checks because it bypasses the explicit
construction of many-body wavefunctions.

\begin{example}[One fermionic mode]
Let
\begin{equation}
    H=\varepsilon\left(c^\dagger c-\frac12\right).
\end{equation}
Then
\begin{equation}
    H=-\frac{\ii\varepsilon}{2}a_1a_2,
    \qquad
    A=
    \begin{pmatrix}
        0&-\varepsilon\\
        \varepsilon&0
    \end{pmatrix}.
\end{equation}
For $\varepsilon>0$, the zero-temperature state is the empty state, and
\begin{equation}
    \Gamma_{12}=1,
    \qquad
    \langle \ii a_1a_2\rangle=1.
\end{equation}
This agrees with $\sigma^z=1-2n$ under the spin-down-as-occupied convention.
\end{example}

\section{Local observables from the covariance matrix}
\label{sec:local-observables-from-covariance}

The simplest observables are local density and magnetization.  From
\cref{eq:complex-fermion-from-majorana},
\begin{equation}
\label{eq:number-from-majoranas}
    n_j=c_j^\dagger c_j
    =\frac12\left(1-\ii a_{2j-1}a_{2j}\right),
\end{equation}
and therefore
\begin{equation}
\label{eq:sigmaz-from-majoranas}
    \sigma_j^z=1-2n_j=\ii a_{2j-1}a_{2j}.
\end{equation}
Taking the expectation value gives
\begin{equation}
\label{eq:sigmaz-expectation-from-gamma}
    \langle \sigma_j^z\rangle
    =
    \Gamma_{2j-1,2j},
    \qquad
    \langle n_j\rangle
    =
    \frac12\left(1-\Gamma_{2j-1,2j}\right).
\end{equation}

For $i\neq j$, the longitudinal spin correlator is a four-Majorana expectation:
\begin{equation}
\label{eq:sigmaz-sigmaz-pfaffian}
    \langle \sigma_i^z\sigma_j^z\rangle
    =
    -\Pf\left(-\ii\Gamma_{\{2i-1,2i,2j-1,2j\}}\right).
\end{equation}
Equivalently, expanding the four-point Pfaffian gives
\begin{align}
\label{eq:sigmaz-sigmaz-expanded}
    \langle \sigma_i^z\sigma_j^z\rangle
    &
    =
    \Gamma_{2i-1,2i}\Gamma_{2j-1,2j}
    -
    \Gamma_{2i-1,2j-1}\Gamma_{2i,2j}
    +
    \Gamma_{2i-1,2j}\Gamma_{2i,2j-1}.
\end{align}
Thus even density-density and longitudinal spin correlations are already
encoded in the same covariance matrix.

\section{Spin correlators as Pfaffians}
\label{sec:spin-correlators-as-pfaffians}

The nontrivial feature of spin correlators is the Jordan--Wigner string.  Define
\begin{equation}
\label{eq:jw-string-this-chapter}
    S_j=\prod_{\ell<j}\sigma_\ell^z
    =
    \prod_{\ell<j}\ii a_{2\ell-1}a_{2\ell}.
\end{equation}
With the convention used in these notes,
\begin{equation}
\label{eq:spin-majorana-local-formulas}
    \sigma_j^x=S_ja_{2j-1},
    \qquad
    \sigma_j^y=-S_ja_{2j},
    \qquad
    \sigma_j^z=\ii a_{2j-1}a_{2j}.
\end{equation}
For $i<j$, multiplying the strings gives finite products of Majorana operators:
\begin{align}
\label{eq:sigmax-sigmax-majorana-string}
    \sigma_i^x\sigma_j^x
    &=
    \ii^{j-i}
    a_{2i}
    \left(\prod_{m=i+1}^{j-1}a_{2m-1}a_{2m}\right)
    a_{2j-1},\\
\label{eq:sigmay-sigmay-majorana-string}
    \sigma_i^y\sigma_j^y
    &=
    -\ii^{j-i}
    a_{2i-1}
    \left(\prod_{m=i+1}^{j-1}a_{2m-1}a_{2m}\right)
    a_{2j},\\
\label{eq:sigmax-sigmay-majorana-string}
    \sigma_i^x\sigma_j^y
    &=
    -\ii^{j-i}
    a_{2i}
    \left(\prod_{m=i+1}^{j-1}a_{2m-1}a_{2m}\right)
    a_{2j}.
\end{align}
The products are ordered from left to right as written.  The string length grows
linearly with the spin separation, but the expectation value is still exactly
computable: it is a Pfaffian of a finite submatrix of $\Gamma$.

For example, define the ordered index list
\begin{equation}
\label{eq:index-list-xx-correlator}
    \mathcal{I}_{xx}(i,j)
    =
    (2i,2i+1,2i+2,\ldots,2j-2,2j-1),
    \qquad i<j.
\end{equation}
Then
\begin{equation}
\label{eq:xx-correlator-pfaffian}
    \langle \sigma_i^x\sigma_j^x\rangle
    =
    \ii^{j-i}
    \Pf\left(-\ii\Gamma_{\mathcal{I}_{xx}(i,j)}\right).
\end{equation}
Similarly, with
\begin{equation}
\label{eq:index-list-yy-correlator}
    \mathcal{I}_{yy}(i,j)
    =
    (2i-1,2i+1,2i+2,\ldots,2j-2,2j),
\end{equation}
one has
\begin{equation}
\label{eq:yy-correlator-pfaffian}
    \langle \sigma_i^y\sigma_j^y\rangle
    =
    -\ii^{j-i}
    \Pf\left(-\ii\Gamma_{\mathcal{I}_{yy}(i,j)}\right).
\end{equation}
These equations are one of the most important practical outputs of the
Jordan--Wigner solution: a highly nonlocal spin correlator has been reduced to a
finite-dimensional linear-algebra problem.

\begin{remark}[Why Pfaffians rather than determinants?]
A determinant appears when the Wick contractions have a bipartite structure, for
example when only contractions of the form $\langle B_mA_n\rangle$ enter.  A
Pfaffian is the more general object.  It remains valid when the state has complex
pairing, broken time-reversal symmetry, finite-size boundary subtleties, or mixed
$AA$, $BB$, and $AB$ contractions.
\end{remark}

\section{Toeplitz determinant form in translation-invariant chains}
\label{sec:toeplitz-determinant-translation-invariant}

For many real XY-chain conventions, symmetries imply a simpler determinant
formula.  Introduce
\begin{equation}
\label{eq:A-B-operators-toeplitz}
    A_j=c_j+c_j^\dagger=a_{2j-1},
    \qquad
    B_j=c_j^\dagger-c_j=\ii a_{2j}.
\end{equation}
These satisfy
\begin{equation}
    A_j^2=1,
    \qquad
    B_j^2=-1,
    \qquad
    \{A_i,B_j\}=0.
\end{equation}
The local parity is
\begin{equation}
    \sigma_j^z=A_jB_j,
\end{equation}
and the transverse correlator for $r>0$ becomes
\begin{equation}
\label{eq:xx-string-A-B-form}
    \sigma_j^x\sigma_{j+r}^x
    =
    B_jA_{j+1}B_{j+1}A_{j+2}\cdots B_{j+r-1}A_{j+r}.
\end{equation}
If the Gaussian state is translation invariant and its only nontrivial
contractions in this basis are
\begin{equation}
\label{eq:G-r-definition-toeplitz}
    G_r=\langle B_jA_{j+r}\rangle,
\end{equation}
then Wick theorem reduces \cref{eq:xx-string-A-B-form} to the Toeplitz
determinant
\begin{equation}
\label{eq:xx-toeplitz-determinant}
    \langle \sigma_j^x\sigma_{j+r}^x\rangle
    =
    \det
    \begin{pmatrix}
        G_1&G_2&\cdots&G_r\\
        G_0&G_1&\cdots&G_{r-1}\\
        \vdots&\vdots&\ddots&\vdots\\
        G_{2-r}&G_{3-r}&\cdots&G_1
    \end{pmatrix}.
\end{equation}
This is a Toeplitz determinant because the matrix entries depend only on the
index difference.  Analogous determinant formulas exist for
$\langle\sigma_j^y\sigma_{j+r}^y\rangle$ and mixed transverse correlators, with
shifted kernels.  In conventions or models where the bipartite contraction
structure is absent, one should return to the Pfaffian formulas
\cref{eq:xx-correlator-pfaffian,eq:yy-correlator-pfaffian}.

\begin{remark}[Thermodynamic limit]
In the thermodynamic limit, $G_r$ is usually represented by a Fourier integral
\begin{equation}
    G_r=\int_{-\pi}^{\pi}\frac{\dd k}{2\pi}\,\ee^{-\ii kr}g(k),
\end{equation}
where $g(k)$ is the symbol determined by the BdG ground state.  Long-distance
spin correlations then reduce to the asymptotic analysis of Toeplitz
determinants.  This is the bridge between free-fermion diagonalization and exact
critical exponents in the Ising and XY chains.
\end{remark}

\section[Entanglement from covariance matrices]{Reduced density matrices and entanglement from covariance matrices}
\label{sec:entanglement-from-covariance-matrix}

Although the present notes mainly use covariance matrices for correlation
functions, the same object also determines entanglement in Gaussian states.  Let
$A$ be a subsystem consisting of $\ell$ lattice sites, hence $2\ell$ Majorana
operators.  Restrict the full covariance matrix to this subsystem:
\begin{equation}
    \Gamma_A=\Gamma|_A.
\end{equation}
The eigenvalues of $\Gamma_A$ come in pairs $\pm\ii\nu_\alpha$, with
$0\leq\nu_\alpha\leq1$, $\alpha=1,\ldots,\ell$.  The reduced density matrix
$\rho_A$ is Gaussian, and its von Neumann entropy is
\begin{equation}
\label{eq:gaussian-entanglement-entropy}
    S_A
    =
    -\Tr \rho_A\log\rho_A
    =
    \sum_{\alpha=1}^{\ell}
    \left[
    -\frac{1+\nu_\alpha}{2}\log\frac{1+\nu_\alpha}{2}
    -\frac{1-\nu_\alpha}{2}\log\frac{1-\nu_\alpha}{2}
    \right].
\end{equation}
Thus the same matrix $\Gamma$ simultaneously encodes spin correlators,
fermionic correlators, and entanglement spectra.

\section{Practical recipe}
\label{sec:correlation-functions-practical-recipe}

For a quadratic spin chain solved by Jordan--Wigner, Fourier transform, and
Bogoliubov transformation, the correlation-function workflow is:
\begin{enumerate}[label=(\roman*)]
    \item Diagonalize the BdG Hamiltonian and obtain the coherence factors
    $u_k,v_k$ or, equivalently, the occupied negative-energy BdG projector.
    \item Construct the normal and anomalous correlators $C_{ij}$ and $F_{ij}$,
    either in momentum space using \cref{eq:translation-invariant-C-r,eq:translation-invariant-F-r}
    or directly in real space.
    \item Convert $C$ and $F$ into the Majorana covariance matrix $\Gamma$ using
    \cref{eq:gamma-block-odd-odd,eq:gamma-block-even-even,eq:gamma-block-odd-even}.
    \item Compute local fermionic observables directly from $C$ and $F$, or from
    the relevant submatrix of $\Gamma$.
    \item Compute spin correlators by rewriting the Jordan--Wigner strings as
    ordered Majorana products and evaluating the Pfaffian
    $\Pf(-\ii\Gamma_{\mathcal{I}})$.
    \item In translation-invariant cases with the appropriate symmetry, simplify
    the Pfaffian to a Toeplitz determinant and use known determinant asymptotics
    for long-distance behavior.
\end{enumerate}

The main conceptual point is that the nonlocality introduced by the
Jordan--Wigner transformation does not destroy exact solvability.  It only
changes the final algebraic object from a two-point function to a Pfaffian of a
finite covariance submatrix.  This Pfaffian structure is the common language
behind the exact correlation functions of the TFIM, the XY chain, Kitaev-chain
spin representations, the two-dimensional Ising model, and more general
free-Majorana solvable systems.


\chapter{Transfer Matrix and Quantum--Classical Correspondence}
\label{chap:transfer-matrix-quantum-classical}

The previous chapters developed the free-fermion machinery used to diagonalize
one-dimensional quantum spin chains: Jordan--Wigner transformation, Fourier
transform, BdG blocks, Bogoliubov quasiparticles, and Gaussian correlation
functions. Transfer matrices provide a complementary route to exact solvability.
They are indispensable for two related reasons:
\begin{enumerate}[label=(\roman*)]
    \item a classical statistical model in one higher dimension can often be
    interpreted as an imaginary-time path integral of a quantum model; and
    \item a family of commuting transfer matrices can generate a quantum
    Hamiltonian by a logarithmic derivative or Hamiltonian limit.
\end{enumerate}
The first point is the quantum--classical correspondence. The second point is the
bridge to vertex-model and Bethe-ansatz integrability in later chapters.

\section{Classical transfer matrices}
\label{sec:classical-transfer-matrices}

Consider a classical statistical model on a strip with $M$ layers. Let
$\bm{s}_m$ denote the entire configuration on the $m$th layer. If the Boltzmann
weight couples only neighbouring layers, the partition function can be written as
\begin{equation}
\label{eq:transfer-matrix-general-partition-function}
    Z
    =
    \sum_{\bm{s}_1,\ldots,\bm{s}_M}
    T_{\bm{s}_1\bm{s}_2}
    T_{\bm{s}_2\bm{s}_3}
    \cdots
    T_{\bm{s}_M\bm{s}_1}
    =
    \Tr T^M .
\end{equation}
The matrix $T$ is the \emph{transfer matrix}. Its matrix element
$T_{\bm{s}\bm{s}'}$ is the Boltzmann weight for two adjacent layers
$\bm{s}$ and $\bm{s}'$, including a chosen assignment of intra-layer weights.
The trace implements periodic boundary conditions in the transfer direction.

If $T$ is finite-dimensional and diagonalizable, with eigenvalues
\begin{equation}
    \lambda_0,\lambda_1,\lambda_2,\ldots,
    \qquad
    \abs{\lambda_0}\geq \abs{\lambda_1}\geq\cdots,
\end{equation}
then
\begin{equation}
\label{eq:transfer-matrix-free-energy}
    Z=\sum_n \lambda_n^M .
\end{equation}
In the thermodynamic limit $M\to\infty$, the largest eigenvalue controls the free
energy per transfer step. If each layer contains $L$ sites, then
\begin{equation}
\label{eq:transfer-matrix-free-energy-density}
    f
    =
    -\frac{1}{\beta_{\rm cl}L}
    \lim_{M\to\infty}\frac{1}{M}\log Z
    =
    -\frac{1}{\beta_{\rm cl}L}\log \lambda_0,
\end{equation}
up to the convention-dependent normalization of the classical inverse temperature
$\beta_{\rm cl}$.

Transfer matrices also encode correlation lengths. For a layer observable $O$,
\begin{equation}
\label{eq:transfer-matrix-correlation-spectral}
    \langle O_m O_0\rangle
    =
    \frac{1}{Z}\Tr\left(O T^m O T^{M-m}\right).
\end{equation}
After inserting the spectral resolution of $T$, the connected correlation
function has the large-distance form
\begin{equation}
\label{eq:transfer-matrix-correlation-length}
    \langle O_m O_0\rangle_c
    \sim
    \left(\frac{\lambda_1}{\lambda_0}\right)^m,
    \qquad
    \xi^{-1}
    =
    \log\frac{\lambda_0}{\abs{\lambda_1}},
\end{equation}
provided $O$ has nonzero matrix element between the leading eigenstate and the
first relevant subleading eigenstate. Thus the spectrum of $T$ plays the same
role for a classical model that the spectrum of a Hamiltonian plays for a quantum
model.

\section{Warm-up: the one-dimensional classical Ising chain}
\label{sec:transfer-matrix-1d-classical-ising}

For the one-dimensional classical Ising model with spins $s_j=\pm1$ and periodic
boundary conditions,
\begin{equation}
\label{eq:classical-ising-1d-partition-function}
    Z_L(K)
    =
    \sum_{s_1,\ldots,s_L=\pm1}
    \exp\left(K\sum_{j=1}^{L}s_js_{j+1}\right),
    \qquad
    s_{L+1}=s_1,
\end{equation}
the transfer matrix is the $2\times2$ matrix
\begin{equation}
\label{eq:1d-ising-transfer-matrix}
    T_{ss'}=\ee^{Kss'},
    \qquad
    T=
    \begin{pmatrix}
        \ee^K & \ee^{-K} \\
        \ee^{-K} & \ee^K
    \end{pmatrix}.
\end{equation}
Its eigenvalues are
\begin{equation}
\label{eq:1d-ising-transfer-eigenvalues}
    \lambda_+=\ee^K+\ee^{-K}=2\cosh K,
    \qquad
    \lambda_- =\ee^K-\ee^{-K}=2\sinh K.
\end{equation}
Therefore
\begin{equation}
\label{eq:1d-ising-exact-partition-function}
    Z_L(K)=\lambda_+^L+\lambda_-^L.
\end{equation}
For finite $K$, the correlation length is finite:
\begin{equation}
\label{eq:1d-ising-correlation-length}
    \xi^{-1}
    =
    \log\frac{\lambda_+}{\lambda_-}
    =
    \log\coth K .
\end{equation}
This elementary example already contains the main lesson: thermodynamics and
correlations are spectral properties of the transfer matrix.

\section{Transfer matrix as imaginary-time evolution}
\label{sec:transfer-matrix-imaginary-time-evolution}

A quantum thermal partition function is
\begin{equation}
\label{eq:quantum-thermal-partition-function}
    Z_{\rm q}(\beta)
    =
    \Tr \ee^{-\beta H} .
\end{equation}
If imaginary time is discretized into $M$ steps,
\begin{equation}
    \beta=M\delta\tau,
\end{equation}
then
\begin{equation}
\label{eq:quantum-partition-transfer-matrix}
    Z_{\rm q}(\beta)
    =
    \Tr\left(\ee^{-\delta\tau H}\right)^M .
\end{equation}
Thus the short imaginary-time evolution operator
\begin{equation}
\label{eq:quantum-transfer-matrix-definition}
    T_{\rm q}(\delta\tau)=\ee^{-\delta\tau H}
\end{equation}
is itself a transfer matrix. Conversely, whenever a positive transfer matrix $T$
can be written as
\begin{equation}
\label{eq:hamiltonian-from-transfer-matrix}
    T=C\ee^{-a_\tau H},
\end{equation}
with $C>0$ and $a_\tau$ interpreted as an imaginary-time lattice spacing, it
defines the quantum Hamiltonian
\begin{equation}
\label{eq:hamiltonian-log-transfer-matrix}
    H
    =
    -\frac{1}{a_\tau}\log T
    +
    \frac{\log C}{a_\tau}\id .
\end{equation}
The additive constant proportional to $\id$ shifts the zero of energy and does not
affect normalized expectation values or energy gaps.

If $T\ket{n}=\lambda_n\ket{n}$ and
$T=C\ee^{-a_\tau H}$, then
\begin{equation}
\label{eq:transfer-eigenvalues-energy-levels}
    \lambda_n=C\ee^{-a_\tau E_n},
    \qquad
    E_n-E_0
    =
    -\frac{1}{a_\tau}
    \log\frac{\lambda_n}{\lambda_0} .
\end{equation}
Therefore transfer-matrix gaps become quantum energy gaps. In particular,
\begin{equation}
\label{eq:correlation-length-quantum-gap}
    \xi_\tau^{-1}
    =
    \log\frac{\lambda_0}{\abs{\lambda_1}}
    =
    a_\tau\,(E_1-E_0)
\end{equation}
when the first subleading eigenvalue controls the decay.

\begin{remark}[Hermiticity and positivity]
A symmetric positive transfer matrix defines a Hermitian Hamiltonian by
$H=-a_\tau^{-1}\log T$, up to a constant. More general transfer matrices may be
non-Hermitian. Such matrices are still useful, for example in stochastic and
non-equilibrium problems, but the direct interpretation as a Hermitian quantum
Hamiltonian then requires additional structure.
\end{remark}

\section{Trotter decomposition and the extra dimension}
\label{sec:trotter-extra-dimension}

The relation between a $d$-dimensional quantum model and a $(d+1)$-dimensional
classical model follows from the Trotter product formula. Suppose
\begin{equation}
\label{eq:hamiltonian-split-A-B}
    H=A+B,
\end{equation}
where $A$ and $B$ do not commute but have simple matrix elements in appropriate
bases. Then
\begin{equation}
\label{eq:trotter-first-order}
    \ee^{-\beta H}
    =
    \lim_{M\to\infty}
    \left(
        \ee^{-\delta\tau A}\ee^{-\delta\tau B}
    \right)^M,
    \qquad
    \delta\tau=\frac{\beta}{M}.
\end{equation}
A symmetric version is
\begin{equation}
\label{eq:trotter-symmetric}
    \ee^{-\delta\tau(A+B)}
    =
    \ee^{-\frac{\delta\tau}{2}A}
    \ee^{-\delta\tau B}
    \ee^{-\frac{\delta\tau}{2}A}
    +O(\delta\tau^3).
\end{equation}
Inserting complete sets of states between the $M$ imaginary-time factors produces
a classical sum over configurations on $M$ time slices. The additional discrete
index is imaginary time. Schematically,
\begin{equation}
\label{eq:quantum-classical-schematic}
    Z_{d\text{-dimensional quantum}}
    =
    \Tr\ee^{-\beta H}
    \quad\longleftrightarrow\quad
    Z_{(d+1)\text{-dimensional classical}} .
\end{equation}
This correspondence is exact in the Trotter limit $M\to\infty$ and becomes a
controlled approximation at finite $M$.

The correspondence is not a statement that the quantum model and the classical
model are identical microscopic objects. Rather, their partition functions and
correlation functions are two representations of the same Euclidean statistical
sum. Quantum noncommutativity is converted into classical coupling along the
imaginary-time direction.

\section{From the anisotropic two-dimensional Ising model to the TFIM}
\label{sec:anisotropic-ising-to-tfim-transfer}

The most important example for these notes is the correspondence between the
anisotropic two-dimensional classical Ising model and the one-dimensional TFIM.
Let
\begin{equation}
\label{eq:anisotropic-2d-ising-partition-function}
    Z_{\rm 2D}
    =
    \sum_{s_{j,m}=\pm1}
    \exp\left[
        K_x\sum_{j,m}s_{j,m}s_{j+1,m}
        +
        K_\tau\sum_{j,m}s_{j,m}s_{j,m+1}
    \right],
\end{equation}
where $j=1,\ldots,L$ is the spatial coordinate and $m=1,\ldots,M$ is the transfer
direction. The couplings $K_x$ and $K_\tau$ are classical dimensionless
couplings.

To match the Pauli convention used in these notes, take the classical variable
$s_j=\pm1$ to label the eigenvalue of $\sigma_j^x$:
\begin{equation}
\label{eq:x-basis-transfer}
    \sigma_j^x\ket{s_j}_x=s_j\ket{s_j}_x .
\end{equation}
The horizontal classical coupling is diagonal in this basis and gives the
operator
\begin{equation}
\label{eq:horizontal-transfer-operator}
    \exp\left(K_x\sum_{j=1}^{L}\sigma_j^x\sigma_{j+1}^x\right).
\end{equation}
The vertical one-site Boltzmann weight is
\begin{equation}
\label{eq:vertical-weight-matrix}
    W_{ss'}=\ee^{K_\tau ss'}
    =
    \begin{pmatrix}
        \ee^{K_\tau} & \ee^{-K_\tau} \\
        \ee^{-K_\tau} & \ee^{K_\tau}
    \end{pmatrix}_{s,s'} .
\end{equation}
In the $\sigma^x$ eigenbasis, the operator $\sigma^z$ flips $s$. Hence
\begin{equation}
\label{eq:vertical-weight-exp-sigmaz}
    W
    =
    A\ee^{\kappa\sigma^z},
    \qquad
    \tanh\kappa=\ee^{-2K_\tau},
    \qquad
    A=\sqrt{2\sinh(2K_\tau)} .
\end{equation}
Indeed,
\begin{equation}
    A\cosh\kappa=\ee^{K_\tau},
    \qquad
    A\sinh\kappa=\ee^{-K_\tau}.
\end{equation}
A convenient symmetric row-to-row transfer matrix is therefore
\begin{equation}
\label{eq:2d-ising-row-transfer-operator}
    T
    =
    A^L
    \exp\left(
        \frac{K_x}{2}\sum_j\sigma_j^x\sigma_{j+1}^x
    \right)
    \exp\left(
        \kappa\sum_j\sigma_j^z
    \right)
    \exp\left(
        \frac{K_x}{2}\sum_j\sigma_j^x\sigma_{j+1}^x
    \right).
\end{equation}
Using the symmetric Trotter formula, the strongly anisotropic limit is obtained
by setting
\begin{equation}
\label{eq:ising-tfim-parameter-scaling}
    K_x=\delta\tau J,
    \qquad
    \kappa=\delta\tau h,
    \qquad
    \tanh(\delta\tau h)=\ee^{-2K_\tau}.
\end{equation}
Then
\begin{equation}
\label{eq:2d-ising-transfer-tfim-limit}
    T
    =
    A^L
    \exp\left[-\delta\tau H_{\rm TFIM}+O(\delta\tau^3)\right],
\end{equation}
with
\begin{equation}
\label{eq:transfer-chapter-tfim-hamiltonian}
    H_{\rm TFIM}
    =
    -J\sum_j\sigma_j^x\sigma_{j+1}^x
    -h\sum_j\sigma_j^z .
\end{equation}
This is precisely the TFIM convention fixed earlier in these notes.

For small $\delta\tau$,
\begin{equation}
\label{eq:large-vertical-coupling}
    \ee^{-2K_\tau}
    =
    \tanh(\delta\tau h)
    =
    \delta\tau h+O(\delta\tau^3),
\end{equation}
so the imaginary-time coupling $K_\tau$ is large. The quantum limit is therefore
a strongly anisotropic limit of the classical two-dimensional model:
\begin{equation}
    K_x\to0,
    \qquad
    K_\tau\to\infty,
    \qquad
    K_x\ee^{2K_\tau}\;\text{fixed}.
\end{equation}
The classical critical condition of the anisotropic two-dimensional Ising model,
\begin{equation}
\label{eq:2d-ising-critical-condition}
    \sinh(2K_x)\sinh(2K_\tau)=1,
\end{equation}
then becomes
\begin{equation}
\label{eq:tfim-critical-from-ising}
    \frac{J}{h}=1,
    \qquad
    \text{or equivalently}\qquad
    h=J.
\end{equation}
In the common normalization $J=1$, this gives the TFIM critical point $h=1$.

\begin{remark}[Why the basis matters]
Many textbooks write the TFIM as
$-J\sum_j\sigma_j^z\sigma_{j+1}^z-\Gamma\sum_j\sigma_j^x$. These notes use the
rotated convention
$-J\sum_j\sigma_j^x\sigma_{j+1}^x-h\sum_j\sigma_j^z$ because it is adapted to the
Jordan--Wigner occupation convention $n_j=(1-\sigma_j^z)/2$. The two forms are
related by a global spin rotation. In the transfer-matrix derivation above, this
rotation is implemented by labeling the classical Ising variable as the eigenvalue
of $\sigma^x$ rather than $\sigma^z$.
\end{remark}

\section{Operator insertions and Euclidean correlators}
\label{sec:operator-insertions-euclidean-correlators}

The transfer-matrix correspondence also maps correlation functions. A diagonal
classical row observable $O(\bm{s})$ becomes an operator $O$ acting on the row
Hilbert space. Its correlation along the transfer direction is
\begin{equation}
\label{eq:transfer-correlator-as-quantum-correlator}
    \langle O_m O_0\rangle_{\rm cl}
    =
    \frac{1}{Z}
    \Tr\left(O T^m O T^{M-m}\right).
\end{equation}
If $T=C\ee^{-a_\tau H}$, then
\begin{equation}
\label{eq:euclidean-quantum-correlator}
    \langle O_m O_0\rangle_{\rm cl}
    =
    \frac{1}{Z_{\rm q}}
    \Tr\left(
        O(\tau)O(0)\ee^{-\beta H}
    \right),
    \qquad
    \tau=m a_\tau,
\end{equation}
where
\begin{equation}
\label{eq:imaginary-time-heisenberg-operator}
    O(\tau)=\ee^{\tau H}O\ee^{-\tau H}.
\end{equation}
At zero temperature, this becomes the ground-state Euclidean correlator. Its
large-$\tau$ decay is controlled by the lowest quantum excitation carrying the
quantum numbers of $O$:
\begin{equation}
\label{eq:euclidean-spectral-decay}
    \langle 0\vert O(\tau)O(0)\vert0\rangle_c
    \sim
    \ee^{-\tau(E_1-E_0)}.
\end{equation}
Thus the transfer-matrix spectrum, the classical correlation length, and the
quantum excitation gap are the same spectral information in different languages.

\section{Quantum-to-classical dictionary}
\label{sec:quantum-classical-dictionary}

For the TFIM/Ising example in the conventions of these notes, the essential
parameter dictionary is
\begin{equation}
\label{eq:tfim-ising-dictionary}
    \boxed{
    K_x=\delta\tau J,
    \qquad
    \ee^{-2K_\tau}=\tanh(\delta\tau h)
    }
\end{equation}
with row transfer matrix
\begin{equation}
\label{eq:tfim-ising-transfer-final}
    T
    \propto
    \exp\left(
        \frac{\delta\tau J}{2}
        \sum_j\sigma_j^x\sigma_{j+1}^x
    \right)
    \exp\left(
        \delta\tau h\sum_j\sigma_j^z
    \right)
    \exp\left(
        \frac{\delta\tau J}{2}
        \sum_j\sigma_j^x\sigma_{j+1}^x
    \right).
\end{equation}
The proportionality constant shifts the quantum free energy by an additive
constant. The operator in the exponential is $-\delta\tau H_{\rm TFIM}$ to
leading order.

The following table summarizes the correspondence.
\begin{center}
\begin{tabularx}{0.95\linewidth}{@{}L{0.31\linewidth}Y Y@{}}
\toprule
Concept & Classical transfer-matrix language & Quantum language \\
\midrule
Transfer direction & Layer index $m$ & Imaginary time $\tau=m a_\tau$ \\
Row configuration & Classical spin row $\bm{s}_m$ & Basis vector of the quantum Hilbert space \\
Transfer matrix & $T_{\bm{s}\bm{s}'}$ & $\ee^{-a_\tau H}$ up to normalization \\
Largest eigenvalue & Dominant Boltzmann weight & Ground-state energy density \\
Eigenvalue ratio & Correlation length $\xi^{-1}=\log(\lambda_0/\abs{\lambda_1})$ & Energy gap $a_\tau(E_1-E_0)$ \\
Classical anisotropy & Different $K_x,K_\tau$ & Spatial velocity and time-lattice spacing \\
Criticality & Divergent correlation length & Closing quantum gap \\
\bottomrule
\end{tabularx}
\end{center}

\section{Hamiltonian limits of commuting transfer matrices}
\label{sec:hamiltonian-limits-commuting-transfer-matrices}

In integrable lattice models one often has a one-parameter family of commuting
transfer matrices,
\begin{equation}
\label{eq:commuting-transfer-matrices}
    [T(u),T(v)]=0,
\end{equation}
where $u$ is a spectral parameter. If $T(u)$ is regular at a point $u=u_0$, the
associated quantum Hamiltonian is commonly obtained from the logarithmic
derivative
\begin{equation}
\label{eq:logarithmic-derivative-transfer-matrix}
    H
    \propto
    \left.
    \frac{\dd}{\dd u}\log T(u)
    \right|_{u=u_0}
    =
    \left.
    T(u)^{-1}\frac{\dd T(u)}{\dd u}
    \right|_{u=u_0},
\end{equation}
plus an additive and multiplicative normalization. Because all $T(u)$ commute,
the Hamiltonian obtained this way commutes with the entire transfer-matrix family:
\begin{equation}
\label{eq:hamiltonian-commutes-transfer-family}
    [H,T(v)]=0.
\end{equation}
This is the transfer-matrix origin of infinitely many conserved quantities in
Bethe-ansatz integrable chains. Later, the six-vertex transfer matrix will yield
the XXZ Hamiltonian in precisely this sense.

\section{Practical summary}
\label{sec:transfer-practical-summary}

The transfer-matrix method should be read as a spectral equivalence principle:
\begin{equation}
\label{eq:transfer-spectral-equivalence-principle}
    \boxed{
    \text{classical layer-to-layer transfer}
    \quad\Longleftrightarrow\quad
    \text{quantum imaginary-time evolution}
    }.
\end{equation}
For the purposes of these notes, the essential workflow is:
\begin{enumerate}[label=(\roman*)]
    \item Write the classical partition function as $Z=\Tr T^M$.
    \item Interpret the layer configuration as a quantum basis state.
    \item If possible, write $T$ as $C\ee^{-a_\tau H}$ or take a controlled
    Hamiltonian limit.
    \item Extract quantum energies from transfer-matrix eigenvalue ratios.
    \item Translate classical correlation lengths into quantum gaps.
    \item Use anisotropic limits to connect two-dimensional classical models to
    one-dimensional quantum Hamiltonians.
\end{enumerate}
The 2D classical Ising model and the 1D TFIM are the canonical example. The same
logic also underlies the relation between the six-vertex model and the XXZ chain,
the eight-vertex model and the XYZ chain, and more general tensor-network
interpretations of Euclidean path integrals.


\chapter{Duality}
\label{chap:duality}

Duality is one of the most efficient organizing principles in exactly solvable
models.  A duality rewrites the same partition function, Hamiltonian, operator
algebra, or low-energy field theory in variables adapted to a different physical
regime.  In the Ising family, the essential effect is the exchange
\begin{equation}
\label{eq:duality-weak-strong-schematic}
\begin{aligned}
    \text{ordered variables}
    &\longleftrightarrow
    \text{domain-wall or disorder variables},\\
    \text{strong coupling}
    &\longleftrightarrow
    \text{weak coupling}.
\end{aligned}
\end{equation}
Thus duality is not merely a computational trick.  It explains why an apparently
nonlocal description can become local after changing variables, why the
transverse-field Ising critical point is located at the self-dual coupling, and
why the two phases of the Ising chain are naturally described as order and
disorder phases.

This chapter develops duality at three levels:
\begin{enumerate}[label=(\roman*)]
    \item Kramers--Wannier duality of the two-dimensional classical Ising model;
    \item the operator Kramers--Wannier duality of the one-dimensional TFIM; and
    \item the precise logical status of self-duality as a criterion for
    criticality.
\end{enumerate}
The conventions are matched to the TFIM Hamiltonian used in these notes,
\begin{equation}
\label{eq:duality-chapter-tfim}
    H_{\rm TFIM}
    =
    -J\sum_j \sigma_j^x\sigma_{j+1}^x
    -h\sum_j \sigma_j^z,
    \qquad
    J,h\geq0 .
\end{equation}
The ordered phase is therefore the phase with long-range order in
$\sigma^x$, while the transverse-field or disordered phase is polarized by
$\sigma^z$.

\section{Classical Kramers--Wannier duality}
\label{sec:classical-kramers-wannier-duality}

Consider the anisotropic two-dimensional classical Ising model on a square
lattice,
\begin{equation}
\label{eq:duality-2d-classical-ising}
    Z(K_x,K_y)
    =
    \sum_{s_{r}=\pm1}
    \exp\left[
        K_x\sum_{\langle r r'\rangle_x}s_rs_{r'}
        +
        K_y\sum_{\langle r r'\rangle_y}s_rs_{r'}
    \right],
\end{equation}
where $K_x$ and $K_y$ are dimensionless couplings along the two lattice
directions.  The high-temperature expansion uses
\begin{equation}
\label{eq:duality-high-temp-bond}
    \ee^{K s_rs_{r'}}
    =
    \cosh K\left(1+s_rs_{r'}\tanh K\right).
\end{equation}
After expanding all bonds and summing over the Ising spins, only configurations
in which an even number of occupied bonds touches each site survive.  These are
closed loops on the original lattice.  Schematically,
\begin{equation}
\label{eq:duality-high-temp-loops}
    Z(K_x,K_y)
    =
    2^N(\cosh K_x)^{E_x}(\cosh K_y)^{E_y}
    \sum_{\mathcal{C}\,\text{closed}}
    (\tanh K_x)^{n_x(\mathcal{C})}
    (\tanh K_y)^{n_y(\mathcal{C})}.
\end{equation}
Here $N$ is the number of spins, $E_x,E_y$ are the numbers of bonds in the two
directions, and $n_x(\mathcal{C}),n_y(\mathcal{C})$ count occupied loop edges.

The low-temperature expansion instead starts from the two ferromagnetic ground
states.  A configuration is specified by domain walls on the dual lattice.  A
domain wall segment crossing a bond of coupling $K$ costs a Boltzmann factor
$\ee^{-2K}$.  Matching the loop weights in the high-temperature expansion to the
domain-wall weights of a dual Ising model gives the dual couplings
\begin{equation}
\label{eq:duality-classical-dual-couplings-exp}
    \ee^{-2K_x^*}=\tanh K_y,
    \qquad
    \ee^{-2K_y^*}=\tanh K_x .
\end{equation}
Equivalently,
\begin{equation}
\label{eq:duality-classical-dual-couplings-sinh}
    \sinh(2K_x^*)\sinh(2K_y)=1,
    \qquad
    \sinh(2K_y^*)\sinh(2K_x)=1 .
\end{equation}
The exchange of directions in \cref{eq:duality-classical-dual-couplings-exp}
comes from the fact that a horizontal bond of the dual lattice crosses a vertical
bond of the original lattice, and conversely.

Up to a non-singular prefactor depending on the couplings and the boundary
conditions, the partition function of the original model equals the partition
function of the dual model,
\begin{equation}
\label{eq:duality-classical-partition-map}
    Z(K_x,K_y)
    \propto
    Z(K_x^*,K_y^*) .
\end{equation}
The proportionality factor affects the analytic part of the free energy but not
the location of a nonanalytic singularity in the thermodynamic limit.

For the isotropic model $K_x=K_y=K$, the self-dual condition is
\begin{equation}
\label{eq:duality-isotropic-ising-critical}
    \ee^{-2K_c}=\tanh K_c,
    \qquad
    \sinh(2K_c)=1,
    \qquad
    K_c=\frac12\log(1+\sqrt2).
\end{equation}
For the anisotropic model, the self-dual curve is
\begin{equation}
\label{eq:duality-anisotropic-ising-critical}
    \sinh(2K_x)\sinh(2K_y)=1 .
\end{equation}
This is the same condition that appeared in the transfer-matrix derivation of
the quantum--classical correspondence in the previous chapter.
In the strongly anisotropic limit, it reduces to the TFIM condition $h=J$.

\section{Operator duality in the transverse-field Ising chain}
\label{sec:operator-duality-tfim}

The quantum version of Kramers--Wannier duality is an exact change of operator
variables.  We first work on an infinite chain or in the bulk of a long open
chain, where boundary subtleties do not obscure the local algebra.  Define dual
Pauli operators on the bonds of the original lattice by
\begin{equation}
\label{eq:duality-dual-spin-definitions}
    \mu_{j+\frac12}^z
    =
    \sigma_j^x\sigma_{j+1}^x,
    \qquad
    \mu_{j+\frac12}^x
    =
    \prod_{\ell\leq j}\sigma_\ell^z .
\end{equation}
The first operator measures whether neighbouring $\sigma^x$ spins are aligned;
the second is a semi-infinite string of transverse-field operators and is called
a \emph{disorder operator}.

The operators in \cref{eq:duality-dual-spin-definitions} obey the Pauli algebra.
Indeed,
\begin{equation}
\label{eq:duality-mu-pauli-algebra}
    (\mu_{j+\frac12}^x)^2=(\mu_{j+\frac12}^z)^2=\id,
    \qquad
    \{\mu_{j+\frac12}^x,\mu_{j+\frac12}^z\}=0,
\end{equation}
and dual operators on distinct dual sites commute.  The anticommutation at the
same dual site occurs because $\mu_{j+1/2}^x$ overlaps with
$\sigma_j^x\sigma_{j+1}^x$ on exactly one site.  For different dual sites, the
number of overlaps is either zero or two, so the operators commute.

The inverse local identities are
\begin{equation}
\label{eq:duality-inverse-local-identities}
    \sigma_j^x\sigma_{j+1}^x=\mu_{j+\frac12}^z,
    \qquad
    \sigma_j^z=\mu_{j-\frac12}^x\mu_{j+\frac12}^x .
\end{equation}
Substituting these into \cref{eq:duality-chapter-tfim} gives
\begin{equation}
\label{eq:duality-tfim-dual-hamiltonian-mu}
    H_{\rm TFIM}
    =
    -J\sum_j\mu_{j+\frac12}^z
    -h\sum_j\mu_{j-\frac12}^x\mu_{j+\frac12}^x .
\end{equation}
If we rotate the dual spin axes by defining
\begin{equation}
\label{eq:duality-rotated-dual-spins}
    \widetilde{\sigma}_{j+\frac12}^x=\mu_{j+\frac12}^z,
    \qquad
    \widetilde{\sigma}_{j+\frac12}^z=\mu_{j+\frac12}^x,
\end{equation}
then the dual Hamiltonian becomes
\begin{equation}
\label{eq:duality-tfim-dual-hamiltonian-tilde}
    H_{\rm TFIM}(J,h)
    \longmapsto
    -h\sum_j
    \widetilde{\sigma}_{j-\frac12}^z
    \widetilde{\sigma}_{j+\frac12}^z
    -J\sum_j
    \widetilde{\sigma}_{j+\frac12}^x .
\end{equation}
This is the same TFIM form after relabeling the axes, with the two couplings
exchanged:
\begin{equation}
\label{eq:duality-tfim-coupling-exchange}
    (J,h)\longleftrightarrow (h,J).
\end{equation}
Equivalently, in terms of the dimensionless ratio $g=h/J$,
\begin{equation}
\label{eq:duality-g-to-inverse-g}
    g\longmapsto g^*=\frac1g,
\end{equation}
up to an overall energy scale.

\begin{remark}[Why the dual lattice is shifted]
The dual spins live on bonds $j+1/2$, not on the original sites $j$.  This shift
is not cosmetic.  It encodes the fact that the elementary dual variable is a
domain-wall variable: it naturally lives between two neighbouring order
variables.  In the continuum limit, this shift becomes invisible, but on the
lattice it is essential for getting the operator algebra and boundary conditions
right.
\end{remark}

\section{Order and disorder operators}
\label{sec:order-disorder-operators}

The duality map makes the distinction between order and disorder completely
explicit.  In the original spin language, the ferromagnetic order parameter is
$\sigma_j^x$.  In the dual language, a two-point order correlator becomes a
string of dual domain-wall variables:
\begin{equation}
\label{eq:duality-order-correlator-string}
    \sigma_i^x\sigma_j^x
    =
    \prod_{m=i}^{j-1}
    \mu_{m+\frac12}^z,
    \qquad i<j .
\end{equation}
Conversely, the dual order operator $\mu^x$ is a disorder string in the original
variables.  Its two-point function is
\begin{equation}
\label{eq:duality-disorder-correlator-string}
    \mu_{i+\frac12}^x\mu_{j+\frac12}^x
    =
    \prod_{\ell=i+1}^{j}\sigma_\ell^z,
    \qquad i<j .
\end{equation}
Thus Kramers--Wannier duality exchanges a local order operator with a nonlocal
string operator.

In the ordered phase $J>h$, the original spins have long-range $\sigma^x$ order,
while the dual variables are disordered.  In the paramagnetic phase $h>J$, the
original order is absent, but the disorder operator has long-range order:
\begin{equation}
\label{eq:duality-order-disorder-phases}
    \begin{array}{ccl}
    J>h &:& \langle \sigma^x\rangle\neq0,
    \qquad \langle \mu^x\rangle=0,\\[2pt]
    h>J &:& \langle \sigma^x\rangle=0,
    \qquad \langle \mu^x\rangle\neq0 .
    \end{array}
\end{equation}
This is the quantum-chain version of the classical statement that the low-
temperature expansion in domain walls is dual to the high-temperature expansion
in closed loops.

\section{Finite chains, constraints, and boundary conditions}
\label{sec:duality-boundary-conditions}

The local formulas above are exact in the bulk.  On a finite chain, duality also
acts on global sectors and boundary conditions.  This point is important because
many solvable-chain calculations in these notes use periodic boundary
conditions, parity sectors, and momentum quantization.

For an open chain with $L$ sites,
\begin{equation}
\label{eq:duality-open-tfim}
    H_{\rm OBC}
    =
    -J\sum_{j=1}^{L-1}\sigma_j^x\sigma_{j+1}^x
    -h\sum_{j=1}^{L}\sigma_j^z,
\end{equation}
define
\begin{equation}
\label{eq:duality-open-definitions}
    \mu_{j+\frac12}^z=\sigma_j^x\sigma_{j+1}^x,
    \qquad
    \mu_{j+\frac12}^x=\prod_{\ell=1}^{j}\sigma_\ell^z,
    \qquad
    j=1,\ldots,L-1.
\end{equation}
Then
\begin{equation}
\label{eq:duality-open-boundary-identities}
    \sigma_j^z
    =
    \mu_{j-\frac12}^x\mu_{j+\frac12}^x,
    \qquad
    \mu_{\frac12}^x\equiv\id,
    \qquad
    \mu_{L+\frac12}^x\equiv P_z,
\end{equation}
where
\begin{equation}
\label{eq:duality-global-z2-charge}
    P_z=\prod_{j=1}^{L}\sigma_j^z
\end{equation}
is the global Ising $\ZZ_2$ symmetry generator.  Therefore the dual open-chain
Hamiltonian contains boundary terms involving the fixed endpoint variables
$\mu_{1/2}^x$ and $\mu_{L+1/2}^x$.  In a fixed $P_z$ sector these are ordinary
boundary conditions for the dual chain.

For a periodic chain, the same issue appears as a relation between symmetry
sectors and boundary twists.  The definitions
\begin{equation}
\label{eq:duality-periodic-definitions}
    \mu_{j+\frac12}^z=\sigma_j^x\sigma_{j+1}^x,
    \qquad
    \mu_{j+\frac12}^x=\prod_{\ell=1}^{j}\sigma_\ell^z
\end{equation}
lead to
\begin{equation}
\label{eq:duality-periodic-constraint}
    \prod_{j=1}^{L}\mu_{j+\frac12}^z
    =
    \prod_{j=1}^{L}\sigma_j^x\sigma_{j+1}^x
    =
    \id,
\end{equation}
and
\begin{equation}
\label{eq:duality-periodic-twist}
    \mu_{L+\frac12}^x=P_z\mu_{\frac12}^x .
\end{equation}
Hence a periodic original chain maps to a dual chain with a constraint and a
boundary condition determined by the original global $\ZZ_2$ charge.  Conversely,
the global symmetry sector of the dual model determines the boundary condition
of the original model.  This is the spin-chain analogue of the parity-dependent
fermion boundary conditions encountered in the Jordan--Wigner transformation.

\begin{remark}[Duality is not always an on-site unitary]
On an infinite chain, Kramers--Wannier duality is a faithful isomorphism of local
operator algebras.  On a finite periodic chain, it is not simply an on-site
unitary map from one unconstrained $2^L$-dimensional spin Hilbert space to
another unconstrained $2^L$-dimensional spin Hilbert space.  One must keep track
of constraints, global sectors, and boundary twists.  Ignoring this point is a
common source of erroneous factors of two in partition functions and spectra.
\end{remark}

\section{Self-duality and the TFIM critical point}
\label{sec:self-duality-tfim-critical-point}

The TFIM is self-dual when $J=h$.  Since duality exchanges $J$ and $h$, the point
$J=h$ is invariant under the transformation.  The exact free-fermion spectrum of
\cref{eq:duality-chapter-tfim} is
\begin{equation}
\label{eq:duality-tfim-spectrum}
    E(k)
    =
    2\sqrt{
        (h-J\cos k)^2
        +
        (J\sin k)^2
    }.
\end{equation}
For $J,h>0$, the excitation gap is obtained at $k=0$:
\begin{equation}
\label{eq:duality-tfim-gap}
    \Delta_{\rm gap}=2\abs{h-J}.
\end{equation}
Thus the self-dual point is indeed the quantum critical point,
\begin{equation}
\label{eq:duality-tfim-critical-point}
    h=J .
\end{equation}
At this point the order and disorder descriptions are equally natural, the
Majorana mass in the continuum limit changes sign, and the scaling limit is the
Ising conformal field theory.

This conclusion can also be obtained from the two-dimensional classical Ising
model.  The transfer-matrix chapter showed that the anisotropic classical
critical condition
\begin{equation}
    \sinh(2K_x)\sinh(2K_\tau)=1
\end{equation}
reduces, in the strongly anisotropic quantum limit, to $h=J$.  Kramers--Wannier
duality and the quantum--classical correspondence therefore locate the same
critical point.

\section{When does self-duality determine a critical point?}
\label{sec:self-duality-criterion-limitations}

Self-duality is powerful, but it is not by itself a proof of criticality.  The
precise statement is conditional.

\begin{principle}[Self-duality criterion]
Suppose a model depends on a coupling ratio $g$, and an exact duality maps
$g\mapsto 1/g$.  If the model has exactly one phase transition separating two
phases exchanged by the duality, and if the transition is not hidden by an
extended critical region or by additional symmetry-breaking/topological
structure, then the transition must occur at the self-dual point $g=1$.
\end{principle}

The assumptions in this principle matter.  A self-dual point can fail to be a
unique continuous critical point for several reasons:
\begin{enumerate}[label=(\roman*)]
    \item there may be more than one transition, appearing in dual pairs away
    from the self-dual point;
    \item there may be a first-order transition at the self-dual point rather
    than a continuous critical point;
    \item there may be an extended critical phase rather than an isolated
    critical point; or
    \item the duality may exchange sectors or boundary conditions rather than
    mapping a single finite-size Hamiltonian strictly to itself.
\end{enumerate}
For the ferromagnetic TFIM with $J,h>0$, the extra spectral information in
\cref{eq:duality-tfim-spectrum} verifies that there is a single gap closing at
$J=h$.  Thus in this model self-duality is not merely suggestive; it correctly
identifies the exact critical point.

\begin{remark}[Sign conventions]
The sign of $h$ in \cref{eq:duality-chapter-tfim} is removable by the global
unitary $\prod_j\sigma_j^x$, which flips $\sigma_j^z\mapsto-\sigma_j^z$ and leaves
$\sigma_j^x\sigma_{j+1}^x$ invariant.  On a bipartite chain, the sign of $J$ can
also be changed by a staggered rotation, up to boundary subtleties for periodic
chains.  Therefore the essential ferromagnetic TFIM phase diagram is usually
drawn in terms of $\abs{h}/\abs{J}$, with the critical point at
$\abs{h}=\abs{J}$.
\end{remark}

\section{Relation to fermions and Majorana mass}
\label{sec:duality-fermions-majorana-mass}

Under the Jordan--Wigner transformation, the TFIM becomes a quadratic Majorana
or BdG Hamiltonian.  In the convention of these notes, the single-particle gap is
controlled by the difference $h-J$.  Near $k=0$ and near $h=J$, the BdG spectrum
has the relativistic form
\begin{equation}
\label{eq:duality-majorana-low-energy-spectrum}
    E(k)
    \simeq
    2\sqrt{(h-J)^2+J h\, k^2}
\end{equation}
up to higher-order lattice corrections.  The continuum theory is therefore a
massive Majorana theory with mass proportional to
\begin{equation}
\label{eq:duality-majorana-mass}
    m\propto h-J.
\end{equation}
Kramers--Wannier duality exchanges the two signs of this mass.  The ordered
Ising phase and the disordered Ising phase are therefore two massive Majorana
phases separated by the massless Majorana point $m=0$.

This observation is also the bridge to the Kitaev-chain interpretation.  The
TFIM ordered/disordered transition corresponds, after fermionization, to the
transition between the topological and trivial phases of a one-dimensional
spinless $p$-wave superconductor.  In the spin language this distinction is
expressed as order versus disorder; in the fermion language it is expressed as a
change in the Majorana mass sign and, for open chains, the presence or absence
of unpaired boundary Majorana modes.  The next chapter will formulate this same
transition in terms of winding numbers and BdG topology.

\section{Summary}
\label{sec:duality-summary}

The essential identities of this chapter are
\begin{equation}
\label{eq:duality-summary-identities}
    \mu_{j+\frac12}^z=\sigma_j^x\sigma_{j+1}^x,
    \qquad
    \mu_{j+\frac12}^x=\prod_{\ell\leq j}\sigma_\ell^z,
    \qquad
    \sigma_j^z=\mu_{j-\frac12}^x\mu_{j+\frac12}^x .
\end{equation}
They imply the bulk Hamiltonian mapping
\begin{equation}
\label{eq:duality-summary-hamiltonian-map}
    -J\sum_j\sigma_j^x\sigma_{j+1}^x
    -h\sum_j\sigma_j^z
    \quad\longleftrightarrow\quad
    -h\sum_j\widetilde{\sigma}_j^x\widetilde{\sigma}_{j+1}^x
    -J\sum_j\widetilde{\sigma}_j^z,
\end{equation}
with the two couplings exchanged.  Consequently the TFIM is self-dual at
$h=J$, and the exact spectrum confirms that this point is the unique continuous
critical point of the ferromagnetic chain.

The conceptual lesson is that duality turns nonlocal degrees of freedom into
local ones.  The same physical theory can be described either by the original
order variables $\sigma^x$, by the disorder variables $\mu^x$, by Jordan--Wigner
fermions, or by Majorana fields.  This is precisely the equivalence-map viewpoint
that underlies the organization of these lecture notes.


\chapter{Topology of Two-Band Hamiltonians}
\chaptermark{Topology of Two-Band Hamiltonians}
\label{chap:topology-two-band-hamiltonians}

The previous chapters developed the algebraic machinery needed to reduce many
one-dimensional spin models to quadratic fermions and BdG momentum-space
blocks.  Once the Hamiltonian has been written as a family of finite-dimensional
matrices depending on momentum, exact solvability is not only a question of
finding the spectrum.  It also becomes meaningful to ask whether the family of
Bloch or BdG Hamiltonians can be continuously deformed into a trivial one while
keeping the bulk gap open and preserving the required symmetries.

This chapter gives the minimal topological toolkit needed for the later chapters
on the Kitaev chain, the SSH chain, and the XY/TFIM mother model.  The central
object is a two-band Hamiltonian
\begin{equation}
\label{eq:topo-two-band-general}
    H(k)=d_0(k)\id+\bm d(k)\cdot\bm\sigma,
    \qquad
    \bm d(k)=(d_x(k),d_y(k),d_z(k)),
\end{equation}
where $k$ belongs to the Brillouin zone.  For BdG Hamiltonians the Pauli
matrices act in Nambu space, while for ordinary band Hamiltonians they may act
in sublattice, orbital, or pseudospin space.  The same formulas therefore apply
to two physically different but mathematically parallel settings:
\begin{equation}
\label{eq:topo-two-realizations}
\begin{array}{ccl}
\text{SSH chain} &:& \text{sublattice two-band Hamiltonian},\\[2pt]
\text{Kitaev chain} &:& \text{BdG two-band Hamiltonian in Nambu space}.
\end{array}
\end{equation}
The distinction between these two realizations is important: the Pauli matrices
have the same algebra, but the symmetries and the meaning of the two components
of the spinor are different.

\section{Two-band Hamiltonians and maps to the Bloch sphere}
\label{sec:topo-two-band-bloch-sphere}

The spectrum of \cref{eq:topo-two-band-general} is
\begin{equation}
\label{eq:topo-two-band-spectrum}
    E_\pm(k)=d_0(k)\pm \abs{\bm d(k)}.
\end{equation}
Thus a two-band model is gapped at momentum $k$ precisely when
\begin{equation}
\label{eq:topo-gap-condition}
    \abs{\bm d(k)}\neq 0.
\end{equation}
For topology the scalar term $d_0(k)\id$ is usually irrelevant: it shifts the two
energies by the same amount but does not change the eigenvectors.  If the model
is gapped over the entire Brillouin zone, one may define the normalized vector
\begin{equation}
\label{eq:topo-normalized-d-vector}
    \widehat{\bm d}(k)=\frac{\bm d(k)}{\abs{\bm d(k)}}\in S^2 .
\end{equation}
Therefore a gapped two-band Hamiltonian defines a map
\begin{equation}
\label{eq:topo-map-bz-to-sphere}
    \widehat{\bm d}:\mathrm{BZ}\longrightarrow S^2.
\end{equation}
A topological phase is a homotopy class of such maps, possibly restricted by
symmetry.

The projectors onto the upper and lower bands of $\bm d\cdot\bm\sigma$ are
\begin{equation}
\label{eq:topo-band-projectors}
    P_+(k)=\frac{1+\widehat{\bm d}(k)\cdot\bm\sigma}{2},
    \qquad
    P_-(k)=\frac{1-\widehat{\bm d}(k)\cdot\bm\sigma}{2}.
\end{equation}
If the lower band is occupied, the occupied-band geometry is completely encoded
in $P_-(k)$.

\begin{principle}[Gap closing is the only way to change topology]
For a two-band Hamiltonian, the topological invariant can change only if there
exists a momentum $k_0$ such that
\begin{equation}
    \bm d(k_0)=0.
\end{equation}
At this point the two bands touch, the normalized vector
$\widehat{\bm d}(k)$ is undefined, and the homotopy class of the map can change.
\end{principle}

This principle is the topological version of the gap-closing conditions already
encountered in the BdG chapter.  For example, the Kitaev-chain spectrum
\begin{equation}
    E(k)=\sqrt{\xi(k)^2+\abs{\Delta(k)}^2}
\end{equation}
can close only when both $\xi(k)$ and $\Delta(k)$ vanish.  Since the nearest-neighbor
$p$-wave gap is proportional to $\sin k$, the possible gap closings occur at
$k=0,\pi$.

\section{Berry connection, Berry curvature, and Chern number}
\label{sec:topo-berry-curvature-chern}

Although much of this chapter focuses on one-dimensional models, the general
two-band language is most transparent if we first state the two-dimensional
Berry-curvature formula.  Let $\ket{u_-(\bm k)}$ be a normalized eigenvector of
the lower band.  The Berry connection and Berry curvature are
\begin{equation}
\label{eq:topo-berry-connection-curvature}
    A_a(\bm k)=\ii\braket{u_-(\bm k)}{\partial_{k_a}u_-(\bm k)},
    \qquad
    F_{xy}=\partial_{k_x}A_y-\partial_{k_y}A_x .
\end{equation}
Equivalently, in a manifestly gauge-invariant form,
\begin{equation}
\label{eq:topo-projector-curvature}
    F_{xy}(\bm k)
    =
    \ii\Tr\Bigl(P_-(\bm k)\bigl[\partial_{k_x}P_-(\bm k),
    \partial_{k_y}P_-(\bm k)\bigr]\Bigr).
\end{equation}
Substituting \cref{eq:topo-band-projectors} gives the standard two-band formula
\begin{equation}
\label{eq:topo-two-band-berry-curvature}
    F_{xy}(\bm k)
    =
    \frac12\widehat{\bm d}\cdot
    \left(
        \partial_{k_x}\widehat{\bm d}
        \times
        \partial_{k_y}\widehat{\bm d}
    \right).
\end{equation}
The first Chern number of the occupied band is
\begin{equation}
\label{eq:topo-chern-number}
    C
    =
    \frac{1}{2\pi}\int_{\mathrm{BZ}} F_{xy}(\bm k)\,\dd^2k
    =
    \frac{1}{4\pi}\int_{\mathrm{BZ}}
    \widehat{\bm d}\cdot
    \left(
        \partial_{k_x}\widehat{\bm d}
        \times
        \partial_{k_y}\widehat{\bm d}
    \right)\dd^2k .
\end{equation}
Geometrically, $C$ counts how many times the Brillouin-zone torus wraps the unit
sphere under the map $\widehat{\bm d}$.

\begin{example}[A standard lattice Chern-insulator block]
\label{ex:topo-qwz-model}
Consider
\begin{equation}
\label{eq:topo-qwz-d-vector}
    \bm d(\bm k)=
    \bigl(
        \sin k_x,
        \sin k_y,
        m+\cos k_x+\cos k_y
    \bigr).
\end{equation}
The gap can close only at the four high-symmetry points
$(0,0)$, $(\pi,0)$, $(0,\pi)$, and $(\pi,\pi)$, giving critical values
$m=-2,0,2$.  With the orientation convention in
\cref{eq:topo-chern-number}, the occupied-band Chern number is
\begin{equation}
\label{eq:topo-qwz-chern-values}
    C=\begin{cases}
    0, & m>2,\\
    -1, & 0<m<2,\\
    +1, & -2<m<0,\\
    0, & m<-2.
    \end{cases}
\end{equation}
This example is not one of the main solvable chains of these notes, but it is a
useful reference point: in two dimensions the full sphere $S^2$ can be wrapped,
whereas in one-dimensional chiral models topology usually comes from restricting
$\widehat{\bm d}$ to an equator.
\end{example}

\section{Why one-dimensional two-band topology needs symmetry}
\label{sec:topo-1d-needs-symmetry}

For a generic one-dimensional two-band Hamiltonian, the map is
\begin{equation}
    S^1_{\rm BZ}\longrightarrow S^2 .
\end{equation}
Since $\pi_1(S^2)=0$, every loop on the sphere can be contracted to a point.
Therefore a generic gapped one-dimensional two-band Hamiltonian has no stable
integer winding number.

A nontrivial one-dimensional invariant appears when a symmetry forces
$\widehat{\bm d}(k)$ to lie on a great circle of the sphere.  The most important
case for these notes is chiral symmetry.

\begin{definition}[Chiral symmetry]
A two-band Hamiltonian has chiral symmetry if there exists a unitary matrix
$S$ with $S^2=\id$ such that
\begin{equation}
\label{eq:topo-chiral-symmetry-definition}
    S H(k) S^{-1}=-H(k).
\end{equation}
For the conventional choice $S=\sigma_z$, chiral symmetry forbids the identity
and $\sigma_z$ terms, so that
\begin{equation}
\label{eq:topo-chiral-two-band}
    H(k)=d_x(k)\sigma_x+d_y(k)\sigma_y.
\end{equation}
\end{definition}

The vector $\bm d(k)$ now lies in the $d_x$--$d_y$ plane.  If the system is
gapped, the complex number
\begin{equation}
\label{eq:topo-q-definition}
    q(k)=d_x(k)+\ii d_y(k)
\end{equation}
never vanishes.  Thus $q(k)$ defines a map
\begin{equation}
\label{eq:topo-q-map}
    S^1_{\rm BZ}\longrightarrow \CC\setminus\{0\}\simeq S^1 .
\end{equation}
This map is classified by an integer winding number.

\section{Winding number of chiral two-band Hamiltonians}
\label{sec:topo-winding-number}

For the chiral Hamiltonian
\begin{equation}
\label{eq:topo-chiral-hamiltonian}
    H(k)=d_x(k)\sigma_x+d_y(k)\sigma_y
    =
    \begin{pmatrix}
        0 & q(k)^*\\
        q(k) & 0
    \end{pmatrix},
    \qquad
    q(k)=d_x(k)+\ii d_y(k),
\end{equation}
where the second equality follows from the standard Pauli matrices.  The winding
number is
\begin{equation}
\label{eq:topo-winding-number}
    \nu
    =
    \frac{1}{2\pi\ii}\int_{\rm BZ}\dd k\,\partial_k\log q(k).
\end{equation}
Equivalently, writing
\begin{equation}
    q(k)=\abs{q(k)}\ee^{\ii\varphi(k)},
\end{equation}
one obtains
\begin{equation}
\label{eq:topo-winding-angle}
    \nu
    =
    \frac{1}{2\pi}\int_{\rm BZ}\dd k\,\partial_k\varphi(k)
    =
    \frac{\varphi(2\pi)-\varphi(0)}{2\pi}\in\ZZ .
\end{equation}
Thus $\nu$ counts how many times the curve
\begin{equation}
    k\longmapsto (d_x(k),d_y(k))
\end{equation}
winds around the origin.

\begin{proposition}[Topological stability of the winding number]
If $q_s(k)$ is a smooth one-parameter family with $q_s(k)\neq 0$ for all
$k$ and all $s\in[0,1]$, then the winding number \cref{eq:topo-winding-number}
is independent of $s$.
\end{proposition}

\begin{proof}
The $s$-derivative of the winding number is
\begin{equation}
    \partial_s\nu
    =
    \frac{1}{2\pi\ii}\int_{\rm BZ}\dd k\,
    \partial_k\left(q^{-1}\partial_s q\right).
\end{equation}
Since the Brillouin zone is a circle and $q^{-1}\partial_s q$ is single-valued
when $q$ never vanishes, the integral of the total derivative is zero.
\end{proof}

The winding number can therefore change only when $q(k_0)=0$ for some momentum.
This is the one-dimensional chiral version of the general gap-closing principle.

\section{Zak phase and its relation to winding}
\label{sec:topo-zak-phase}

For a one-dimensional occupied band, the Berry phase over the Brillouin zone is
called the Zak phase:
\begin{equation}
\label{eq:topo-zak-phase-definition}
    \gamma_{\rm Zak}
    =
    \oint_{\rm BZ}\dd k\,
    \ii\braket{u_-(k)}{\partial_k u_-(k)}
    \quad\mod 2\pi .
\end{equation}
For a chiral two-band Hamiltonian, one may choose the lower-band eigenvector as
\begin{equation}
\label{eq:topo-chiral-eigenvector}
    \ket{u_-(k)}
    =
    \frac{1}{\sqrt{2}}
    \begin{pmatrix}
        -\ee^{-\ii\varphi(k)}\\
        1
    \end{pmatrix},
    \qquad
    q(k)=\abs{q(k)}\ee^{\ii\varphi(k)}.
\end{equation}
This gives
\begin{equation}
\label{eq:topo-zak-winding-relation}
    \gamma_{\rm Zak}
    =
    +\frac12\int_{\rm BZ}\dd k\,\partial_k\varphi(k)
    =
    +\pi\nu
    \quad\mod 2\pi .
\end{equation}
The sign depends on the eigenvector and Pauli-matrix convention, but the
quantized statement is invariant:
\begin{equation}
\label{eq:topo-zak-quantized-statement}
    \gamma_{\rm Zak}=0\ \text{or}\ \pi\quad\mod 2\pi
\end{equation}
for a chiral or inversion-symmetric two-band chain.  Thus the parity of the
winding number is visible as a Berry phase.

\section{SSH chain as a chiral two-band model}
\label{sec:topo-ssh-chain}

The Su--Schrieffer--Heeger chain is the canonical ordinary-band example.  It has
two sublattices $A$ and $B$ in each unit cell, with intracell hopping $t_1$ and
intercell hopping $t_2$:
\begin{equation}
\label{eq:topo-ssh-real-space}
    H_{\rm SSH}
    =
    \sum_j
    \left(
        t_1 a_j^\dagger b_j
        +t_2 a_{j+1}^\dagger b_j
        +\text{h.c.}
    \right).
\end{equation}
Using the Bloch spinor
\begin{equation}
    \Psi_k=\begin{pmatrix}a_k\\ b_k\end{pmatrix},
\end{equation}
the Hamiltonian becomes
\begin{equation}
\label{eq:topo-ssh-bloch-hamiltonian}
    H_{\rm SSH}
    =
    \sum_k\Psi_k^\dagger
    \begin{pmatrix}
        0 & q_{\rm SSH}(k)^*\\
        q_{\rm SSH}(k) & 0
    \end{pmatrix}
    \Psi_k,
    \qquad
    q_{\rm SSH}(k)=t_1+t_2\ee^{\ii k}.
\end{equation}
Equivalently,
\begin{equation}
\label{eq:topo-ssh-d-vector}
    d_x(k)=t_1+t_2\cos k,
    \qquad
    d_y(k)=t_2\sin k,
    \qquad
    d_z(k)=0.
\end{equation}
The curve $q_{\rm SSH}(k)$ is a circle of radius $\abs{t_2}$ centered at $t_1$ on
the real axis.  It winds around the origin if and only if the origin lies inside
this circle.  For $t_1,t_2>0$ and the convention in
\cref{eq:topo-ssh-bloch-hamiltonian},
\begin{equation}
\label{eq:topo-ssh-winding-result}
    \nu_{\rm SSH}
    =
    \begin{cases}
    0, & \abs{t_1}>\abs{t_2},\\
    1, & \abs{t_2}>\abs{t_1}.
    \end{cases}
\end{equation}
The transition occurs at $\abs{t_1}=\abs{t_2}$, where the gap closes at
$k=\pi$ for $t_1=t_2$.

\begin{remark}[Unit-cell convention and the sign of winding]
Changing the unit-cell convention multiplies $q(k)$ by a phase
$\ee^{\ii n k}$.  This shifts the integer assigned to a particular Bloch
basis.  The robust physical statement for the SSH chain with a fixed boundary
termination is whether the dimerization leaves unpaired edge sites.  In the
convention above, $\abs{t_2}>\abs{t_1}$ is the topological phase for the usual
termination with weak intracell bonds at the boundary.
\end{remark}

\subsection{Edge zero modes}
\label{subsec:topo-ssh-edge-zero-modes}

For an open chain in the topological regime $\abs{t_2}>\abs{t_1}$, the left-edge
zero mode has support only on one sublattice.  A zero-energy ansatz of the form
\begin{equation}
\label{eq:topo-ssh-left-edge-ansatz}
    \ket{\psi_L}=\sum_{j\geq 1}\alpha_j a_j^\dagger\ket{0}
\end{equation}
obeys the recursion
\begin{equation}
\label{eq:topo-ssh-edge-recursion}
    t_1\alpha_j+t_2\alpha_{j+1}=0,
    \qquad
    \alpha_{j+1}=-\frac{t_1}{t_2}\alpha_j .
\end{equation}
Therefore
\begin{equation}
\label{eq:topo-ssh-edge-wavefunction}
    \alpha_j=\left(-\frac{t_1}{t_2}\right)^{j-1}\alpha_1,
\end{equation}
which is normalizable precisely when $\abs{t_1/t_2}<1$.  The bulk winding and
the boundary zero mode therefore encode the same phase distinction.

\section{BdG topology: Kitaev chain and Majorana zero modes}
\label{sec:topo-bdg-kitaev}

The Kitaev chain is also a two-band Hamiltonian, but the two components of the
spinor are not sublattices.  They are particle and hole components of a Nambu
spinor.  In the convention of the BdG chapter,
\begin{equation}
\label{eq:topo-kitaev-bdg-block}
    \mathcal{H}_{\rm K}(k)
    =
    \begin{pmatrix}
        \xi(k) & \Delta_\phi(k)\\
        \Delta_\phi(k)^* & -\xi(k)
    \end{pmatrix},
    \qquad
    \xi(k)=-\mu-2t\cos k,
\end{equation}
with
\begin{equation}
\label{eq:topo-kitaev-gap-function}
    \Delta_\phi(k)
    =
    2\ii\abs{\Delta}\ee^{\ii(\theta_\Delta-2\phi)}\sin k .
\end{equation}
The uniform phase $\theta_\Delta$ and the Fourier/Nambu gauge phase $\phi$ enter
only through the combination $\theta_\Delta-2\phi$.  For a single uniform chain,
one may choose
\begin{equation}
    \phi=\frac{\theta_\Delta}{2}
\end{equation}
so that
\begin{equation}
\label{eq:topo-kitaev-real-gap-gauge}
    \Delta_\phi(k)=2\ii\abs{\Delta}\sin k.
\end{equation}
In this real pairing gauge,
\begin{equation}
\label{eq:topo-kitaev-d-vector}
    \mathcal{H}_{\rm K}(k)
    =
    d_y(k)\tau_y+d_z(k)\tau_z,
    \qquad
    d_y(k)=-2\abs{\Delta}\sin k,
    \qquad
    d_z(k)=\xi(k).
\end{equation}
Here the Pauli matrices $\tau_a$ act in Nambu space.  Since the vector lies in a
plane, the real Kitaev chain has a winding-number description analogous to a
chiral two-band model.

A convenient complex function is
\begin{equation}
\label{eq:topo-kitaev-q-function}
    q_{\rm K}(k)=\xi(k)+2\ii\abs{\Delta}\sin k
    =-\mu-2t\cos k+2\ii\abs{\Delta}\sin k.
\end{equation}
Its winding number is
\begin{equation}
\label{eq:topo-kitaev-bdi-winding}
    \nu_{\rm BDI}
    =
    \frac{1}{2\pi\ii}\int_{\rm BZ}\dd k\,
    \partial_k\log q_{\rm K}(k).
\end{equation}
Geometrically, $q_{\rm K}(k)$ traces an ellipse centered at $-\mu$ on the real
axis.  The origin lies inside this ellipse if and only if
\begin{equation}
\label{eq:topo-kitaev-topological-condition}
    \abs{\mu}<2\abs{t}.
\end{equation}
Thus the Kitaev chain is topological for $\abs{\mu}<2\abs{t}$ and trivial for
$\abs{\mu}>2\abs{t}$, assuming $\Delta\neq0$.

\subsection{Class D Pfaffian invariant}
\label{subsec:topo-class-d-pfaffian}

The real gauge above makes an additional chiral symmetry manifest and places the
nearest-neighbor Kitaev chain in class BDI.  More generally, a superconducting
BdG Hamiltonian always has particle-hole symmetry.  In the Nambu convention
used earlier,
\begin{equation}
\label{eq:topo-phs-condition}
    \tau_x\mathcal{H}_{\rm BdG}(k)^*\tau_x=-\mathcal{H}_{\rm BdG}(-k).
\end{equation}
Without a real gauge or time-reversal-like symmetry, the one-dimensional
superconductor is in class D and has a $\ZZ_2$ invariant rather than an integer
winding number.

For a two-band Kitaev chain the class-D invariant reduces to the sign of the
normal-state term at the particle-hole-invariant momenta $k=0,\pi$, because the
odd-parity pairing vanishes there:
\begin{equation}
\label{eq:topo-kitaev-class-d-invariant}
    \mathcal{M}
    =
    \sgn\bigl[\xi(0)\xi(\pi)\bigr]
    =
    \sgn\bigl[(-\mu-2t)(-\mu+2t)\bigr]
    =
    \sgn(\mu^2-4t^2).
\end{equation}
The nontrivial phase is
\begin{equation}
\label{eq:topo-class-d-nontrivial-condition}
    \mathcal{M}=-1
    \quad\Longleftrightarrow\quad
    \abs{\mu}<2\abs{t}.
\end{equation}
This is the $\ZZ_2$ remnant of the integer BDI winding.  Uniform complex phases
of $\Delta$ do not change \cref{eq:topo-kitaev-class-d-invariant}; they can be
removed by a global Nambu gauge choice.  Relative superconducting phases,
spatially varying phases, or fluxes that cannot be removed globally can reduce
the manifest symmetry to class D, but the $\ZZ_2$ invariant remains well-defined
as long as particle-hole symmetry and the bulk gap persist.

\subsection{Exactly solvable Majorana limit}
\label{subsec:topo-majorana-edge-limit}

The boundary meaning of \cref{eq:topo-kitaev-topological-condition} is clearest
at the exactly solvable point
\begin{equation}
\label{eq:topo-kitaev-sweet-spot}
    \mu=0,
    \qquad
    t=\abs{\Delta}>0,
\end{equation}
with open boundary conditions.  Using Majorana operators
\begin{equation}
\label{eq:topo-majorana-definitions}
    a_{2j-1}=c_j+c_j^\dagger,
    \qquad
    a_{2j}=\ii(c_j-c_j^\dagger),
\end{equation}
the Hamiltonian becomes, up to an overall sign convention fixed by the phase of
$\Delta$,
\begin{equation}
\label{eq:topo-kitaev-majorana-sweet-spot}
    H_{\rm K}
    =
    -\ii t\sum_{j=1}^{L-1}a_{2j}a_{2j+1}.
\end{equation}
All bulk Majoranas are paired on bonds, while $a_1$ and $a_{2L}$ are absent from
the Hamiltonian.  These two unpaired Majorana operators form a nonlocal boundary
fermion
\begin{equation}
\label{eq:topo-edge-fermion}
    f_{\rm edge}=\frac{a_1+\ii a_{2L}}{2}.
\end{equation}
The two occupations of $f_{\rm edge}$ are degenerate in the thermodynamic limit.
Away from the exactly solvable point but within the same gapped topological
phase, the Majorana zero modes remain exponentially localized at the edges.

\begin{principle}[Bulk-boundary correspondence in one-dimensional two-band models]
A nontrivial bulk winding or class-D $\ZZ_2$ invariant predicts protected
boundary zero modes, provided the boundary termination respects the protecting
symmetry and the bulk gap remains open.
\end{principle}

For the SSH chain the boundary zero modes are ordinary complex fermion states
localized on sublattices.  For the Kitaev chain they are Majorana zero modes in
Nambu space.  The algebraic similarity of the two-band Hamiltonians should not
obscure this physical difference.

\section{Relation to the TFIM and the Majorana mass sign}
\label{sec:topo-relation-to-tfim}

The duality chapter showed that the transverse-field Ising chain separates two
phases exchanged by Kramers--Wannier duality.  The Jordan--Wigner and BdG
chapters showed that the same model is equivalent to a quadratic Majorana chain.
Topology supplies a third language for the same transition.

Near a gap-closing momentum, a one-dimensional BdG Hamiltonian takes the
continuum Dirac/Majorana form
\begin{equation}
\label{eq:topo-continuum-majorana-dirac}
    \mathcal{H}(q)=v q\,\tau_y+m\,\tau_z+\text{higher-order terms},
\end{equation}
where $q$ measures momentum relative to the gap-closing point.  A topological
transition occurs when the mass $m$ changes sign.  For the Kitaev chain,
\begin{equation}
\label{eq:topo-kitaev-masses}
    m_0=\xi(0)=-\mu-2t,
    \qquad
    m_\pi=\xi(\pi)=-\mu+2t.
\end{equation}
The class-D invariant is simply the relative sign of these two masses:
\begin{equation}
\label{eq:topo-mass-sign-invariant}
    \mathcal{M}=\sgn(m_0m_\pi).
\end{equation}
The nontrivial phase has opposite signs at $k=0$ and $k=\pi$.

In the TFIM image, the topological Majorana phase corresponds to the ordered
Ising phase for open chains.  The edge Majoranas encode the same twofold
structure that appears as symmetry breaking in the spin variables.  Thus the
following descriptions are equivalent ways of describing the same physics:
\begin{equation}
\label{eq:topo-tfim-equivalent-descriptions}
\begin{array}{ccl}
\text{spin language} &:& \text{ordered versus disordered Ising phase},\\[2pt]
\text{duality language} &:& \text{order variable versus disorder variable},\\[2pt]
\text{fermion language} &:& \text{trivial versus topological Majorana chain},\\[2pt]
\text{continuum language} &:& \text{sign of a Majorana mass}.
\end{array}
\end{equation}

\section{Symmetry classes relevant to the mother models}
\label{sec:topo-symmetry-classes}

The full tenfold classification is beyond the scope of these notes, but the
following restricted table is sufficient for the solvable models treated here.

\begin{center}
\small
\begin{tabular}{L{0.15\textwidth}L{0.22\textwidth}L{0.21\textwidth}L{0.28\textwidth}}
\toprule
Class & Symmetry content & 1D invariant & Representative model in these notes \\
\midrule
AIII & Chiral symmetry only & $\ZZ$ winding & SSH chain with chiral sublattice symmetry \\
BDI & Particle-hole plus real time-reversal-like symmetry, hence chiral symmetry & $\ZZ$ winding & Real Kitaev chain / JW image of XY-type chains in a real gauge \\
D & Particle-hole symmetry only & $\ZZ_2$ Pfaffian invariant & Kitaev chain with only BdG particle-hole symmetry manifest \\
A & No protecting symmetry & No stable 1D two-band invariant; $\ZZ$ Chern number in 2D & Two-dimensional Chern-insulator block \\
\bottomrule
\end{tabular}
\end{center}

The table should be read operationally.  The same matrix expression may belong
to different classes depending on what physical structure the Pauli matrices
represent and which perturbations are allowed.  For example, adding a sublattice
staggering term $m\sigma_z$ to the SSH chain breaks chiral symmetry and destroys
the integer winding classification.  Adding a perturbation to a BdG Hamiltonian
must still respect particle-hole redundancy, so the allowed deformations are
different.

\section{Gauge choices, phases, and what is physically invariant}
\label{sec:topo-gauge-phases}

The Fourier and BdG chapters emphasized that a constant phase in the Fourier
transform,
\begin{equation}
\label{eq:topo-fourier-gauge-reminder}
    c_k=\frac{\ee^{-\ii\phi}}{\sqrt{L}}\sum_j\ee^{-\ii k j}c_j,
\end{equation}
is a gauge convention.  In BdG language it rotates the Nambu spinor and changes
the phase of the gap function as
\begin{equation}
\label{eq:topo-gap-gauge-rotation}
    \Delta_\phi(k)=\ee^{-2\ii\phi}\Delta_0(k).
\end{equation}
Consequently the components $(d_x,d_y)$ of a BdG vector may rotate in the
transverse plane under a gauge change.  This rotation does not change the energy
spectrum, the existence of a gap, or the class-D invariant
\cref{eq:topo-kitaev-class-d-invariant}.

For an ordinary chiral band Hamiltonian such as the SSH chain, a change of Bloch
basis can also multiply $q(k)$ by a phase.  If the phase is periodic and smooth,
it does not change the physically meaningful distinction between phases with and
without protected boundary states for a fixed termination.  If the phase has the
form $\ee^{\ii n k}$, it reflects a shift of unit-cell convention and changes the
integer assigned to the winding by a convention-dependent offset.  Boundary
physics is recovered only after the unit cell and termination are specified.

The correct hierarchy is therefore
\begin{equation}
\label{eq:topo-invariant-hierarchy}
\begin{array}{ccl}
\text{basis-dependent} &:& d_x(k),d_y(k),d_z(k)\ \text{as components},\\[2pt]
\text{gauge-covariant} &:& q(k)\ \text{and its phase},\\[2pt]
\text{physically invariant} &: &
\begin{gathered}
\text{gap closings, Berry phases modulo }2\pi,\\
\text{winding differences, }\ZZ_2\text{ parity, and edge modes}
\end{gathered}.
\end{array}
\end{equation}
This distinction is especially important for the Kitaev chain with a complex
superconducting gap: a uniform complex phase is removable by a global gauge
choice, while a relative phase across a junction or a spatial phase twist can
represent physical superconducting phase bias or flux.

\section{Practical diagnostic formulas}
\label{sec:topo-practical-diagnostics}

For later use, we collect the most useful diagnostic formulas.

\subsection*{Chiral two-band chain}
For
\begin{equation}
    H(k)=d_x(k)\sigma_x+d_y(k)\sigma_y,
    \qquad
    q(k)=d_x(k)+\ii d_y(k),
\end{equation}
the integer invariant is
\begin{equation}
\label{eq:topo-summary-chiral-winding}
    \nu=\frac{1}{2\pi\ii}\oint_{\rm BZ}\dd k\,\partial_k\log q(k).
\end{equation}
A gap closing occurs when $q(k)=0$.

\subsection*{SSH chain}
For
\begin{equation}
    q_{\rm SSH}(k)=t_1+t_2\ee^{\ii k},
\end{equation}
the usual topological regime for the standard open-chain termination is
\begin{equation}
\label{eq:topo-summary-ssh}
    \abs{t_2}>\abs{t_1}.
\end{equation}

\subsection*{Kitaev chain}
For
\begin{equation}
    \xi(k)=-\mu-2t\cos k,
    \qquad
    \Delta(k)=2\ii\abs{\Delta}\sin k
\end{equation}
in a real pairing gauge, the BDI winding may be computed from
\begin{equation}
\label{eq:topo-summary-kitaev-winding}
    q_{\rm K}(k)=\xi(k)+2\ii\abs{\Delta}\sin k.
\end{equation}
The class-D invariant is
\begin{equation}
\label{eq:topo-summary-kitaev-z2}
    \mathcal{M}=\sgn\bigl[\xi(0)\xi(\pi)\bigr]
    =\sgn(\mu^2-4t^2).
\end{equation}
The topological Majorana phase is
\begin{equation}
\label{eq:topo-summary-kitaev-topological-phase}
    \mathcal{M}=-1
    \quad\Longleftrightarrow\quad
    \abs{\mu}<2\abs{t}.
\end{equation}

\subsection*{Two-dimensional two-band block}
For a two-dimensional occupied lower band,
\begin{equation}
\label{eq:topo-summary-chern}
    C
    =
    \frac{1}{4\pi}\int_{\rm BZ}
    \widehat{\bm d}\cdot
    \left(
        \partial_{k_x}\widehat{\bm d}
        \times
        \partial_{k_y}\widehat{\bm d}
    \right)\dd^2k.
\end{equation}

\section{Summary}
\label{sec:topo-summary}

The essential lesson of this chapter is that a solvable two-band Hamiltonian
contains more information than its dispersion.  The eigenvectors define a map
from momentum space to a target space, and the topology of this map constrains
possible phases and boundary states.

The main results are:
\begin{enumerate}[label=(\roman*)]
    \item A gapped two-band Hamiltonian defines a normalized vector
    $\widehat{\bm d}(k)$ and can change topology only when $\bm d(k)=0$.
    \item In one dimension, a generic two-band Hamiltonian is topologically
    trivial unless a symmetry restricts the target space.
    \item Chiral two-band chains have an integer winding number
    $\nu=(2\pi\ii)^{-1}\int\partial_k\log q(k)$.
    \item The SSH chain realizes this winding in sublattice space.
    \item The Kitaev chain realizes an analogous structure in Nambu space; in
    class D the robust invariant is the Pfaffian sign
    $\mathcal{M}=\sgn[\xi(0)\xi(\pi)]$.
    \item Majorana edge modes are the boundary manifestation of the nontrivial
    Kitaev-chain invariant.
    \item The TFIM transition, the Kitaev-chain topological transition, and the
    sign change of a continuum Majorana mass are different languages for the
    same solvable critical structure.
\end{enumerate}


\chapter[Bethe Ansatz and Yang--Baxter Integrability]{Preview of Bethe Ansatz and Yang--Baxter Integrability}
\label{chap:bethe-yang-baxter-preview}

The preceding chapters developed two main solution mechanisms.  First, transfer
matrices convert a classical statistical model into a spectral problem.  Second,
Jordan--Wigner and Bogoliubov transformations solve a large class of spin chains
because their fermionic images are quadratic.  The next integrability mechanism
is different.  Bethe-ansatz integrability solves genuinely interacting models,
not by reducing them to independent one-particle modes, but by exploiting the
factorization of many-body scattering and the existence of an infinite commuting
family of transfer matrices.

This chapter is a conceptual and notational preview.  Its purpose is to explain
what the Bethe ansatz solves, why the Yang--Baxter equation appears, and how the
row-to-row transfer matrices of vertex models become quantum Hamiltonians such as
the XXZ chain.  The complete coordinate and algebraic Bethe-ansatz derivations
are postponed to later chapters.

\section{From free particles to factorized scattering}
\label{sec:ba-factorized-scattering}

A free-fermion Hamiltonian is solvable because it can be diagonalized into
independent quasiparticles,
\begin{equation}
    H=E_0+\sum_k E_k\gamma_k^\dagger\gamma_k .
\end{equation}
The many-body spectrum is then obtained by filling one-particle levels.  Bethe
ansatz uses a weaker but still powerful property: particles interact, but the
scattering is elastic and factorized.

\begin{definition}[Factorized scattering]
Consider a one-dimensional many-body system with stable asymptotic quasiparticles
labelled by momenta or rapidities.  The scattering is called \emph{factorized} if
an $M$-body collision can be decomposed into an ordered product of two-body
scattering processes, with no particle production and no diffraction.
\end{definition}

In one spatial dimension the ordering of particles is meaningful.  For
coordinates
\begin{equation}
    x_1<x_2<\cdots<x_M,
\end{equation}
a coordinate Bethe wave function takes the form
\begin{equation}
\label{eq:coordinate-ba-wavefunction-preview}
    \Psi(x_1,\ldots,x_M)
    =
    \sum_{P\in S_M} A(P)
    \exp\left(\ii\sum_{a=1}^{M}k_{P_a}x_a\right),
\end{equation}
where $P$ permutes the momenta among the ordered particles.  Interactions do not
change the set of momenta.  They only relate amplitudes associated with different
permutations.  For neighboring exchanges one writes schematically
\begin{equation}
\label{eq:two-body-scattering-amplitude-preview}
    A(P\circ s_a)
    =S(k_{P_a},k_{P_{a+1}})A(P),
\end{equation}
where $s_a$ exchanges the $a$th and $(a+1)$st momenta.  For a scalar scattering
phase, factorization is almost automatic because complex numbers commute.  For
particles with internal spin, color, or auxiliary indices, $S$ becomes a matrix,
and consistency is highly nontrivial.

The basic physical meaning is the following.  In a generic interacting model, a
many-particle collision can redistribute momenta in ways that cannot be described
as a sequence of independent two-body events.  In an integrable one-dimensional
model, the infinite family of conservation laws is strong enough to forbid such
diffraction.  The only allowed effect of scattering is to permute the rapidities
and multiply the wave function by two-body scattering matrices.

\begin{principle}[Bethe-ansatz solvability]
A one-dimensional interacting model is Bethe-ansatz solvable when its eigenstates
can be assembled from plane waves whose amplitudes are related by consistent
factorized two-body scattering.
\end{principle}

\section{A two-magnon preview}
\label{sec:two-magnon-preview}

The simplest way to see the Bethe idea is to consider a spin chain above a
ferromagnetic reference state.  Let
\begin{equation}
    \ket{0}=\ket{\uparrow\uparrow\cdots\uparrow}
\end{equation}
and let a down spin be interpreted as a magnon.  A one-magnon state is
\begin{equation}
    \ket{k}=\sum_{x=1}^{L}\ee^{\ii kx}\sigma_x^-\ket{0} .
\end{equation}
For a translation-invariant chain, it is an eigenstate in the one-magnon sector.
The nontrivial part begins in the two-magnon sector, where the wave function has
one expression for $x_1<x_2$:
\begin{equation}
\label{eq:two-magnon-wavefunction-preview}
    \Psi(x_1,x_2)
    =A_{12}\ee^{\ii(k_1x_1+k_2x_2)}
    +A_{21}\ee^{\ii(k_2x_1+k_1x_2)} .
\end{equation}
The interaction only matters when the two magnons approach each other.  The
Schrodinger equation at contact fixes the ratio
\begin{equation}
\label{eq:two-magnon-S-preview}
    A_{21}=S(k_1,k_2)A_{12} .
\end{equation}
For a periodic chain, moving one magnon around the ring makes it scatter once
with every other magnon.  This gives the schematic Bethe equations
\begin{equation}
\label{eq:schematic-bethe-equations-preview}
    \ee^{\ii k_j L}
    =
    \prod_{\ell\neq j}S(k_j,k_\ell),
    \qquad j=1,\ldots,M .
\end{equation}
Thus periodic boundary conditions quantize the momenta in an interaction-dependent
way.  Free particles would have $S=\pm1$ and hence ordinary integer or
half-integer momenta.  Interacting Bethe-ansatz particles acquire phase shifts
from all other particles.

\begin{example}[The XXZ two-body scattering phase]
For one common Pauli-matrix convention, the XXZ chain may be written as
\begin{equation}
\label{eq:xxz-preview-pauli-convention}
    H_{\rm XXZ}
    =
    J\sum_{j=1}^{L}
    \left(
    \sigma_j^x\sigma_{j+1}^x
    +\sigma_j^y\sigma_{j+1}^y
    +\Delta(\sigma_j^z\sigma_{j+1}^z-1)
    \right).
\end{equation}
In the two-magnon coordinate Bethe ansatz, the contact equation gives a scattering
phase of the form
\begin{equation}
\label{eq:xxz-momentum-S-preview}
    S(k_1,k_2)
    =-
    \frac{1+\ee^{\ii(k_1+k_2)}-2\Delta\ee^{\ii k_1}}
         {1+\ee^{\ii(k_1+k_2)}-2\Delta\ee^{\ii k_2}} .
\end{equation}
The precise overall sign and energy convention depend on the normalization of
\cref{eq:xxz-preview-pauli-convention}, but the structural message is universal:
the interaction is encoded in a two-body phase shift, and the many-body momenta
are fixed by \cref{eq:schematic-bethe-equations-preview}.
\end{example}

It is often more convenient to replace momenta by rapidities.  In the XXX limit,
the Bethe equations are commonly written as
\begin{equation}
\label{eq:xxx-bethe-equations-preview}
    \left(
    \frac{\lambda_j+\ii/2}{\lambda_j-\ii/2}
    \right)^L
    =
    \prod_{\ell\neq j}
    \frac{\lambda_j-\lambda_\ell+\ii}
         {\lambda_j-\lambda_\ell-\ii} .
\end{equation}
For the gapless XXZ chain with $\Delta=\cos\gamma$, a standard rapidity form is
\begin{equation}
\label{eq:xxz-bethe-equations-preview}
    \left[
    \frac{\sinh(\lambda_j+\ii\gamma/2)}
         {\sinh(\lambda_j-\ii\gamma/2)}
    \right]^L
    =
    \prod_{\ell\neq j}
    \frac{\sinh(\lambda_j-\lambda_\ell+\ii\gamma)}
         {\sinh(\lambda_j-\lambda_\ell-\ii\gamma)} .
\end{equation}
The left-hand side is free propagation around the ring.  The right-hand side is
the accumulated product of two-body scattering phases.

\section{Why a three-body consistency condition is needed}
\label{sec:three-body-consistency-preview}

For three particles there are multiple ways to exchange the same initial and
final orderings.  For instance, the permutation reversing $(1,2,3)$ to $(3,2,1)$
can be decomposed as
\begin{equation}
    s_1s_2s_1=s_2s_1s_2 .
\end{equation}
If the scattering amplitudes are scalar phases, this identity causes no serious
constraint.  If the two-body scatterings are matrices acting on internal degrees
of freedom, the two possible sequences must give the same operator.  This is the
origin of the Yang--Baxter equation.

Let $R_{ab}(u)$ denote a two-body scattering or vertex operator acting nontrivially
on spaces $a$ and $b$, with spectral parameter difference $u$.  The Yang--Baxter
equation is
\begin{equation}
\label{eq:yang-baxter-equation-preview}
    R_{12}(u-v)R_{13}(u-w)R_{23}(v-w)
    =
    R_{23}(v-w)R_{13}(u-w)R_{12}(u-v) .
\end{equation}
It is an equality of operators on a triple tensor product.  Diagrammatically it
states that two different orders of resolving a three-line crossing are
equivalent.  Algebraically it is precisely the consistency condition required for
factorized scattering with internal degrees of freedom.

\begin{remark}[Yang--Baxter versus ordinary commutativity]
The Yang--Baxter equation is not the statement that two matrices commute.  It is
a cubic relation among three two-body operators acting on three tensor factors.
Its consequence, after constructing a monodromy matrix and tracing over an
auxiliary space, is the commutativity of transfer matrices.
\end{remark}

\section{R-matrix as vertex weight and as scattering matrix}
\label{sec:r-matrix-two-interpretations}

The same $R$-matrix has two complementary meanings.

First, in a two-dimensional classical vertex model, $R$ assigns a local Boltzmann
weight to a vertex.  The two tensor indices may be viewed as spin or arrow states
living on the lines entering and leaving the vertex.  The partition function is
built by multiplying local $R$-matrix weights and summing over internal indices.

Second, in a one-dimensional quantum chain, $R$ is the algebraic object that
controls two-body scattering and the commuting transfer matrices.  The spectral
parameter $u$ plays the role of an auxiliary rapidity.  Taking a logarithmic
derivative of the transfer matrix produces a quantum Hamiltonian.

For the six-vertex model, a standard trigonometric or hyperbolic $R$-matrix is
\begin{equation}
\label{eq:six-vertex-R-preview}
    R(u)=
    \begin{pmatrix}
    a(u)&0&0&0\\
    0&b(u)&c(u)&0\\
    0&c(u)&b(u)&0\\
    0&0&0&a(u)
    \end{pmatrix},
\end{equation}
where, for the hyperbolic parametrization,
\begin{equation}
\label{eq:six-vertex-weights-hyperbolic-preview}
    a(u)=\sinh(u+\eta),
    \qquad
    b(u)=\sinh u,
    \qquad
    c(u)=\sinh\eta .
\end{equation}
The nonzero entries encode the ice rule: the number of incoming arrows equals the
number of outgoing arrows at each vertex.  In the associated XXZ chain, the
anisotropy parameter is controlled by
\begin{equation}
\label{eq:xxz-anisotropy-preview}
    \Delta=\cosh\eta
\end{equation}
for this hyperbolic convention.  In the gapless convention one uses sines instead
of hyperbolic sines and obtains $\Delta=\cos\gamma$.

\section{Monodromy matrix and commuting transfer matrices}
\label{sec:monodromy-transfer-preview}

To build a length-$L$ quantum chain, introduce an auxiliary space $a$ and quantum
spaces $1,\ldots,L$.  The monodromy matrix is the ordered product
\begin{equation}
\label{eq:monodromy-matrix-preview}
    M_a(u)
    =
    R_{aL}(u-\xi_L)R_{a,L-1}(u-\xi_{L-1})\cdots R_{a1}(u-\xi_1),
\end{equation}
where $\xi_j$ are inhomogeneity parameters.  The homogeneous chain is obtained by
setting all $\xi_j=0$.  The transfer matrix is the auxiliary trace
\begin{equation}
\label{eq:transfer-matrix-from-monodromy-preview}
    \tau(u)=\Tr_a M_a(u).
\end{equation}

The Yang--Baxter equation implies the RTT relation
\begin{equation}
\label{eq:rtt-relation-preview}
    R_{ab}(u-v)M_a(u)M_b(v)
    =
    M_b(v)M_a(u)R_{ab}(u-v) .
\end{equation}
Taking the trace over the two auxiliary spaces $a$ and $b$ gives
\begin{equation}
\label{eq:commuting-transfer-matrices-preview}
    [\tau(u),\tau(v)]=0,
    \qquad
    \text{for all }u,v .
\end{equation}
This is the central algebraic output of Yang--Baxter integrability.

\begin{theorem}[Commuting transfer matrices]
If the $R$-matrix satisfies the Yang--Baxter equation and the monodromy matrix is
constructed as in \cref{eq:monodromy-matrix-preview}, then the transfer matrices
\cref{eq:transfer-matrix-from-monodromy-preview} commute for all values of the
spectral parameter:
\begin{equation}
    [\tau(u),\tau(v)]=0 .
\end{equation}
\end{theorem}

\begin{proof}
The Yang--Baxter equation implies the RTT relation
\cref{eq:rtt-relation-preview}.  Taking $\Tr_a\Tr_b$ of both sides and using
cyclicity of the trace in the auxiliary spaces gives
\begin{equation}
    \Tr_a\Tr_b\left[R_{ab}(u-v)M_a(u)M_b(v)\right]
    =
    \Tr_a\Tr_b\left[M_b(v)M_a(u)R_{ab}(u-v)\right].
\end{equation}
Cyclicity moves $R_{ab}(u-v)$ through the auxiliary trace, leaving
\begin{equation}
    \Tr_a M_a(u)\,\Tr_b M_b(v)
    =
    \Tr_b M_b(v)\,\Tr_a M_a(u).
\end{equation}
This is exactly $\tau(u)\tau(v)=\tau(v)\tau(u)$.
\end{proof}

The physical meaning of \cref{eq:commuting-transfer-matrices-preview} is that
$\tau(u)$ is a generating function for conserved quantities.  Expanding the
logarithm around a regular point gives
\begin{equation}
\label{eq:charges-from-log-transfer-preview}
    \log\tau(u)
    =
    \sum_{n=0}^{\infty} (u-u_0)^n Q_n,
    \qquad
    [Q_m,Q_n]=0 .
\end{equation}
One of these charges is the Hamiltonian; the rest are higher conserved charges.
This is the transfer-matrix form of quantum integrability.

\section{Hamiltonian limit}
\label{sec:hamiltonian-limit-preview}

A local quantum Hamiltonian is obtained from a regular transfer matrix.  The key
regularity property is
\begin{equation}
\label{eq:R-regularity-preview}
    R(0)=\rho P,
\end{equation}
where $P$ is the permutation operator on two neighboring spin spaces and $\rho$ is
a scalar normalization.  Then the logarithmic derivative of the transfer matrix
at $u=0$ produces a sum of local nearest-neighbor terms:
\begin{equation}
\label{eq:hamiltonian-log-transfer-preview}
    H
    =
    C\left.\frac{\dd}{\dd u}\log\tau(u)\right|_{u=0}
    +C_0
    =
    \sum_{j=1}^{L}h_{j,j+1}+C_0 .
\end{equation}
For the six-vertex $R$-matrix, this Hamiltonian is the spin-$1/2$ XXZ chain,
up to normalization, additive constants, and boundary conventions:
\begin{equation}
\label{eq:xxz-from-six-vertex-preview}
    H_{\rm XXZ}
    =
    J\sum_{j=1}^{L}
    \left(
    S_j^xS_{j+1}^x
    +S_j^yS_{j+1}^y
    +\Delta S_j^zS_{j+1}^z
    \right)
    +\text{constant} .
\end{equation}
This is the basic bridge between the six-vertex model and the XXZ quantum chain.
It is the interacting analogue of the quantum--classical correspondence discussed
in the transfer-matrix chapter.

\section{Why commuting transfer matrices imply many conserved quantities}
\label{sec:commuting-transfer-conserved-quantities}

Suppose $\tau(u)$ is analytic near $u_0$ and $[\tau(u),\tau(v)]=0$ for all $u,v$.
Then every coefficient in the expansion of $\log\tau(u)$ commutes with every other
coefficient.  More explicitly, define
\begin{equation}
    Q_n
    =
    \left.
    \frac{\dd^n}{\dd u^n}\log\tau(u)
    \right|_{u=u_0} .
\end{equation}
Since the whole family $\tau(u)$ commutes, one obtains
\begin{equation}
    [Q_m,Q_n]=0,
    \qquad
    [Q_n,H]=0 .
\end{equation}
The first nontrivial local charge is usually the Hamiltonian.  Higher charges
have increasing spatial range.  In a local integrable chain one may schematically
write
\begin{equation}
    Q_2=H,
    \qquad
    Q_3=\sum_j q_{j,j+1,j+2},
    \qquad
    Q_4=\sum_j q_{j,j+1,j+2,j+3},
    \quad \ldots
\end{equation}
The existence of these charges is the algebraic reason why scattering is
elastic and factorized.

\begin{remark}[Integrability is stronger than one conservation law]
Translation invariance gives momentum conservation, and time-translation
invariance gives energy conservation.  These two conservation laws are not enough
to solve an interacting many-body problem.  Bethe-ansatz integrability requires an
infinite tower of mutually commuting charges, generated by a commuting family of
transfer matrices.
\end{remark}

\section{Coordinate Bethe ansatz versus algebraic Bethe ansatz}
\label{sec:coordinate-vs-algebraic-ba-preview}

There are two complementary ways to use this structure.

\subsection{Coordinate Bethe ansatz}

The coordinate Bethe ansatz works directly with wave functions such as
\cref{eq:coordinate-ba-wavefunction-preview}.  Its main steps are:
\begin{enumerate}[label=(\roman*)]
    \item choose a simple reference state;
    \item construct $M$-particle wave functions in ordered coordinate sectors;
    \item impose contact conditions when particles become adjacent;
    \item extract two-body scattering phases;
    \item impose periodic boundary conditions to obtain Bethe equations.
\end{enumerate}
This method is concrete and physically transparent.  It is especially useful for
the XXX and XXZ chains and will be developed in the coordinate Bethe-ansatz
chapter.

\subsection{Algebraic Bethe ansatz}

The algebraic Bethe ansatz starts from the monodromy matrix.  For a spin-$1/2$
chain the monodromy matrix is a $2\times2$ matrix in the auxiliary space,
\begin{equation}
\label{eq:monodromy-ABCD-preview}
    M_a(u)
    =
    \begin{pmatrix}
    A(u)&B(u)\\
    C(u)&D(u)
    \end{pmatrix}_a,
\end{equation}
where $A,B,C,D$ are operators acting on the quantum spin chain.  The transfer
matrix is
\begin{equation}
    \tau(u)=A(u)+D(u).
\end{equation}
One chooses a reference state $\ket{0}$ such that
\begin{equation}
    C(u)\ket{0}=0,
    \qquad
    A(u)\ket{0}=a(u)\ket{0},
    \qquad
    D(u)\ket{0}=d(u)\ket{0} .
\end{equation}
The operator $B(u)$ creates Bethe excitations:
\begin{equation}
\label{eq:algebraic-bethe-state-preview}
    \ket{\lambda_1,\ldots,\lambda_M}
    =B(\lambda_1)\cdots B(\lambda_M)\ket{0} .
\end{equation}
The RTT relation determines how $A(u)$ and $D(u)$ commute through the $B$'s.  The
unwanted terms vanish exactly when the rapidities satisfy the Bethe equations.
Then \cref{eq:algebraic-bethe-state-preview} is an eigenstate of $\tau(u)$ and
therefore also of the Hamiltonian extracted from $\log\tau(u)$.

\begin{center}
\begin{tabularx}{\textwidth}{L{0.28\textwidth}Y Y}
\toprule
Approach & Main object & Best use \\
\midrule
Coordinate Bethe ansatz & Piecewise plane-wave many-body wave function & Direct derivation of scattering phases and Bethe equations \\
Algebraic Bethe ansatz & Monodromy matrix and RTT algebra & Systematic construction of eigenstates from $B(u)$ operators \\
Transfer-matrix method & Commuting family $\tau(u)$ & Conserved charges and Hamiltonian limits of vertex models \\
\bottomrule
\end{tabularx}
\end{center}

\section{Relation to free-fermion solvability}
\label{sec:bethe-vs-free-fermion-preview}

It is useful to compare the free-fermion and Bethe-ansatz mechanisms.

\begin{center}
\begin{tabularx}{\textwidth}{L{0.26\textwidth}Y Y}
\toprule
Feature & Free-fermion/BdG solvability & Bethe/Yang--Baxter integrability \\
\midrule
Elementary objects & Independent fermionic modes or quasiparticles & Interacting Bethe quasiparticles with rapidities \\
Many-body construction & Fill independent modes & Solve coupled Bethe equations \\
Scattering & Trivial exchange signs or BdG mixing & Nontrivial two-body phase or matrix scattering \\
Conserved quantities & Quadratic mode occupations & Infinite tower from $\log\tau(u)$ \\
Representative models & XY, TFIM, Kitaev, SSH & XXX, XXZ, XYZ, six-vertex, Hubbard \\
Core algebra & Canonical anticommutation and Bogoliubov rotation & Yang--Baxter equation and RTT relation \\
\bottomrule
\end{tabularx}
\end{center}

The XX chain sits at the overlap between these viewpoints.  It is the
$\Delta=0$ limit of the XXZ chain and is also free after the Jordan--Wigner
transformation.  From the Bethe-ansatz perspective, its scattering phase reduces
to the fermionic exchange sign.  From the free-fermion perspective, it is simply a
number-conserving tight-binding chain.  This overlap is a useful consistency
check between the two solution languages.

\section{Physical interpretation of rapidities}
\label{sec:rapidities-preview}

In relativistic field theory a rapidity directly parametrizes energy and momentum.
In lattice Bethe ansatz, the word is used more broadly: a rapidity is a spectral
parameter that uniformizes the dispersion and scattering.  The physical momentum
is a function of rapidity,
\begin{equation}
    p=p(\lambda),
\end{equation}
and the two-body scattering depends only on rapidity differences in many
homogeneous integrable models,
\begin{equation}
    S(\lambda_i,\lambda_j)=S(\lambda_i-\lambda_j).
\end{equation}
This difference property is one reason why the Yang--Baxter equation becomes a
tractable functional equation for $R(u)$.

For the XXX chain one often uses
\begin{equation}
\label{eq:xxx-momentum-rapidity-preview}
    \ee^{\ii p(\lambda)}
    =
    \frac{\lambda+\ii/2}{\lambda-\ii/2} .
\end{equation}
Then \cref{eq:xxx-bethe-equations-preview} has the direct interpretation
\begin{equation}
    \ee^{\ii p(\lambda_j)L}
    =
    \prod_{\ell\neq j}S(\lambda_j-\lambda_\ell),
    \qquad
    S(\lambda)=\frac{\lambda+\ii}{\lambda-\ii} .
\end{equation}

\section{What Yang--Baxter integrability does not mean}
\label{sec:what-integrability-does-not-mean}

The term ``integrable'' is used in several related but distinct senses.  In these
notes, Yang--Baxter integrability will mean the existence of an $R$-matrix
satisfying the Yang--Baxter equation, a commuting transfer-matrix family, and a
Hamiltonian obtained from a logarithmic derivative.  This is stronger than merely
knowing many exact eigenstates or finding one special solvable point.

There are several useful cautions.
\begin{enumerate}[label=(\roman*)]
    \item A model may be exactly solvable by free fermions without requiring a
    nontrivial Yang--Baxter structure.
    \item A model may have a transfer matrix but not a commuting family
    $\tau(u)$ for all spectral parameters.
    \item Self-duality may identify a candidate critical point, but it is not the
    same as Bethe-ansatz integrability.
    \item Topological two-band classification concerns eigenvector winding or
    Chern structure; Yang--Baxter integrability concerns factorized scattering
    and commuting transfer matrices.
\end{enumerate}
The value of the mother-model perspective is that these solution mechanisms can
be compared without conflating them.

\section{Roadmap to the later integrable chapters}
\label{sec:integrability-roadmap-preview}

The material in this preview will be used later in three directions.

\begin{enumerate}[label=(\roman*)]
    \item The six-vertex model will provide the canonical example of an
    $R$-matrix satisfying the Yang--Baxter equation.
    \item The Hamiltonian limit of the six-vertex transfer matrix will give the
    XXZ spin chain.
    \item The coordinate Bethe ansatz will solve the XXZ chain by deriving the
    many-body Bethe equations from factorized scattering.
    \item The algebraic Bethe ansatz will solve the same chain from the RTT
    relation and the operator algebra of $A(u),B(u),C(u),D(u)$.
    \item The Hubbard chain will require a nested Bethe ansatz because its
    quasiparticles carry more than one kind of internal degree of freedom.
\end{enumerate}

\section{Summary}
\label{sec:bethe-yang-baxter-summary}

The essential ideas are:
\begin{enumerate}[label=(\roman*)]
    \item Bethe ansatz replaces independent free modes by interacting rapidities
    whose many-body scattering factorizes into two-body scattering.
    \item The coordinate Bethe wave function is a sum of plane waves over
    momentum permutations, with amplitudes related by two-body scattering
    matrices or phases.
    \item Periodic boundary conditions give Bethe equations of the form
    $\ee^{\ii k_jL}=\prod_{\ell\neq j}S(k_j,k_\ell)$.
    \item For matrix-valued scattering, three-body consistency is the
    Yang--Baxter equation.
    \item The $R$-matrix is both a local vertex weight in a two-dimensional
    classical model and the algebraic scattering/intertwining object of a quantum
    chain.
    \item The Yang--Baxter equation implies the RTT relation, which implies
    $[\tau(u),\tau(v)]=0$.
    \item Expanding $\log\tau(u)$ produces mutually commuting conserved charges;
    the quantum Hamiltonian appears as one of these charges.
    \item The six-vertex model and the XXZ chain are the basic interacting
    examples connecting transfer matrices, Bethe ansatz, and Yang--Baxter
    integrability.
\end{enumerate}


% -----------------------------------------------------------------------------
\part{Free-Fermion Mother Models}
% -----------------------------------------------------------------------------

\chapter{The Transverse-Field XY Chain as a Mother Model}
\label{chap:xy-chain-mother-model}

The transverse-field XY chain is the first genuine mother model in these notes.
It is not merely one more spin chain in a list of examples.  It contains, as
parameter limits or equivalent fermionic descriptions, the transverse-field
Ising model, the XX chain, and the spinless Kitaev chain.  It is also the
simplest setting in which all technical tools developed in Part~I meet in one
calculation: Jordan--Wigner fermionization, parity-dependent boundary sectors,
Fourier transformation, BdG blocks, Bogoliubov quasiparticles, Gaussian
correlators, duality, and two-band topology.

The purpose of this chapter is therefore twofold.  First, we solve the XY chain
completely as a quadratic fermion problem.  Second, we organize its special
limits so that the later chapters can be read as controlled reductions of one
solved mother model.

\section{Hamiltonian and conventions}
\label{sec:xy-hamiltonian-conventions}

The transverse-field XY chain is defined by
\begin{equation}
\label{eq:xy-hamiltonian}
    H_{\rm XY}
    =-
    \sum_{j=1}^{L}
    \left[
        \frac{J(1+\gamma)}{2}\sigma_j^x\sigma_{j+1}^x
        +
        \frac{J(1-\gamma)}{2}\sigma_j^y\sigma_{j+1}^y
    \right]
    -h\sum_{j=1}^{L}\sigma_j^z .
\end{equation}
Here $J$ sets the exchange scale, $h$ is the transverse field, and $\gamma$ is the
XY anisotropy.  The isotropic XX chain is obtained at $\gamma=0$, while the
transverse-field Ising chain is obtained at $\gamma=1$.

Unless otherwise stated, the bulk formulas below are independent of whether one
starts from open or periodic spin boundary conditions.  Boundary conditions only
enter when the Jordan--Wigner string wraps around the chain.  This distinction is
central enough that we keep it explicit throughout the derivation.

The spin and fermion conventions are those fixed in Part~I:
\begin{equation}
\label{eq:xy-mother-jw-convention-reminder}
    c_j
    =
    \left(\prod_{\ell=1}^{j-1}\sigma_\ell^z\right)\sigma_j^+,
    \qquad
    c_j^\dagger
    =
    \left(\prod_{\ell=1}^{j-1}\sigma_\ell^z\right)\sigma_j^- ,
\end{equation}
with
\begin{equation}
    n_j=c_j^\dagger c_j=\frac{1-\sigma_j^z}{2},
    \qquad
    \sigma_j^z=1-2n_j.
\end{equation}
Thus a spin-down state is interpreted as an occupied fermion state.  This is the
convention responsible for the sign of the chemical potential in the Kitaev-chain
dictionary below.

It is useful to rewrite the exchange part of \cref{eq:xy-hamiltonian}
as
\begin{equation}
\label{eq:xy-mother-isotropic-anisotropic-split}
    \frac{J(1+\gamma)}{2}\sigma_j^x\sigma_{j+1}^x
    +
    \frac{J(1-\gamma)}{2}\sigma_j^y\sigma_{j+1}^y
    =
    \frac{J}{2}
    \left(
        \sigma_j^x\sigma_{j+1}^x
        +
        \sigma_j^y\sigma_{j+1}^y
    \right)
    +
    \frac{J\gamma}{2}
    \left(
        \sigma_j^x\sigma_{j+1}^x
        -
        \sigma_j^y\sigma_{j+1}^y
    \right).
\end{equation}
The first bracket becomes fermion hopping; the second becomes $p$-wave pairing.
This is the algebraic reason why the XY chain is a BdG mother model.

\section{Special limits and the mother-model viewpoint}
\label{sec:xy-special-limits-mother-model}

The most important reductions of \cref{eq:xy-hamiltonian} are
summarized in \cref{tab:xy-mother-special-limits}.

\begin{table}[h]
\centering
\small
\setlength{\tabcolsep}{3pt}
\caption{Special limits and equivalent descriptions of the transverse-field XY chain.}
\label{tab:xy-mother-special-limits}
\begin{tabularx}{\textwidth}{L{0.23\textwidth}L{0.29\textwidth}Y}
\toprule
Parameter choice or map & Model obtained & Interpretation \\
\midrule
$\gamma=1$ & Transverse-field Ising model & Maximal anisotropy; Ising criticality at $h=\pm J$ in the convention of this chapter \\
$\gamma=0$ & XX chain in a transverse field & Number-conserving free fermions; gapless for $\abs{h}<\abs{J}$ \\
$h=0$ & Zero-field anisotropic XY chain & Pure competition between hopping and pairing after Jordan--Wigner transformation \\
General $h,\gamma$ & Spinless BdG chain & Independent $(k,-k)$ two-level systems after Fourier transformation \\
Jordan--Wigner image & Kitaev $p$-wave chain & $t=J$, $\mu=-2h$, and pairing amplitude fixed by $J\gamma$ up to pairing convention \\
Near $h=\pm J$, $\gamma\neq0$ & Massive Majorana theory & Ising universality class with a sign-changing Majorana mass \\
Near $\gamma=0$, $\abs{h}<\abs{J}$ & Two Fermi points & XX/Luttinger-liquid universality class rather than Ising criticality \\
\bottomrule
\end{tabularx}
\end{table}

The role of the XY chain as a mother model can be summarized as
\begin{equation}
\label{eq:xy-mother-model-map}
\begin{aligned}
    \text{XY spin chain}
    &\xrightarrow{\rm JW}
    \text{spinless BdG chain}\\
    &\xrightarrow{\rm Fourier}
    \prod_{k>0}\text{two-level BdG blocks}
    \xrightarrow{\rm Bogoliubov}
    \text{free quasiparticles}.
\end{aligned}
\end{equation}
All exact results in the later TFIM, XX, and Kitaev-chain chapters are obtained
by specializing this chain of maps.

\section{Jordan--Wigner image}
\label{sec:xy-jordan-wigner-image}

We now derive the fermionic Hamiltonian.  The nearest-neighbor Jordan--Wigner
identities established earlier give
\begin{align}
\label{eq:xy-mother-jw-xx-plus-yy}
    \sigma_j^x\sigma_{j+1}^x+
    \sigma_j^y\sigma_{j+1}^y
    &=
    2\left(
        c_j^\dagger c_{j+1}+c_{j+1}^\dagger c_j
    \right),
    \\
\label{eq:xy-mother-jw-xx-minus-yy}
    \sigma_j^x\sigma_{j+1}^x-
    \sigma_j^y\sigma_{j+1}^y
    &=
    2\left(
        c_j^\dagger c_{j+1}^\dagger+c_{j+1}c_j
    \right),
    \\
\label{eq:xy-mother-jw-z}
    \sigma_j^z&=1-2c_j^\dagger c_j.
\end{align}
Substituting these identities into \cref{eq:xy-hamiltonian} gives the
bulk Jordan--Wigner image
\begin{equation}
\label{eq:xy-mother-bulk-fermion-hamiltonian}
    H_{\rm XY}^{\rm bulk}
    =
    -J\sum_j
    \left(
        c_j^\dagger c_{j+1}+c_{j+1}^\dagger c_j
    \right)
    -J\gamma\sum_j
    \left(
        c_j^\dagger c_{j+1}^\dagger+c_{j+1}c_j
    \right)
    -h\sum_j(1-2n_j).
\end{equation}
Equivalently,
\begin{equation}
\label{eq:xy-mother-bulk-fermion-expanded}
    H_{\rm XY}^{\rm bulk}
    =
    -hL
    +2h\sum_j n_j
    -J\sum_j
    \left(
        c_j^\dagger c_{j+1}+c_{j+1}^\dagger c_j
    \right)
    -J\gamma\sum_j
    \left(
        c_j^\dagger c_{j+1}^\dagger+c_{j+1}c_j
    \right).
\end{equation}
The field term therefore plays the role of a chemical-potential term for the
Jordan--Wigner fermions.

\begin{remark}[Why the anisotropy becomes pairing]
The isotropic part
$\sigma^x\sigma^x+\sigma^y\sigma^y$ conserves the number of spin flips and becomes
fermion hopping.  The anisotropic part
$\sigma^x\sigma^x-\sigma^y\sigma^y$ creates or annihilates spin flips in pairs and
therefore becomes a superconducting pairing term.  Fermion number is not
conserved for $\gamma\neq0$, but fermion parity is conserved.
\end{remark}

For an open chain, \cref{eq:xy-mother-bulk-fermion-hamiltonian} is the full
fermionic Hamiltonian, with the sums running over $j=1,\ldots,L-1$.

For a periodic spin chain, the boundary term carries the Jordan--Wigner parity
sign.  Let
\begin{equation}
\label{eq:xy-mother-fermion-parity}
    P_F=(-1)^N,
    \qquad
    N=\sum_{j=1}^{L}n_j,
\end{equation}
and fix a parity sector
\begin{equation}
    P_F=p,
    \qquad
    p=\pm1.
\end{equation}
As shown in the Jordan--Wigner chapter, the periodic spin chain is represented in
that sector by the effective fermionic boundary condition
\begin{equation}
\label{eq:xy-mother-effective-fermion-boundary}
    c_{L+1}=-p\,c_1.
\end{equation}
Thus
\begin{equation}
\label{eq:xy-mother-pbc-apbc-sector}
    p=+1
    \quad\Longleftrightarrow\quad
    c_{L+1}=-c_1
    \quad\text{(anti-periodic fermions)},
\end{equation}
whereas
\begin{equation}
\label{eq:xy-mother-pbc-pbc-sector}
    p=-1
    \quad\Longleftrightarrow\quad
    c_{L+1}=c_1
    \quad\text{(periodic fermions)}.
\end{equation}
In a fixed sector, one may use
\cref{eq:xy-mother-bulk-fermion-hamiltonian} with $j=1,\ldots,L$ provided the
boundary condition \cref{eq:xy-mother-effective-fermion-boundary} is imposed.  The
exact finite-size spin partition function or spectrum is obtained only after
combining the two parity sectors with the correct parity projection.  In the
thermodynamic bulk spectrum, this distinction affects only finite-size momentum
quantization.

\section{Momentum quantization and Fourier transform}
\label{sec:xy-momentum-quantization-fourier}

In a fixed parity sector $p$, the allowed momenta are
\begin{equation}
\label{eq:xy-mother-allowed-momenta}
    k_m=\frac{2\pi m+\vartheta_p}{L},
    \qquad
    m=0,1,\ldots,L-1,
\end{equation}
where
\begin{equation}
\label{eq:xy-mother-boundary-twist-sector}
    \vartheta_p=
    \begin{cases}
        \pi, & p=+1 \quad \text{(even sector, APBC)},\\
        0, & p=-1 \quad \text{(odd sector, PBC)}.
    \end{cases}
\end{equation}
We write the corresponding momentum set as $\mathcal{K}_{\vartheta_p}$.

We use the Fourier convention developed earlier,
\begin{equation}
\label{eq:xy-mother-fourier-convention}
    c_k
    =
    \frac{\ee^{-\ii\phi}}{\sqrt{L}}
    \sum_{j=1}^{L}\ee^{-\ii kj}c_j,
    \qquad
    c_j
    =
    \frac{\ee^{\ii\phi}}{\sqrt{L}}
    \sum_{k\in\mathcal{K}_{\vartheta_p}}\ee^{\ii kj}c_k .
\end{equation}
The twist $\vartheta_p$ is physical boundary data.  The constant phase $\phi$ is a
Fourier-gauge convention; it rotates the phase of the pairing gap but not the
spectrum.

The hopping and field terms give
\begin{equation}
\label{eq:xy-mother-normal-dispersion}
    -J\sum_j
    \left(
        c_j^\dagger c_{j+1}+c_{j+1}^\dagger c_j
    \right)
    +2h\sum_j n_j
    =
    \sum_{k\in\mathcal{K}_{\vartheta_p}}
    \xi(k)c_k^\dagger c_k,
\end{equation}
with
\begin{equation}
\label{eq:xy-mother-xi-k}
    \xi(k)=2(h-J\cos k).
\end{equation}
The pairing term gives an odd gap function.  With the real-space sign in
\cref{eq:xy-mother-bulk-fermion-hamiltonian}, the direct Fourier transform gives
\begin{equation}
\label{eq:xy-mother-raw-gap}
    \Delta_\phi^{\rm raw}(k)
    =
    -2\ii J\gamma\,\ee^{-2\ii\phi}\sin k.
\end{equation}
This sign is a convention for the orientation of the Nambu pairing plane.  Many
BdG summaries, including some compact formulas in these notes, use the equivalent
orientation
\begin{equation}
\label{eq:xy-mother-standard-gap}
    \Delta_\phi(k)
    =
    +2\ii J\gamma\,\ee^{-2\ii\phi}\sin k.
\end{equation}
The two choices are related by a harmless reversal of the pairing-plane
orientation, for example by $\gamma\mapsto-\gamma$ or by an equivalent Nambu-gauge
relabelling.  They have the same $\abs{\Delta(k)}$, the same quasiparticle
spectrum, the same critical lines, and the same absolute winding number.  Only
orientation-sensitive signs of winding numbers change.  In the rest of this
chapter we display formulas with the convention
\cref{eq:xy-mother-standard-gap}, because it matches the BdG-vector convention
used in the preceding chapters.

With this convention, the momentum-space Hamiltonian is
\begin{equation}
\label{eq:xy-mother-momentum-hamiltonian}
    H_{\rm XY}^{(p)}
    =
    -hL
    +
    \sum_{k\in\mathcal{K}_{\vartheta_p}}
    \xi(k)c_k^\dagger c_k
    +
    \frac12
    \sum_{k\in\mathcal{K}_{\vartheta_p}}
    \left[
        \Delta_\phi(k)c_k^\dagger c_{-k}^\dagger
        +
        \Delta_\phi(k)^*c_{-k}c_k
    \right],
\end{equation}
up to the finite-size parity projection already discussed.

\section{BdG block and quasiparticle spectrum}
\label{sec:xy-bdg-block-spectrum}

Introduce the Nambu spinor
\begin{equation}
\label{eq:xy-mother-nambu-spinor}
    \Psi_k=
    \begin{pmatrix}
        c_k\\ c_{-k}^\dagger
    \end{pmatrix}.
\end{equation}
For $k\not\equiv -k$ modulo $2\pi$, the pair $(k,-k)$ is governed by the BdG
block
\begin{equation}
\label{eq:xy-mother-bdg-block}
    \mathcal{H}_{\rm XY}(k)
    =
    \begin{pmatrix}
        \xi(k)&\Delta_\phi(k)\\
        \Delta_\phi(k)^*&-\xi(k)
    \end{pmatrix}.
\end{equation}
In the gauge $\phi=0$ and with \cref{eq:xy-mother-standard-gap}, this becomes
\begin{equation}
\label{eq:xy-mother-bdg-block-explicit}
    \mathcal{H}_{\rm XY}(k)
    =
    \begin{pmatrix}
        2(h-J\cos k)&2\ii J\gamma\sin k\\
        -2\ii J\gamma\sin k&-2(h-J\cos k)
    \end{pmatrix}.
\end{equation}
Equivalently,
\begin{equation}
\label{eq:xy-mother-d-vector}
    \mathcal{H}_{\rm XY}(k)
    =
    \bm d_{\rm XY}(k)\cdot\bm\tau,
    \qquad
    \bm d_{\rm XY}(k)
    =
    \bigl(0,-2J\gamma\sin k,2(h-J\cos k)\bigr),
\end{equation}
where $\tau_x,\tau_y,\tau_z$ act in Nambu space.  The sign of the $d_y$
component follows from the convention $d_y=-\operatorname{Im}\Delta(k)$.

The quasiparticle energies are
\begin{equation}
\label{eq:xy-mother-quasiparticle-spectrum}
    E_\pm(k)=\pm E(k),
    \qquad
    E(k)
    =
    2\sqrt{(h-J\cos k)^2+(J\gamma\sin k)^2}.
\end{equation}
This is the central exact result of the chapter.  All phase boundaries and low-energy limits follow from the zeros of this dispersion.

\begin{remark}[BdG eigenvalues versus many-body excitations]
The negative branch $-E(k)$ is the particle--hole partner of the positive branch.
It is not an independent physical particle.  After the Bogoliubov
transformation, the many-body Hamiltonian is written in terms of positive-energy
quasiparticles, with a ground-state energy shift.
\end{remark}

\section{Bogoliubov angle and diagonal form}
\label{sec:xy-bogoliubov-angle-diagonal-form}

For each independent non-self-conjugate pair $(k,-k)$, define
\begin{equation}
\label{eq:xy-mother-pair-xi-delta}
    \xi_k=2(h-J\cos k),
    \qquad
    \Delta_k=2\ii J\gamma\,\ee^{-2\ii\phi}\sin k,
    \qquad
    E_k=\sqrt{\xi_k^2+\abs{\Delta_k}^2}.
\end{equation}
The Bogoliubov polar angle is
\begin{equation}
\label{eq:xy-mother-bogoliubov-angle}
    \cos\theta_k=\frac{\xi_k}{E_k},
    \qquad
    \sin\theta_k=\frac{\abs{\Delta_k}}{E_k}.
\end{equation}
Writing
\begin{equation}
    \Delta_k=\abs{\Delta_k}\ee^{\ii\alpha_k},
\end{equation}
the coherence factors may be chosen as
\begin{equation}
\label{eq:xy-mother-coherence-factors}
    u_k=
    \sqrt{\frac{E_k+\xi_k}{2E_k}},
    \qquad
    v_k=
    -\ee^{\ii\alpha_k}
    \sqrt{\frac{E_k-\xi_k}{2E_k}}.
\end{equation}
The quasiparticle operators are
\begin{equation}
\label{eq:xy-mother-quasiparticle-operators}
    \gamma_k=u_kc_k-v_kc_{-k}^\dagger,
    \qquad
    \gamma_{-k}=u_kc_{-k}+v_kc_k^\dagger.
\end{equation}
They satisfy canonical anticommutation relations and diagonalize the pair block:
\begin{equation}
\label{eq:xy-mother-pair-diagonal}
    H_k
    =
    E_k
    \left(
        \gamma_k^\dagger\gamma_k
        +
        \gamma_{-k}^\dagger\gamma_{-k}
        -1
    \right).
\end{equation}
The normalized even-parity ground state of the pair block is
\begin{equation}
\label{eq:xy-mother-pair-ground-state}
    \ket{\Omega_k}
    =
    \left(
        u_k+v_k c_k^\dagger c_{-k}^\dagger
    \right)
    \ket{0}_{k,-k}.
\end{equation}
Thus, ignoring possible self-conjugate modes for the moment, the BCS ground state
in a fixed boundary sector is
\begin{equation}
\label{eq:xy-mother-bcs-ground-state}
    \ket{\Omega}
    =
    \prod_{k\in\mathcal{K}_+}
    \left(
        u_k+v_k c_k^\dagger c_{-k}^\dagger
    \right)
    \ket{0},
\end{equation}
where $\mathcal{K}_+$ contains one representative from each pair $\{k,-k\}$.

The corresponding ground-state energy in the thermodynamic limit is
\begin{equation}
\label{eq:xy-mother-ground-state-energy-density}
    e_0
    =
    \lim_{L\to\infty}\frac{E_0}{L}
    =
    -\frac{1}{2\pi}\int_0^\pi E(k)\,\dd k,
\end{equation}
with $E(k)$ given in \cref{eq:xy-mother-quasiparticle-spectrum}.  This formula
includes the constant $-hL$ in the original spin Hamiltonian.  For example, in
the large-field limit $h\gg J$, it gives $e_0\simeq -h$, as expected for the
fully $+z$ polarized spin state.

\section{Self-conjugate momenta and finite-size parity bookkeeping}
\label{sec:xy-self-conjugate-momenta}

The momenta $k=0$ and $k=\pi$ satisfy $k\equiv-k$ modulo $2\pi$.  They occur only
in the periodic fermion sector and only for compatible system sizes.  Since the
spinless pairing gap is odd in momentum,
\begin{equation}
    \Delta(0)=0,
    \qquad
    \Delta(\pi)=0,
\end{equation}
these modes are not paired.  They must be treated as ordinary fermionic modes.
Their normal energies are
\begin{equation}
\label{eq:xy-mother-self-conjugate-xi}
    \xi(0)=2(h-J),
    \qquad
    \xi(\pi)=2(h+J).
\end{equation}
Occupying such a mode changes fermion parity and therefore affects the exact
finite-size matching between the spin Hilbert space and the fermionic sector.

For bulk thermodynamics, the contribution of these isolated modes is negligible
as $L\to\infty$.  For exact diagonalization, finite-size spectra, and parity
selection, they are essential.  A safe finite-size procedure is:
\begin{enumerate}[label=(\roman*)]
    \item choose a spin boundary condition;
    \item split the fermionic problem into $P_F=+1$ and $P_F=-1$ sectors;
    \item impose APBC in the even sector and PBC in the odd sector;
    \item diagonalize each quadratic Hamiltonian;
    \item keep only states whose quasiparticle occupations have the required
    fermion parity.
\end{enumerate}

\section{Excitation gap and phase diagram}
\label{sec:xy-phase-diagram}

The bulk gap closes when
\begin{equation}
\label{eq:xy-mother-gap-closing-condition}
    h-J\cos k=0,
    \qquad
    J\gamma\sin k=0.
\end{equation}
Assume $J\neq0$.  For $\gamma\neq0$, the second equation forces
\begin{equation}
    k=0
    \quad\text{or}\quad
    k=\pi.
\end{equation}
The first equation then gives the two Ising critical lines
\begin{equation}
\label{eq:xy-mother-ising-critical-lines}
    h=J
    \quad(k=0),
    \qquad
    h=-J
    \quad(k=\pi).
\end{equation}
Near $h=J$, set $k\ll1$.  Then
\begin{equation}
\label{eq:xy-mother-near-h-plus-J}
    E(k)
    \simeq
    2\sqrt{
        \left(h-J+\frac{Jk^2}{2}\right)^2
        +(J\gamma k)^2
    }.
\end{equation}
At the critical point $h=J$ with $\gamma\neq0$, the leading dispersion is
linear,
\begin{equation}
\label{eq:xy-mother-ising-linear-dispersion}
    E(k)\simeq 2\abs{J\gamma}\,\abs{k},
\end{equation}
which is the lattice signature of an Ising/Majorana critical point with dynamical
exponent $z=1$.

Similarly, near $h=-J$ one expands around $k=\pi$.  Writing $q=k-\pi$, one finds
another Ising critical point with the same low-energy structure.

The situation changes qualitatively when $\gamma=0$.  The pairing term vanishes
and the Hamiltonian is number conserving.  The dispersion reduces to
\begin{equation}
\label{eq:xy-mother-xx-dispersion}
    \xi(k)=2(h-J\cos k).
\end{equation}
The system is gapless whenever there are Fermi points satisfying
\begin{equation}
\label{eq:xy-mother-fermi-points}
    \cos k_F=\frac{h}{J}.
\end{equation}
Thus, for $J>0$,
\begin{equation}
\label{eq:xy-mother-xx-critical-region}
    \gamma=0,
    \qquad
    \abs{h}<J
\end{equation}
defines the gapless XX region.  This is not an Ising critical line.  It is a
number-conserving free-fermion critical phase with two Fermi points and central
charge $c=1$ in the continuum limit.  By contrast, the Ising lines at
$\gamma\neq0$ have a single Majorana critical mode and central charge $c=1/2$.

The phase structure is therefore:
\begin{equation}
\label{eq:xy-mother-phase-summary}
\begin{array}{ccl}
\gamma\neq0,\ h=J &:& \text{Ising critical line at } k=0,\\[2pt]
\gamma\neq0,\ h=-J &:& \text{Ising critical line at } k=\pi,\\[2pt]
\gamma=0,\ \abs{h}<J &:& \text{gapless XX/free-fermion critical phase},\\[2pt]
\abs{h}>J &:& \text{gapped polarized phase},\\[2pt]
\abs{h}<J,\ \gamma\neq0 &:& \text{gapped paired/ordered phase}.
\end{array}
\end{equation}
At the multicritical points $(h,\gamma)=(\pm J,0)$, the linear pairing velocity
vanishes and the dispersion along the $\gamma=0$ direction becomes quadratic.
This is why the multicritical point is not described by the same scaling limit as
a generic Ising point.

\section{Relation to the transverse-field Ising model}
\label{sec:xy-relation-to-tfim}

Setting $\gamma=1$ in \cref{eq:xy-hamiltonian} gives
\begin{equation}
\label{eq:xy-mother-tfim-limit}
    H_{\rm TFIM}
    =
    -J\sum_j\sigma_j^x\sigma_{j+1}^x
    -h\sum_j\sigma_j^z.
\end{equation}
The quasiparticle spectrum becomes
\begin{equation}
\label{eq:xy-mother-tfim-spectrum}
    E_{\rm TFIM}(k)
    =
    2\sqrt{(h-J\cos k)^2+J^2\sin^2 k}
    =
    2\sqrt{h^2+J^2-2hJ\cos k}.
\end{equation}
For $J>0$ and $h\ge0$, the physically standard transition is at
\begin{equation}
\label{eq:xy-mother-tfim-critical-h-positive}
    h=J.
\end{equation}
The excitation gap is
\begin{equation}
\label{eq:xy-mother-tfim-gap}
    \Delta_{\rm gap}=2\abs{h-J}.
\end{equation}
The sign of $h-J$ becomes the sign of the continuum Majorana mass.  This is the
fermionic counterpart of the order-disorder transition of the Ising chain.

The second line $h=-J$ is present in the full XY parameter plane.  For the usual
TFIM convention one often restricts to $J>0$ and $h\ge0$, in which case this
second line is outside the standard positive-field discussion.  Keeping it in the
mother-model phase diagram is useful because it is required by the full lattice
BdG structure.

\section{Relation to the XX chain}
\label{sec:xy-relation-to-xx}

Setting $\gamma=0$ removes the pairing term in
\cref{eq:xy-mother-bulk-fermion-hamiltonian}.  The Hamiltonian becomes
\begin{equation}
\label{eq:xy-mother-xx-fermion-hamiltonian}
    H_{\rm XX}
    =
    -J\sum_j
    \left(
        c_j^\dagger c_{j+1}+c_{j+1}^\dagger c_j
    \right)
    -hL+2h\sum_j n_j.
\end{equation}
This is a tight-binding model with dispersion
\begin{equation}
\label{eq:xy-mother-xx-tight-binding-dispersion}
    \xi(k)=2(h-J\cos k).
\end{equation}
The ground state is a Fermi sea, not a BCS paired state.  For $J>0$, modes with
\begin{equation}
    \xi(k)<0
    \quad\Longleftrightarrow\quad
    \cos k>\frac{h}{J}
\end{equation}
are occupied.  When $\abs{h}<J$, the two Fermi points are
\begin{equation}
\label{eq:xy-mother-xx-fermi-momenta}
    k=\pm k_F,
    \qquad
    k_F=\arccos\left(\frac{h}{J}\right).
\end{equation}
Linearizing near the Fermi points gives two chiral fermion modes, or equivalently
a compact boson/Luttinger-liquid continuum theory.  This is why the XX critical
region has central charge $c=1$, not the Ising value $c=1/2$.

\section{Relation to the Kitaev chain}
\label{sec:xy-relation-to-kitaev-chain}

The bulk fermionic image of the XY chain is a spinless $p$-wave superconducting
chain.  To compare conventions cleanly, write the Kitaev chain as
\begin{equation}
\label{eq:xy-mother-kitaev-convention-annihilation}
    H_{\rm K}
    =
    -t\sum_j
    \left(
        c_j^\dagger c_{j+1}+c_{j+1}^\dagger c_j
    \right)
    -\mu\sum_j\left(n_j-\frac12\right)
    +\Delta\sum_j
    \left(
        c_jc_{j+1}+c_{j+1}^\dagger c_j^\dagger
    \right),
\end{equation}
with real $\Delta$ for the moment.  Since
\begin{equation}
    c_{j+1}^\dagger c_j^\dagger=-c_j^\dagger c_{j+1}^\dagger,
    \qquad
    c_jc_{j+1}=-c_{j+1}c_j,
\end{equation}
comparison with \cref{eq:xy-mother-bulk-fermion-expanded} gives the bulk
dictionary
\begin{equation}
\label{eq:xy-mother-kitaev-dictionary}
    t=J,
    \qquad
    \Delta=J\gamma,
    \qquad
    \mu=-2h.
\end{equation}
The constant term also matches:
\begin{equation}
    -\mu\sum_j\left(n_j-\frac12\right)
    =
    2h\sum_j n_j-hL
    \qquad
    \text{when } \mu=-2h.
\end{equation}

Some references instead define the pairing term as
\begin{equation}
    \Delta_{\rm cr}\sum_j
    \left(
        c_j^\dagger c_{j+1}^\dagger+c_{j+1}c_j
    \right).
\end{equation}
In that convention the same XY Hamiltonian corresponds to
\begin{equation}
\label{eq:xy-mother-kitaev-creation-dictionary}
    \Delta_{\rm cr}=-J\gamma.
\end{equation}
This is only a convention for how the Hermitian conjugate pair is written.  The
physical spectrum depends on $\abs{\Delta}$ and on the relative orientation of the
BdG vector, not on this bookkeeping sign alone.

Using \cref{eq:xy-mother-kitaev-dictionary}, the Kitaev-chain topological
condition
\begin{equation}
    \abs{\mu}<2\abs{t}
\end{equation}
becomes
\begin{equation}
\label{eq:xy-mother-kitaev-topological-condition-spin}
    \abs{h}<\abs{J}.
\end{equation}
This is precisely the region between the two Ising critical lines of the XY
mother model when $\gamma\neq0$.  In the spin language it is the ordered region
for the Ising-anisotropic chain; in the fermion language it is the topological
superconducting phase with Majorana edge modes under open boundary conditions.

\section{Topological interpretation of the BdG vector}
\label{sec:xy-bdg-vector-topological-interpretation}

For real $J$, $h$, and $\gamma$, and in a real pairing gauge, the BdG vector lies
in a two-dimensional plane:
\begin{equation}
\label{eq:xy-mother-bdg-plane-vector}
    \bm d(k)=
    \bigl(0,-2J\gamma\sin k,2(h-J\cos k)\bigr).
\end{equation}
Equivalently, one may encode the loop by the complex function
\begin{equation}
\label{eq:xy-mother-q-function}
    q(k)=d_z(k)+\ii d_y(k)
    =
    2(h-J\cos k)-2\ii J\gamma\sin k.
\end{equation}
As $k$ runs around the Brillouin zone, $q(k)$ traces an ellipse centered at
$2h$ on the real axis.  The origin is enclosed if and only if
\begin{equation}
    \abs{h}<\abs{J}
    \qquad
    \text{and}
    \qquad
    \gamma\neq0.
\end{equation}
The corresponding winding number may be written as
\begin{equation}
\label{eq:xy-mother-winding-number}
    \nu
    =
    \frac{1}{2\pi\ii}
    \int_{-\pi}^{\pi}\dd k\,\partial_k\log q(k).
\end{equation}
With the orientation convention in \cref{eq:xy-mother-q-function},
\begin{equation}
\label{eq:xy-mother-winding-result}
    \nu
    =
    \begin{cases}
        -\operatorname{sgn}(J\gamma), & \abs{h}<\abs{J},\\
        0, & \abs{h}>\abs{J}.
    \end{cases}
\end{equation}
Changing the sign convention for the pairing reverses the sign of $\nu$ but does
not change whether the phase is topological.  The robust class-D statement is
the corresponding $\ZZ_2$ distinction:
\begin{equation}
\label{eq:xy-mother-z2-topological-condition}
    \abs{h}<\abs{J}
    \quad\Longleftrightarrow\quad
    \text{nontrivial Kitaev phase}.
\end{equation}

\section{Continuum limits}
\label{sec:xy-continuum-limits}

The same exact lattice dispersion produces different continuum theories in
different regions of parameter space.

\subsection*{Ising/Majorana scaling near $h=J$}

Take $J>0$, $\gamma\neq0$, and expand near $k=0$ and $h=J$.  Keeping the leading
terms gives
\begin{equation}
\label{eq:xy-mother-majorana-expansion-h-plus}
    \mathcal{H}_{\rm XY}(k)
    \simeq
    2(h-J)\tau_z
    -2J\gamma k\,\tau_y,
\end{equation}
up to the sign convention of the pairing plane.  This is the momentum-space form
of a massive Majorana Hamiltonian.  The mass is proportional to
\begin{equation}
\label{eq:xy-mother-majorana-mass-h-plus}
    m_+=2(h-J),
\end{equation}
and the velocity is proportional to
\begin{equation}
    v=2J\gamma.
\end{equation}
The critical theory at $m_+=0$ has central charge $c=1/2$.

Near $h=-J$, one expands around $k=\pi$.  Setting $k=\pi+q$ gives another
Majorana theory with mass proportional to
\begin{equation}
\label{eq:xy-mother-majorana-mass-h-minus}
    m_-=2(h+J).
\end{equation}

\subsection*{XX scaling for $\gamma=0$}

For $\gamma=0$ and $\abs{h}<J$, expand around the two Fermi points
$\pm k_F$.  The Fermi velocity is
\begin{equation}
\label{eq:xy-mother-xx-fermi-velocity}
    v_F
    =
    \left.\abs{\frac{\dd \xi(k)}{\dd k}}\right|_{k=k_F}
    =
    2J\sin k_F
    =
    2J\sqrt{1-\left(\frac{h}{J}\right)^2}.
\end{equation}
The resulting low-energy theory is a massless Dirac fermion, or equivalently a
compact boson.  Adding a small nonzero $\gamma$ pairs the two Fermi points and
gaps the XX critical phase, except at the Ising critical endpoints.

\section{Summary}
\label{sec:xy-mother-summary}

The transverse-field XY chain is solved by the following sequence:
\begin{enumerate}[label=(\roman*)]
    \item Jordan--Wigner transformation maps the spin chain to a quadratic
    spinless fermion Hamiltonian with hopping, chemical potential, and $p$-wave
    pairing.
    \item Periodic spin boundary conditions split the fermionic problem into
    parity sectors: even parity gives APBC, odd parity gives PBC.
    \item Fourier transformation decomposes the Hamiltonian into independent
    $(k,-k)$ BdG blocks.
    \item Each block has
    \begin{equation}
        \xi(k)=2(h-J\cos k),
        \qquad
        \abs{\Delta(k)}=2\abs{J\gamma\sin k}.
    \end{equation}
    \item The exact quasiparticle dispersion is
    \begin{equation}
        E(k)=2\sqrt{(h-J\cos k)^2+(J\gamma\sin k)^2}.
    \end{equation}
    \item For $\gamma\neq0$, the bulk gap closes at $h=\pm J$, giving Ising
    critical lines.
    \item For $\gamma=0$, the chain is number conserving and is gapless for
    $\abs{h}<\abs{J}$, giving the XX/free-fermion critical phase.
    \item Under the Kitaev-chain dictionary $t=J$, $\mu=-2h$, and
    $\Delta=J\gamma$ up to pairing convention, the region $\abs{h}<\abs{J}$ is
    the topological superconducting phase.
\end{enumerate}

This single solution therefore supplies the exact solution of the TFIM, the XX
chain, and the Kitaev chain as special cases or equivalent descriptions.


\chapter{The Transverse-Field Ising Model}
\label{chap:tfim}

The transverse-field Ising model is the most important one-parameter reduction of
the transverse-field XY mother model.  It is simple enough to display the full
solution explicitly, but rich enough to exhibit the main structures of the
lecture notes: Jordan--Wigner fermionization, BdG quasiparticles, parity-sector
bookkeeping, Kramers--Wannier duality, order--disorder variables, and the
Majorana continuum limit.

In the present convention the Hamiltonian is
\begin{equation}
\label{eq:tfim-hamiltonian}
    H_{\rm TFIM}
    =
    -J\sum_{j=1}^{L}\sigma_j^x\sigma_{j+1}^x
    -h\sum_{j=1}^{L}\sigma_j^z .
\end{equation}
It is obtained from the XY chain by setting $\gamma=1$.  In most physical
applications one takes $J>0$ and $h\geq0$.  The full lattice solution, however,
also makes the second critical line at $h=-J$ visible.  We therefore keep signs
explicit when doing so is useful.

\section{Reduction from the XY mother model}
\label{sec:tfim-from-xy-mother}

Setting $\gamma=1$ in the XY-chain Hamiltonian gives
\begin{equation}
    H_{\rm XY}\big|_{\gamma=1}
    =
    -J\sum_j\sigma_j^x\sigma_{j+1}^x
    -h\sum_j\sigma_j^z,
\end{equation}
which is exactly \cref{eq:tfim-hamiltonian}.  The exact solution is therefore
obtained by specializing the XY-chain BdG data:
\begin{equation}
\label{eq:tfim-xi-delta-from-xy}
    \xi(k)=2(h-J\cos k),
    \qquad
    \Delta(k)=2\ii J\sin k
\end{equation}
in the standard BdG orientation used in the preceding chapter.  The
quasiparticle dispersion is
\begin{equation}
\label{eq:tfim-spectrum}
    E_{\rm TFIM}(k)
    =
    2\sqrt{(h-J\cos k)^2+J^2\sin^2 k}
    =
    2\sqrt{h^2+J^2-2hJ\cos k}.
\end{equation}
For $J>0$ and $h\geq0$, the minimum occurs at $k=0$ and the excitation gap is
\begin{equation}
\label{eq:tfim-gap-positive-field}
    \Delta_{\rm gap}=2\abs{h-J}.
\end{equation}
Thus the bulk gap closes at
\begin{equation}
\label{eq:tfim-critical-point}
    h=J.
\end{equation}
If negative fields are also retained, a second gap closing occurs at $h=-J$ and
$k=\pi$.

\begin{remark}[Why the TFIM is not merely a classical Ising chain]
The exchange term in \cref{eq:tfim-hamiltonian} is diagonal in the $\sigma^x$
basis, whereas the field term is diagonal in the $\sigma^z$ basis.  These two
parts do not commute.  The transition at $h=J$ is therefore a genuine quantum
phase transition controlled by the competition between ordering in the $x$
direction and quantum fluctuations generated by the transverse field.
\end{remark}

\section{Jordan--Wigner form and boundary sectors}
\label{sec:tfim-jw-form}

Using the Jordan--Wigner convention fixed earlier,
\begin{equation}
    c_j=
    \left(\prod_{\ell<j}\sigma_\ell^z\right)\sigma_j^+,
    \qquad
    c_j^\dagger=
    \left(\prod_{\ell<j}\sigma_\ell^z\right)\sigma_j^-,
\end{equation}
with $n_j=(1-\sigma_j^z)/2$, the open-chain or bulk fermionic Hamiltonian is
\begin{equation}
\label{eq:tfim-jw-bulk}
    H_{\rm TFIM}^{\rm bulk}
    =
    -J\sum_j
    \left(
        c_j^\dagger c_{j+1}+c_{j+1}^\dagger c_j
    \right)
    -J\sum_j
    \left(
        c_j^\dagger c_{j+1}^\dagger+c_{j+1}c_j
    \right)
    -h\sum_j(1-2n_j).
\end{equation}
This is a spinless $p$-wave superconducting chain with hopping and pairing of
equal magnitude.  Fermion number is not conserved, but fermion parity
\begin{equation}
\label{eq:tfim-fermion-parity}
    P_F=(-1)^N,
    \qquad
    N=\sum_j n_j,
\end{equation}
is conserved.

For a periodic spin chain, the Jordan--Wigner string turns the boundary term
into a parity-dependent fermionic boundary condition.  In a fixed parity sector
$P_F=p$, one imposes
\begin{equation}
\label{eq:tfim-parity-boundary-condition}
    c_{L+1}=-p\,c_1.
\end{equation}
Thus the even sector has anti-periodic fermions, while the odd sector has
periodic fermions:
\begin{equation}
\label{eq:tfim-sector-momenta}
    p=+1:
    \quad
    k=\frac{2\pi m+\pi}{L},
    \qquad
    p=-1:
    \quad
    k=\frac{2\pi m}{L}.
\end{equation}
The thermodynamic spectrum is insensitive to this distinction, but the exact
finite-size spectrum and the spin partition function require the correct parity
projection.

\section{BdG block and Bogoliubov diagonalization}
\label{sec:tfim-bdg-block}

For each non-self-conjugate momentum pair $(k,-k)$, define the Nambu spinor
\begin{equation}
    \Psi_k=
    \begin{pmatrix}
        c_k\\ c_{-k}^\dagger
    \end{pmatrix}.
\end{equation}
The BdG block is
\begin{equation}
\label{eq:tfim-bdg-block}
    \mathcal{H}_{\rm TFIM}(k)
    =
    \begin{pmatrix}
        2(h-J\cos k)&2\ii J\sin k\\
        -2\ii J\sin k&-2(h-J\cos k)
    \end{pmatrix}.
\end{equation}
Equivalently,
\begin{equation}
\label{eq:tfim-bdg-vector}
    \mathcal{H}_{\rm TFIM}(k)
    =
    \bm d_{\rm TFIM}(k)\cdot\bm\tau,
    \qquad
    \bm d_{\rm TFIM}(k)=
    \bigl(0,-2J\sin k,2(h-J\cos k)\bigr).
\end{equation}
The Bogoliubov angle may be chosen as
\begin{equation}
\label{eq:tfim-bogoliubov-angle}
    \cos\theta_k=\frac{h-J\cos k}{
    \sqrt{(h-J\cos k)^2+J^2\sin^2 k}},
    \qquad
    \sin\theta_k=\frac{J\abs{\sin k}}{
    \sqrt{(h-J\cos k)^2+J^2\sin^2 k}}.
\end{equation}
The diagonal many-body form is
\begin{equation}
\label{eq:tfim-diagonal-form}
    H_{\rm TFIM}^{(p)}
    =
    E_0^{(p)}+
    \sum_{k\in\mathcal{K}_{\vartheta_p}}
    E_{\rm TFIM}(k)\,\gamma_k^\dagger\gamma_k,
\end{equation}
where the allowed momentum set depends on the parity sector as in
\cref{eq:tfim-sector-momenta}.  Equivalently, after choosing one representative
from each pair $\{k,-k\}$,
\begin{equation}
\label{eq:tfim-pair-diagonal-form}
    H_k
    =
    E_{\rm TFIM}(k)
    \left(
        \gamma_k^\dagger\gamma_k
        +
        \gamma_{-k}^\dagger\gamma_{-k}
        -1
    \right).
\end{equation}
The BCS ground state in a fixed sector is the $\gamma=1$ specialization of the
XY-chain ground state.

\section{Phases and order parameters}
\label{sec:tfim-phases-order-parameters}

For $J>0$ and $h\geq0$, the two phases are:
\begin{equation}
\label{eq:tfim-phase-summary}
\begin{array}{ccl}
    h>J &:& \text{paramagnetic phase polarized along } z,\\[2pt]
    0\leq h<J &:& \text{ferromagnetic Ising phase ordered along } x.
\end{array}
\end{equation}
In the paramagnetic phase the field dominates and the ground state is smoothly
connected to
\begin{equation}
    \ket{\uparrow_z\uparrow_z\cdots\uparrow_z}.
\end{equation}
In the ferromagnetic phase, the thermodynamic ground space breaks the global
Ising symmetry
\begin{equation}
\label{eq:tfim-z2-symmetry}
    U_Z=\prod_{j=1}^{L}\sigma_j^z,
    \qquad
    U_Z\sigma_j^xU_Z^{-1}=-\sigma_j^x.
\end{equation}
The two symmetry-broken states are approximately the two $x$-polarized states at
$h=0$.

The exact spontaneous magnetization for $J>0$ and $0\leq h<J$ is
\begin{equation}
\label{eq:tfim-exact-order-parameter}
    m_x
    =
    \lim_{r\to\infty}
    \sqrt{\langle \sigma_j^x\sigma_{j+r}^x\rangle}
    =
    \left[1-\left(\frac{h}{J}\right)^2\right]^{1/8}.
\end{equation}
It vanishes for $h\geq J$.  The exponent $1/8$ is the two-dimensional Ising order-parameter exponent, reflecting the quantum--classical correspondence between the
one-dimensional TFIM and the two-dimensional classical Ising model.

\section{Kramers--Wannier duality}
\label{sec:tfim-kramers-wannier-duality}

The operator duality is defined on the dual lattice sites $j+1/2$ by
\begin{equation}
\label{eq:tfim-dual-operators}
    \mu_{j+\frac12}^z
    =
    \sigma_j^x\sigma_{j+1}^x,
    \qquad
    \mu_{j+\frac12}^x
    =
    \prod_{\ell\leq j}\sigma_\ell^z.
\end{equation}
These operators satisfy the same Pauli algebra on the dual lattice.  In
particular,
\begin{equation}
\label{eq:tfim-dual-field-term}
    \sigma_j^z
    =
    \mu_{j-\frac12}^x\mu_{j+\frac12}^x.
\end{equation}
Ignoring boundary subtleties for the moment, the Hamiltonian becomes
\begin{equation}
\label{eq:tfim-dual-hamiltonian}
    H_{\rm TFIM}(J,h)
    =
    -J\sum_j\mu_{j+\frac12}^z
    -h\sum_j\mu_{j-\frac12}^x\mu_{j+\frac12}^x.
\end{equation}
After a relabelling of the dual lattice and a spin-axis relabelling, this has the
same form as the original Hamiltonian with the couplings exchanged:
\begin{equation}
\label{eq:tfim-coupling-exchange}
    H_{\rm TFIM}(J,h)
    \longleftrightarrow
    H_{\rm TFIM}(h,J).
\end{equation}
The self-dual point is therefore
\begin{equation}
\label{eq:tfim-self-dual-point}
    h=J.
\end{equation}
Together with the exact solution, this confirms that the self-dual point is the
actual critical point for $J,h>0$.

The dual order parameter is a disorder operator in the original spin variables:
\begin{equation}
\label{eq:tfim-disorder-correlator}
    \langle \mu_{j+\frac12}^x\mu_{j+r+\frac12}^x\rangle
    =
    \left\langle
    \prod_{\ell=j+1}^{j+r}\sigma_\ell^z
    \right\rangle.
\end{equation}
Thus the ferromagnetic phase of the original spins maps to the disordered phase
of the dual spins, and the paramagnetic phase maps to the ordered phase of the
dual variables.

\section{Continuum Majorana limit}
\label{sec:tfim-continuum-majorana-limit}

For $J>0$ and $h\simeq J$, the gap closes near $k=0$.  Expanding
\cref{eq:tfim-bdg-block} at small $k$ gives
\begin{equation}
\label{eq:tfim-majorana-bdg-expansion}
    \mathcal{H}_{\rm TFIM}(k)
    \simeq
    2(h-J)\tau_z
    -2Jk\,\tau_y.
\end{equation}
The continuum mass and velocity are therefore
\begin{equation}
\label{eq:tfim-majorana-mass-velocity}
    m=2(h-J),
    \qquad
    v=2J.
\end{equation}
In real space, the low-energy Hamiltonian may be written in terms of two real
Majorana fields $\chi_R$ and $\chi_L$ as
\begin{equation}
\label{eq:tfim-majorana-field-theory}
    H_{\rm cont}
    =
    \frac{\ii v}{2}\int \dd x
    \left(
        \chi_R\partial_x\chi_R
        -
        \chi_L\partial_x\chi_L
    \right)
    +
    \ii m\int \dd x\,\chi_R\chi_L.
\end{equation}
The critical point $m=0$ is the Ising conformal field theory with central charge
\begin{equation}
\label{eq:tfim-central-charge}
    c=\frac12.
\end{equation}
The two signs of $m$ correspond to the two gapped phases.  In the fermionic
Kitaev-chain description, the sign change of the Majorana mass is the continuum
version of the topological transition at $\abs{\mu}=2\abs{t}$.

\section{Correlation functions from the Gaussian solution}
\label{sec:tfim-correlators-gaussian}

The fermionic ground state is Gaussian.  Therefore all fermionic correlators are
obtained from the two-point covariance matrix, and all spin correlators follow
from Wick's theorem after expressing Jordan--Wigner strings in terms of Majorana
operators.

For example, the transverse magnetization is local in fermions:
\begin{equation}
\label{eq:tfim-z-magnetization-fermions}
    \sigma_j^z=1-2c_j^\dagger c_j.
\end{equation}
By contrast, the Ising correlator contains a Jordan--Wigner string:
\begin{equation}
\label{eq:tfim-x-correlator-string}
    \sigma_i^x\sigma_j^x
    =
    \left(c_i^\dagger-c_i\right)
    \left[\prod_{\ell=i+1}^{j-1}(1-2n_\ell)\right]
    \left(c_j+c_j^\dagger\right),
    \qquad i<j.
\end{equation}
In Majorana language this becomes a Pfaffian of a finite antisymmetric matrix.
For translation-invariant infinite chains, the Pfaffian often reduces to a
Toeplitz determinant.  This is the computational bridge between the exact
Bogoliubov solution and the exact long-distance order parameter
\cref{eq:tfim-exact-order-parameter}.

\section{Summary}
\label{sec:tfim-summary}

The TFIM is the $\gamma=1$ reduction of the XY mother model.  Its exact solution
is controlled by the BdG dispersion
\begin{equation}
    E(k)=2\sqrt{h^2+J^2-2hJ\cos k}.
\end{equation}
For $J,h>0$, the gap closes at $h=J$, the same point selected by
Kramers--Wannier self-duality.  The phase $h>J$ is paramagnetic, while
$0\leq h<J$ has Ising ferromagnetic order in the $x$ direction.  The continuum
critical theory is a massless Majorana fermion, and the two gapped phases are
distinguished by the sign of the Majorana mass.


\chapter{The XX Chain and Free Fermions}
\label{chap:xx-chain-free-fermions}

The XX chain is the number-conserving limit of the XY mother model.  It is the
cleanest example of how a spin chain can become an ordinary tight-binding
fermion problem after the Jordan--Wigner transformation.  Unlike the TFIM and
the generic XY chain, the XX chain has a conserved fermion number, a Fermi sea,
and a gapless critical phase with two Fermi points.

\section{The \texorpdfstring{$\gamma=0$}{gamma=0} limit of the XY chain}
\label{sec:xx-gamma-zero-limit}

Setting $\gamma=0$ in the transverse-field XY chain gives
\begin{equation}
\label{eq:xx-spin-hamiltonian}
    H_{\rm XX}
    =
    -\frac{J}{2}\sum_{j=1}^{L}
    \left(
        \sigma_j^x\sigma_{j+1}^x
        +
        \sigma_j^y\sigma_{j+1}^y
    \right)
    -h\sum_{j=1}^{L}\sigma_j^z.
\end{equation}
Equivalently, in terms of spin raising and lowering operators,
\begin{equation}
\label{eq:xx-spin-flip-form}
    H_{\rm XX}
    =
    -J\sum_j
    \left(
        \sigma_j^+\sigma_{j+1}^-
        +
        \sigma_j^-\sigma_{j+1}^+
    \right)
    -h\sum_j\sigma_j^z.
\end{equation}
The first term moves a spin flip from one site to a neighboring site.  It does
not create or annihilate spin flips in pairs.  This is the spin-language origin
of fermion-number conservation.

Using the Jordan--Wigner convention of these notes, the bulk fermionic
Hamiltonian is
\begin{equation}
\label{eq:xx-fermion-hamiltonian-real-space}
    H_{\rm XX}^{\rm bulk}
    =
    -J\sum_j
    \left(
        c_j^\dagger c_{j+1}+c_{j+1}^\dagger c_j
    \right)
    -hL+2h\sum_j n_j .
\end{equation}
There is no pairing term.  Therefore
\begin{equation}
\label{eq:xx-number-conservation}
    [H_{\rm XX},N]=0,
    \qquad
    N=\sum_j n_j.
\end{equation}
This is stronger than the parity conservation of the generic XY chain.

For a periodic spin chain, the same parity-sector boundary condition as before
is present:
\begin{equation}
\label{eq:xx-parity-boundary-condition}
    P_F=p
    \quad\Longrightarrow\quad
    c_{L+1}=-p\,c_1.
\end{equation}
In the thermodynamic limit this affects only the quantization of allowed momenta.

\section{Fourier diagonalization}
\label{sec:xx-fourier-diagonalization}

In a fixed boundary sector, use
\begin{equation}
    c_j=\frac{1}{\sqrt{L}}\sum_{k\in\mathcal{K}_\vartheta}
    \ee^{\ii kj}c_k,
\end{equation}
where $\vartheta=0$ or $\pi$ depending on the sector.  The Hamiltonian becomes
\begin{equation}
\label{eq:xx-momentum-hamiltonian}
    H_{\rm XX}
    =
    -hL+
    \sum_{k\in\mathcal{K}_\vartheta}
    \xi(k)c_k^\dagger c_k,
\end{equation}
with single-particle dispersion
\begin{equation}
\label{eq:xx-dispersion}
    \xi(k)=2(h-J\cos k).
\end{equation}
The many-body eigenstates are occupation-number states
\begin{equation}
\label{eq:xx-occupation-basis}
    \ket{\{n_k\}}
    =
    \prod_{k}(c_k^\dagger)^{n_k}\ket{0},
    \qquad
    n_k=0,1,
\end{equation}
with energy
\begin{equation}
\label{eq:xx-many-body-energy}
    E[\{n_k\}]
    =
    -hL+
    \sum_k \xi(k)n_k.
\end{equation}
The diagonalization is therefore simpler than the BdG problem: no Bogoliubov
rotation is required.

\section{Ground-state Fermi sea}
\label{sec:xx-ground-state-fermi-sea}

At zero temperature the ground state fills all modes with negative
single-particle energy:
\begin{equation}
\label{eq:xx-fill-negative-energy}
    \xi(k)<0
    \quad\Longleftrightarrow\quad
    \cos k>\frac{h}{J}
    \qquad (J>0).
\end{equation}
There are three regimes:
\begin{equation}
\label{eq:xx-ground-state-regimes}
\begin{array}{ccl}
    h>J &:& \text{no fermions are occupied},\\[2pt]
    \abs{h}<J &:& \text{a Fermi sea is present},\\[2pt]
    h<-J &:& \text{all fermion modes are occupied}.
\end{array}
\end{equation}
For $\abs{h}<J$, the Fermi points are
\begin{equation}
\label{eq:xx-fermi-momentum}
    k=\pm k_F,
    \qquad
    k_F=\arccos\left(\frac{h}{J}\right).
\end{equation}
The filled region is $\abs{k}<k_F$ in the convention of
\cref{eq:xx-dispersion}.  The fermion density in the thermodynamic limit is
\begin{equation}
\label{eq:xx-fermion-density}
    n
    =
    \langle c_j^\dagger c_j\rangle
    =
    \frac{k_F}{\pi},
    \qquad
    \abs{h}<J.
\end{equation}
Since $\sigma_j^z=1-2n_j$, the magnetization is
\begin{equation}
\label{eq:xx-z-magnetization}
    m_z
    =
    \langle \sigma_j^z\rangle
    =
    1-\frac{2k_F}{\pi},
    \qquad
    k_F=\arccos\left(\frac{h}{J}\right).
\end{equation}
This interpolates from $m_z=-1$ at $h=-J$ to $m_z=+1$ at $h=J$.

The ground-state energy density for $\abs{h}<J$ is
\begin{equation}
\label{eq:xx-ground-state-energy-density}
    e_0
    =
    -h+
    \frac{1}{2\pi}
    \int_{-k_F}^{k_F}2(h-J\cos k)\,\dd k
    =
    -h+
    \frac{2hk_F}{\pi}
    -
    \frac{2J\sin k_F}{\pi}.
\end{equation}
Outside the partially filled regime this expression must be replaced by the
empty or completely filled Fermi sea.

\section{Low-energy Dirac and Luttinger-liquid limit}
\label{sec:xx-low-energy-limit}

In the gapless region $\abs{h}<J$, expand near the two Fermi points:
\begin{equation}
\label{eq:xx-linearization}
    k=\pm k_F+q,
    \qquad
    \abs{q}\ll1.
\end{equation}
The dispersion linearizes as
\begin{equation}
\label{eq:xx-linear-dispersion}
    \xi(+k_F+q)\simeq v_F q,
    \qquad
    \xi(-k_F+q)\simeq -v_F q,
\end{equation}
where
\begin{equation}
\label{eq:xx-fermi-velocity}
    v_F
    =
    \left.\abs{\frac{\dd \xi}{\dd k}}\right|_{k=k_F}
    =
    2J\sin k_F
    =
    2J\sqrt{1-\left(\frac{h}{J}\right)^2}.
\end{equation}
The continuum fermion is decomposed as
\begin{equation}
\label{eq:xx-right-left-decomposition}
    c_j
    \sim
    \ee^{\ii k_F x}\psi_R(x)
    +
    \ee^{-\ii k_F x}\psi_L(x),
    \qquad
    x=ja.
\end{equation}
The low-energy Hamiltonian is the massless Dirac Hamiltonian
\begin{equation}
\label{eq:xx-dirac-continuum-hamiltonian}
    H_{\rm cont}
    =
    -\ii v_F\int\dd x
    \left(
        \psi_R^\dagger\partial_x\psi_R
        -
        \psi_L^\dagger\partial_x\psi_L
    \right).
\end{equation}
In spin-chain language this is the free-fermion point of a Luttinger liquid.  Its
central charge is
\begin{equation}
\label{eq:xx-central-charge}
    c=1,
\end{equation}
which distinguishes the XX critical region from the Ising critical point of the
TFIM, where $c=1/2$.

\section{Correlation functions}
\label{sec:xx-correlation-functions}

Because the Hamiltonian is number conserving and diagonal in momentum space, the
basic equal-time fermion correlator is elementary.  In the thermodynamic ground
state with $\abs{h}<J$,
\begin{equation}
\label{eq:xx-fermion-two-point}
    \langle c_i^\dagger c_j\rangle
    =
    \frac{1}{2\pi}\int_{-k_F}^{k_F}
    \ee^{\ii k(j-i)}\,\dd k
    =
    \begin{cases}
        \dfrac{k_F}{\pi}, & i=j,\\[8pt]
        \dfrac{\sin[k_F(j-i)]}{\pi(j-i)}, & i\neq j.
    \end{cases}
\end{equation}
This $1/r$ decay is the fermionic signature of gaplessness.

The longitudinal spin correlator is local in fermion density:
\begin{equation}
\label{eq:xx-longitudinal-correlator-density}
    \sigma_j^z=1-2n_j.
\end{equation}
The transverse spin correlators, such as
$\langle\sigma_i^x\sigma_j^x\rangle$, contain Jordan--Wigner strings and are
computed as determinants or Pfaffians.  Their long-distance behavior in the
critical region is algebraic, rather than exponential.  At half filling
($h=0$, $k_F=\pi/2$), the leading transverse spin correlations decay as a power
law with the characteristic free-fermion/Luttinger-liquid exponent.

\section{Comparison with the generic XY chain}
\label{sec:xx-comparison-generic-xy}

The XX chain is obtained from the generic XY chain by turning off the pairing
amplitude.  The distinction can be summarized as follows.

\begin{table}[h]
\centering
\small
\setlength{\tabcolsep}{3pt}
\caption{Generic XY chain versus the XX limit.}
\label{tab:xx-vs-xy}
\begin{tabularx}{\textwidth}{L{0.25\textwidth}L{0.33\textwidth}Y}
\toprule
Feature & Generic XY chain $(\gamma\neq0)$ & XX chain $(\gamma=0)$ \\
\midrule
Fermion number & Not conserved & Conserved \\
Fermion parity & Conserved & Conserved, with full $U(1)$ enhanced symmetry \\
Momentum-space block & BdG $(k,-k)$ block & Single-particle dispersion $\xi(k)$ \\
Ground state & BCS paired state & Fermi sea \\
Generic criticality & Ising lines at $h=\pm J$ & Gapless region $\abs{h}<J$ \\
Continuum theory & Majorana, $c=1/2$ & Dirac/Luttinger liquid, $c=1$ \\
\bottomrule
\end{tabularx}
\end{table}

Adding a small nonzero anisotropy $\gamma$ to the gapless XX region introduces a
$p$-wave pairing perturbation.  In the fermionic language this pairs the right-
and left-moving modes and opens a gap except where the lattice mass also closes
at the Ising transition.

\section{Summary}
\label{sec:xx-summary}

The XX chain is the $\gamma=0$ limit of the XY mother model.  It is exactly a
number-conserving tight-binding model with dispersion
\begin{equation}
    \xi(k)=2(h-J\cos k).
\end{equation}
For $J>0$, the ground state is empty for $h>J$, completely filled for $h<-J$,
and a partially filled Fermi sea for $\abs{h}<J$.  The latter regime has two
Fermi points, a linear low-energy Dirac description, and central charge $c=1$.
This chapter therefore illustrates the difference between free-fermion
solvability by a Fermi sea and free-fermion solvability by BdG pairing.


\chapter{The Kitaev Chain as the Fermionic Image of the XY Mother Model}
\label{chap:kitaev-chain-fermionic-image}

The Kitaev chain is the fermionic image of the transverse-field XY chain.  From
the spin point of view it is obtained by the Jordan--Wigner transformation; from
the fermion point of view it is a one-dimensional spinless $p$-wave
superconductor.  This chapter isolates the fermionic interpretation of the XY
solution and connects it to Majorana edge modes and the BdG topological
invariant.

\section{Spinless \texorpdfstring{$p$}{p}-wave superconducting chain}
\label{sec:kitaev-p-wave-chain}

We use the convention
\begin{equation}
\label{eq:kitaev-hamiltonian}
    H_{\rm K}
    =
    -t\sum_{j}
    \left(
        c_j^\dagger c_{j+1}+c_{j+1}^\dagger c_j
    \right)
    -\mu\sum_j\left(n_j-\frac12\right)
    +\Delta\sum_j
    \left(
        c_jc_{j+1}+c_{j+1}^\dagger c_j^\dagger
    \right),
\end{equation}
where $n_j=c_j^\dagger c_j$.  The hopping $t$ and chemical potential $\mu$ are
ordinary number-conserving terms, while $\Delta$ is a nearest-neighbor odd-parity
pairing amplitude.  For the moment $\Delta$ is taken real.  A uniform complex
phase will be discussed below.

The pairing term breaks fermion-number conservation but preserves fermion parity:
\begin{equation}
\label{eq:kitaev-parity-conservation}
    [H_{\rm K},(-1)^N]=0,
    \qquad
    N=\sum_j n_j.
\end{equation}
Thus the Kitaev chain belongs naturally to the same BdG framework used for the
XY chain.

\section{Dictionary from the XY mother model}
\label{sec:kitaev-xy-dictionary}

The bulk Jordan--Wigner image of the XY chain in the conventions of these notes
is
\begin{equation}
\label{eq:kitaev-xy-bulk-image}
    H_{\rm XY}^{\rm bulk}
    =
    -J\sum_j
    \left(
        c_j^\dagger c_{j+1}+c_{j+1}^\dagger c_j
    \right)
    -J\gamma\sum_j
    \left(
        c_j^\dagger c_{j+1}^\dagger+c_{j+1}c_j
    \right)
    -h\sum_j(1-2n_j).
\end{equation}
Using
\begin{equation}
    c_{j+1}^\dagger c_j^\dagger=-c_j^\dagger c_{j+1}^\dagger,
    \qquad
    c_jc_{j+1}=-c_{j+1}c_j,
\end{equation}
comparison between \cref{eq:kitaev-hamiltonian} and
\cref{eq:kitaev-xy-bulk-image} gives the bulk dictionary
\begin{equation}
\label{eq:kitaev-xy-dictionary}
    t=J,
    \qquad
    \Delta=J\gamma,
    \qquad
    \mu=-2h.
\end{equation}
The constant term matches because
\begin{equation}
    -\mu\sum_j\left(n_j-\frac12\right)
    =
    2h\sum_j n_j-hL
    \qquad
    (\mu=-2h).
\end{equation}

This dictionary is convention-dependent.  If one writes the pairing term as
$\Delta_{\rm cr}\sum_j(c_j^\dagger c_{j+1}^\dagger+c_{j+1}c_j)$ instead, then
$\Delta_{\rm cr}=-J\gamma$ in the present Jordan--Wigner convention.  The
spectrum and phase diagram are unchanged; only the orientation convention for
the BdG pairing plane is changed.

For a periodic spin chain, the boundary term is also sector dependent.  In a
fixed fermion-parity sector,
\begin{equation}
\label{eq:kitaev-boundary-from-spin}
    P_F=p
    \quad\Longrightarrow\quad
    c_{L+1}=-p\,c_1.
\end{equation}
For an intrinsically fermionic Kitaev chain, one may instead specify periodic,
anti-periodic, or open boundary conditions directly.  The bulk topological phase
diagram is the same, but exact finite-size spectra depend on this choice.

\section{BdG block and spectrum}
\label{sec:kitaev-bdg-block-spectrum}

For a translation-invariant chain, the normal-state dispersion is
\begin{equation}
\label{eq:kitaev-xi-k}
    \xi(k)=-\mu-2t\cos k.
\end{equation}
With the Fourier and BdG convention used in the earlier chapters, the odd-parity
pairing gap for real $\Delta$ may be written as
\begin{equation}
\label{eq:kitaev-gap-real}
    \Delta(k)=2\ii\Delta\sin k.
\end{equation}
The BdG block is therefore
\begin{equation}
\label{eq:kitaev-bdg-block}
    \mathcal{H}_{\rm K}(k)
    =
    \begin{pmatrix}
        -\mu-2t\cos k & 2\ii\Delta\sin k\\
        -2\ii\Delta\sin k & \mu+2t\cos k
    \end{pmatrix}.
\end{equation}
Equivalently,
\begin{equation}
\label{eq:kitaev-d-vector}
    \mathcal{H}_{\rm K}(k)
    =
    \bm d_{\rm K}(k)\cdot\bm\tau,
    \qquad
    \bm d_{\rm K}(k)
    =
    \bigl(0,-2\Delta\sin k,-\mu-2t\cos k\bigr).
\end{equation}
The quasiparticle spectrum is
\begin{equation}
\label{eq:kitaev-spectrum}
    E(k)
    =
    \sqrt{(-\mu-2t\cos k)^2+4\Delta^2\sin^2 k}.
\end{equation}
Under \cref{eq:kitaev-xy-dictionary}, this reproduces the XY-chain spectrum
\begin{equation}
    E(k)
    =
    2\sqrt{(h-J\cos k)^2+(J\gamma\sin k)^2}.
\end{equation}

For $\Delta\neq0$, the bulk gap can close only at $k=0$ or $k=\pi$.  The two
conditions are
\begin{equation}
\label{eq:kitaev-gap-closing}
    k=0:\quad \mu=-2t,
    \qquad
    k=\pi:
    \quad \mu=2t.
\end{equation}
Thus the gapped phases are separated by the lines
\begin{equation}
\label{eq:kitaev-critical-lines}
    \mu=\pm 2t.
\end{equation}

\section{Complex pairing phase and Nambu gauge}
\label{sec:kitaev-complex-gap-phase}

Let the pairing amplitude be complex:
\begin{equation}
\label{eq:kitaev-complex-delta}
    \Delta=\abs{\Delta}\ee^{\ii\theta_\Delta}.
\end{equation}
With the Fourier convention containing a constant phase $\phi$,
\begin{equation}
    c_k=\frac{\ee^{-\ii\phi}}{\sqrt{L}}
    \sum_j\ee^{-\ii kj}c_j,
\end{equation}
the momentum-space gap becomes
\begin{equation}
\label{eq:kitaev-complex-gap-fourier-phase}
    \Delta_\phi(k)
    =
    2\ii\abs{\Delta}\,
    \ee^{\ii(\theta_\Delta-2\phi)}\sin k.
\end{equation}
The phase $\phi$ is a Nambu/Fourier gauge convention.  Choosing
\begin{equation}
\label{eq:kitaev-gauge-remove-uniform-phase}
    \phi=\frac{\theta_\Delta}{2}
\end{equation}
removes the uniform superconducting phase from the bulk BdG block.  Therefore a
spatially uniform pairing phase is not by itself a new bulk phase of matter.  It
rotates the BdG vector in the pairing plane but leaves $E(k)$ and the topological
criterion unchanged.

A phase difference, phase winding, or boundary flux is different.  If the pairing
phase cannot be removed by a single-valued global gauge transformation, it may
represent physical supercurrent or flux data and can affect boundary spectra.

\section{Majorana representation}
\label{sec:kitaev-majorana-representation}

Define site Majorana operators using the convention of the covariance-matrix
chapter:
\begin{equation}
\label{eq:kitaev-majorana-definitions}
    A_j=c_j+c_j^\dagger,
    \qquad
    B_j=\ii(c_j-c_j^\dagger).
\end{equation}
They satisfy
\begin{equation}
    A_j^\dagger=A_j,
    \qquad
    B_j^\dagger=B_j,
    \qquad
    \{A_i,A_j\}=2\delta_{ij},
    \qquad
    \{B_i,B_j\}=2\delta_{ij},
    \qquad
    \{A_i,B_j\}=0.
\end{equation}
Also
\begin{equation}
\label{eq:kitaev-density-majorana}
    n_j-\frac12=-\frac{\ii}{2}A_jB_j.
\end{equation}
For real $\Delta$, the Hamiltonian becomes
\begin{equation}
\label{eq:kitaev-majorana-hamiltonian}
    H_{\rm K}
    =
    \frac{\ii}{2}\sum_j
    \left[
        \mu A_jB_j
        +(t-\Delta)A_jB_{j+1}
        -(t+\Delta)B_jA_{j+1}
    \right],
\end{equation}
up to boundary terms fixed by the chosen boundary condition.

Two limiting cases reveal the phase structure immediately.

\paragraph{Trivial atomic limit.}
For $t=\Delta=0$ and $\mu<0$ or $\mu>0$, Majoranas on the same site are paired by
$\ii\mu A_jB_j/2$.  There are no protected edge Majoranas.

\paragraph{Topological dimerized limit.}
For $\mu=0$ and $\Delta=t>0$,
\begin{equation}
\label{eq:kitaev-topological-special-point}
    H_{\rm K}
    =
    -\ii t\sum_j B_jA_{j+1}.
\end{equation}
On an open chain this pairs $B_j$ with $A_{j+1}$ on each bond.  The operators
$A_1$ and $B_L$ are left unpaired and become zero-energy Majorana edge modes.
The two edge Majoranas combine into a nonlocal fermion
\begin{equation}
\label{eq:kitaev-edge-fermion}
    f_{\rm edge}=\frac{1}{2}(A_1+\ii B_L),
\end{equation}
producing two nearly degenerate ground states of opposite fermion parity in a
long open chain.

\section{Topological invariant}
\label{sec:kitaev-topological-invariant}

For real parameters, the BdG vector lies in a plane.  Define
\begin{equation}
\label{eq:kitaev-q-function}
    q(k)
    =
    d_z(k)+\ii d_y(k)
    =
    -\mu-2t\cos k-2\ii\Delta\sin k.
\end{equation}
The integer winding number in the real-pairing BDI convention is
\begin{equation}
\label{eq:kitaev-winding-number}
    \nu
    =
    \frac{1}{2\pi\ii}
    \int_{-\pi}^{\pi}\dd k\,\partial_k\log q(k).
\end{equation}
For $t\Delta\neq0$, the loop winds around the origin if and only if
\begin{equation}
\label{eq:kitaev-topological-condition}
    \abs{\mu}<2\abs{t}.
\end{equation}
The orientation of the winding depends on the sign of $t\Delta$ and on the
chosen definition of $q(k)$, but the distinction between zero and nonzero winding
is robust.

If only particle--hole symmetry is kept, the class-D invariant is the
corresponding $\ZZ_2$ quantity.  In the present convention it can be written as
\begin{equation}
\label{eq:kitaev-pfaffian-invariant}
    \mathcal{M}
    =
    \sgn\{\xi(0)\xi(\pi)\}
    =
    \sgn\{(-\mu-2t)(-\mu+2t)\}
    =
    \sgn(\mu^2-4t^2).
\end{equation}
The topological phase has
\begin{equation}
\label{eq:kitaev-class-d-nontrivial}
    \mathcal{M}=-1
    \quad\Longleftrightarrow\quad
    \abs{\mu}<2\abs{t}.
\end{equation}
Using the XY dictionary $\mu=-2h$ and $t=J$, this becomes
\begin{equation}
\label{eq:kitaev-xy-topological-condition}
    \abs{h}<\abs{J},
\end{equation}
which is the region between the two Ising gap-closing lines of the XY mother
model.

\section{Relation to the TFIM and the XY chain}
\label{sec:kitaev-relation-to-tfim-xy}

The transverse-field Ising model corresponds to the Kitaev chain with
\begin{equation}
\label{eq:kitaev-tfim-dictionary}
    \gamma=1
    \quad\Longleftrightarrow\quad
    \Delta=t=J,
    \qquad
    \mu=-2h.
\end{equation}
For $J,h>0$, the TFIM transition at $h=J$ is the Kitaev transition at
\begin{equation}
    \mu=-2t.
\end{equation}
The spin ferromagnetic phase $h<J$ is the topological superconducting phase
$\abs{\mu}<2\abs{t}$.  The spin paramagnet $h>J$ is the trivial superconducting
phase.

This identification is conceptually useful but should not be overinterpreted.
The fermions in the Kitaev-chain description are Jordan--Wigner fermions, not
microscopic electrons of the original spin chain.  The Majorana edge modes of
the open Kitaev chain correspond, in the spin language, to the near degeneracy of
the symmetry-broken Ising phase and to nonlocal string operators.

\section{Summary}
\label{sec:kitaev-summary}

The Kitaev chain is the fermionic form of the XY mother model.  The convention
used in these notes gives the bulk dictionary
\begin{equation}
    t=J,
    \qquad
    \Delta=J\gamma,
    \qquad
    \mu=-2h.
\end{equation}
Its BdG spectrum is
\begin{equation}
    E(k)=\sqrt{(-\mu-2t\cos k)^2+4\abs{\Delta}^2\sin^2 k}.
\end{equation}
For nonzero pairing, the bulk transitions occur at $\mu=\pm2t$.  The phase
$\abs{\mu}<2\abs{t}$ is topological and supports Majorana edge modes on an open
chain.  A uniform complex phase of $\Delta$ rotates the BdG pairing plane and can
be absorbed by a Nambu gauge choice, while phase winding or boundary flux may
carry physical information.


\chapter{The SSH Chain and Two-Band Chiral Fermions}
\label{chap:ssh-chain-two-band-chiral-fermions}

The Su--Schrieffer--Heeger (SSH) chain is not a Jordan--Wigner image of the XY
mother model.  It is nevertheless an essential companion example because it is
the simplest number-conserving two-band topological Hamiltonian.  The Kitaev
chain uses a Nambu particle--hole basis; the SSH chain uses a sublattice basis.
Both reduce topology to the winding of a two-component vector around the origin,
but the Hilbert-space meaning of the two components is different.

\section{SSH Hamiltonian}
\label{sec:ssh-hamiltonian}

Consider a one-dimensional chain with two sites per unit cell, labelled $A$ and
$B$.  Let $a_j$ annihilate a fermion on sublattice $A$ in unit cell $j$, and let
$b_j$ annihilate a fermion on sublattice $B$.  The SSH Hamiltonian is
\begin{equation}
\label{eq:ssh-real-space-hamiltonian}
    H_{\rm SSH}
    =
    \sum_{j}
    \left(
        t_1 a_j^\dagger b_j
        +
        t_2 b_j^\dagger a_{j+1}
        +
        \text{h.c.}
    \right).
\end{equation}
Here $t_1$ is the intracell hopping and $t_2$ is the intercell hopping.  The
model conserves fermion number:
\begin{equation}
\label{eq:ssh-number-conservation}
    [H_{\rm SSH},N]=0,
    \qquad
    N=\sum_j(a_j^\dagger a_j+b_j^\dagger b_j).
\end{equation}
Thus it is a Bloch-band problem rather than a BdG problem.

The two exactly solvable limits are instructive.  If $t_2=0$, the chain is a set
of isolated intracell dimers.  If $t_1=0$, the chain is a set of intercell dimers;
for an open chain, the leftmost $A$ site and rightmost $B$ site are unpaired edge
states.  The topological distinction between these two dimerizations is the
central content of the SSH model.

\section{Bloch form}
\label{sec:ssh-bloch-form}

Use the Fourier convention
\begin{equation}
    a_j=\frac{1}{\sqrt{L}}\sum_k\ee^{\ii kj}a_k,
    \qquad
    b_j=\frac{1}{\sqrt{L}}\sum_k\ee^{\ii kj}b_k.
\end{equation}
In the basis
\begin{equation}
\label{eq:ssh-bloch-spinor}
    \psi_k=
    \begin{pmatrix}
        a_k\\ b_k
    \end{pmatrix},
\end{equation}
the Hamiltonian becomes
\begin{equation}
\label{eq:ssh-momentum-hamiltonian}
    H_{\rm SSH}
    =
    \sum_k\psi_k^\dagger h_{\rm SSH}(k)\psi_k,
\end{equation}
where
\begin{equation}
\label{eq:ssh-bloch-matrix}
    h_{\rm SSH}(k)
    =
    \begin{pmatrix}
        0&q(k)\\
        q(k)^*&0
    \end{pmatrix},
    \qquad
    q(k)=t_1+t_2\ee^{-\ii k}.
\end{equation}
Equivalently,
\begin{equation}
\label{eq:ssh-d-vector}
    h_{\rm SSH}(k)
    =
    d_x(k)\sigma_x+d_y(k)\sigma_y,
\end{equation}
with
\begin{equation}
\label{eq:ssh-d-components}
    d_x(k)=t_1+t_2\cos k,
    \qquad
    d_y(k)=t_2\sin k.
\end{equation}
The energy bands are
\begin{equation}
\label{eq:ssh-band-energies}
    E_\pm(k)=\pm\abs{q(k)}
    =
    \pm\sqrt{t_1^2+t_2^2+2t_1t_2\cos k}.
\end{equation}
The gap closes when
\begin{equation}
\label{eq:ssh-gap-closing}
    q(k)=0.
\end{equation}
For real positive hoppings this occurs at
\begin{equation}
\label{eq:ssh-critical-point}
    t_1=t_2,
    \qquad
    k=\pi.
\end{equation}

\section{Chiral symmetry and winding number}
\label{sec:ssh-chiral-symmetry-winding}

The SSH Bloch Hamiltonian anticommutes with $\sigma_z$:
\begin{equation}
\label{eq:ssh-chiral-symmetry}
    \sigma_z h_{\rm SSH}(k)\sigma_z=-h_{\rm SSH}(k).
\end{equation}
This is sublattice, or chiral, symmetry.  It forces the Hamiltonian to be
off-diagonal in the $(A,B)$ basis and pins the spectrum symmetrically around zero
energy.

When the bulk is gapped, $q(k)$ never vanishes, so the map
\begin{equation}
\label{eq:ssh-q-map}
    k\in S^1_{\rm BZ}
    \longmapsto
    q(k)\in\CC\setminus\{0\}
\end{equation}
has an integer winding number.  With the off-diagonal convention
\cref{eq:ssh-bloch-matrix}, define
\begin{equation}
\label{eq:ssh-winding-number}
    \nu_{\rm SSH}
    =
    \frac{1}{2\pi\ii}
    \int_{-\pi}^{\pi}\dd k\,
    \partial_k\log q(k).
\end{equation}
For $t_1,t_2>0$ and $q(k)=t_1+t_2\ee^{-\ii k}$, the result is
\begin{equation}
\label{eq:ssh-winding-result}
    \nu_{\rm SSH}
    =
    \begin{cases}
        -1, & t_2>t_1,\\
        0, & t_1>t_2.
    \end{cases}
\end{equation}
The sign is an orientation convention coming from the choice $\ee^{-\ii k}$ in
$q(k)$.  Replacing $q(k)$ by $t_1+t_2\ee^{\ii k}$ reverses the sign.  The robust
physical distinction is whether the winding is zero or nonzero.

Geometrically, $q(k)$ traces a circle of radius $\abs{t_2}$ centered at $t_1$ in
the complex plane.  The nontrivial phase is precisely the case in which this
circle encloses the origin:
\begin{equation}
\label{eq:ssh-topological-condition}
    \abs{t_2}>\abs{t_1}.
\end{equation}

\section{Zak phase}
\label{sec:ssh-zak-phase}

Let
\begin{equation}
    q(k)=\abs{q(k)}\ee^{\ii\varphi(k)}.
\end{equation}
A normalized occupied-band eigenstate may be chosen locally as
\begin{equation}
\label{eq:ssh-occupied-eigenstate}
    \ket{u_-(k)}
    =
    \frac{1}{\sqrt{2}}
    \begin{pmatrix}
        -\ee^{\ii\varphi(k)}\\ 1
    \end{pmatrix}.
\end{equation}
Its Berry connection is
\begin{equation}
\label{eq:ssh-berry-connection}
    \mathcal{A}(k)
    =
    \ii\braket{u_-(k)}{\partial_k u_-(k)}
    =
    -\frac12\partial_k\varphi(k),
\end{equation}
up to gauge conventions.  The Zak phase is therefore
\begin{equation}
\label{eq:ssh-zak-phase}
    \Phi_{\rm Zak}
    =
    \int_{-\pi}^{\pi}\mathcal{A}(k)\,\dd k
    =
    -\pi\nu_{\rm SSH}
    \quad
    \text{mod }2\pi.
\end{equation}
Thus the nontrivial SSH phase has Zak phase $\pi$ modulo $2\pi$, provided the
unit-cell convention is fixed.  Changing the unit cell shifts the Zak phase by an
integer multiple of $2\pi$ or changes the representative convention, but it does
not remove the protected boundary-state distinction for a specified physical
termination.

\section{Open chain and edge modes}
\label{sec:ssh-edge-modes}

For an open chain with $L$ unit cells, the zero-energy eigenvalue equation at
energy $E=0$ decouples between sublattices.  A left edge mode supported on the
$A$ sublattice has the form
\begin{equation}
\label{eq:ssh-left-edge-ansatz}
    \ket{\psi_L}
    =
    \sum_{j=1}^{L}\alpha_j a_j^\dagger\ket{0}.
\end{equation}
The Schrödinger equation gives the recursion
\begin{equation}
\label{eq:ssh-edge-recursion}
    t_1\alpha_j+t_2\alpha_{j+1}=0,
\end{equation}
so that
\begin{equation}
\label{eq:ssh-edge-wavefunction}
    \alpha_j
    =
    \left(-\frac{t_1}{t_2}\right)^{j-1}\alpha_1.
\end{equation}
This mode is normalizable at the left edge in the thermodynamic limit if and only
if
\begin{equation}
\label{eq:ssh-edge-normalizability}
    \abs{t_1}<\abs{t_2}.
\end{equation}
The right edge mode is similarly localized on the $B$ sublattice.  Thus the
bulk winding condition and the edge-state condition agree.

At the exactly dimerized point $t_1=0$, the edge modes are simply the unpaired
operators $a_1^\dagger$ and $b_L^\dagger$.  This is the SSH analogue of the
unpaired Majorana operators in the topological limit of the Kitaev chain.

\section{Relation and distinction from the Kitaev chain}
\label{sec:ssh-relation-distinction-kitaev}

Both the SSH and Kitaev chains are two-band topological models, but the meanings
of the two-component spinors are different:
\begin{equation}
\label{eq:ssh-kitaev-spinor-comparison}
    \psi_k^{\rm SSH}
    =
    \begin{pmatrix}
        a_k\\ b_k
    \end{pmatrix},
    \qquad
    \Psi_k^{\rm Kitaev}
    =
    \begin{pmatrix}
        c_k\\ c_{-k}^\dagger
    \end{pmatrix}.
\end{equation}
The SSH Pauli matrices act in sublattice space.  The Kitaev Pauli matrices act in
Nambu particle--hole space.  Consequently:
\begin{itemize}[leftmargin=2.2em]
    \item the SSH chain conserves fermion number, while the Kitaev chain
    conserves only fermion parity;
    \item the SSH spectral symmetry comes from sublattice chiral symmetry,
    while the Kitaev BdG spectral symmetry comes from particle--hole redundancy;
    \item SSH edge states are ordinary complex-fermion boundary orbitals, while
    Kitaev edge states are Majorana zero modes;
    \item the SSH winding is naturally a chiral-band winding, while the Kitaev
    winding is a BdG/Nambu winding, reducible to a class-D $\ZZ_2$ invariant when
    the real-pairing chiral symmetry is not imposed.
\end{itemize}

This comparison is useful because the same mathematical form
\begin{equation}
\label{eq:ssh-kitaev-two-band-common-form}
    H(k)=d_x(k)\sigma_x+d_y(k)\sigma_y+d_z(k)\sigma_z
\end{equation}
can describe physically distinct objects.  A two-band matrix by itself does not
say whether the Pauli matrices refer to sublattice, spin, orbital, or Nambu
space.  The symmetry interpretation must always be specified.

\section{SSH as a two-band chiral mother example}
\label{sec:ssh-two-band-chiral-mother}

The SSH chain represents the simplest member of a broader family of one-dimensional chiral two-band Hamiltonians:
\begin{equation}
\label{eq:ssh-general-chiral-two-band}
    h(k)=
    \begin{pmatrix}
        0&q(k)\\ q(k)^*&0
    \end{pmatrix}.
\end{equation}
Every such gapped Hamiltonian has a winding number
\begin{equation}
\label{eq:ssh-general-winding}
    \nu
    =
    \frac{1}{2\pi\ii}\int_{\rm BZ}\dd k\,
    \partial_k\log q(k).
\end{equation}
Longer-range hoppings simply replace $q(k)=t_1+t_2\ee^{-\ii k}$ by a finite
Laurent polynomial in $\ee^{\ii k}$.  The topology is still controlled by whether
and how this complex loop winds around the origin.  The SSH model is the minimal
case in which this winding can be nonzero.

\section{Summary}
\label{sec:ssh-summary}

The SSH chain is a number-conserving two-band model with Bloch Hamiltonian
\begin{equation}
    h_{\rm SSH}(k)=
    \begin{pmatrix}
        0&t_1+t_2\ee^{-\ii k}\\
        t_1+t_2\ee^{\ii k}&0
    \end{pmatrix}.
\end{equation}
It has chiral symmetry and an integer winding number.  For real positive
hoppings, the nontrivial phase occurs when $t_2>t_1$ and supports boundary zero
modes on an open chain.  It parallels the Kitaev chain at the level of two-band
topology, but differs physically because its two-component structure is
sublattice space rather than Nambu space.


% -----------------------------------------------------------------------------
\part{Transfer Matrices and Quantum--Classical Correspondence}
% -----------------------------------------------------------------------------

\chapter{The Two-Dimensional Classical Ising Model}
\label{chap:two-dimensional-classical-ising}

The two-dimensional classical Ising model is the first genuinely nontrivial
exactly solved lattice model in these notes.  It is also the central bridge
between two themes developed earlier: transfer matrices and free fermions.  In
one dimension, the classical Ising transfer matrix is only a $2\times2$ matrix;
in two dimensions, a row-to-row transfer matrix acts on a $2^L$-dimensional row
Hilbert space.  The remarkable fact, discovered by Onsager and reformulated by
Kaufman, is that this exponentially large transfer matrix is still exactly
solvable because it is a Gaussian operator in Majorana variables.

This chapter has three goals.  First, we define the anisotropic zero-field model
and construct its row-to-row transfer matrix in the same convention used in the
quantum--classical correspondence chapter.  Second, we explain why the transfer
matrix belongs to a free-fermion algebra.  Third, we extract the Onsager critical
condition
\begin{equation}
\label{eq:2d-classical-ising-onsager-critical-preview}
    \sinh(2K_x)\sinh(2K_y)=1
\end{equation}
from the closing of the transfer-matrix gap.

\begin{remark}[Zero field is essential]
The exact free-fermion solution discussed here applies to the square-lattice
Ising model at zero magnetic field.  Adding a field term
$H\sum_{j,m}s_{j,m}$ destroys the Gaussian structure of the transfer matrix.
The two-dimensional Ising model in a nonzero magnetic field is not known to have
a comparable closed-form exact solution.
\end{remark}

\section{Classical partition function}
\label{sec:2d-classical-ising-partition-function}

Consider an $L\times M$ square lattice with periodic boundary conditions in both
directions.  A classical spin variable
\begin{equation}
    s_{j,m}=\pm1,
    \qquad
    j=1,\ldots,L,
    \qquad
    m=1,\ldots,M,
\end{equation}
lives at each lattice site.  We identify
\begin{equation}
    s_{L+1,m}=s_{1,m},
    \qquad
    s_{j,M+1}=s_{j,1}.
\end{equation}
The anisotropic zero-field Ising partition function is
\begin{equation}
\label{eq:2d-classical-ising-partition-function}
    Z_{L,M}(K_x,K_y)
    =
    \sum_{\{s_{j,m}=\pm1\}}
    \exp\left[
        K_x\sum_{j,m}s_{j,m}s_{j+1,m}
        +
        K_y\sum_{j,m}s_{j,m}s_{j,m+1}
    \right].
\end{equation}
Here $K_x$ and $K_y$ are dimensionless couplings.  If the physical Hamiltonian is
\begin{equation}
    E[s]
    =
    -J_x\sum_{j,m}s_{j,m}s_{j+1,m}
    -J_y\sum_{j,m}s_{j,m}s_{j,m+1},
\end{equation}
then
\begin{equation}
    K_x=\beta J_x,
    \qquad
    K_y=\beta J_y .
\end{equation}
For most of this chapter we assume $K_x,K_y>0$, corresponding to the
ferromagnetic model.  Antiferromagnetic signs can often be moved by sublattice
spin flips on a bipartite lattice, but the ferromagnetic case is the canonical
one for discussing the thermal order--disorder transition.

Let
\begin{equation}
    \bm{s}_m=(s_{1,m},\ldots,s_{L,m})
\end{equation}
denote the $m$th row configuration.  The Boltzmann weight between two adjacent
rows can be distributed symmetrically as
\begin{equation}
\label{eq:2d-classical-ising-symmetric-row-weight}
    T_{\bm{s},\bm{s}'}
    =
    \exp\left[
        \frac{K_x}{2}\sum_{j=1}^{L}s_js_{j+1}
        +
        K_y\sum_{j=1}^{L}s_js_j'
        +
        \frac{K_x}{2}\sum_{j=1}^{L}s_j's_{j+1}'
    \right].
\end{equation}
Then
\begin{equation}
\label{eq:2d-classical-ising-z-tr-tm}
    Z_{L,M}(K_x,K_y)=\Tr T^M .
\end{equation}
The symmetrization in \cref{eq:2d-classical-ising-symmetric-row-weight} is not
physically necessary, but it makes $T$ real symmetric and therefore convenient
for spectral analysis.

\section{Operator form of the row-to-row transfer matrix}
\label{sec:2d-classical-ising-row-transfer-operator}

We now rewrite $T$ as an operator acting on a row Hilbert space
\begin{equation}
    \calH_{\rm row}=(\CC^2)^{\otimes L}.
\end{equation}
To remain consistent with the transverse-field Ising convention used earlier in
these notes, we label a classical row configuration by eigenvalues of
$\sigma_j^x$:
\begin{equation}
\label{eq:2d-classical-ising-x-basis}
    \sigma_j^x\ket{s_1,\ldots,s_L}_x
    =
    s_j\ket{s_1,\ldots,s_L}_x .
\end{equation}
In this basis, the intra-row interaction is diagonal:
\begin{equation}
\label{eq:2d-classical-ising-horizontal-operator}
    \exp\left(
        \frac{K_x}{2}\sum_{j=1}^{L}\sigma_j^x\sigma_{j+1}^x
    \right).
\end{equation}
The vertical one-site Boltzmann matrix is
\begin{equation}
\label{eq:2d-classical-ising-vertical-weight}
    W_{ss'}=\ee^{K_yss'}
    =
    \begin{pmatrix}
        \ee^{K_y} & \ee^{-K_y} \\
        \ee^{-K_y} & \ee^{K_y}
    \end{pmatrix}_{s,s'} .
\end{equation}
Since $\sigma^z$ flips the $\sigma^x$ eigenvalue, this matrix can be written as
\begin{equation}
\label{eq:2d-classical-ising-vertical-exp}
    W=A\ee^{\kappa\sigma^z},
    \qquad
    \tanh\kappa=\ee^{-2K_y},
    \qquad
    A=\sqrt{2\sinh(2K_y)} .
\end{equation}
Indeed,
\begin{equation}
    A\cosh\kappa=\ee^{K_y},
    \qquad
    A\sinh\kappa=\ee^{-K_y}.
\end{equation}
The parameter $\kappa$ is the dual coupling associated with $K_y$:
\begin{equation}
\label{eq:2d-classical-ising-kappa-dual}
    \ee^{-2K_y}=\tanh\kappa
    \quad\Longleftrightarrow\quad
    \sinh(2K_y)\sinh(2\kappa)=1 .
\end{equation}
Thus the symmetric row-to-row transfer matrix is
\begin{equation}
\label{eq:2d-classical-ising-symmetric-transfer-operator}
    T
    =
    A^L
    \exp\left(
        \frac{K_x}{2}\sum_{j=1}^{L}\sigma_j^x\sigma_{j+1}^x
    \right)
    \exp\left(
        \kappa\sum_{j=1}^{L}\sigma_j^z
    \right)
    \exp\left(
        \frac{K_x}{2}\sum_{j=1}^{L}\sigma_j^x\sigma_{j+1}^x
    \right).
\end{equation}
This is the precise point at which the two-dimensional classical model starts to
look like a one-dimensional quantum spin chain: $T$ is an exponential product of
operators acting on a one-dimensional row Hilbert space.

\begin{remark}[Relation to the notation $K_\tau$]
In the quantum--classical correspondence chapter, the transfer direction was
written as an imaginary-time direction and the vertical coupling was denoted
$K_\tau$.  In the present chapter we write the two classical directions as
$x$ and $y$.  Thus $K_y$ here plays the role of $K_\tau$ there.
\end{remark}

\section{Kaufman's free-fermion structure}
\label{sec:2d-classical-ising-kaufman-free-fermion}

At first sight, \cref{eq:2d-classical-ising-symmetric-transfer-operator} is not
easier than the original classical problem, because it acts on a
$2^L$-dimensional vector space.  The key observation is that both exponential
factors are Gaussian after the Jordan--Wigner transformation.

Using the Jordan--Wigner convention fixed earlier in these notes, define
Majorana operators
\begin{equation}
\label{eq:2d-classical-ising-majorana-definitions}
    a_{2j-1}
    =
    \left(\prod_{\ell<j}\sigma_\ell^z\right)\sigma_j^x,
    \qquad
    a_{2j}
    =
    \left(\prod_{\ell<j}\sigma_\ell^z\right)\sigma_j^y .
\end{equation}
They satisfy
\begin{equation}
\label{eq:2d-classical-ising-majorana-algebra}
    a_m^\dagger=a_m,
    \qquad
    \{a_m,a_n\}=2\delta_{mn} .
\end{equation}
In the bulk one has the local identities
\begin{equation}
\label{eq:2d-classical-ising-majorana-spin-identities}
    \sigma_j^z=-\ii a_{2j-1}a_{2j},
    \qquad
    \sigma_j^x\sigma_{j+1}^x=-\ii a_{2j}a_{2j+1} .
\end{equation}
Therefore the two building blocks of $T$ become
\begin{align}
\label{eq:2d-classical-ising-vertical-majorana-gaussian}
    \exp\left(\kappa\sum_j\sigma_j^z\right)
    &=
    \exp\left(-\ii\kappa\sum_j a_{2j-1}a_{2j}\right),
    \\
\label{eq:2d-classical-ising-horizontal-majorana-gaussian}
    \exp\left(\frac{K_x}{2}\sum_j\sigma_j^x\sigma_{j+1}^x\right)
    &=
    \exp\left(-\frac{\ii K_x}{2}\sum_j a_{2j}a_{2j+1}\right),
\end{align}
up to the usual boundary-sector term for a periodic spin chain.

Equations \cref{eq:2d-classical-ising-vertical-majorana-gaussian,eq:2d-classical-ising-horizontal-majorana-gaussian}
show that $T$ is a product of exponentials of Majorana bilinears.  Such operators
are fermionic Gaussian operators.  Their adjoint action maps Majoranas linearly
to Majoranas:
\begin{equation}
\label{eq:2d-classical-ising-gaussian-adjoint-action}
    T^{-1}a_mT
    =
    \sum_{n=1}^{2L}R_{mn}a_n,
\end{equation}
where $R$ is a $2L\times2L$ matrix preserving the Majorana anticommutation
relations.  The exponentially large transfer-matrix problem is therefore reduced
to the diagonalization of a one-particle Majorana transformation.

\begin{remark}[Boundary sectors]
For periodic spin boundary conditions, the boundary term in
$\sigma_L^x\sigma_1^x$ contains the fermion-parity operator
$P_F=\prod_{j=1}^{L}\sigma_j^z$.  Consequently, the exact finite-$L$ partition
function is obtained by combining parity-projected fermion traces with periodic
and antiperiodic momenta, in the same spirit as the TFIM parity-sector
bookkeeping.  In the thermodynamic limit these sector distinctions affect only
subleading finite-size terms, not the bulk free energy or the critical curve.
\end{remark}

Because the model is translation invariant, the one-particle Majorana problem
block-diagonalizes in momentum.  For each momentum $k$, the relevant
$2\times2$ block has eigenvalues $\ee^{\pm\varepsilon(k)}$, where
\begin{equation}
\label{eq:2d-classical-ising-onsager-dispersion}
    \cosh\varepsilon(k)
    =
    \cosh(2\kappa)\cosh(2K_x)
    -
    \sinh(2\kappa)\sinh(2K_x)\cos k .
\end{equation}
The quantity $\varepsilon(k)$ is the transfer-matrix analogue of a quasiparticle
energy.  More precisely, it is the single-particle gap in the logarithm of the
transfer matrix.  Its minimum occurs at $k=0$ for ferromagnetic couplings, and
\begin{equation}
\label{eq:2d-classical-ising-transfer-gap}
    \varepsilon(0)
    =
    2\abs{\kappa-K_x} .
\end{equation}
Thus the transfer-matrix gap closes when
\begin{equation}
\label{eq:2d-classical-ising-kappa-equals-kx}
    \kappa=K_x .
\end{equation}
Using \cref{eq:2d-classical-ising-kappa-dual}, this is equivalent to
\begin{equation}
\label{eq:2d-classical-ising-onsager-critical}
    \boxed{\sinh(2K_x)\sinh(2K_y)=1.}
\end{equation}
This is the Onsager critical line for the anisotropic two-dimensional Ising
model.

\section{Largest eigenvalue and Onsager free energy}
\label{sec:2d-classical-ising-onsager-free-energy}

Let
\begin{equation}
\label{eq:2d-classical-ising-phi-definition}
    \Phi(K_x,K_y)
    =
    \lim_{L,M\to\infty}\frac{1}{LM}\log Z_{L,M}(K_x,K_y)
\end{equation}
be the dimensionless free-energy density with the sign convention appropriate to
$\log Z$.  The physical free energy per site is
\begin{equation}
    f=-\beta^{-1}\Phi
\end{equation}
when $K_\mu=\beta J_\mu$.  Since $Z=\Tr T^M$, $\Phi$ is the logarithm of the
largest transfer-matrix eigenvalue per row site.

From the Gaussian diagonalization one obtains
\begin{equation}
\label{eq:2d-classical-ising-onsager-free-energy-kappa}
    \Phi(K_x,K_y)
    =
    \frac12\log\!\bigl(2\sinh(2K_y)\bigr)
    +
    \frac12\int_{-\pi}^{\pi}\frac{\dd k}{2\pi}\,\varepsilon(k),
\end{equation}
where $\varepsilon(k)$ is defined by
\cref{eq:2d-classical-ising-onsager-dispersion}.  Equivalently,
\begin{align}
\label{eq:2d-classical-ising-onsager-free-energy-standard}
    \Phi(K_x,K_y)
    &=
    \frac12\log\!\bigl(2\sinh(2K_y)\bigr)
    +
    \frac{1}{2\pi}
    \int_0^\pi \dd k\,\operatorname{arcosh} X(k),
    \\
    X(k)
    &=
    \cosh(2\kappa)\cosh(2K_x)
    -
    \sinh(2\kappa)\sinh(2K_x)\cos k .
\end{align}
The appearance of $\operatorname{arcosh}$ is the transfer-matrix version of
``taking a quasiparticle energy from a one-particle evolution eigenvalue.''

A useful check is the decoupled-row limit $K_x=0$.  Then
\begin{equation}
    \varepsilon(k)=2\kappa,
\end{equation}
and \cref{eq:2d-classical-ising-onsager-free-energy-kappa} gives
\begin{equation}
    \Phi(0,K_y)
    =
    \frac12\log\!\bigl(2\sinh(2K_y)\bigr)+\kappa
    =
    \log(2\cosh K_y),
\end{equation}
which is exactly the free-energy density of independent vertical Ising chains
viewed through the chosen transfer direction.

For the isotropic model $K_x=K_y=K$, the critical point follows from
\begin{equation}
    \sinh^2(2K_c)=1,
\end{equation}
and therefore
\begin{equation}
\label{eq:2d-classical-ising-isotropic-kc}
    \boxed{K_c=\frac12\log(1+\sqrt2).}
\end{equation}
If $K=\beta J$, this gives
\begin{equation}
\label{eq:2d-classical-ising-isotropic-tc}
    k_B T_c
    =
    \frac{2J}{\log(1+\sqrt2)} .
\end{equation}

\section{Duality interpretation of the critical curve}
\label{sec:2d-classical-ising-duality-interpretation}

The same critical condition follows from Kramers--Wannier duality.  The
high-temperature expansion writes each bond factor as
\begin{equation}
\label{eq:2d-classical-ising-high-temp-bond}
    \ee^{Ks_rs_{r'}}
    =
    \cosh K\left(1+s_rs_{r'}\tanh K\right).
\end{equation}
After summing over spins, only closed even-degree bond configurations survive.
The low-temperature expansion of a dual Ising model is instead a sum over domain
walls.  Matching high-temperature loops of the original lattice with
domain-wall loops of the dual lattice gives
\begin{equation}
\label{eq:2d-classical-ising-dual-couplings}
    \ee^{-2K_x^*}=\tanh K_y,
    \qquad
    \ee^{-2K_y^*}=\tanh K_x .
\end{equation}
Equivalently,
\begin{equation}
\label{eq:2d-classical-ising-dual-couplings-sinh}
    \sinh(2K_x^*)\sinh(2K_y)=1,
    \qquad
    \sinh(2K_y^*)\sinh(2K_x)=1 .
\end{equation}
A self-dual point satisfies
\begin{equation}
    K_x^*=K_x,
    \qquad
    K_y^*=K_y,
\end{equation}
and hence obeys
\begin{equation}
    \sinh(2K_x)\sinh(2K_y)=1 .
\end{equation}
Duality alone identifies the possible transition curve if there is a unique
ferromagnetic critical line.  The transfer-matrix solution gives the sharper
statement: the one-particle transfer gap actually closes on this curve.

\section{Order parameter and critical behavior}
\label{sec:2d-classical-ising-order-parameter-critical-behavior}

The finite-volume partition function at zero field has a global $\ZZ_2$ symmetry
$s_{j,m}\mapsto -s_{j,m}$, so the finite-volume magnetization vanishes exactly.
The spontaneous magnetization is defined by taking the thermodynamic limit before
removing an infinitesimal symmetry-breaking field:
\begin{equation}
\label{eq:2d-classical-ising-spontaneous-magnetization-definition}
    M
    =
    \lim_{H\to0^+}\lim_{L,M\to\infty}
    \frac{1}{LM}\sum_{j,m}\langle s_{j,m}\rangle_H .
\end{equation}
Yang's exact result for the anisotropic square-lattice model is
\begin{equation}
\label{eq:2d-classical-ising-yang-magnetization}
    M(K_x,K_y)
    =
    \begin{cases}
    \displaystyle
    \left[
        1-\frac{1}{\sinh^2(2K_x)\sinh^2(2K_y)}
    \right]^{1/8},
    &
    \sinh(2K_x)\sinh(2K_y)>1,\\[12pt]
    0,
    &
    \sinh(2K_x)\sinh(2K_y)\leq1 .
    \end{cases}
\end{equation}
For the isotropic model this reduces to
\begin{equation}
\label{eq:2d-classical-ising-isotropic-magnetization}
    M(K)
    =
    \left[1-\sinh^{-4}(2K)\right]^{1/8},
    \qquad
    K>K_c .
\end{equation}
Thus the magnetization exponent is
\begin{equation}
\label{eq:2d-classical-ising-beta-exponent}
    M\sim (T_c-T)^{1/8},
    \qquad
    \beta_{\rm Ising}=\frac18 .
\end{equation}

The transfer-matrix dispersion also reveals the divergence of the correlation
length.  Along the transfer direction, the inverse correlation length is given by
the relevant transfer gap.  Near the ferromagnetic critical line,
\begin{equation}
\label{eq:2d-classical-ising-correlation-length-gap}
    \xi_y^{-1}\sim\varepsilon(0)=2\abs{\kappa-K_x} .
\end{equation}
Because $\kappa-K_x$ is linear in the distance from the critical curve, the
correlation-length exponent is
\begin{equation}
\label{eq:2d-classical-ising-nu-exponent}
    \xi\sim\abs{T-T_c}^{-1},
    \qquad
    \nu=1 .
\end{equation}
The full critical behavior of the two-dimensional Ising universality class is
summarized by
\begin{equation}
\label{eq:2d-classical-ising-critical-exponents}
\begin{gathered}
    \alpha=0\quad(\text{logarithmic specific-heat divergence}),
    \qquad
    \beta=\frac18,
    \qquad
    \gamma=\frac74,
    \\
    \nu=1,
    \qquad
    \eta=\frac14,
    \qquad
    \delta=15 .
\end{gathered}
\end{equation}
In these notes, the most important point is not the list of exponents itself,
but their origin: the classical thermal transition is controlled by a massless
free-fermion transfer-matrix theory at criticality.

\section{Connection to the one-dimensional TFIM}
\label{sec:2d-classical-ising-connection-to-tfim}

The transfer matrix in
\cref{eq:2d-classical-ising-symmetric-transfer-operator} becomes the imaginary-
time evolution operator of the one-dimensional TFIM in a strongly anisotropic
limit.  Set
\begin{equation}
\label{eq:2d-classical-ising-tfim-scaling}
    K_x=\delta\tau J,
    \qquad
    \kappa=\delta\tau h,
    \qquad
    \tanh(\delta\tau h)=\ee^{-2K_y} .
\end{equation}
Then, using the symmetric Trotter formula,
\begin{equation}
\label{eq:2d-classical-ising-tfim-transfer-limit}
    T
    =
    A^L\exp\left[-\delta\tau H_{\rm TFIM}+O(\delta\tau^3)\right],
\end{equation}
with
\begin{equation}
\label{eq:2d-classical-ising-tfim-hamiltonian}
    H_{\rm TFIM}
    =
    -J\sum_{j=1}^{L}\sigma_j^x\sigma_{j+1}^x
    -h\sum_{j=1}^{L}\sigma_j^z .
\end{equation}
The Onsager condition becomes
\begin{equation}
    \sinh(2\delta\tau J)\sinh(2K_y)=1.
\end{equation}
Using
\begin{equation}
    \ee^{-2K_y}=\tanh(\delta\tau h)
    =
    \delta\tau h+O(\delta\tau^3),
\end{equation}
one finds
\begin{equation}
\label{eq:2d-classical-ising-tfim-critical-limit}
    h=J .
\end{equation}
Thus the thermal critical point of the anisotropic two-dimensional classical
Ising model maps to the zero-temperature quantum critical point of the
one-dimensional transverse-field Ising model.

\begin{center}
\begin{tabular}{@{}L{0.31\linewidth}L{0.31\linewidth}L{0.28\linewidth}@{}}
\toprule
Classical 2D Ising object & Transfer-matrix object & 1D quantum interpretation \\
\midrule
Row configuration $\bm{s}_m$ & Basis vector in $\calH_{\rm row}$ & Spin-chain configuration \\
Row-to-row transfer matrix $T$ & Positive operator on $\calH_{\rm row}$ & Imaginary-time step $\ee^{-\delta\tau H}$ \\
Largest eigenvalue of $T$ & Ground-state sector of $-\log T$ & Free-energy density / ground-state energy density \\
Transfer gap $\varepsilon(0)$ & Mass gap of $-\log T$ & Quantum excitation gap \\
Onsager curve & Gapless transfer matrix & TFIM critical point $h=J$ \\
\bottomrule
\end{tabular}
\end{center}

\section{Summary}
\label{sec:2d-classical-ising-summary}

The two-dimensional zero-field Ising model is exactly solvable because its
row-to-row transfer matrix is a fermionic Gaussian operator.  The essential chain
of ideas is
\begin{equation}
\label{eq:2d-classical-ising-solution-chain}
\begin{gathered}
    \text{2D classical partition function}
    \longrightarrow
    \text{row transfer matrix}
    \longrightarrow
    \text{Majorana Gaussian operator}
    \\
    \longrightarrow
    \text{one-particle dispersion }\varepsilon(k).
\end{gathered}
\end{equation}
The Onsager critical point is the point where the transfer-matrix quasiparticle
gap closes:
\begin{equation}
    \varepsilon(0)=0
    \quad\Longleftrightarrow\quad
    \kappa=K_x
    \quad\Longleftrightarrow\quad
    \sinh(2K_x)\sinh(2K_y)=1 .
\end{equation}
This is why the model belongs naturally in a lecture note on solvable quantum
many-body systems: its exact solution is not merely a classical combinatorial
miracle, but a free-fermion solution in transfer-matrix language.


\chapter{From 2D Classical Ising to 1D Quantum TFIM}
\label{chap:2d-ising-to-1d-tfim}

The previous chapter solved the square-lattice Ising model by rewriting its
row-to-row transfer matrix as a free-fermion operator.  We now sharpen one
particular consequence of that construction: in a strongly anisotropic limit,
the two-dimensional classical Ising transfer matrix becomes the short
imaginary-time evolution operator of the one-dimensional transverse-field Ising
model (TFIM).  This is the cleanest example of the general principle
\begin{equation}
\label{eq:ising-tfim-main-correspondence-schematic}
    d\text{-dimensional quantum system}
    \quad\longleftrightarrow\quad
    (d+1)\text{-dimensional classical statistical system}.
\end{equation}
For the present case, the correspondence is
\begin{equation}
    \text{1D quantum TFIM}
    \quad\longleftrightarrow\quad
    \text{anisotropic 2D classical Ising model}.
\end{equation}
The additional classical direction is Euclidean imaginary time.

The discussion in this chapter has three complementary goals.  First, we derive
precisely how the row-to-row transfer matrix becomes
$\ee^{-\delta\tau H_{\rm TFIM}}$.  Second, we give the explicit dictionary between
classical couplings and quantum Hamiltonian parameters in the conventions of
these notes.  Third, we explain the physical interpretation of the map: spatial
Ising bonds become the quantum Ising interaction, while rare domain walls in the
transfer direction become transverse-field spin flips.

\section{Row transfer matrix and basis convention}
\label{sec:ising-tfim-row-transfer-review}

Consider the anisotropic zero-field classical Ising model on an $L\times M$
square lattice,
\begin{equation}
\label{eq:ising-tfim-classical-partition-function}
    Z_{L,M}(K_x,K_y)
    =
    \sum_{\{s_{j,m}=\pm1\}}
    \exp\left[
        K_x\sum_{j,m}s_{j,m}s_{j+1,m}
        +
        K_y\sum_{j,m}s_{j,m}s_{j,m+1}
    \right],
\end{equation}
with periodic boundary conditions in both directions.  The index $j$ labels the
one-dimensional spatial direction, while $m$ labels the transfer direction.  In
the quantum interpretation, $m$ will become a discretized imaginary-time index.

Following the convention used throughout these notes, a classical row
configuration is identified with a simultaneous eigenstate of the Pauli
operators $\sigma_j^x$:
\begin{equation}
\label{eq:ising-tfim-x-basis-row-state}
    \sigma_j^x\ket{s_1,\ldots,s_L}_x
    =
    s_j\ket{s_1,\ldots,s_L}_x,
    \qquad
    s_j=\pm1 .
\end{equation}
This convention is chosen so that the resulting quantum Hamiltonian is written
in the same form as the TFIM chapter,
\begin{equation}
\label{eq:ising-tfim-quantum-hamiltonian-target}
    H_{\rm TFIM}
    =
    -J\sum_{j=1}^{L}\sigma_j^x\sigma_{j+1}^x
    -h\sum_{j=1}^{L}\sigma_j^z .
\end{equation}
Many books use the rotated convention
$-J\sum_j\sigma_j^z\sigma_{j+1}^z-\Gamma\sum_j\sigma_j^x$.  The two are related
by a global spin rotation.  In these notes the $x$-basis convention is more
natural because the Jordan--Wigner occupation convention uses
$n_j=(1-\sigma_j^z)/2$.

The symmetrized row-to-row transfer matrix is
\begin{equation}
\label{eq:ising-tfim-symmetric-transfer-matrix}
    T
    =
    A_y^L
    \exp\left(
        \frac{K_x}{2}\sum_{j=1}^{L}\sigma_j^x\sigma_{j+1}^x
    \right)
    \exp\left(
        \kappa_y\sum_{j=1}^{L}\sigma_j^z
    \right)
    \exp\left(
        \frac{K_x}{2}\sum_{j=1}^{L}\sigma_j^x\sigma_{j+1}^x
    \right),
\end{equation}
where
\begin{equation}
\label{eq:ising-tfim-kappa-definition}
    \tanh\kappa_y=\ee^{-2K_y},
    \qquad
    A_y=\sqrt{2\sinh(2K_y)} .
\end{equation}
Equivalently,
\begin{equation}
\label{eq:ising-tfim-kappa-log-coth}
    \kappa_y
    =
    \operatorname{arctanh}(\ee^{-2K_y})
    =
    \frac12\log\coth K_y .
\end{equation}
The normalization $A_y^L$ is important for the absolute classical free energy,
but it becomes an additive constant in the quantum Hamiltonian and can be ignored
when discussing normalized correlators and excitation gaps.

\begin{remark}[Why $\sigma^z$ appears in the vertical transfer operator]
In the $\sigma^x$ basis, $\sigma^z$ flips the sign of $s_j$.  Therefore the
one-site matrix
\begin{equation}
    W_{ss'}=\ee^{K_yss'}
\end{equation}
can be represented as $W=A_y\ee^{\kappa_y\sigma^z}$.  A temporal mismatch
$s\neq s'$ is a domain wall in the transfer direction; in the quantum model it is
generated by the transverse-field operator $\sigma^z$.
\end{remark}

\section{Anisotropic Hamiltonian limit}
\label{sec:ising-tfim-anisotropic-limit}

The classical transfer matrix defines an exact operator
\begin{equation}
\label{eq:ising-tfim-exact-effective-hamiltonian}
    H_{\rm eff}
    =
    -\frac{1}{a_\tau}\log T,
\end{equation}
where $a_\tau$ is an arbitrary microscopic time spacing.  For generic finite
$K_x$ and $K_y$, however, $H_{\rm eff}$ is not simply the local nearest-neighbour
TFIM Hamiltonian.  The logarithm of a product of noncommuting many-body
operators generally contains longer-range and higher-body terms.

The local TFIM is obtained in the \emph{strongly anisotropic limit}.  Introduce a
small imaginary-time lattice spacing $\delta\tau$ and set
\begin{equation}
\label{eq:ising-tfim-anisotropic-scaling}
    K_x=\delta\tau J,
    \qquad
    \kappa_y=\delta\tau h,
    \qquad
    \tanh(\delta\tau h)=\ee^{-2K_y} .
\end{equation}
Thus
\begin{equation}
\label{eq:ising-tfim-strong-anisotropy}
    K_x\to0,
    \qquad
    K_y\to\infty,
    \qquad
    K_x\ee^{2K_y}\;\text{fixed}.
\end{equation}
Indeed, for small $\delta\tau$,
\begin{equation}
\label{eq:ising-tfim-ky-asymptotic}
    \ee^{-2K_y}
    =
    \tanh(\delta\tau h)
    =
    \delta\tau h+O(\delta\tau^3),
    \qquad
    K_y
    =
    -\frac12\log(\delta\tau h)+O(\delta\tau^2).
\end{equation}
The horizontal classical coupling becomes weak, while the vertical coupling
becomes strong.  This is exactly what one expects from an imaginary-time path
integral: neighbouring time slices are almost identical, and spin flips in time
are rare events with probability proportional to $\delta\tau h$.

Define
\begin{equation}
\label{eq:ising-tfim-ax-az-definitions}
    X
    =
    J\sum_{j=1}^{L}\sigma_j^x\sigma_{j+1}^x,
    \qquad
    Z
    =
    h\sum_{j=1}^{L}\sigma_j^z .
\end{equation}
Then \cref{eq:ising-tfim-symmetric-transfer-matrix} becomes
\begin{equation}
\label{eq:ising-tfim-symmetric-strang-product}
    T
    =
    A_y^L
    \ee^{\frac{\delta\tau}{2}X}
    \ee^{\delta\tau Z}
    \ee^{\frac{\delta\tau}{2}X} .
\end{equation}
The symmetric Trotter formula gives
\begin{equation}
\label{eq:ising-tfim-strang-expansion}
    \ee^{\frac{\delta\tau}{2}X}
    \ee^{\delta\tau Z}
    \ee^{\frac{\delta\tau}{2}X}
    =
    \exp\left[
        \delta\tau(X+Z)+O(\delta\tau^3)
    \right].
\end{equation}
Since $X+Z=-H_{\rm TFIM}$ in the convention of
\cref{eq:ising-tfim-quantum-hamiltonian-target}, we obtain
\begin{equation}
\label{eq:ising-tfim-transfer-hamiltonian-limit}
    T
    =
    A_y^L
    \exp\left[
        -\delta\tau H_{\rm TFIM}
        +O(\delta\tau^3)
    \right].
\end{equation}
Equivalently,
\begin{equation}
\label{eq:ising-tfim-effective-hamiltonian-expansion}
    -\frac{1}{\delta\tau}\log T
    =
    H_{\rm TFIM}
    -\frac{L}{\delta\tau}\log A_y\,\id
    +O(\delta\tau^2).
\end{equation}
The second term is an additive constant.  It affects the absolute free energy
but not the eigenvectors, normalized expectation values, or excitation gaps.

\begin{proposition}[Hamiltonian limit of the anisotropic Ising transfer matrix]
\label{prop:ising-tfim-hamiltonian-limit}
With the scaling
\begin{equation}
    K_x=\delta\tau J,
    \qquad
    \tanh(\delta\tau h)=\ee^{-2K_y},
\end{equation}
the row-to-row transfer matrix of the anisotropic two-dimensional classical Ising
model satisfies
\begin{equation}
    T
    =
    A_y^L\ee^{-\delta\tau H_{\rm TFIM}+O(\delta\tau^3)},
\end{equation}
where
\begin{equation}
    H_{\rm TFIM}
    =
    -J\sum_j\sigma_j^x\sigma_{j+1}^x
    -h\sum_j\sigma_j^z .
\end{equation}
Therefore the local quantum Hamiltonian is recovered from the logarithm of the
classical transfer matrix up to an additive constant and errors of order
$O(\delta\tau^2)$ in the Hamiltonian density.
\end{proposition}

\section{Partition functions and the continuum-time limit}
\label{sec:ising-tfim-partition-functions}

The classical partition function is
\begin{equation}
\label{eq:ising-tfim-z-tr-tm}
    Z_{L,M}(K_x,K_y)=\Tr T^M .
\end{equation}
Using \cref{eq:ising-tfim-transfer-hamiltonian-limit}, one obtains
\begin{equation}
\label{eq:ising-tfim-z-to-zq-finite-dtau}
    Z_{L,M}(K_x,K_y)
    =
    A_y^{LM}
    \Tr\exp\left[
        -M\delta\tau H_{\rm TFIM}
        +O(M\delta\tau^3)
    \right].
\end{equation}
If
\begin{equation}
\label{eq:ising-tfim-beta-M-dtau}
    \beta=M\delta\tau
\end{equation}
is held fixed while $M\to\infty$ and $\delta\tau\to0$, then
\begin{equation}
\label{eq:ising-tfim-quantum-classical-partition-function}
    Z_{L,M}(K_x,K_y)
    =
    A_y^{LM}
    Z_{\rm TFIM}(\beta)
    \left[1+\text{controlled Trotter corrections}\right],
\end{equation}
where
\begin{equation}
\label{eq:ising-tfim-quantum-partition-function}
    Z_{\rm TFIM}(\beta)=\Tr\ee^{-\beta H_{\rm TFIM}} .
\end{equation}
For local thermodynamic quantities, the logarithm of the correction is of order
$O(\beta L\delta\tau^2)$ for the symmetric decomposition.  Thus the limit is
controlled at fixed system size and fixed inverse temperature, and remains
controlled for thermodynamic densities when the usual local Trotter estimates are
used.

The normalization factor $A_y^{LM}$ gives an additive contribution to the
classical free energy:
\begin{equation}
\label{eq:ising-tfim-free-energy-shift}
    -\log Z_{L,M}
    =
    -LM\log A_y
    -\log Z_{\rm TFIM}(\beta)
    +\cdots .
\end{equation}
This term has no effect on normalized quantum correlation functions.  It is,
however, part of the absolute classical free energy.  One should therefore keep
it when comparing the full classical free-energy density, but drop it when
extracting the quantum Hamiltonian spectrum.

\section{Dictionary of couplings}
\label{sec:ising-tfim-coupling-dictionary}

There are two equivalent ways to present the coupling dictionary.

\subsection{From quantum TFIM to classical Ising}
\label{subsec:ising-tfim-quantum-to-classical}

Given a TFIM Hamiltonian
\begin{equation}
    H_{\rm TFIM}
    =
    -J\sum_j\sigma_j^x\sigma_{j+1}^x
    -h\sum_j\sigma_j^z,
\end{equation}
and a small imaginary-time spacing $\delta\tau$, the corresponding classical
couplings are
\begin{equation}
\label{eq:ising-tfim-quantum-to-classical-dictionary}
    \boxed{
    K_x=\delta\tau J,
    \qquad
    K_y=-\frac12\log\tanh(\delta\tau h)
    }.
\end{equation}
Equivalently,
\begin{equation}
\label{eq:ising-tfim-quantum-to-classical-kappa}
    \kappa_y=\delta\tau h,
    \qquad
    \tanh\kappa_y=\ee^{-2K_y} .
\end{equation}
The classical lattice has $M=\beta/\delta\tau$ rows in the imaginary-time
direction when representing a quantum thermal partition function at inverse
temperature $\beta$.

\subsection{From classical Ising to quantum TFIM}
\label{subsec:ising-tfim-classical-to-quantum}

Given classical couplings $K_x$ and $K_y$, define the dual transfer-direction
coupling
\begin{equation}
\label{eq:ising-tfim-classical-to-quantum-kappa}
    \kappa_y
    =
    \frac12\log\coth K_y .
\end{equation}
If the transfer matrix is close to the identity in the sense that
$K_x\ll1$ and $\kappa_y\ll1$, then it is naturally interpreted as a short-time
quantum evolution operator.  Choosing an imaginary-time spacing $\delta\tau$,
\begin{equation}
\label{eq:ising-tfim-classical-to-quantum-dictionary}
    \boxed{
    J=\frac{K_x}{\delta\tau},
    \qquad
    h=\frac{\kappa_y}{\delta\tau}
    =
    \frac{1}{2\delta\tau}\log\coth K_y
    }.
\end{equation}
Only the dimensionless ratio
\begin{equation}
\label{eq:ising-tfim-ratio-dictionary}
    \boxed{
    \frac{h}{J}
    =
    \frac{\kappa_y}{K_x}
    =
    \frac{1}{2K_x}\log\coth K_y
    }
\end{equation}
is independent of the arbitrary choice of $\delta\tau$.  The overall scale of
$J$ and $h$ is fixed only after one decides how the classical transfer step is to
be measured in imaginary-time units.

\begin{center}
\begin{tabularx}{0.95\linewidth}{@{}L{0.30\linewidth}Y Y@{}}
\toprule
Object & Classical Ising transfer-matrix language & Quantum TFIM language \\
\midrule
Spatial coordinate & Column index $j$ & Spin-chain site $j$ \\
Transfer coordinate & Row index $m$ & Imaginary time $\tau=m\delta\tau$ \\
Row spin $s_{j,m}$ & Eigenvalue of $\sigma_j^x$ & Order-parameter basis variable \\
Horizontal coupling & $K_x=\delta\tau J$ & Ising exchange $J$ \\
Vertical coupling & $\tanh\kappa_y=\ee^{-2K_y}$ & Transverse field $h=\kappa_y/\delta\tau$ \\
Transfer matrix & $T$ & $A_y^L\ee^{-\delta\tau H_{\rm TFIM}}$ \\
Vertical domain wall & $s_{j,m}\neq s_{j,m+1}$ & Spin flip generated by $\sigma_j^z$ \\
\bottomrule
\end{tabularx}
\end{center}

\section{Critical line and self-duality}
\label{sec:ising-tfim-critical-line-self-duality}

The anisotropic square-lattice Ising model is critical on the Onsager curve
\begin{equation}
\label{eq:ising-tfim-onsager-critical-line}
    \sinh(2K_x)\sinh(2K_y)=1 .
\end{equation}
Using
\begin{equation}
\label{eq:ising-tfim-kappa-duality-identity}
    \sinh(2K_y)\sinh(2\kappa_y)=1,
\end{equation}
this condition is equivalent, for positive couplings, to
\begin{equation}
\label{eq:ising-tfim-critical-kx-kappa}
    \kappa_y=K_x .
\end{equation}
Under the Hamiltonian-limit dictionary $K_x=\delta\tau J$ and
$\kappa_y=\delta\tau h$, this becomes
\begin{equation}
\label{eq:ising-tfim-critical-h-equals-j}
    h=J .
\end{equation}
Thus the thermal critical line of the anisotropic two-dimensional classical Ising
model reduces to the zero-temperature quantum critical point of the
one-dimensional TFIM.

This result is also the transfer-matrix form of Kramers--Wannier duality.  The
coupling $\kappa_y$ is the dual of $K_y$, and the critical line is the self-dual
condition
\begin{equation}
\label{eq:ising-tfim-self-dual-condition}
    K_x=\kappa_y .
\end{equation}
The TFIM inherits this as the quantum self-dual condition $J=h$.  On one side of
this point, $J>h$, the quantum ground state is ferromagnetically ordered in the
$x$ direction.  On the other side, $h>J$, the ground state is paramagnetic.  The
classical high-temperature and low-temperature phases become the two quantum
phases separated by a zero-temperature transition.

\begin{remark}[Thermal versus quantum criticality]
The two-dimensional classical transition is driven by thermal fluctuations.  The
one-dimensional TFIM transition at $T=0$ is driven by quantum fluctuations from
the transverse field.  The transfer-matrix correspondence identifies their
Euclidean statistical weights, and therefore their universal critical theory, not
their microscopic dynamical origin.
\end{remark}

\section{Transfer-matrix spectrum and the TFIM dispersion}
\label{sec:ising-tfim-spectrum-dispersion}

The correspondence is sharper than a relation between critical points.  It maps
the full transfer-matrix quasiparticle spectrum to the quantum excitation
spectrum.

In the free-fermion solution of the two-dimensional Ising transfer matrix, the
one-particle transfer dispersion may be written as
\begin{equation}
\label{eq:ising-tfim-onsager-transfer-dispersion}
    \cosh\varepsilon_{\rm tr}(k)
    =
    \cosh(2K_x)\cosh(2\kappa_y)
    -
    \sinh(2K_x)\sinh(2\kappa_y)\cos k .
\end{equation}
The transfer gap closes at $k=0$ when $K_x=\kappa_y$, reproducing the critical
condition above.

Now impose the Hamiltonian-limit scaling
\begin{equation}
    K_x=\delta\tau J,
    \qquad
    \kappa_y=\delta\tau h .
\end{equation}
For small $\delta\tau$,
\begin{align}
\label{eq:ising-tfim-dispersion-expansion}
    \cosh\varepsilon_{\rm tr}(k)
    & =
    1
    +2\delta\tau^2
    \left(J^2+h^2-2Jh\cos k\right)
    +O(\delta\tau^4).
\end{align}
Therefore
\begin{equation}
\label{eq:ising-tfim-transfer-to-quantum-dispersion}
    \varepsilon_{\rm tr}(k)
    =
    \delta\tau\,E_{\rm TFIM}(k)+O(\delta\tau^3),
\end{equation}
where
\begin{equation}
\label{eq:ising-tfim-quantum-dispersion}
    E_{\rm TFIM}(k)
    =
    2\sqrt{h^2+J^2-2Jh\cos k}
    =
    2\sqrt{(h-J\cos k)^2+J^2\sin^2 k}
\end{equation}
is exactly the Bogoliubov quasiparticle dispersion of the TFIM in the convention
of \cref{eq:ising-tfim-quantum-hamiltonian-target}.

Thus transfer-matrix eigenvalue ratios become quantum energy gaps:
\begin{equation}
\label{eq:ising-tfim-eigenvalue-gap-map}
    -\log\frac{\lambda_n}{\lambda_0}
    =
    \delta\tau\,(E_n-E_0)+O(\delta\tau^3).
\end{equation}
At the ferromagnetic critical point $h=J>0$,
\begin{equation}
\label{eq:ising-tfim-critical-linear-dispersion}
    E_{\rm TFIM}(k)
    =
    4J\abs{\sin(k/2)}
    \simeq
    2J\abs{k}
    \qquad
    (k\to0).
\end{equation}
The continuum limit is therefore a massless Majorana theory.  Away from the
critical point, the quantum gap is
\begin{equation}
\label{eq:ising-tfim-quantum-gap}
    \Delta
    =
    E_{\rm TFIM}(0)
    =
    2\abs{h-J}
\end{equation}
for $J,h>0$.

\section{Operator dictionary and Euclidean correlators}
\label{sec:ising-tfim-operator-dictionary}

The transfer-matrix map also identifies correlation functions.  A classical spin
on row $m$ becomes the quantum operator $\sigma_j^x$ inserted at imaginary time
\begin{equation}
    \tau=m\delta\tau .
\end{equation}
Thus, in the Hamiltonian limit,
\begin{equation}
\label{eq:ising-tfim-spin-correlator-map}
    \langle s_{j,m}s_{j',0}\rangle_{\rm cl}
    \longrightarrow
    \frac{1}{Z_{\rm TFIM}}
    \Tr\left[
        \ee^{-\beta H_{\rm TFIM}}
        \sigma_j^x(\tau)\sigma_{j'}^x(0)
    \right],
\end{equation}
where
\begin{equation}
\label{eq:ising-tfim-imaginary-time-heisenberg}
    \sigma_j^x(\tau)
    =
    \ee^{\tau H_{\rm TFIM}}\sigma_j^x\ee^{-\tau H_{\rm TFIM}} .
\end{equation}
At zero quantum temperature, $\beta\to\infty$, the trace becomes a ground-state
expectation value:
\begin{equation}
\label{eq:ising-tfim-zero-temperature-correlator}
    \langle s_{j,m}s_{j',0}\rangle_{\rm cl}
    \longrightarrow
    \bra{0}\sigma_j^x(\tau)\sigma_{j'}^x(0)\ket{0} .
\end{equation}

Spatial classical bonds map to the Ising energy operator,
\begin{equation}
\label{eq:ising-tfim-spatial-bond-map}
    s_{j,m}s_{j+1,m}
    \quad\longleftrightarrow\quad
    \sigma_j^x\sigma_{j+1}^x .
\end{equation}
Temporal bonds are represented by insertions obtained by differentiating the
vertical transfer matrix with respect to $K_y$.  Since
\begin{equation}
    W=A_y\ee^{\kappa_y\sigma^z},
\end{equation}
vertical-bond observables contain an identity contribution and a contribution
proportional to $\sigma^z$.  In the quantum Hamiltonian limit, the nontrivial
operator associated with temporal fluctuations is therefore the transverse-field
operator
\begin{equation}
\label{eq:ising-tfim-temporal-bond-map}
    \text{temporal domain-wall fluctuations}
    \quad\longleftrightarrow\quad
    \sigma_j^z .
\end{equation}
This is the operator-level reason why the vertical classical coupling becomes
the transverse field rather than another Ising exchange term.

\section{Physical interpretation of the anisotropic limit}
\label{sec:ising-tfim-physical-interpretation}

The anisotropic limit has a simple worldline interpretation.  In the $\sigma^x$
basis, a quantum spin configuration at imaginary time $\tau$ is a row
configuration
\begin{equation}
    \bm{s}(\tau)=(s_1(\tau),\ldots,s_L(\tau)).
\end{equation}
The diagonal part of the Hamiltonian,
\begin{equation}
    -J\sum_j\sigma_j^x\sigma_{j+1}^x,
\end{equation}
assigns an energy to each row configuration.  Over a short imaginary-time
interval $\delta\tau$, it contributes the Boltzmann factor
\begin{equation}
\label{eq:ising-tfim-diagonal-row-weight}
    \exp\left[
        \delta\tau J\sum_j s_j(\tau)s_{j+1}(\tau)
    \right],
\end{equation}
which is precisely the horizontal classical Ising weight with
$K_x=\delta\tau J$.

The off-diagonal part of the Hamiltonian,
\begin{equation}
    -h\sum_j\sigma_j^z,
\end{equation}
flips $\sigma^x$ eigenvalues.  For one spin,
\begin{equation}
\label{eq:ising-tfim-one-site-short-time-matrix}
    \bra{s'}_x\ee^{\delta\tau h\sigma^z}\ket{s}_x
    =
    \begin{cases}
    \cosh(\delta\tau h), & s'=s,\\
    \sinh(\delta\tau h), & s'=-s.
    \end{cases}
\end{equation}
The relative weight of a flip is therefore
\begin{equation}
\label{eq:ising-tfim-flip-weight}
    \frac{\sinh(\delta\tau h)}{\cosh(\delta\tau h)}
    =
    \tanh(\delta\tau h)
    =
    \ee^{-2K_y}.
\end{equation}
A vertical classical domain wall costs a factor $\ee^{-2K_y}$ relative to a
non-domain-wall segment.  Hence the transverse field is the rate per unit
imaginary time at which a worldline flips:
\begin{equation}
\label{eq:ising-tfim-transverse-field-as-flip-rate}
    \ee^{-2K_y}\simeq\delta\tau h .
\end{equation}
This is the microscopic origin of the dictionary.

The ordered and disordered phases can be read from this worldline picture:
\begin{itemize}[leftmargin=2.0em]
    \item If $J\gg h$, spatial alignment dominates and spin worldlines remain
    ordered in the $x$ direction.  This is the ferromagnetic phase of the TFIM.
    \item If $h\gg J$, spin flips proliferate in imaginary time and destroy
    long-range $x$ order.  This is the quantum paramagnet.
    \item At $h=J$, domain-wall proliferation is critical.  The transfer-matrix
    gap closes and the Euclidean theory becomes scale invariant.
\end{itemize}

\section{Finite temperature, finite size, and the extra dimension}
\label{sec:ising-tfim-finite-temperature-finite-size}

The number of rows $M$ in the transfer direction determines the quantum inverse
temperature:
\begin{equation}
    \beta=M\delta\tau .
\end{equation}
Therefore:
\begin{itemize}[leftmargin=2.0em]
    \item A classical cylinder with finite circumference $M$ in the transfer
    direction represents a quantum TFIM at finite temperature $T_{\rm q}=1/\beta$.
    \item The limit $M\to\infty$ at fixed $\delta\tau$ projects onto the largest
    eigenvalue of the transfer matrix.  In the quantum language this is the
    zero-temperature or ground-state limit.
    \item The continuum-time quantum limit requires $M\to\infty$ and
    $\delta\tau\to0$ with $\beta=M\delta\tau$ fixed.
\end{itemize}
The spatial size $L$ is the number of sites in the quantum spin chain.  Thus the
classical thermodynamic limit $L,M\to\infty$ contains both the quantum
thermodynamic limit and the zero-temperature limit.

The transfer-matrix correspondence also explains why a quantum critical point in
$d$ spatial dimensions is naturally described by a $(d+1)$-dimensional Euclidean
field theory.  The quantum partition function
\begin{equation}
    Z_{\rm q}=\Tr\ee^{-\beta H}
\end{equation}
is a sum over histories in imaginary time.  When the quantum gap closes, the
correlation length in space and the correlation length in imaginary time both
diverge.  The relative scaling of the two is controlled by the dynamical exponent
$z$.  For the TFIM critical point, $z=1$, and the long-distance continuum theory
is effectively isotropic after a nonuniversal rescaling of time.

\section{Common pitfalls}
\label{sec:ising-tfim-common-pitfalls}

There are several points that are easy to misstate.

\begin{enumerate}[label=(\roman*)]
    \item \textbf{The generic transfer matrix is not automatically a local
    TFIM Hamiltonian.}  For finite $K_x$ and $K_y$, the exact logarithm
    $-\log T$ is a well-defined free-fermion Hamiltonian, but it is generally not
    the nearest-neighbour TFIM.  Locality in the simple TFIM form emerges in the
    strongly anisotropic limit.

    \item \textbf{The normalization $A_y^L$ is not a physical transverse field.}
    It shifts the absolute free energy or the zero of quantum energy.  The
    transverse field is encoded in $\kappa_y$, not in $A_y$.

    \item \textbf{The choice of spin basis matters.}  In these notes the
    classical Ising variable is the eigenvalue of $\sigma^x$.  If one instead
    labels classical spins by $\sigma^z$, the same physics is obtained after a
    global spin rotation, but the symbols in the Hamiltonian are interchanged.

    \item \textbf{Classical temperature and quantum temperature are different
    concepts.}  The classical couplings $K_x,K_y$ already include the classical
    inverse temperature.  The quantum inverse temperature is the length of the
    transfer direction, $\beta=M\delta\tau$.

    \item \textbf{Boundary conditions introduce fermion parity sectors.}  After
    Jordan--Wigner transformation, periodic spin boundary conditions split the
    free-fermion solution into parity-dependent boundary conditions.  This
    sector bookkeeping is essential for exact finite-size spectra, although it
    does not change the local Hamiltonian-limit dictionary above.
\end{enumerate}

\section{Summary}
\label{sec:ising-tfim-summary}

The anisotropic two-dimensional classical Ising model contains the
one-dimensional quantum TFIM as its Hamiltonian limit.  The essential chain of
relations is
\begin{equation}
\label{eq:ising-tfim-summary-chain}
\begin{gathered}
    Z_{\rm 2D\ Ising}=\Tr T^M,
    \qquad
    T
    =
    A_y^L
    \ee^{\frac{K_x}{2}\sum_j\sigma_j^x\sigma_{j+1}^x}
    \ee^{\kappa_y\sum_j\sigma_j^z}
    \ee^{\frac{K_x}{2}\sum_j\sigma_j^x\sigma_{j+1}^x},
    \\
    K_x=\delta\tau J,
    \qquad
    \kappa_y=\delta\tau h,
    \qquad
    \tanh\kappa_y=\ee^{-2K_y},
    \\
    T
    =
    A_y^L\ee^{-\delta\tau H_{\rm TFIM}+O(\delta\tau^3)},
    \qquad
    H_{\rm TFIM}
    =
    -J\sum_j\sigma_j^x\sigma_{j+1}^x
    -h\sum_j\sigma_j^z .
\end{gathered}
\end{equation}
The coupling ratio is
\begin{equation}
\label{eq:ising-tfim-summary-ratio}
    \frac{h}{J}
    =
    \frac{\kappa_y}{K_x}
    =
    \frac{1}{2K_x}\log\coth K_y .
\end{equation}
The Onsager critical curve becomes the TFIM quantum critical point:
\begin{equation}
\label{eq:ising-tfim-summary-critical}
    \sinh(2K_x)\sinh(2K_y)=1
    \quad\Longleftrightarrow\quad
    K_x=\kappa_y
    \quad\Longleftrightarrow\quad
    J=h .
\end{equation}
At the spectral level,
\begin{equation}
\label{eq:ising-tfim-summary-dispersion}
    \varepsilon_{\rm tr}(k)
    =
    \delta\tau\,2\sqrt{h^2+J^2-2Jh\cos k}
    +O(\delta\tau^3),
\end{equation}
so the transfer-matrix gap is the quantum excitation gap measured in imaginary-
time units.  This chapter therefore gives the canonical solvable example of how
classical transfer matrices, quantum Hamiltonians, duality, and free-fermion
criticality are different representations of the same underlying structure.


\chapter{Six-Vertex Model and the XXZ Chain}
\label{chap:six-vertex-xxz}

The two-dimensional Ising model was solvable because its transfer matrix could be
rewritten in terms of free fermions.  The six-vertex model is solvable for a
different reason.  It is generally not a free-fermion model.  Its solvability
comes from the existence of an $R$-matrix satisfying the Yang--Baxter equation.
This equation implies a commuting family of transfer matrices, and the
logarithmic derivative of this family produces the one-dimensional XXZ quantum
spin chain.

This chapter has three goals:
\begin{enumerate}[label=(\roman*)]
    \item define the six allowed local vertices and the row-to-row transfer
    matrix;
    \item introduce the six-vertex $R$-matrix and explain how the Yang--Baxter
    equation gives commuting transfer matrices; and
    \item take the Hamiltonian limit and obtain the XXZ chain with anisotropy
    parameter $\Delta$.
\end{enumerate}
The coordinate and algebraic Bethe-ansatz diagonalizations are postponed to the
following chapters.  Here the main point is the exact map
\[
    \text{six-vertex transfer matrix}
    \quad\Longrightarrow\quad
    \text{commuting transfer matrices}
    \quad\Longrightarrow\quad
    \text{XXZ Hamiltonian}.
\]

\section{Arrows, the ice rule, and six local weights}
\label{sec:six-vertex-weights-ice-rule}

The six-vertex model is a classical statistical model on a square lattice.  The
degrees of freedom are arrows living on the edges.  At each vertex four arrows
meet.  The allowed configurations obey the \emph{ice rule}: exactly two arrows
point into the vertex and exactly two arrows point out of the vertex.

Equivalently, each local vertex conserves a two-state charge.  We identify an
arrow state with a spin-$1/2$ basis vector
\begin{equation}
    \ket{\uparrow}=\begin{pmatrix}1\\0\end{pmatrix},
    \qquad
    \ket{\downarrow}=\begin{pmatrix}0\\1\end{pmatrix},
    \qquad
    V\simeq \CC^2 .
\end{equation}
A local vertex weight is then an operator on $V\otimes V$.  In the ordered basis
\begin{equation}
    \ket{\uparrow\uparrow},
    \quad
    \ket{\uparrow\downarrow},
    \quad
    \ket{\downarrow\uparrow},
    \quad
    \ket{\downarrow\downarrow},
\end{equation}
the most common six-vertex weight matrix has the block form
\begin{equation}
\label{eq:six-vertex-constant-R-matrix}
    R=
    \begin{pmatrix}
        a&0&0&0\\
        0&b&c&0\\
        0&c&b&0\\
        0&0&0&a
    \end{pmatrix}.
\end{equation}
The six nonzero matrix elements are
\begin{equation}
\label{eq:six-vertex-six-nonzero-weights}
\begin{aligned}
    &\ket{\uparrow\uparrow}\longrightarrow\ket{\uparrow\uparrow}
    &&\text{with weight }a,\\
    &\ket{\downarrow\downarrow}\longrightarrow\ket{\downarrow\downarrow}
    &&\text{with weight }a,\\
    &\ket{\uparrow\downarrow}\longrightarrow\ket{\uparrow\downarrow}
    &&\text{with weight }b,\\
    &\ket{\downarrow\uparrow}\longrightarrow\ket{\downarrow\uparrow}
    &&\text{with weight }b,\\
    &\ket{\uparrow\downarrow}\longrightarrow\ket{\downarrow\uparrow}
    &&\text{with weight }c,\\
    &\ket{\downarrow\uparrow}\longrightarrow\ket{\uparrow\downarrow}
    &&\text{with weight }c.
\end{aligned}
\end{equation}
All other matrix elements vanish.  Thus $R$ conserves the total $z$-spin on the
two lines:
\begin{equation}
\label{eq:six-vertex-U1-conservation-local}
    \left[ R,\sigma^z\otimes\id+\id\otimes\sigma^z\right]=0 .
\end{equation}
This is the algebraic form of the ice rule.

\begin{definition}[Six-vertex anisotropy parameter]
For nonzero $a$ and $b$, the standard six-vertex anisotropy parameter is
\begin{equation}
\label{eq:six-vertex-Delta-from-abc}
    \Delta
    =
    \frac{a^2+b^2-c^2}{2ab} .
\end{equation}
Multiplying all weights by a common scalar changes the normalization of the
partition function but not $\Delta$.
\end{definition}

The parameter $\Delta$ is the same anisotropy parameter that appears in the XXZ
spin chain obtained from the Hamiltonian limit.  The local weights $(a,b,c)$
control a two-dimensional classical model, whereas $\Delta$ controls the relative
strength of the $z$ interaction in the associated one-dimensional quantum chain.

\begin{remark}[Why there are six vertices]
Without the ice rule there would be $2^4=16$ possible arrow configurations at a
four-valent vertex.  The constraint ``two in and two out'' leaves six.  In the
spin notation this is exactly the statement that the local weight preserves the
number of up spins, or equivalently the total $S^z$ charge, on the two lines.
\end{remark}

\section{Row-to-row transfer matrix}
\label{sec:six-vertex-row-transfer-matrix}

Consider a square lattice with $L$ columns and $M$ rows.  We impose periodic
boundary conditions in the horizontal direction.  A row configuration is specified
by the $L$ vertical arrows crossing a horizontal cut, hence the row Hilbert space
is
\begin{equation}
\label{eq:six-vertex-row-space}
    \calH_L=V_1\otimes V_2\otimes\cdots\otimes V_L,
    \qquad
    V_j\simeq\CC^2 .
\end{equation}
A row-to-row transfer matrix maps one row of vertical arrows to the next row,
summing over the internal horizontal arrows.

The efficient way to write this object is to introduce an auxiliary space
$V_a\simeq\CC^2$, which represents the horizontal arrow passing through the row.
Let $R_{aj}(u)$ act on the auxiliary space $a$ and the quantum space $j$.  The
monodromy matrix is
\begin{equation}
\label{eq:six-vertex-monodromy-matrix}
    M_a(u)
    =
    R_{aL}(u)R_{a,L-1}(u)\cdots R_{a1}(u),
\end{equation}
and the transfer matrix is the auxiliary trace
\begin{equation}
\label{eq:six-vertex-transfer-matrix}
    \tau(u)=\Tr_a M_a(u)
    =
    \Tr_a\!\bigl[R_{aL}(u)R_{a,L-1}(u)\cdots R_{a1}(u)\bigr] .
\end{equation}
The trace over $a$ closes the horizontal arrow into a periodic loop.  The
partition function on an $L\times M$ torus is then
\begin{equation}
\label{eq:six-vertex-partition-function-transfer}
    Z_{L,M}(u)=\Tr_{\calH_L}\tau(u)^M .
\end{equation}
Thus the thermodynamics of the two-dimensional classical model is controlled by
the spectrum of $\tau(u)$.

\begin{remark}[Two meanings of the same space]
In the classical model, $\calH_L$ is a convenient vector space for row
configurations.  In the quantum Hamiltonian limit, the same tensor product becomes
the Hilbert space of a length-$L$ spin chain.  This is the same conceptual move as
in the Ising transfer-matrix construction: a classical transfer direction becomes
imaginary time for a quantum problem.
\end{remark}

The ice rule implies a global $U(1)$ symmetry of the transfer matrix.  Let
\begin{equation}
\label{eq:six-vertex-total-magnetization}
    S^z_{\rm tot}=\frac12\sum_{j=1}^{L}\sigma_j^z .
\end{equation}
Since every local $R_{aj}(u)$ conserves the sum of the auxiliary and quantum
$z$-spins, tracing over the auxiliary space gives
\begin{equation}
\label{eq:six-vertex-transfer-U1-conservation}
    [\tau(u),S^z_{\rm tot}]=0 .
\end{equation}
The corresponding XXZ Hamiltonian therefore conserves total magnetization.

\section{The trigonometric six-vertex \texorpdfstring{$R$}{R}-matrix}
\label{sec:six-vertex-r-matrix}

A single transfer matrix is not yet enough to imply exact solvability.  The
integrable structure appears when the vertex weights are embedded in a
spectral-parameter family $R(u)$ satisfying the Yang--Baxter equation.  A standard
hyperbolic parametrization is
\begin{equation}
\label{eq:six-vertex-R-hyperbolic}
    R(u)=
    \begin{pmatrix}
        \sinh(u+\eta)&0&0&0\\
        0&\sinh u&\sinh\eta&0\\
        0&\sinh\eta&\sinh u&0\\
        0&0&0&\sinh(u+\eta)
    \end{pmatrix}.
\end{equation}
Here $u$ is the spectral parameter and $\eta$ is the anisotropy parameter.  In
this convention
\begin{equation}
\label{eq:six-vertex-hyperbolic-weights}
    a(u)=\sinh(u+\eta),
    \qquad
    b(u)=\sinh u,
    \qquad
    c(u)=\sinh\eta,
\end{equation}
and \cref{eq:six-vertex-Delta-from-abc} gives
\begin{equation}
\label{eq:xxz-Delta-cosh-eta}
    \Delta=\cosh\eta .
\end{equation}
For the critical regime one often uses the trigonometric parametrization
\begin{equation}
\label{eq:six-vertex-R-trigonometric}
    R(u)=
    \begin{pmatrix}
        \sin(u+\gamma)&0&0&0\\
        0&\sin u&\sin\gamma&0\\
        0&\sin\gamma&\sin u&0\\
        0&0&0&\sin(u+\gamma)
    \end{pmatrix},
    \qquad
    \Delta=\cos\gamma .
\end{equation}
The two forms are related by analytic continuation $\eta\leftrightarrow \ii\gamma$,
up to harmless scalar and phase normalizations.  The hyperbolic form is natural
for $\Delta>1$, while the trigonometric form is natural for $\abs{\Delta}<1$.

\begin{remark}[Scalar normalization]
The transformation
\begin{equation}
    R(u)\longmapsto f(u)R(u)
\end{equation}
with scalar $f(u)\neq0$ preserves the Yang--Baxter equation.  It multiplies the
transfer matrix by $f(u)^L$ and therefore changes only the scalar part of
$\log\tau(u)$.  In the Hamiltonian limit this produces an additive constant in
$H$, not a change of the interacting physics.
\end{remark}

The $R$-matrix in \cref{eq:six-vertex-R-hyperbolic} is \emph{regular}:
\begin{equation}
\label{eq:six-vertex-regularity}
    R(0)=\sinh\eta\,P,
\end{equation}
where $P:V\otimes V\to V\otimes V$ is the permutation operator
\begin{equation}
\label{eq:permutation-operator-definition}
    P(\ket{\alpha}\otimes\ket{\beta})
    =
    \ket{\beta}\otimes\ket{\alpha} .
\end{equation}
In Pauli matrices,
\begin{equation}
\label{eq:permutation-operator-pauli}
    P=\frac12\left(
        \id\otimes\id
        +\sigma^x\otimes\sigma^x
        +\sigma^y\otimes\sigma^y
        +\sigma^z\otimes\sigma^z
    \right).
\end{equation}
Regularity is the key property that makes the logarithmic derivative of
$\tau(u)$ local.

\section{Yang--Baxter equation and commuting transfer matrices}
\label{sec:six-vertex-yang-baxter}

Let $R_{ij}(u)$ denote $R(u)$ acting on the $i$th and $j$th tensor factors and as
identity on all others.  The Yang--Baxter equation is
\begin{equation}
\label{eq:six-vertex-YBE}
    R_{12}(u-v)R_{13}(u)R_{23}(v)
    =
    R_{23}(v)R_{13}(u)R_{12}(u-v) .
\end{equation}
It is an operator identity on $V_1\otimes V_2\otimes V_3$.  The six-vertex
$R$-matrix \cref{eq:six-vertex-R-hyperbolic} satisfies this equation.  In the
classical vertex-model language, \cref{eq:six-vertex-YBE} is a local
star--triangle move.  In the quantum scattering language, it is the consistency
condition for factorized three-body scattering.

\begin{theorem}[Commuting six-vertex transfer matrices]
\label{thm:six-vertex-commuting-transfer}
Let $R(u)$ satisfy the Yang--Baxter equation \cref{eq:six-vertex-YBE}.  Define
$\tau(u)$ by \cref{eq:six-vertex-transfer-matrix}.  Then
\begin{equation}
\label{eq:six-vertex-transfer-commutativity}
    [\tau(u),\tau(v)]=0,
    \qquad
    \text{for all }u,v .
\end{equation}
\end{theorem}

\begin{proof}
Introduce two auxiliary spaces $a$ and $b$.  The Yang--Baxter equation implies
the local relation
\begin{equation}
\label{eq:six-vertex-local-RLL}
    R_{ab}(u-v)R_{aj}(u)R_{bj}(v)
    =
    R_{bj}(v)R_{aj}(u)R_{ab}(u-v)
\end{equation}
for each quantum site $j$.  Multiplying this identity over all sites gives the
RTT relation
\begin{equation}
\label{eq:six-vertex-RTT}
    R_{ab}(u-v)M_a(u)M_b(v)
    =
    M_b(v)M_a(u)R_{ab}(u-v) .
\end{equation}
Taking $\Tr_a\Tr_b$ and using cyclicity of the trace in the auxiliary spaces,
$R_{ab}(u-v)$ can be moved through the trace and canceled.  The result is
\begin{equation}
    \Tr_a M_a(u)\,\Tr_b M_b(v)
    =
    \Tr_b M_b(v)\,\Tr_a M_a(u),
\end{equation}
which is exactly $\tau(u)\tau(v)=\tau(v)\tau(u)$.
\end{proof}

The implication is substantial.  Expanding the logarithm of $\tau(u)$ around a
regular point gives an infinite sequence of mutually commuting conserved charges:
\begin{equation}
\label{eq:six-vertex-log-transfer-charges}
    \log\tau(u)
    =
    \sum_{n=0}^{\infty}(u-u_0)^n Q_n,
    \qquad
    [Q_m,Q_n]=0 .
\end{equation}
The first nontrivial local charge is the XXZ Hamiltonian.  Higher charges have
larger spatial range and are responsible for integrability.

\begin{principle}[Yang--Baxter integrability]
The Yang--Baxter equation is the local algebraic identity that upgrades one
transfer matrix into a commuting one-parameter family.  The logarithm of this
family generates the Hamiltonian and the higher conserved charges.
\end{principle}

\section{Hamiltonian limit and the XXZ chain}
\label{sec:six-vertex-hamiltonian-limit}

We now derive the quantum Hamiltonian associated with the regular six-vertex
transfer matrix.  For a homogeneous periodic chain, regularity gives
\begin{equation}
\label{eq:six-vertex-transfer-at-zero-shift}
    \tau(0)
    =
    (\sinh\eta)^L\,U,
\end{equation}
where $U$ is the one-site translation operator on the spin chain.  Taking the
logarithmic derivative removes this translation and produces a nearest-neighbor
operator:
\begin{equation}
\label{eq:six-vertex-log-derivative-locality}
    \left.\frac{\dd}{\dd u}\log\tau(u)\right|_{u=0}
    =
    \sum_{j=1}^{L}
    \frac{1}{\sinh\eta}
    P_{j,j+1}R'_{j,j+1}(0),
\end{equation}
with $j+1$ understood modulo $L$.  The local operator can be computed explicitly.
From \cref{eq:six-vertex-R-hyperbolic},
\begin{equation}
\label{eq:six-vertex-R-prime-zero}
    R'(0)=
    \begin{pmatrix}
        \cosh\eta&0&0&0\\
        0&1&0&0\\
        0&0&1&0\\
        0&0&0&\cosh\eta
    \end{pmatrix}.
\end{equation}
Therefore
\begin{equation}
\label{eq:six-vertex-local-log-derivative-pauli}
    \frac{1}{\sinh\eta}P R'(0)
    =
    \frac{1}{2\sinh\eta}
    \left(
        \sigma^x\otimes\sigma^x
        +\sigma^y\otimes\sigma^y
        +\cosh\eta\,\sigma^z\otimes\sigma^z
        +\cosh\eta\,\id\otimes\id
    \right).
\end{equation}
Multiplying the logarithmic derivative by $2\sinh\eta$ and subtracting an
additive constant gives
\begin{equation}
\label{eq:xxz-from-six-vertex-integrability-normalization}
    H_{\rm XXZ}^{(+)}
    =
    J\sum_{j=1}^{L}
    \left(
        \sigma_j^x\sigma_{j+1}^x
        +\sigma_j^y\sigma_{j+1}^y
        +\Delta\bigl(\sigma_j^z\sigma_{j+1}^z-1\bigr)
    \right),
    \qquad
    \Delta=\cosh\eta .
\end{equation}
Equivalently,
\begin{equation}
\label{eq:xxz-hamiltonian-log-derivative-final}
    H_{\rm XXZ}^{(+)}
    =
    2J\sinh\eta
    \left.\frac{\dd}{\dd u}\log\tau(u)\right|_{u=0}
    -2J\Delta L .
\end{equation}
The superscript $(+)$ indicates the standard antiferromagnetic integrability
normalization for $J>0$ and $\Delta>0$.  The constant $-2J\Delta L$ fixes the
energy of the fully polarized reference state to zero; it has no effect on
eigenvectors, gaps, or correlation functions.

In spin-operator normalization, $S_j^\alpha=\sigma_j^\alpha/2$, the same
Hamiltonian is
\begin{equation}
\label{eq:xxz-spin-operator-normalization-from-six-vertex}
    H_{\rm XXZ}^{(+)}
    =
    4J\sum_{j=1}^{L}
    \left(
        S_j^xS_{j+1}^x
        +S_j^yS_{j+1}^y
        +\Delta S_j^zS_{j+1}^z
    \right)
    -J\Delta L .
\end{equation}
Only the overall energy scale and additive constant differ from other common XXZ
conventions.

\begin{remark}[Relation to the earlier Pauli convention]
Earlier chapters used the ferromagnetic Pauli convention
\begin{equation}
    H_{\rm XXZ}^{\rm ferro}
    =
    -\frac{J_{\rm f}}{2}\sum_j
    \left(
        \sigma_j^x\sigma_{j+1}^x
        +\sigma_j^y\sigma_{j+1}^y
        +\Delta\sigma_j^z\sigma_{j+1}^z
    \right)
    -h\sum_j\sigma_j^z .
\end{equation}
The Hamiltonian generated directly by the six-vertex $R$-matrix is instead the
integrability convention \cref{eq:xxz-from-six-vertex-integrability-normalization}.
The two are related by an overall sign and energy-scale choice, plus an additive
constant.  Whenever Bethe-ansatz formulas are quoted, the sign convention must be
stated explicitly.
\end{remark}

The trigonometric $R$-matrix gives the same formula with
\begin{equation}
\label{eq:xxz-gapless-Delta-cos-gamma}
    \Delta=\cos\gamma,
    \qquad
    \sinh\eta\longrightarrow\sin\gamma .
\end{equation}
Thus the critical XXZ regime $\abs{\Delta}<1$ is naturally parametrized by
$0<\gamma<\pi$.

\section{Why the Hamiltonian is local}
\label{sec:six-vertex-why-local}

It is useful to isolate the mechanism behind locality.  The transfer matrix
\begin{equation}
    \tau(u)=\Tr_a R_{aL}(u)\cdots R_{a1}(u)
\end{equation}
is a nonlocal operator on the whole chain.  Nevertheless, at the regular point
$u=0$ every $R_{aj}(0)$ is a permutation.  The product of permutations simply
moves the auxiliary spin through the chain and becomes the translation operator
after the auxiliary trace.  Differentiating with respect to $u$ replaces exactly
one local permutation by $R'(0)$, while all other factors remain permutations.
After multiplying by $\tau(0)^{-1}$, this single defect becomes a local
nearest-neighbor density.

This is the vertex-model analogue of the Trotter logic in the Ising--TFIM
correspondence.  The difference is that the present transfer matrix is not merely
one short-time evolution operator.  It belongs to a commuting spectral-parameter
family, and this additional structure produces infinitely many conserved charges.

\begin{proposition}[Local charge from regularity]
\label{prop:regular-transfer-local-charge}
Let $R(u)$ satisfy $R(0)=\rho P$ with $\rho\neq0$, and define the homogeneous
periodic transfer matrix by \cref{eq:six-vertex-transfer-matrix}.  Then
\begin{equation}
    \left.\frac{\dd}{\dd u}\log\tau(u)\right|_{u=0}
    =
    \sum_{j=1}^{L}h_{j,j+1},
    \qquad
    h_{j,j+1}=\rho^{-1}P_{j,j+1}R'_{j,j+1}(0),
\end{equation}
up to the conventional orientation of the one-site translation operator.
\end{proposition}

\section{Conserved magnetization and interacting character}
\label{sec:six-vertex-conserved-magnetization-interaction}

The XXZ Hamiltonian obtained above satisfies
\begin{equation}
\label{eq:xxz-U1-symmetry}
    [H_{\rm XXZ}^{(+)},S^z_{\rm tot}]=0 .
\end{equation}
Using $\sigma_j^\pm=(\sigma_j^x\pm\ii\sigma_j^y)/2$, one may rewrite the exchange
part as
\begin{equation}
\label{eq:xxz-ladder-form}
    \sigma_j^x\sigma_{j+1}^x+
    \sigma_j^y\sigma_{j+1}^y
    =
    2\left(
        \sigma_j^+\sigma_{j+1}^-
        +
        \sigma_j^-\sigma_{j+1}^+
    \right).
\end{equation}
Thus the $XY$ part hops a down spin through an up-spin background, while the
$\Delta\sigma_j^z\sigma_{j+1}^z$ term gives a nearest-neighbor interaction between
magnons.  At $\Delta=0$ the model is the XX chain and becomes free after the
Jordan--Wigner transformation.  At generic $\Delta$, the Jordan--Wigner image
contains a density--density interaction and is not a quadratic fermion model.

Explicitly, under Jordan--Wigner fermionization, the $XY$ part becomes hopping,
whereas
\begin{equation}
\label{eq:xxz-density-density-interaction}
    \sigma_j^z\sigma_{j+1}^z
    =
    (1-2n_j)(1-2n_{j+1})
    =
    1-2(n_j+n_{j+1})+4n_jn_{j+1} .
\end{equation}
The quartic term $n_jn_{j+1}$ is the interacting part.  Therefore the XXZ chain
is a genuinely interacting integrable model, except at the free-fermion point
$\Delta=0$.

\begin{remark}[Why this chapter leaves the spectrum unfinished]
For the Ising, XY, Kitaev, and SSH chains, exact solvability meant reducing the
Hamiltonian to independent BdG blocks.  For the XXZ chain this is impossible at
generic $\Delta$ because of the quartic Jordan--Wigner interaction.  Its solution
requires Bethe ansatz: many-body eigenstates are built from interacting magnons
whose scattering factorizes.  The six-vertex transfer matrix supplies the
Yang--Baxter reason why that factorization is consistent.
\end{remark}

\section{Anisotropy regimes}
\label{sec:six-vertex-xxz-regimes}

For the antiferromagnetic XXZ convention
\begin{equation}
\label{eq:xxz-standard-spin-regime-convention}
    H_{\rm XXZ}^{S}
    =
    \mathcal{J}\sum_j
    \left(
        S_j^xS_{j+1}^x
        +S_j^yS_{j+1}^y
        +\Delta S_j^zS_{j+1}^z
    \right),
    \qquad
    \mathcal{J}>0,
\end{equation}
the familiar zero-field phase structure is summarized as follows:
\begin{center}
\begin{tabular}{@{}L{0.22\linewidth}L{0.30\linewidth}L{0.34\linewidth}@{}}
\toprule
Parameter regime & Useful parametrization & Quantum-chain behavior \\
\midrule
$\Delta>1$ & $\Delta=\cosh\eta$ & gapped antiferromagnetic, or N\'{e}el, regime \\
$-1<\Delta<1$ & $\Delta=\cos\gamma$ & gapless critical Luttinger-liquid regime \\
$\Delta=1$ & isotropic XXX point & critical with full $SU(2)$ symmetry and marginal corrections \\
$\Delta<-1$ & analytic continuation or sign-twisted convention & gapped ferromagnetic regime \\
\bottomrule
\end{tabular}
\end{center}
For the classical six-vertex model with positive weights, one also classifies
regions using
\begin{equation}
    \Delta=\frac{a^2+b^2-c^2}{2ab} .
\end{equation}
The regime $\abs{\Delta}<1$ is the critical disordered phase of the classical
vertex model.  The regions $\Delta>1$ and $\Delta<-1$ correspond to ordered
ferroelectric and antiferroelectric regimes, respectively, depending on which
weights dominate.  The exact terminology is classical-model dependent, but the
same invariant $\Delta$ controls the integrable family.

\begin{remark}[Classical phases versus quantum phases]
The two-dimensional six-vertex partition function at a fixed spectral parameter
and the one-dimensional XXZ Hamiltonian obtained by a logarithmic derivative are
not the same physical object.  They are two different uses of the same commuting
transfer-matrix family.  The shared parameter $\Delta$ explains why their phase
diagrams are closely related, but one should not identify a classical ordered
region with a quantum ground-state phase without specifying the Hamiltonian limit
and the sign convention.
\end{remark}

\section{What the transfer matrix solution will become}
\label{sec:six-vertex-what-comes-next}

The six-vertex transfer matrix is diagonalized by Bethe ansatz.  In the algebraic
form, one writes the monodromy matrix as a $2\times2$ matrix in auxiliary space,
\begin{equation}
\label{eq:six-vertex-ABA-monodromy-entries}
    M_a(u)=
    \begin{pmatrix}
        A(u)&B(u)\\
        C(u)&D(u)
    \end{pmatrix}_a,
    \qquad
    \tau(u)=A(u)+D(u).
\end{equation}
The operator $B(u)$ creates Bethe quasiparticles above a reference state, and the
RTT relation generated by the Yang--Baxter equation gives the algebra needed to
move $A(u)$ and $D(u)$ through a product of $B$'s.  Requiring unwanted terms to
vanish gives the Bethe equations.

In the coordinate form, the same result appears as factorized magnon scattering.
In a sector with $M$ down spins, a wave function is a sum over plane waves in each
ordered coordinate region,
\begin{equation}
    \Psi(x_1,\ldots,x_M)
    =
    \sum_{P\in S_M}A(P)
    \exp\left(\ii\sum_{a=1}^{M}k_{P_a}x_a\right),
    \qquad
    x_1<\cdots<x_M .
\end{equation}
The two-body scattering phase determines how the amplitudes $A(P)$ change when
neighboring momenta are exchanged.  Periodic boundary conditions then lead to
Bethe equations.  The following chapters derive these statements explicitly.

\section{Summary}
\label{sec:six-vertex-summary}

The essential chain of ideas is
\begin{equation}
\label{eq:six-vertex-summary-chain}
\boxed{
\begin{gathered}
    \text{ice-rule vertex weights}
    \;\Longrightarrow\;
    R(u)\text{ with }U(1)\text{ conservation}
    \\
    \Longrightarrow\;
    \text{Yang--Baxter equation}
    \;\Longrightarrow\;
    [\tau(u),\tau(v)]=0
    \\
    \Longrightarrow\;
    H_{\rm XXZ}\propto
    \left.\frac{\dd}{\dd u}\log\tau(u)\right|_{u=0}
    \;\Longrightarrow\;
    \text{interacting Bethe-ansatz integrability.}
\end{gathered}}
\end{equation}
The six-vertex model is therefore the canonical bridge between two-dimensional
classical vertex models and one-dimensional interacting quantum spin chains.
Unlike the Ising--TFIM correspondence, the resulting quantum chain is not solved
by free fermions at generic anisotropy.  Its solvability is instead controlled by
Yang--Baxter integrability and factorized scattering.


\chapter{Eight-Vertex Model and the XYZ Chain}
\label{chap:eight-vertex-xyz}

The six-vertex model is the vertex-model counterpart of the XXZ chain.  Its ice
rule gives a local $U(1)$ conservation law, and the trigonometric $R$-matrix gives
an interacting but integrable quantum spin chain.  The eight-vertex model is the
next member in the same hierarchy.  It relaxes the ice rule while preserving an
even-arrow, or $\ZZ_2$, constraint.  Algebraically this replaces the six-vertex
$R$-matrix by an elliptic $R$-matrix.  The Hamiltonian limit is no longer XXZ but
the fully anisotropic XYZ spin chain.

The exact correspondence is
\[
\begin{aligned}
    \text{eight-vertex transfer matrix}
    &\quad\Longrightarrow\quad
    \text{elliptic Yang--Baxter equation}\\
    &\quad\Longrightarrow\quad
    \text{XYZ Hamiltonian}.
\end{aligned}
\]
The important conceptual change from the previous chapter is that the model is
still Yang--Baxter integrable, but it is no longer organized by total
magnetization sectors.  The missing $U(1)$ symmetry is the main reason why the
solution of the eight-vertex model is technically harder than the six-vertex
case.

This chapter does not derive Baxter's full solution.  Its goal is to set up the
model, state the elliptic $R$-matrix structure, and derive the Hamiltonian limit
that produces the XYZ chain.  Baxter's $TQ$ relation and the full spectral theory
belong to a more advanced treatment of elliptic integrability.

\section{Eight allowed vertices and local weights}
\label{sec:eight-vertex-local-weights}

The eight-vertex model is a classical statistical model on the square lattice.
As in the six-vertex model, arrows live on edges and a Boltzmann weight is
assigned to each vertex.  The local constraint is weaker than the ice rule:
\begin{equation}
\label{eq:eight-vertex-even-arrow-rule}
    \text{an allowed vertex has an even number of arrows pointing in.}
\end{equation}
Equivalently, the allowed numbers of incoming arrows are $0$, $2$, and $4$.
There are therefore eight allowed local configurations.  The six ice-rule
vertices are included, and the two additional vertices have all arrows pointing
in or all arrows pointing out.

In spin notation a local vertex weight is an operator on
$V\otimes V$, where $V\simeq\CC^2$.  We use the same ordered basis as in the
six-vertex chapter,
\begin{equation}
    \ket{\uparrow\uparrow},\quad
    \ket{\uparrow\downarrow},\quad
    \ket{\downarrow\uparrow},\quad
    \ket{\downarrow\downarrow}.
\end{equation}
A standard symmetric eight-vertex weight matrix has the form
\begin{equation}
\label{eq:eight-vertex-R-abcd}
    R=
    \begin{pmatrix}
        a&0&0&d\\
        0&b&c&0\\
        0&c&b&0\\
        d&0&0&a
    \end{pmatrix}.
\end{equation}
The six-vertex model is recovered when $d=0$.  The extra $d$-terms are the new
eight-vertex processes
\begin{equation}
\label{eq:eight-vertex-pair-flip-processes}
    \ket{\uparrow\uparrow}\longleftrightarrow\ket{\downarrow\downarrow}.
\end{equation}
They change the total $S^z$ of the two lines by two units.  Thus the model no
longer conserves total $S^z$, but it conserves its parity.

\begin{proposition}[$\ZZ_2$ conservation]
Let
\begin{equation}
    Q_2=\sigma^z\otimes\sigma^z .
\end{equation}
For the eight-vertex weight \cref{eq:eight-vertex-R-abcd},
\begin{equation}
\label{eq:eight-vertex-local-Z2-conservation}
    [R,Q_2]=0 .
\end{equation}
However, for $d\neq0$,
\begin{equation}
    \left[R,\sigma^z\otimes\id+\id\otimes\sigma^z\right]\neq0 .
\end{equation}
Thus the local $U(1)$ conservation of the six-vertex model is reduced to a
$\ZZ_2$ parity conservation.
\end{proposition}

\begin{proof}
The diagonal $a$ and $b$ processes clearly preserve both the total $z$-spin and
its parity.  The $c$ processes exchange
$\ket{\uparrow\downarrow}$ and $\ket{\downarrow\uparrow}$, so they also preserve
total $z$-spin.  The $d$ processes exchange
$\ket{\uparrow\uparrow}$ and $\ket{\downarrow\downarrow}$, changing total $S^z$
by $\pm2$ but preserving the eigenvalue of $\sigma^z\otimes\sigma^z$.  This is
exactly \cref{eq:eight-vertex-local-Z2-conservation}.
\end{proof}

\begin{remark}[Why the eight-vertex model is harder]
The six-vertex transfer matrix decomposes into fixed-magnetization sectors.  This
sector decomposition is the basic input behind the usual algebraic Bethe ansatz
with a simple ferromagnetic reference state.  In the eight-vertex model the
$d$-processes destroy $S^z$ conservation, so the standard highest-weight
construction no longer applies directly.  Baxter's solution replaces it with an
elliptic functional-relation method.
\end{remark}

\section{Row-to-row transfer matrix}
\label{sec:eight-vertex-transfer-matrix}

The transfer-matrix construction is formally the same as in the six-vertex
model.  On a row of length $L$, the row space is
\begin{equation}
\label{eq:eight-vertex-row-space}
    \calH_L=V_1\otimes V_2\otimes\cdots\otimes V_L,
    \qquad V_j\simeq\CC^2 .
\end{equation}
Introduce an auxiliary space $V_a\simeq\CC^2$.  Given a spectral-parameter family
$R(u)$ of eight-vertex weights, define the monodromy matrix
\begin{equation}
\label{eq:eight-vertex-monodromy}
    M_a(u)=R_{aL}(u)R_{a,L-1}(u)\cdots R_{a1}(u),
\end{equation}
and the row-to-row transfer matrix
\begin{equation}
\label{eq:eight-vertex-transfer-matrix}
    \tau(u)=\Tr_a M_a(u)
    =\Tr_a\!\left[R_{aL}(u)R_{a,L-1}(u)\cdots R_{a1}(u)\right] .
\end{equation}
For an $L\times M$ torus,
\begin{equation}
\label{eq:eight-vertex-partition-function}
    Z_{L,M}(u)=\Tr_{\calH_L}\tau(u)^M .
\end{equation}
The spectrum of $\tau(u)$ controls the thermodynamic limit of the classical
model.

Because of the local parity conservation in
\cref{eq:eight-vertex-local-Z2-conservation}, the transfer matrix commutes with
the global spin-parity operator
\begin{equation}
\label{eq:eight-vertex-global-parity}
    \mathcal{P}=\prod_{j=1}^{L}\sigma_j^z,
    \qquad
    [\tau(u),\mathcal{P}]=0 .
\end{equation}
For $d\neq0$ one does not generally have
\begin{equation}
    [\tau(u),S^z_{\rm tot}]=0,
    \qquad
    S^z_{\rm tot}=\frac12\sum_{j=1}^{L}\sigma_j^z .
\end{equation}
This is the transfer-matrix reflection of the fact that the corresponding XYZ
Hamiltonian has only a discrete spin-flip/parity symmetry, not a continuous
$U(1)$ symmetry.

\section{Baxter's elliptic parametrization}
\label{sec:baxter-elliptic-parametrization}

A generic matrix of the form \cref{eq:eight-vertex-R-abcd} does not automatically
produce commuting transfer matrices.  Integrability requires that the weights be
embedded in an elliptic spectral-parameter family satisfying the Yang--Baxter
equation.

Let $\vartheta_1(u)$ and $\vartheta_4(u)$ denote Jacobi theta functions with a
fixed nome $q$, $\abs{q}<1$.  Let $\eta$ be the crossing parameter.  One common
Baxter parametrization, shifted so that the regular point is at $u=0$, is
\begin{equation}
\label{eq:eight-vertex-elliptic-weights}
\begin{aligned}
    a(u)&=\rho\,\vartheta_4(2\eta)\,\vartheta_4(u)\,\vartheta_1(u+2\eta),\\
    b(u)&=\rho\,\vartheta_4(2\eta)\,\vartheta_1(u)\,\vartheta_4(u+2\eta),\\
    c(u)&=\rho\,\vartheta_1(2\eta)\,\vartheta_4(u)\,\vartheta_4(u+2\eta),\\
    d(u)&=\rho\,\vartheta_1(2\eta)\,\vartheta_1(u)\,\vartheta_1(u+2\eta).
\end{aligned}
\end{equation}
Here $\rho$ is an arbitrary nonzero scalar normalization.  The corresponding
$R$-matrix is
\begin{equation}
\label{eq:eight-vertex-R-u}
    R(u)=
    \begin{pmatrix}
        a(u)&0&0&d(u)\\
        0&b(u)&c(u)&0\\
        0&c(u)&b(u)&0\\
        d(u)&0&0&a(u)
    \end{pmatrix}.
\end{equation}
Since $\vartheta_1(0)=0$ and $\vartheta_4'(0)=0$, the regular point satisfies
\begin{equation}
\label{eq:eight-vertex-regularity}
    b(0)=d(0)=0,
    \qquad
    a(0)=c(0)=\rho_0,
\end{equation}
where
\begin{equation}
\label{eq:eight-vertex-rho0}
    \rho_0
    =\rho\,\vartheta_4(0)\vartheta_4(2\eta)\vartheta_1(2\eta).
\end{equation}
Therefore
\begin{equation}
\label{eq:eight-vertex-R0-P}
    R(0)=\rho_0 P,
\end{equation}
where $P$ is the two-spin permutation operator,
\begin{equation}
\label{eq:eight-vertex-permutation-operator}
    P\ket{\alpha}\otimes\ket{\beta}
    =\ket{\beta}\otimes\ket{\alpha}.
\end{equation}

\begin{remark}[Normalization and shift conventions]
The eight-vertex literature contains several equivalent theta-function
conventions.  Some place the regular point at $u=\eta$ rather than $u=0$, and
some use theta functions denoted by $H$ and $\Theta$.  The convention in
\cref{eq:eight-vertex-elliptic-weights} is chosen to match the Hamiltonian-limit
convention used in these notes: $R(0)$ is proportional to the permutation
operator.
\end{remark}

For the family \cref{eq:eight-vertex-R-u}, the Yang--Baxter equation takes the
same formal form as in the six-vertex case:
\begin{equation}
\label{eq:eight-vertex-YBE}
    R_{12}(u-v)R_{13}(u-w)R_{23}(v-w)
    =
    R_{23}(v-w)R_{13}(u-w)R_{12}(u-v).
\end{equation}
This identity is a nontrivial theta-function identity.  Its consequence is the
commutativity of transfer matrices:
\begin{equation}
\label{eq:eight-vertex-commuting-transfer-matrices}
    [\tau(u),\tau(v)]=0
    \qquad
    \text{for all }u,v .
\end{equation}
Thus the eight-vertex model is integrable in the same structural sense as the
six-vertex model: there is a one-parameter commuting family of transfer
matrices.

\begin{principle}[Elliptic integrability]
The hierarchy
\begin{equation}
    \text{rational }R\text{-matrix}
    \longrightarrow
    \text{trigonometric }R\text{-matrix}
    \longrightarrow
    \text{elliptic }R\text{-matrix}
\end{equation}
corresponds, respectively, to the XXX, XXZ, and XYZ spin chains.  The eight-vertex
model is the elliptic vertex model whose Hamiltonian limit is the XYZ chain.
\end{principle}

A useful way to see the difference from the six-vertex model is the Pauli-basis
form.  Since
\begin{align}
    \id\otimes\id
    &\mapsto
    \begin{pmatrix}
        1&0&0&0\\0&1&0&0\\0&0&1&0\\0&0&0&1
    \end{pmatrix},\\
    \sigma^x\otimes\sigma^x
    &\mapsto
    \begin{pmatrix}
        0&0&0&1\\0&0&1&0\\0&1&0&0\\1&0&0&0
    \end{pmatrix},\\
    \sigma^y\otimes\sigma^y
    &\mapsto
    \begin{pmatrix}
        0&0&0&-1\\0&0&1&0\\0&1&0&0\\-1&0&0&0
    \end{pmatrix},\\
    \sigma^z\otimes\sigma^z
    &\mapsto
    \begin{pmatrix}
        1&0&0&0\\0&-1&0&0\\0&0&-1&0\\0&0&0&1
    \end{pmatrix},
\end{align}
we may rewrite \cref{eq:eight-vertex-R-u} as
\begin{equation}
\label{eq:eight-vertex-R-Pauli-basis}
    R(u)=
    A(u)\,\id\otimes\id
    +B(u)\,\sigma^x\otimes\sigma^x
    +C(u)\,\sigma^y\otimes\sigma^y
    +D(u)\,\sigma^z\otimes\sigma^z,
\end{equation}
with
\begin{equation}
\label{eq:eight-vertex-ABCD}
    A(u)=\frac{a(u)+b(u)}{2},
    \qquad
    D(u)=\frac{a(u)-b(u)}{2},
\end{equation}
\begin{equation}
\label{eq:eight-vertex-BCD}
    B(u)=\frac{c(u)+d(u)}{2},
    \qquad
    C(u)=\frac{c(u)-d(u)}{2}.
\end{equation}
The distinction between $\sigma^x\sigma^x$ and $\sigma^y\sigma^y$ is controlled
by the $d$-weight.  When $d=0$, the coefficients of
$\sigma^x\otimes\sigma^x$ and $\sigma^y\otimes\sigma^y$ coincide, and the model
reduces to the six-vertex/XXZ type.

\section{Hamiltonian limit and the XYZ chain}
\label{sec:eight-vertex-xyz-hamiltonian-limit}

The Hamiltonian limit is obtained from the logarithmic derivative of the transfer
matrix at the regular point.  Because $R(0)=\rho_0P$, the transfer matrix at
$u=0$ is proportional to the one-site translation operator.  The first derivative
therefore gives a sum of local nearest-neighbour terms.

\begin{proposition}[Hamiltonian from the eight-vertex transfer matrix]
Let $\tau(u)$ be the periodic transfer matrix built from the regular
$R$-matrix \cref{eq:eight-vertex-R-u}.  Define
\begin{equation}
\label{eq:eight-vertex-H-log-derivative}
    H_{\rm XYZ}
    =\left.\frac{\dd}{\dd u}\log\tau(u)\right|_{u=0}.
\end{equation}
After dropping an additive scalar multiple of the identity and allowing an
overall energy-scale normalization, one obtains
\begin{equation}
\label{eq:XYZ-Hamiltonian-main}
    H_{\rm XYZ}
    =
    \sum_{j=1}^{L}
    \left(
        J_x\,\sigma_j^x\sigma_{j+1}^x
        +J_y\,\sigma_j^y\sigma_{j+1}^y
        +J_z\,\sigma_j^z\sigma_{j+1}^z
    \right),
    \qquad
    \sigma_{L+1}^{\alpha}\equiv\sigma_1^{\alpha}.
\end{equation}
The three couplings are determined by the first derivatives of the vertex
weights at the regular point:
\begin{equation}
\label{eq:eight-vertex-Jxyz-from-weight-derivatives}
    J_x=\frac{b'(0)+d'(0)}{2\rho_0},
    \qquad
    J_y=\frac{b'(0)-d'(0)}{2\rho_0},
    \qquad
    J_z=\frac{a'(0)-c'(0)}{2\rho_0}.
\end{equation}
The discarded additive constant is
\begin{equation}
\label{eq:eight-vertex-additive-constant}
    E_0=\frac{a'(0)+c'(0)}{2\rho_0}
\end{equation}
per bond.
\end{proposition}

\begin{proof}
By regularity, $R(0)=\rho_0P$.  The standard transfer-matrix argument gives a
local logarithmic derivative
\begin{equation}
\label{eq:eight-vertex-local-h-P-Rprime}
    h_{j,j+1}=\frac{1}{\rho_0}P_{j,j+1}R'_{j,j+1}(0).
\end{equation}
For the matrix form \cref{eq:eight-vertex-R-u},
\begin{equation}
\label{eq:eight-vertex-PRprime-matrix}
    \frac{1}{\rho_0}PR'(0)
    =\frac{1}{\rho_0}
    \begin{pmatrix}
        a'(0)&0&0&d'(0)\\
        0&c'(0)&b'(0)&0\\
        0&b'(0)&c'(0)&0\\
        d'(0)&0&0&a'(0)
    \end{pmatrix}.
\end{equation}
Comparing this matrix with
\begin{equation}
\label{eq:eight-vertex-local-XYZ-comparison}
    E_0\id
    +J_x\sigma^x\otimes\sigma^x
    +J_y\sigma^y\otimes\sigma^y
    +J_z\sigma^z\otimes\sigma^z
\end{equation}
gives exactly \cref{eq:eight-vertex-Jxyz-from-weight-derivatives} and
\cref{eq:eight-vertex-additive-constant}.
\end{proof}

For the theta-function weights in \cref{eq:eight-vertex-elliptic-weights}, these
formulas give
\begin{equation}
\label{eq:eight-vertex-JxJy-theta}
    J_x
    =
    \frac{\vartheta_1'(0)}{2\vartheta_4(0)\vartheta_4(2\eta)\vartheta_1(2\eta)}
    \left[\vartheta_4(2\eta)^2+\vartheta_1(2\eta)^2\right],
\end{equation}
\begin{equation}
\label{eq:eight-vertex-Jy-theta}
    J_y
    =
    \frac{\vartheta_1'(0)}{2\vartheta_4(0)\vartheta_4(2\eta)\vartheta_1(2\eta)}
    \left[\vartheta_4(2\eta)^2-\vartheta_1(2\eta)^2\right],
\end{equation}
and
\begin{equation}
\label{eq:eight-vertex-Jz-theta}
    J_z
    =
    \frac12
    \left[
        \frac{\vartheta_1'(2\eta)}{\vartheta_1(2\eta)}
        -
        \frac{\vartheta_4'(2\eta)}{\vartheta_4(2\eta)}
    \right].
\end{equation}
These expressions depend on the chosen normalization of the spectral parameter.
A rescaling $u\mapsto\lambda u$ rescales all three couplings by the common factor
$\lambda$.  The physically meaningful anisotropies are the coupling ratios.

\begin{remark}[Sign convention]
Some condensed-matter texts write
\begin{equation}
    H=-\sum_j
    \left(
        J_x\sigma_j^x\sigma_{j+1}^x
        +J_y\sigma_j^y\sigma_{j+1}^y
        +J_z\sigma_j^z\sigma_{j+1}^z
    \right).
\end{equation}
This differs from \cref{eq:XYZ-Hamiltonian-main} by redefining
$J_{\alpha}\mapsto -J_{\alpha}$.  The integrable structure is unaffected by this
overall convention.
\end{remark}

\section{Relation to the XXZ chain and special limits}
\label{sec:eight-vertex-special-limits}

The XYZ chain is the most general nearest-neighbour spin-$1/2$ chain with
diagonal exchange couplings and no magnetic field:
\begin{equation}
\label{eq:XYZ-general-H}
    H_{\rm XYZ}
    =
    \sum_j
    \left(
        J_x\sigma_j^x\sigma_{j+1}^x
        +J_y\sigma_j^y\sigma_{j+1}^y
        +J_z\sigma_j^z\sigma_{j+1}^z
    \right).
\end{equation}
The XXZ chain is the special case
\begin{equation}
\label{eq:XXZ-as-XYZ-special-case}
    J_x=J_y,
    \qquad
    J_z=\Delta J_x .
\end{equation}
In the vertex-model language this corresponds to the degeneration
\begin{equation}
\label{eq:eight-to-six-d-zero}
    d(u)\to0,
\end{equation}
so that the eight-vertex $R$-matrix becomes a six-vertex $R$-matrix and the
elliptic parametrization degenerates to a trigonometric one.  In theta-function
language this occurs in the trigonometric limit of the elliptic nome.  In that
limit the $d$-weight disappears, $J_x-J_y\to0$, and the $U(1)$ symmetry of the
XXZ chain is restored.

\begin{center}
\begin{tabular}{@{}lll@{}}
\toprule
Vertex model & Quantum chain & Symmetry \\
\midrule
Six-vertex & XXZ & $U(1)$ generated by $S^z_{\rm tot}$ \\
Eight-vertex & XYZ & $\ZZ_2$ generated by $\prod_j\sigma_j^z$ \\
Trigonometric $R$ & XXZ anisotropy & $J_x=J_y$ \\
Elliptic $R$ & XYZ anisotropy & $J_x\neq J_y\neq J_z$ generically \\
\bottomrule
\end{tabular}
\end{center}

The transverse-field Ising and XY chains studied earlier in these notes are also
anisotropic spin chains, but their solvability mechanism is different.  They
become quadratic fermion models after the Jordan--Wigner transformation.  The
generic XYZ chain does not become a free-fermion Hamiltonian under
Jordan--Wigner; it becomes an interacting fermion model.  Its exact solvability
comes instead from the commuting transfer matrices inherited from the elliptic
Yang--Baxter equation.

\section{Fermionic image and loss of free-fermion structure}
\label{sec:XYZ-JW-image}

It is useful to compare the XYZ chain with the XY mother model.  Write
\begin{equation}
\label{eq:XYZ-in-ladder-operators}
\begin{aligned}
    J_x\sigma_j^x\sigma_{j+1}^x+J_y\sigma_j^y\sigma_{j+1}^y
    &=(J_x+J_y)
    \left(\sigma_j^+\sigma_{j+1}^-+\sigma_j^-\sigma_{j+1}^+\right)\\
    &\quad
    +(J_x-J_y)
    \left(\sigma_j^+\sigma_{j+1}^++\sigma_j^-\sigma_{j+1}^-\right).
\end{aligned}
\end{equation}
The first line gives fermion hopping after Jordan--Wigner, while the second line
gives fermion pairing.  The additional $J_z$ term gives
\begin{equation}
\label{eq:ZZ-to-density-interaction}
    \sigma_j^z\sigma_{j+1}^z
    =(1-2n_j)(1-2n_{j+1})
    =1-2(n_j+n_{j+1})+4n_jn_{j+1},
\end{equation}
up to the convention for which spin state is identified with an occupied
fermion.  Therefore the XYZ chain maps schematically to
\begin{equation}
\label{eq:XYZ-JW-schematic}
    H_{\rm XYZ}
    \longmapsto
    H_{\rm hopping}+H_{\rm pairing}+4J_z\sum_j n_jn_{j+1}
    +\text{chemical-potential and boundary terms}.
\end{equation}
The density-density interaction is absent only when $J_z=0$.  Hence the generic
XYZ chain is not a quadratic BdG problem.

\begin{remark}[Two different meanings of anisotropy]
In the free-fermion XY chain, anisotropy between $x$ and $y$ exchange produces
pairing but does not destroy quadratic solvability.  In the XYZ chain, the
$J_z\sigma^z\sigma^z$ term produces a genuine fermion interaction.  Therefore
``anisotropic spin chain'' does not by itself specify the solution method.
\end{remark}

\section{Baxter functional relations}
\label{sec:eight-vertex-functional-relations}

The six-vertex model can be diagonalized by coordinate or algebraic Bethe ansatz
because the transfer matrix has a simple reference state, such as the all-up
state.  The eight-vertex transfer matrix does not preserve the number of down
spins, so that reference-state construction is not directly available.
Baxter's solution instead uses functional relations.

The central object is a second commuting family of operators $Q(u)$ satisfying a
relation of the schematic form
\begin{equation}
\label{eq:eight-vertex-TQ-schematic}
    \tau(u)Q(u)
    =
    \phi_+(u)Q(u-2\eta)+\phi_-(u)Q(u+2\eta),
\end{equation}
where $\phi_{\pm}(u)$ are known scalar functions determined by the elliptic
weights and the system size.  On a common eigenvector of $\tau(u)$ and $Q(u)$,
this becomes a scalar finite-difference equation for the transfer-matrix
eigenvalue.

\begin{principle}[Baxter's replacement for the Bethe ansatz]
For the eight-vertex model, the role of Bethe roots is played by zeros of the
$Q$-eigenvalue.  The Bethe equations arise from requiring the apparent poles or
zeros in the $TQ$ relation to be consistent.  Thus the solution is still an
exact spectral solution, but it is naturally formulated as an elliptic functional
relation rather than as a magnon scattering problem.
\end{principle}

This is why the eight-vertex and XYZ models are often regarded as more advanced
than the six-vertex and XXZ models.  They are not merely more anisotropic; they
belong to the elliptic level of Yang--Baxter integrability.

\section{Summary of the model hierarchy}
\label{sec:eight-vertex-summary-hierarchy}

The hierarchy developed in the last two chapters can be summarized as follows:
\begin{equation}
\label{eq:R-matrix-hierarchy-summary}
\begin{array}{ccccc}
\text{XXX chain}
&\longleftrightarrow&
\text{rational }R\text{-matrix}
&\longleftrightarrow&
\Delta=1,\\[3pt]
\text{XXZ chain}
&\longleftrightarrow&
\text{trigonometric }R\text{-matrix}
&\longleftrightarrow&
J_x=J_y\neq J_z,\\[3pt]
\text{XYZ chain}
&\longleftrightarrow&
\text{elliptic }R\text{-matrix}
&\longleftrightarrow&
J_x\neq J_y\neq J_z.
\end{array}
\end{equation}
The same progression appears on the classical two-dimensional side:
\begin{equation}
\label{eq:vertex-model-hierarchy-summary}
    \text{six-vertex model}
    \subset
    \text{eight-vertex model},
\end{equation}
where the inclusion is obtained by setting the pair-flip weight $d$ to zero.

\boxedpar{%
The eight-vertex model is the elliptic vertex model whose regular transfer
matrix generates the XYZ chain.  Its solvability is not free-fermionic and not
based on $U(1)$ magnon-number conservation.  It is Yang--Baxter integrability at
the elliptic level, realized through commuting transfer matrices and Baxter
functional relations.}


% -----------------------------------------------------------------------------
\part{Interacting Integrable Quantum Chains}
% -----------------------------------------------------------------------------

\chapter{Coordinate Bethe Ansatz for the XXZ Chain}
\label{chap:coordinate-bethe-ansatz-xxz}

The preceding chapters introduced the XXZ chain as the Hamiltonian limit of the
six-vertex transfer matrix.  We now solve the same quantum chain directly in real
space.  This is the coordinate Bethe ansatz.  Its logic is deliberately close to
ordinary few-body quantum mechanics: write a plane-wave superposition in each
ordered coordinate sector, impose the Schrodinger equation at particle contact,
and finally impose periodic boundary conditions.

Throughout this chapter we use the integrability normalization obtained from the
six-vertex logarithmic derivative,
\begin{equation}
\label{eq:cba-xxz-hamiltonian}
    H_{\rm XXZ}^{(+)}
    =
    J\sum_{j=1}^{L}
    \left(
        \sigma_j^x\sigma_{j+1}^x
        +\sigma_j^y\sigma_{j+1}^y
        +\Delta\bigl(\sigma_j^z\sigma_{j+1}^z-1\bigr)
    \right),
\end{equation}
with periodic boundary conditions $\sigma_{L+1}^\alpha=\sigma_1^\alpha$.  The
fully polarized state
\begin{equation}
\label{eq:cba-ferromagnetic-reference}
    \ket{0}=\ket{\uparrow\uparrow\cdots\uparrow}
\end{equation}
is an eigenstate with zero energy.  It is not necessarily the ground state of the
antiferromagnetic chain, but it is the natural pseudovacuum for Bethe ansatz.
Down spins above this reference state will be called magnons.

The three goals of the chapter are:
\begin{enumerate}[label=(\roman*)]
    \item solve the one-magnon problem and identify the bare dispersion;
    \item derive the two-magnon scattering amplitude from the contact equation;
    \item assemble the $M$-magnon Bethe wave function and obtain the Bethe
    equations.
\end{enumerate}
The result is the first genuinely interacting exact solution in these notes: the
energy is additive in Bethe momenta, while the allowed momenta are determined by
many-body quantization conditions containing all two-body phase shifts.

\section{Conserved magnon number and coordinate basis}
\label{sec:cba-coordinate-basis}

Using
\begin{equation}
    \sigma_j^x\sigma_{j+1}^x+\sigma_j^y\sigma_{j+1}^y
    =
    2\left(\sigma_j^+\sigma_{j+1}^-+\sigma_j^-\sigma_{j+1}^+\right),
\end{equation}
one sees that the $XY$ part only exchanges neighboring up and down spins.  The
total magnetization
\begin{equation}
\label{eq:cba-total-magnetization}
    S^z_{\rm tot}=\frac12\sum_{j=1}^{L}\sigma_j^z
\end{equation}
therefore commutes with $H_{\rm XXZ}^{(+)}$.  Equivalently, the number of down
spins is conserved.

In the sector with $M$ down spins, a convenient basis is
\begin{equation}
\label{eq:cba-magnon-coordinate-basis}
    \ket{x_1,\ldots,x_M}
    =
    \sigma_{x_1}^-\cdots\sigma_{x_M}^-\ket{0},
    \qquad
    1\leq x_1<x_2<\cdots<x_M\leq L .
\end{equation}
A general state in this sector is
\begin{equation}
\label{eq:cba-general-M-magnon-state}
    \ket{\Psi_M}
    =
    \sum_{1\leq x_1<\cdots<x_M\leq L}
    \Psi(x_1,\ldots,x_M)\ket{x_1,\ldots,x_M} .
\end{equation}
The coordinate Bethe ansatz is an ansatz for the amplitudes
$\Psi(x_1,\ldots,x_M)$.

It is useful to separate two kinds of equations.  When no two magnons are
neighbors, the Schrodinger equation is just a discrete free-particle equation.
When $x_{a+1}=x_a+1$, the two magnons touch, and the nearest-neighbor Ising term
produces a contact condition.  This contact condition is the source of the
scattering phase.

\section{One-magnon sector}
\label{sec:cba-one-magnon-sector}

A one-magnon state has the form
\begin{equation}
\label{eq:cba-one-magnon-state}
    \ket{\Psi_1}
    =
    \sum_{x=1}^{L}\psi(x)\ket{x},
    \qquad
    \ket{x}=\sigma_x^-\ket{0} .
\end{equation}
Acting with \cref{eq:cba-xxz-hamiltonian}, one obtains
\begin{equation}
\label{eq:cba-one-magnon-difference-equation}
    E\psi(x)
    =
    2J\psi(x-1)+2J\psi(x+1)-4J\Delta\psi(x),
\end{equation}
where $x$ is understood modulo $L$.  The hopping amplitude is $2J$, and the
potential term $-4J\Delta$ comes from the two domain walls adjacent to the down
spin.

The plane wave
\begin{equation}
\label{eq:cba-one-magnon-plane-wave}
    \psi_k(x)=\ee^{\ii kx}
\end{equation}
is therefore an eigenfunction with dispersion
\begin{equation}
\label{eq:cba-one-magnon-dispersion}
    \varepsilon(k)
    =
    4J(\cos k-\Delta).
\end{equation}
Periodic boundary conditions give
\begin{equation}
\label{eq:cba-one-magnon-momentum-quantization}
    \ee^{\ii kL}=1,
    \qquad
    k=\frac{2\pi n}{L},
    \qquad
    n\in\ZZ .
\end{equation}

\begin{remark}[Energy convention]
The dispersion in \cref{eq:cba-one-magnon-dispersion} is written relative to the
fully polarized reference state, whose energy was fixed to zero by the constant
in \cref{eq:cba-xxz-hamiltonian}.  Reversing the overall sign of the Hamiltonian
or removing the constant shifts and rescales the formulas, but does not change
the Bethe wave functions or the scattering phase.
\end{remark}

\section{Two-magnon scattering}
\label{sec:cba-two-magnon-scattering}

The first nontrivial part of the Bethe ansatz appears for two magnons.  For
coordinates $x_1<x_2$, take
\begin{equation}
\label{eq:cba-two-magnon-wavefunction}
    \Psi(x_1,x_2)
    =
    A_{12}\,\ee^{\ii(k_1x_1+k_2x_2)}
    +
    A_{21}\,\ee^{\ii(k_2x_1+k_1x_2)} .
\end{equation}
The two terms describe the same two momenta assigned to the two ordered
coordinates in two different ways.  Away from contact, namely for
$x_2>x_1+1$, the wave function satisfies the free difference equation and has
energy
\begin{equation}
\label{eq:cba-two-magnon-energy}
    E(k_1,k_2)
    =
    4J(\cos k_1+\cos k_2-2\Delta).
\end{equation}
The energy is additive in the two bare one-magnon energies.

At contact, $x_2=x_1+1$, two of the four domain walls of separated magnons have
annihilated.  The Schrodinger equation becomes
\begin{equation}
\label{eq:cba-two-magnon-contact-equation}
    (E+4J\Delta)\Psi(x,x+1)
    =
    2J\Psi(x-1,x+1)+2J\Psi(x,x+2).
\end{equation}
Substituting \cref{eq:cba-two-magnon-wavefunction} into
\cref{eq:cba-two-magnon-contact-equation} fixes the ratio of amplitudes:
\begin{equation}
\label{eq:cba-two-magnon-scattering-amplitude}
    \frac{A_{21}}{A_{12}}
    =
    S(k_1,k_2)
    =
    -
    \frac{1+\ee^{\ii(k_1+k_2)}-2\Delta\ee^{\ii k_1}}
         {1+\ee^{\ii(k_1+k_2)}-2\Delta\ee^{\ii k_2}} .
\end{equation}
This is the two-magnon scattering amplitude of the XXZ chain in the convention
of \cref{eq:cba-xxz-hamiltonian}.

Several checks are immediate.  First,
\begin{equation}
\label{eq:cba-S-unitarity}
    S(k_1,k_2)S(k_2,k_1)=1,
\end{equation}
so exchanging two magnons twice gives back the original amplitude.  Second, at
$\Delta=0$,
\begin{equation}
    S(k_1,k_2)=-1,
\end{equation}
which is the free-fermion sign of the XX chain.  Thus the XX point is recovered
as the noninteracting limit of the coordinate Bethe ansatz.

\begin{principle}[Interaction as a boundary condition]
In the coordinate Bethe ansatz, the interaction does not destroy plane waves in
each ordered region.  Instead, it appears as a matching condition at contact,
which relates different plane-wave amplitudes by a two-body scattering phase.
\end{principle}

\section{The many-magnon coordinate Bethe ansatz}
\label{sec:cba-many-magnon-wavefunction}

For $M$ magnons, the ordered coordinate region is
\begin{equation}
    1\leq x_1<x_2<\cdots<x_M\leq L .
\end{equation}
The coordinate Bethe ansatz is
\begin{equation}
\label{eq:cba-M-magnon-wavefunction}
    \Psi(x_1,\ldots,x_M)
    =
    \sum_{P\in S_M}A(P)
    \exp\left(\ii\sum_{a=1}^{M}k_{P_a}x_a\right),
\end{equation}
where $S_M$ is the permutation group of $M$ objects.  The amplitude $A(P)$ is
attached to the assignment of momentum $k_{P_a}$ to the ordered coordinate
$x_a$.

Whenever all magnons are separated, the Schrodinger equation gives the additive
energy
\begin{equation}
\label{eq:cba-M-magnon-energy}
    E(k_1,
    \ldots,k_M)
    =
    \sum_{a=1}^{M}\varepsilon(k_a)
    =
    4J\sum_{a=1}^{M}(\cos k_a-\Delta).
\end{equation}
Whenever two adjacent coordinates touch, the same two-body contact equation as in
the two-magnon problem applies.  Therefore adjacent permutations obey
\begin{equation}
\label{eq:cba-adjacent-permutation-relation}
    A(Ps_a)
    =
    S(k_{P_a},k_{P_{a+1}})A(P),
    \qquad
    a=1,\ldots,M-1,
\end{equation}
where $s_a$ exchanges the entries in positions $a$ and $a+1$.

Because the scattering amplitude is a scalar phase, consistency of different
ways of reducing a permutation follows from ordinary multiplication.  In models
with internal degrees of freedom, the analogous consistency condition becomes the
Yang--Baxter equation.  This is why the coordinate Bethe ansatz here is the
real-space shadow of the six-vertex Yang--Baxter structure.

\section{Periodic boundary conditions and Bethe equations}
\label{sec:cba-bethe-equations-momentum-form}

The remaining step is to impose periodic boundary conditions.  Move magnon $j$
once around the ring.  It accumulates a free propagation phase $\ee^{\ii k_jL}$
and scatters once with every other magnon.  With the scattering convention of
\cref{eq:cba-two-magnon-scattering-amplitude}, the result is
\begin{equation}
\label{eq:cba-momentum-bethe-equations}
    \ee^{\ii k_jL}
    =
    \prod_{\substack{\ell=1\\ \ell\neq j}}^{M}
    S(k_j,k_\ell),
    \qquad
    j=1,\ldots,M .
\end{equation}
Together with the energy formula \cref{eq:cba-M-magnon-energy}, these equations
solve the finite periodic XXZ chain in each fixed-magnon sector, up to the usual
questions of completeness and string organization.

Writing
\begin{equation}
    S(k_j,k_\ell)=\ee^{\ii\theta(k_j,k_\ell)},
\end{equation}
one obtains the logarithmic form
\begin{equation}
\label{eq:cba-logarithmic-bethe-equations}
    Lk_j
    =
    2\pi I_j
    +
    \sum_{\substack{\ell=1\\ \ell\neq j}}^{M}
    \theta(k_j,k_\ell),
    \qquad
    I_j\in\ZZ\ \text{or}\ \ZZ+\frac12 .
\end{equation}
The integer versus half-integer choice depends on $M$, $L$, the branch choice of
$\theta$, and the boundary-condition convention.  The exponential form
\cref{eq:cba-momentum-bethe-equations} is the unambiguous statement.

\begin{remark}[What is solved by the Bethe equations]
The Bethe equations do not say that the interacting chain has free momenta.  They
say that eigenstates are still labelled by rapidities or momenta, but those
labels are quantized collectively.  Each $k_j$ is shifted by the phase shifts
against all other magnons.
\end{remark}

\section{Rapidity parametrizations}
\label{sec:cba-rapidity-parametrizations}

The momentum-space scattering amplitude is explicit but not the most useful form
for analysis.  The standard Bethe equations become simpler after introducing a
rapidity variable.  The appropriate parametrization depends on the anisotropy
regime.

\subsection*{Gapless regime: $-1<\Delta<1$}

Let
\begin{equation}
\label{eq:cba-gapless-parametrization}
    \Delta=\cos\gamma,
    \qquad
    0<\gamma<\pi .
\end{equation}
Define the rapidity $\lambda$ by
\begin{equation}
\label{eq:cba-gapless-momentum-rapidity-map}
    \ee^{\ii k(\lambda)}
    =
    \frac{\sinh(\lambda+\ii\gamma/2)}
         {\sinh(\lambda-\ii\gamma/2)} .
\end{equation}
Then \cref{eq:cba-momentum-bethe-equations} becomes
\begin{equation}
\label{eq:cba-gapless-xxz-bethe-equations}
    \left[
    \frac{\sinh(\lambda_j+\ii\gamma/2)}
         {\sinh(\lambda_j-\ii\gamma/2)}
    \right]^L
    =
    \prod_{\substack{\ell=1\\ \ell\neq j}}^{M}
    \frac{\sinh(\lambda_j-\lambda_\ell+\ii\gamma)}
         {\sinh(\lambda_j-\lambda_\ell-\ii\gamma)} .
\end{equation}
The energy becomes
\begin{equation}
\label{eq:cba-gapless-energy-rapidity}
    E
    =
    -4J\sum_{j=1}^{M}
    \frac{\sin^2\gamma}{\cosh(2\lambda_j)-\cos\gamma} .
\end{equation}
The total lattice momentum is
\begin{equation}
\label{eq:cba-total-momentum}
    P
    =
    \sum_{j=1}^{M}k(\lambda_j)
    \quad (\operatorname{mod} 2\pi).
\end{equation}

\subsection*{Massive antiferromagnetic regime: $\Delta>1$}

Let
\begin{equation}
\label{eq:cba-massive-parametrization}
    \Delta=\cosh\eta,
    \qquad
    \eta>0 .
\end{equation}
A convenient rapidity map is
\begin{equation}
\label{eq:cba-massive-momentum-rapidity-map}
    \ee^{\ii k(\lambda)}
    =
    \frac{\sin(\lambda+\ii\eta/2)}
         {\sin(\lambda-\ii\eta/2)} .
\end{equation}
The Bethe equations are
\begin{equation}
\label{eq:cba-massive-xxz-bethe-equations}
    \left[
    \frac{\sin(\lambda_j+\ii\eta/2)}
         {\sin(\lambda_j-\ii\eta/2)}
    \right]^L
    =
    \prod_{\substack{\ell=1\\ \ell\neq j}}^{M}
    \frac{\sin(\lambda_j-\lambda_\ell+\ii\eta)}
         {\sin(\lambda_j-\lambda_\ell-\ii\eta)} .
\end{equation}
The energy is
\begin{equation}
\label{eq:cba-massive-energy-rapidity}
    E
    =
    -4J\sum_{j=1}^{M}
    \frac{\sinh^2\eta}{\cosh\eta-\cos(2\lambda_j)} .
\end{equation}
In this regime complex rapidities organize into string patterns in the
thermodynamic limit.  Strings describe bound states of magnons and are essential
for thermodynamics, but the finite-size Bethe equations above are the starting
point.

\subsection*{Isotropic XXX point}

At $\Delta=1$, one may use the rational rapidity variable
\begin{equation}
\label{eq:cba-xxx-momentum-rapidity-map}
    \ee^{\ii k(\lambda)}
    =
    \frac{\lambda+\ii/2}{\lambda-\ii/2} .
\end{equation}
The Bethe equations reduce to
\begin{equation}
\label{eq:cba-xxx-bethe-equations}
    \left(
    \frac{\lambda_j+\ii/2}{\lambda_j-\ii/2}
    \right)^L
    =
    \prod_{\substack{\ell=1\\ \ell\neq j}}^{M}
    \frac{\lambda_j-\lambda_\ell+\ii}
         {\lambda_j-\lambda_\ell-\ii} ,
\end{equation}
and the energy is
\begin{equation}
\label{eq:cba-xxx-energy}
    E
    =
    -2J\sum_{j=1}^{M}
    \frac{1}{\lambda_j^2+1/4} .
\end{equation}
This is the rational limit of the trigonometric XXZ equations.

\section{Special limits and physical interpretation}
\label{sec:cba-special-limits-interpretation}

\begin{example}[XX point]
At $\Delta=0$, the scattering amplitude becomes $S=-1$.  The Bethe equations are
therefore
\begin{equation}
    \ee^{\ii k_jL}=(-1)^{M-1} .
\end{equation}
This is precisely the parity-dependent boundary condition of Jordan--Wigner
fermions in the $M$-particle sector.  Thus the coordinate Bethe ansatz reduces to
free fermions at the XX point.
\end{example}

\begin{example}[Two-magnon bound states]
For $\Delta>1$, the scattering amplitude has poles in the complex momentum or
rapidity plane.  Normalizable two-magnon wave functions can then have complex
conjugate momenta, producing an exponentially localized relative coordinate.
These states are the simplest examples of Bethe strings.  In the many-body
problem, strings encode bound-state quasiparticles.
\end{example}

\begin{remark}[Pseudovacuum versus physical ground state]
The coordinate Bethe ansatz above builds states by flipping spins above
$\ket{0}=\ket{\uparrow\cdots\uparrow}$.  For the antiferromagnetic convention
with $J>0$ and $\Delta>0$, this reference state is not the physical ground state
in the zero-magnetization sector.  The true ground state is obtained by solving
the Bethe equations with an extensive number of rapidities.  The pseudovacuum is
used because it makes the scattering construction algebraically simple.
\end{remark}

\section{Relation to the algebraic Bethe ansatz}
\label{sec:cba-relation-to-aba}

The coordinate Bethe ansatz diagonalizes the Hamiltonian by solving the real-space
Schrodinger equation.  The algebraic Bethe ansatz diagonalizes the whole transfer
matrix family using the $R$-matrix algebra.  The two methods produce the same
Bethe equations, but they organize the calculation differently:
\begin{center}
\begin{tabular}{@{}L{0.30\linewidth}L{0.30\linewidth}L{0.30\linewidth}@{}}
\toprule
Method & Basic object & Origin of Bethe equations \\
\midrule
Coordinate Bethe ansatz & Wave function
$\Psi(x_1,\ldots,x_M)$ & Periodic boundary conditions plus two-body scattering \\
Algebraic Bethe ansatz & Bethe vector
$B(\lambda_1)\cdots B(\lambda_M)\ket{0}$ & Cancellation of unwanted terms in the RTT algebra \\
\bottomrule
\end{tabular}
\end{center}
The coordinate method is more physically transparent: it shows directly that the
XXZ chain is solvable because many-body scattering factorizes into two-body
phases.  The algebraic method is more structural: it explains why this
factorization is guaranteed by the Yang--Baxter equation and gives direct access
to transfer-matrix eigenvalues and higher conserved charges.

\section{Summary}
\label{sec:cba-summary}

The coordinate Bethe ansatz for the XXZ chain can be summarized as
\begin{equation}
\label{eq:cba-summary-chain}
\boxed{
\begin{gathered}
    H_{\rm XXZ}^{(+)}
    \quad\Longrightarrow\quad
    M\text{-magnon coordinate sectors}
    \\
    \Longrightarrow\quad
    \Psi(x_1,\ldots,x_M)
    =
    \sum_{P\in S_M}A(P)
    \exp\left(\ii\sum_a k_{P_a}x_a\right)
    \\
    \Longrightarrow\quad
    A(Ps_a)=S(k_{P_a},k_{P_{a+1}})A(P)
    \\
    \Longrightarrow\quad
    \ee^{\ii k_jL}=\prod_{\ell\neq j}S(k_j,k_\ell),
    \qquad
    E=4J\sum_j(\cos k_j-\Delta).
\end{gathered}}
\end{equation}
The one-magnon sector gives the bare dispersion.  The two-magnon contact equation
gives the scattering phase.  The many-magnon solution works because all
collisions factorize into the same two-body scattering amplitude.  This is the
coordinate-space form of the Yang--Baxter integrability already encountered in
the six-vertex model.


\chapter{Algebraic Bethe Ansatz}
\label{chap:algebraic-bethe-ansatz}

The coordinate Bethe ansatz solves the XXZ chain by constructing wave functions in
ordered magnon-coordinate sectors.  The algebraic Bethe ansatz solves the same
problem from the transfer-matrix side.  Instead of writing a many-body wave
function explicitly, one uses the Yang--Baxter equation to build an operator
algebra.  Bethe states are then obtained by applying a creation operator
$B(u)$, extracted from the monodromy matrix, to a simple reference state.

The conceptual chain is
\[
\begin{gathered}
    R\text{-matrix satisfying YBE}
    \Longrightarrow
    \text{RTT algebra for }A(u),B(u),C(u),D(u) \\
    \Longrightarrow
    \text{Bethe vectors }B(\lambda_1)\cdots B(\lambda_M)\ket{0} \\
    \Longrightarrow
    \text{transfer-matrix eigenvalues and Bethe equations}.
\end{gathered}
\]
Thus the algebraic Bethe ansatz diagonalizes not only one Hamiltonian, but the
whole commuting transfer-matrix family.  The Hamiltonian eigenvalues are obtained
afterwards by taking the logarithmic derivative of the transfer-matrix
eigenvalue.

Throughout this chapter we use the homogeneous six-vertex convention of
\cref{chap:six-vertex-xxz}.  The local $R$-matrix is
\begin{equation}
\label{eq:aba-six-vertex-R}
    R(u)=
    \begin{pmatrix}
        \sinh(u+\eta)&0&0&0\\
        0&\sinh u&\sinh\eta&0\\
        0&\sinh\eta&\sinh u&0\\
        0&0&0&\sinh(u+\eta)
    \end{pmatrix},
    \qquad
    \Delta=\cosh\eta .
\end{equation}
The trigonometric critical regime follows by analytic continuation
$\eta\mapsto \ii\gamma$, giving $\Delta=\cos\gamma$.  We keep the hyperbolic
notation because it makes the algebraic formulas compact.

\section{Monodromy matrix}
\label{sec:aba-monodromy-matrix}

Let $V_a\simeq\CC^2$ be an auxiliary spin-$1/2$ space and
\begin{equation}
\label{eq:aba-quantum-space}
    \calH_L=V_1\otimes\cdots\otimes V_L,
    \qquad V_j\simeq\CC^2
\end{equation}
be the quantum spin-chain Hilbert space.  The monodromy matrix is the ordered
product
\begin{equation}
\label{eq:aba-monodromy-definition}
    M_a(u)
    =R_{aL}(u)R_{a,L-1}(u)\cdots R_{a1}(u).
\end{equation}
It is a $2\times2$ matrix in the auxiliary space, whose entries are operators on
$\calH_L$:
\begin{equation}
\label{eq:aba-monodromy-ABCD}
    M_a(u)=
    \begin{pmatrix}
        A(u)&B(u)\\
        C(u)&D(u)
    \end{pmatrix}_a .
\end{equation}
The transfer matrix is the auxiliary trace
\begin{equation}
\label{eq:aba-transfer-matrix}
    \tau(u)=\Tr_a M_a(u)=A(u)+D(u).
\end{equation}
The entries $A,B,C,D$ are not four independent Hamiltonians.  They are
noncommuting operator-valued functions constrained by the RTT algebra below.

It is useful to write the local $R$-operator as a matrix in the auxiliary space
with quantum-site operators as entries.  In the convention of
\cref{eq:aba-six-vertex-R},
\begin{equation}
\label{eq:aba-local-L-operator}
    R_{aj}(u)=
    \begin{pmatrix}
        a_R(u)e^{11}_j+b_R(u)e^{22}_j
        &c_R e^{21}_j\\
        c_R e^{12}_j
        &b_R(u)e^{11}_j+a_R(u)e^{22}_j
    \end{pmatrix}_a,
\end{equation}
where
\begin{equation}
\label{eq:aba-local-weights}
    a_R(u)=\sinh(u+\eta),
    \qquad
    b_R(u)=\sinh u,
    \qquad
    c_R=\sinh\eta,
\end{equation}
and $e^{\alpha\beta}_j=\ket{\alpha}_j\bra{\beta}_j$.  With
$\ket{\uparrow}=\ket{1}$ and $\ket{\downarrow}=\ket{2}$, one has
$e^{21}=\sigma^-$ and $e^{12}=\sigma^+$.

Choose the ferromagnetic reference state
\begin{equation}
\label{eq:aba-reference-state}
    \ket{0}=\ket{\uparrow\uparrow\cdots\uparrow} .
\end{equation}
On this state every local $R_{aj}(u)$ becomes upper triangular in auxiliary
space:
\begin{equation}
\label{eq:aba-local-triangular-action}
    R_{aj}(u)\ket{\uparrow}_j
    =
    \begin{pmatrix}
        a_R(u)\ket{\uparrow}_j & c_R\ket{\downarrow}_j\\
        0&b_R(u)\ket{\uparrow}_j
    \end{pmatrix}_a .
\end{equation}
Multiplying over all sites preserves triangularity.  Hence
\begin{equation}
\label{eq:aba-pseudovacuum-action}
    C(u)\ket{0}=0,
    \qquad
    A(u)\ket{0}=\alpha(u)\ket{0},
    \qquad
    D(u)\ket{0}=\delta(u)\ket{0},
\end{equation}
with vacuum eigenvalues
\begin{equation}
\label{eq:aba-vacuum-eigenvalues}
    \alpha(u)=a_R(u)^L=\sinh^L(u+\eta),
    \qquad
    \delta(u)=b_R(u)^L=\sinh^L u .
\end{equation}
The operator $B(u)$ lowers one or more spins when expanded as an operator on the
chain.  Its leading action on the reference state is a one-magnon superposition,
so $B(u)$ is the algebraic creation operator for Bethe quasiparticles.

\begin{remark}[Pseudovacuum]
The word ``vacuum'' here is algebraic.  For the antiferromagnetic XXZ chain the
physical ground state is usually not $\ket{0}$.  The state $\ket{0}$ is used
because it is a highest-weight state for the RTT algebra: $C(u)$ annihilates it,
while $A(u)$ and $D(u)$ act diagonally.
\end{remark}

\section{RTT relation and commuting transfer matrices}
\label{sec:aba-rtt-commuting-transfer}

The Yang--Baxter equation for the six-vertex $R$-matrix implies the local
intertwining relation
\begin{equation}
\label{eq:aba-local-RLL}
    R_{ab}(u-v)R_{aj}(u)R_{bj}(v)
    =
    R_{bj}(v)R_{aj}(u)R_{ab}(u-v).
\end{equation}
Multiplying this identity over all sites gives the RTT relation
\begin{equation}
\label{eq:aba-RTT-relation}
    R_{ab}(u-v)M_a(u)M_b(v)
    =
    M_b(v)M_a(u)R_{ab}(u-v).
\end{equation}
This is the defining algebraic relation of the algebraic Bethe ansatz.  It is an
operator identity on
\begin{equation}
    V_a\otimes V_b\otimes\calH_L .
\end{equation}
The two auxiliary spaces $a$ and $b$ are bookkeeping devices; the nontrivial
operators act on the same quantum chain $\calH_L$.

\begin{theorem}[Commuting transfer matrices]
\label{thm:aba-commuting-transfer-matrices}
The RTT relation implies
\begin{equation}
\label{eq:aba-transfer-commutativity}
    [\tau(u),\tau(v)]=0,
    \qquad
    \text{for all }u,v .
\end{equation}
\end{theorem}

\begin{proof}
Take $\Tr_a\Tr_b$ of \cref{eq:aba-RTT-relation}.  By cyclicity of the trace in
the auxiliary spaces, $R_{ab}(u-v)$ can be moved through the trace and canceled.
This gives
\begin{equation}
    \Tr_a M_a(u)\,\Tr_b M_b(v)
    =
    \Tr_b M_b(v)\,\Tr_a M_a(u),
\end{equation}
which is exactly \cref{eq:aba-transfer-commutativity}.
\end{proof}

The commutativity is the spectral reason for integrability.  Expanding around the
regular point gives mutually commuting charges,
\begin{equation}
\label{eq:aba-conserved-charges-from-transfer}
    \log\tau(u)=\sum_{n\geq0}u^n Q_n,
    \qquad
    [Q_m,Q_n]=0 .
\end{equation}
The XXZ Hamiltonian is one of these charges:
\begin{equation}
\label{eq:aba-xxz-from-log-transfer}
    H_{\rm XXZ}^{(+)}
    =
    2J\sinh\eta
    \left.\frac{\dd}{\dd u}\log\tau(u)\right|_{u=0}
    -2J\cosh\eta\,L,
\end{equation}
which is the same normalization as
\begin{equation}
\label{eq:aba-xxz-hamiltonian}
    H_{\rm XXZ}^{(+)}
    =
    J\sum_{j=1}^{L}
    \left(
        \sigma_j^x\sigma_{j+1}^x
        +\sigma_j^y\sigma_{j+1}^y
        +\Delta(\sigma_j^z\sigma_{j+1}^z-1)
    \right),
    \qquad
    \Delta=\cosh\eta .
\end{equation}
The constant fixes the energy of the fully polarized reference state to zero.

\section{The fundamental commutation relations}
\label{sec:aba-fundamental-commutation-relations}

The RTT relation contains many component equations.  The algebraic Bethe ansatz
only needs the relations that move $A(u)$ and $D(u)$ through products of $B$'s.
Introduce
\begin{equation}
\label{eq:aba-f-and-g-functions}
    f(u)=\frac{\sinh(u+\eta)}{\sinh u},
    \qquad
    g(u)=\frac{\sinh\eta}{\sinh u} .
\end{equation}
With the monodromy order \cref{eq:aba-monodromy-definition}, the needed RTT
relations are
\begin{align}
\label{eq:aba-BB-commutation}
    B(u)B(v)&=B(v)B(u),\\
\label{eq:aba-AB-commutation}
    A(u)B(v)
    &=
    f(v-u)B(v)A(u)-g(v-u)B(u)A(v),\\
\label{eq:aba-DB-commutation}
    D(u)B(v)
    &=
    f(u-v)B(v)D(u)-g(u-v)B(u)D(v).
\end{align}
The first identity says that Bethe creation operators commute.  The second and
third identities have two parts.  The first terms preserve the set of Bethe
rapidities and will produce the transfer-matrix eigenvalue.  The second terms
replace one rapidity $v$ by the spectral parameter $u$; these are the unwanted
terms whose cancellation gives the Bethe equations.

\begin{remark}[Convention dependence]
If one reverses the monodromy ordering, changes $R(u-v)$ to $R(v-u)$, or shifts
the spectral parameter by $\eta/2$, the displayed formulas may be inverted or
shifted.  The invariant content is the same: the ratio of vacuum eigenvalues is
balanced against the product of two-body scattering factors.
\end{remark}

\section{Bethe vectors}
\label{sec:aba-bethe-vectors}

For rapidities $\lambda_1,
\ldots,\lambda_M$, define the Bethe vector
\begin{equation}
\label{eq:aba-bethe-vector}
    \ket{\lambda_1,\ldots,\lambda_M}
    =
    B(\lambda_1)B(\lambda_2)\cdots B(\lambda_M)\ket{0} .
\end{equation}
Because the $B$'s commute, this vector is symmetric under permutations of the
rapidities.  It lies in the sector with $M$ down spins:
\begin{equation}
\label{eq:aba-magnetization-sector}
    S^z_{\rm tot}\ket{\lambda_1,\ldots,\lambda_M}
    =
    \left(\frac{L}{2}-M\right)
    \ket{\lambda_1,\ldots,\lambda_M} .
\end{equation}
This matches the coordinate Bethe ansatz, where $M$ was the number of magnons.

\begin{example}[One Bethe excitation]
For $M=1$,
\begin{equation}
    \ket{\lambda}=B(\lambda)\ket{0} .
\end{equation}
Using \cref{eq:aba-AB-commutation,eq:aba-DB-commutation},
\begin{equation}
\label{eq:aba-one-particle-action}
\begin{aligned}
    \tau(u)B(\lambda)\ket{0}
    ={}&
    \left[
        \alpha(u)f(\lambda-u)+\delta(u)f(u-\lambda)
    \right]B(\lambda)\ket{0}
    \\
    &
    +g(\lambda-u)
    \left[
        \delta(\lambda)-\alpha(\lambda)
    \right]B(u)\ket{0} .
\end{aligned}
\end{equation}
The first line is proportional to the original vector.  The second line is
unwanted.  It vanishes exactly when
\begin{equation}
\label{eq:aba-one-particle-bethe-equation}
    \alpha(\lambda)=\delta(\lambda),
    \qquad
    \left[\frac{\sinh(\lambda+\eta)}{\sinh\lambda}\right]^L=1 .
\end{equation}
This is the one-magnon momentum quantization condition in algebraic form.
\end{example}

The many-particle calculation is the repeated version of the same mechanism.
Move $A(u)$ and $D(u)$ through the string of $B$'s until they reach the
pseudovacuum.  The wanted part keeps all rapidities fixed.  Each unwanted term
replaces one $\lambda_j$ by $u$.

\section{Transfer-matrix eigenvalue and Bethe equations}
\label{sec:aba-transfer-eigenvalue-bethe-equations}

Acting with $\tau(u)=A(u)+D(u)$ on the Bethe vector gives the off-shell form
\begin{equation}
\label{eq:aba-off-shell-action}
    \tau(u)\ket{\lambda_1,
\ldots,\lambda_M}
    =
    \Lambda(u\mid\boldsymbol{\lambda})
    \ket{\lambda_1,
\ldots,\lambda_M}
    +
    \sum_{j=1}^{M}
    \mathcal{F}_j(u\mid\boldsymbol{\lambda})
    B(u)\prod_{\ell\neq j}B(\lambda_\ell)\ket{0} .
\end{equation}
Here the wanted coefficient is
\begin{equation}
\label{eq:aba-transfer-eigenvalue}
\boxed{
\begin{aligned}
    \Lambda(u\mid\boldsymbol{\lambda})
    ={}&
    \alpha(u)
    \prod_{j=1}^{M}f(\lambda_j-u)
    +
    \delta(u)
    \prod_{j=1}^{M}f(u-\lambda_j)
    \\
    ={}&
    \sinh^L(u+\eta)
    \prod_{j=1}^{M}
    \frac{\sinh(\lambda_j-u+\eta)}{\sinh(\lambda_j-u)}
    \\
    &+
    \sinh^L u
    \prod_{j=1}^{M}
    \frac{\sinh(u-\lambda_j+\eta)}{\sinh(u-\lambda_j)} .
\end{aligned}}
\end{equation}
Up to a nonzero scalar factor, the unwanted coefficient is
\begin{equation}
\label{eq:aba-unwanted-coefficient}
    \mathcal{F}_j(u\mid\boldsymbol{\lambda})
    \propto
    \delta(\lambda_j)
    \prod_{\ell\neq j}f(\lambda_j-\lambda_\ell)
    -
    \alpha(\lambda_j)
    \prod_{\ell\neq j}f(\lambda_\ell-\lambda_j) .
\end{equation}
Therefore the Bethe vector is an eigenvector of $\tau(u)$ for all $u$ if
\begin{equation}
\label{eq:aba-bethe-equations-abstract}
    \frac{\alpha(\lambda_j)}{\delta(\lambda_j)}
    =
    \prod_{\substack{\ell=1\\ \ell\neq j}}^{M}
    \frac{f(\lambda_j-\lambda_\ell)}{f(\lambda_\ell-\lambda_j)},
    \qquad
    j=1,
\ldots,M .
\end{equation}
Substituting the six-vertex functions gives the algebraic Bethe equations
\begin{equation}
\label{eq:aba-xxz-bethe-equations}
\boxed{
    \left[
    \frac{\sinh(\lambda_j+\eta)}{\sinh\lambda_j}
    \right]^L
    =
    \prod_{\substack{\ell=1\\ \ell\neq j}}^{M}
    \frac{\sinh(\lambda_j-\lambda_\ell+\eta)}
         {\sinh(\lambda_j-\lambda_\ell-\eta)} ,
    \qquad
    j=1,
\ldots,M .
}
\end{equation}
When these equations hold, all unwanted terms in \cref{eq:aba-off-shell-action}
vanish, and
\begin{equation}
\label{eq:aba-on-shell-transfer-action}
    \tau(u)\ket{\lambda_1,
\ldots,\lambda_M}
    =
    \Lambda(u\mid\boldsymbol{\lambda})
    \ket{\lambda_1,
\ldots,\lambda_M} .
\end{equation}

\begin{principle}[Bethe equations as analyticity conditions]
The transfer-matrix eigenvalue \cref{eq:aba-transfer-eigenvalue} appears to have
simple poles at $u=\lambda_j$.  The Bethe equations are exactly the conditions
that the residues cancel.  Thus the same equations can be read either as
cancellation of unwanted vectors or as analyticity of the transfer-matrix
eigenvalue.
\end{principle}

\section{Hamiltonian eigenvalues}
\label{sec:aba-hamiltonian-eigenvalues}

The transfer-matrix eigenvalue contains the eigenvalues of all commuting charges.
Using \cref{eq:aba-xxz-from-log-transfer}, the XXZ energy is
\begin{equation}
\label{eq:aba-energy-from-transfer-eigenvalue}
    E(\boldsymbol{\lambda})
    =
    2J\sinh\eta
    \left.
    \frac{\dd}{\dd u}
    \log\Lambda(u\mid\boldsymbol{\lambda})
    \right|_{u=0}
    -2J\cosh\eta\,L .
\end{equation}
For the eigenvalue \cref{eq:aba-transfer-eigenvalue}, this gives
\begin{equation}
\label{eq:aba-energy-rapidity-form}
\boxed{
    E(\boldsymbol{\lambda})
    =
    2J\sinh^2\eta
    \sum_{j=1}^{M}
    \frac{1}{\sinh\lambda_j\,\sinh(\lambda_j+\eta)} .
}
\end{equation}
The empty Bethe state $M=0$ has $E=0$, as required by the constant chosen in
\cref{eq:aba-xxz-hamiltonian}.

The lattice momentum is obtained from the regular-point transfer matrix.  Since
$R(0)=\sinh\eta\,P$, one has
\begin{equation}
\label{eq:aba-transfer-at-zero-translation}
    \tau(0)=(\sinh\eta)^L U,
\end{equation}
where $U$ is the one-site translation operator, with the orientation determined
by the monodromy ordering.  Therefore
\begin{equation}
\label{eq:aba-momentum-from-roots}
    U\ket{\boldsymbol{\lambda}}
    =
    \left[
    \prod_{j=1}^{M}
    \frac{\sinh(\lambda_j+\eta)}{\sinh\lambda_j}
    \right]
    \ket{\boldsymbol{\lambda}} .
\end{equation}
Depending on whether one defines $U$ as a left or right shift, this eigenvalue is
$\ee^{\ii P}$ or $\ee^{-\ii P}$.  This is only an orientation convention.

If one introduces a one-particle momentum by
\begin{equation}
\label{eq:aba-momentum-rapidity-map}
    \ee^{\ii k(\lambda)}
    =
    \frac{\sinh(\lambda+\eta)}{\sinh\lambda},
\end{equation}
then
\begin{equation}
\label{eq:aba-energy-momentum-form}
    \frac{2J\sinh^2\eta}{\sinh\lambda\,\sinh(\lambda+\eta)}
    =
    4J\bigl(\cos k(\lambda)-\cosh\eta\bigr),
\end{equation}
so \cref{eq:aba-energy-rapidity-form} reproduces the coordinate Bethe-ansatz
energy
\begin{equation}
\label{eq:aba-coordinate-energy-recovered}
    E=4J\sum_{j=1}^{M}\bigl(\cos k_j-\Delta\bigr).
\end{equation}
If the opposite translation orientation is used, $k$ is replaced by $-k$ and the
energy is unchanged.

\section[Symmetric rapidities and CBA]{Symmetric rapidities and comparison with the coordinate Bethe ansatz}
\label{sec:aba-comparison-coordinate-bethe-ansatz}

The rapidity in \cref{eq:aba-xxz-bethe-equations} is tied to the unshifted
six-vertex $R$-matrix.  A more symmetric variable is
\begin{equation}
\label{eq:aba-symmetric-rapidity}
    \mu=\lambda+\frac{\eta}{2} .
\end{equation}
Then
\begin{equation}
\label{eq:aba-symmetric-momentum-map}
    \ee^{\ii k(\mu)}
    =
    \frac{\sinh(\mu+\eta/2)}{\sinh(\mu-\eta/2)},
\end{equation}
and the Bethe equations become
\begin{equation}
\label{eq:aba-symmetric-bethe-equations}
    \left[
    \frac{\sinh(\mu_j+\eta/2)}{\sinh(\mu_j-\eta/2)}
    \right]^L
    =
    \prod_{\substack{\ell=1\\ \ell\neq j}}^{M}
    \frac{\sinh(\mu_j-\mu_\ell+\eta)}
         {\sinh(\mu_j-\mu_\ell-\eta)} .
\end{equation}
After analytic continuation $\eta=\ii\gamma$, this is the standard critical XXZ
form
\begin{equation}
\label{eq:aba-critical-xxz-bethe-equations}
    \left[
    \frac{\sinh(\mu_j+\ii\gamma/2)}{\sinh(\mu_j-\ii\gamma/2)}
    \right]^L
    =
    \prod_{\substack{\ell=1\\ \ell\neq j}}^{M}
    \frac{\sinh(\mu_j-\mu_\ell+\ii\gamma)}
         {\sinh(\mu_j-\mu_\ell-\ii\gamma)} .
\end{equation}
This is the same equation obtained in the coordinate Bethe ansatz, up to the
orientation convention for lattice momentum.

The two approaches differ in what they make manifest:
\begin{center}
\begin{tabular}{@{}L{0.30\linewidth}L{0.30\linewidth}L{0.30\linewidth}@{}}
\toprule
Method & Main calculation & Output \\
\midrule
Coordinate Bethe ansatz & Solve contact conditions in real-space wave functions & Two-body scattering phase and momentum quantization \\
Algebraic Bethe ansatz & Commute $A(u)$ and $D(u)$ through products of $B$'s & Transfer-matrix eigenvalue and Bethe equations \\
\bottomrule
\end{tabular}
\end{center}
The coordinate method emphasizes factorized scattering.  The algebraic method
explains why factorized scattering exists: the RTT relation is the operator form
of the Yang--Baxter equation.

\section{Higher conserved charges}
\label{sec:aba-higher-conserved-charges}

Because the Bethe vector is an eigenvector of $\tau(u)$ for all $u$, it is also a
simultaneous eigenvector of all charges generated by $\log\tau(u)$.  If
\begin{equation}
\label{eq:aba-log-transfer-charge-expansion}
    \log\tau(u)=\sum_{n\geq0}u^n Q_n,
\end{equation}
then
\begin{equation}
\label{eq:aba-charge-eigenvalues}
    Q_n\ket{\boldsymbol{\lambda}}
    =q_n(\boldsymbol{\lambda})\ket{\boldsymbol{\lambda}},
    \qquad
    q_n(\boldsymbol{\lambda})
    =
    \left.
    \frac{1}{n!}
    \frac{\dd^n}{\dd u^n}
    \log\Lambda(u\mid\boldsymbol{\lambda})
    \right|_{u=0}
\end{equation}
up to the same normalization convention used in defining the charges.  The
Hamiltonian is the first nontrivial local charge.  Higher $Q_n$ have larger
spatial support and provide the extensive set of conservation laws characteristic
of an integrable chain.

\begin{remark}[Why this is stronger than diagonalizing one Hamiltonian]
A generic many-body Hamiltonian may be diagonalized in principle, but its
eigenvectors do not usually diagonalize a whole analytic family of local and
quasilocal conserved charges.  The algebraic Bethe ansatz gives simultaneous
eigenvectors of the transfer matrix for all spectral parameters.  This is the
finite-size algebraic signature of integrability.
\end{remark}

\section{Trigonometric and rational limits}
\label{sec:aba-trigonometric-rational-limits}

For $\abs{\Delta}<1$, set
\begin{equation}
\label{eq:aba-eta-to-gamma}
    \eta=\ii\gamma,
    \qquad
    0<\gamma<\pi,
    \qquad
    \Delta=\cos\gamma .
\end{equation}
The same algebraic derivation applies, with $\sinh$ functions analytically
continued.  In the symmetric rapidity convention, the Bethe equations are exactly
\cref{eq:aba-critical-xxz-bethe-equations}.

The isotropic XXX chain is the rational degeneration.  Scaling the spectral
parameter and anisotropy together produces the rational $R$-matrix
\begin{equation}
\label{eq:aba-rational-R-schematic}
    R_{\rm XXX}(u)=u\,\id+\ii P,
    \qquad
    \text{up to an overall scalar normalization.}
\end{equation}
In the usual XXX rapidity variable, the Bethe equations become
\begin{equation}
\label{eq:aba-xxx-bethe-equations}
    \left(
    \frac{\lambda_j+\ii/2}{\lambda_j-\ii/2}
    \right)^L
    =
    \prod_{\substack{\ell=1\\ \ell\neq j}}^{M}
    \frac{\lambda_j-\lambda_\ell+\ii}
         {\lambda_j-\lambda_\ell-\ii} .
\end{equation}
This rational limit will reappear when the XXX chain is discussed as a special
case of the XXZ chain.

\section{Common pitfalls}
\label{sec:aba-common-pitfalls}

\begin{remark}[The operators $A,B,C,D$ are nonlocal]
Although $R_{aj}(u)$ is local in the quantum site $j$, the monodromy entries
$A(u),B(u),C(u),D(u)$ are products over the whole chain.  The operator $B(u)$ is
therefore not a local spin-lowering operator.  It is a spectral-parameter
dependent, nonlocal creation operator for Bethe quasiparticles.
\end{remark}

\begin{remark}[Bethe roots cannot sit at generic poles]
The formulas above contain denominators such as $\sinh\lambda_j$ and
$\sinh(\lambda_j-\lambda_\ell)$.  Ordinary Bethe roots are assumed to avoid these
poles.  Singular solutions can occur at special points and require a limiting
procedure.  They are not part of the basic algebraic derivation.
\end{remark}

\begin{remark}[Completeness is a separate question]
The algebraic Bethe ansatz constructs eigenvectors whenever the Bethe equations
hold and the resulting vector is nonzero.  Showing that all eigenvectors are
obtained requires a separate completeness analysis.  For the finite XXZ chain
this is well understood, but it is conceptually distinct from deriving the Bethe
equations.
\end{remark}

\section{Summary}
\label{sec:aba-summary}

The algebraic Bethe ansatz for the XXZ chain can be compressed into the following
sequence:
\begin{equation}
\label{eq:aba-summary-chain}
\boxed{
\begin{gathered}
    R(u)\text{ satisfies YBE}
    \Longrightarrow
    R_{ab}(u-v)M_a(u)M_b(v)=M_b(v)M_a(u)R_{ab}(u-v),
    \\
    M_a(u)=
    \begin{pmatrix}A(u)&B(u)\\ C(u)&D(u)\end{pmatrix}_a,
    \qquad
    \tau(u)=A(u)+D(u),
    \\
    \ket{\boldsymbol{\lambda}}
    =B(\lambda_1)\cdots B(\lambda_M)\ket{0},
    \\
    \tau(u)\ket{\boldsymbol{\lambda}}
    =\Lambda(u\mid\boldsymbol{\lambda})\ket{\boldsymbol{\lambda}}
    \quad
    \text{if }\boldsymbol{\lambda}\text{ satisfies the Bethe equations},
    \\
    \left[
    \frac{\sinh(\lambda_j+\eta)}{\sinh\lambda_j}
    \right]^L
    =
    \prod_{\ell\neq j}
    \frac{\sinh(\lambda_j-\lambda_\ell+\eta)}
         {\sinh(\lambda_j-\lambda_\ell-\eta)} .
\end{gathered}}
\end{equation}
The Bethe equations arise from cancellation of unwanted terms in the RTT algebra.
The transfer-matrix eigenvalue then gives the complete set of commuting-charge
eigenvalues, including the XXZ Hamiltonian energy.  This is the algebraic
counterpart of the factorized scattering picture developed in the coordinate
Bethe-ansatz chapter.


\chapter{Further Integrable Models}
\label{chap:further-integrable-models}

The previous two chapters solved the XXZ chain by two complementary methods.
The coordinate Bethe ansatz emphasized factorized scattering of magnons, while
the algebraic Bethe ansatz emphasized the Yang--Baxter algebra of the six-vertex
transfer matrix.  This chapter explains how the same ideas extend beyond the
nearest-neighbor XXZ chain.  The goal is not to give a second complete solution
of every model.  Instead, the purpose is to identify the structural data that
make each model solvable and to locate it within the mother-model organization
of these notes.

We discuss three classes of examples:
\begin{enumerate}[label=(\roman*)]
    \item the isotropic XXX chain, obtained as the rational limit of the XXZ
    chain;
    \item the one-dimensional Hubbard model, whose exact solution requires a
    nested Bethe ansatz because particles carry both charge and spin;
    \item long-range integrable models, represented by the Haldane--Shastry spin
    chain and the Calogero--Sutherland model.
\end{enumerate}
The recurring theme is that integrability is not synonymous with free
fermions.  In these models the many-body spectrum is controlled by nontrivial
quantization conditions, hidden symmetries, or special long-range interactions.
The excitations are still strongly constrained, but they need not be free
Bogoliubov quasiparticles.

\section{The XXX limit}
\label{sec:further-xxx-limit}

The XXZ Hamiltonian used in \cref{chap:coordinate-bethe-ansatz-xxz,chap:algebraic-bethe-ansatz}
was
\begin{equation}
\label{eq:further-xxz-hamiltonian-recall}
    H_{\rm XXZ}^{(+)}
    =
    J\sum_{j=1}^{L}
    \left(
        \sigma_j^x\sigma_{j+1}^x
        +\sigma_j^y\sigma_{j+1}^y
        +\Delta(\sigma_j^z\sigma_{j+1}^z-1)
    \right),
\end{equation}
with periodic boundary conditions.  The isotropic XXX chain is the special point
\begin{equation}
\label{eq:further-xxx-limit}
    \Delta\to1,
    \qquad
    H_{\rm XXX}^{(+)}
    =
    J\sum_{j=1}^{L}
    \left(
        \bm{\sigma}_j\cdot\bm{\sigma}_{j+1}-1
    \right).
\end{equation}
Equivalently, in terms of spin operators $\bm{S}_j=\bm{\sigma}_j/2$,
\begin{equation}
\label{eq:further-xxx-spin-form}
    H_{\rm XXX}^{(+)}
    =
    4J\sum_{j=1}^{L}
    \left(
        \bm{S}_j\cdot\bm{S}_{j+1}-\frac14
    \right).
\end{equation}
The constant is chosen so that the fully polarized reference state has zero
energy, as in the XXZ Bethe ansatz chapters.

\begin{remark}[Sign convention]
With $J>0$ in \cref{eq:further-xxx-limit}, a one-magnon state has
$\varepsilon(k)=4J(\cos k-1)\leq0$ relative to the ferromagnetic reference
state.  Many condensed-matter texts instead write the antiferromagnetic XXX
Hamiltonian as
$J_{\rm AF}\sum_j\bm{S}_j\cdot\bm{S}_{j+1}$ with $J_{\rm AF}>0$.
The two conventions differ by an overall sign and an additive constant.  The
Bethe equations are convention-independent; only the energy formula changes by
this normalization.
\end{remark}

\subsection{Rational R-matrix}
\label{subsec:further-rational-r-matrix}

The XXX chain is also obtained from the six-vertex $R$-matrix by taking the
rational limit.  In the hyperbolic convention of the algebraic Bethe ansatz
chapter,
\begin{equation}
    R_{\rm XXZ}(u)
    =
    \begin{pmatrix}
        \sinh(u+\eta)&0&0&0\\
        0&\sinh u&\sinh\eta&0\\
        0&\sinh\eta&\sinh u&0\\
        0&0&0&\sinh(u+\eta)
    \end{pmatrix},
    \qquad
    \Delta=\cosh\eta .
\end{equation}
Sending $u\mapsto\eta u$ and then taking $\eta\to0$ gives, up to an irrelevant
scalar normalization,
\begin{equation}
\label{eq:further-rational-R-matrix}
    R_{\rm XXX}(u)
    =u\,\id+P,
\end{equation}
where $P$ permutes the two spin-$1/2$ spaces:
\begin{equation}
\label{eq:further-permutation-operator}
    P\ket{\alpha}\otimes\ket{\beta}
    =
    \ket{\beta}\otimes\ket{\alpha}.
\end{equation}
This rational $R$-matrix satisfies the Yang--Baxter equation
\begin{equation}
\label{eq:further-rational-ybe}
    R_{12}(u-v)R_{13}(u-w)R_{23}(v-w)
    =
    R_{23}(v-w)R_{13}(u-w)R_{12}(u-v).
\end{equation}
It is regular because
\begin{equation}
\label{eq:further-rational-regularity}
    R_{\rm XXX}(0)=P.
\end{equation}
Therefore the logarithmic derivative of the rational transfer matrix gives a
nearest-neighbor Hamiltonian.  Using
\begin{equation}
    P_{j,j+1}=\frac12\left(\id+\bm{\sigma}_j\cdot\bm{\sigma}_{j+1}\right),
\end{equation}
one obtains the isotropic exchange interaction.

\begin{principle}[Trigonometric versus rational integrability]
The XXZ chain is controlled by the trigonometric six-vertex $R$-matrix.  The
XXX chain is controlled by its rational degeneration.  The solution method is
unchanged: monodromy matrix, transfer matrix, RTT algebra, Bethe vectors, Bethe
equations.  What changes is the analytic form of the scattering phase.
\end{principle}

\subsection{Bethe equations in rapidity form}
\label{subsec:further-xxx-bethe-equations}

The rational rapidity variable $\lambda$ is related to the one-magnon momentum
by
\begin{equation}
\label{eq:further-xxx-momentum-rapidity}
    \ee^{\ii k}
    =
    \frac{\lambda+\frac{\ii}{2}}{\lambda-\frac{\ii}{2}}.
\end{equation}
For $M$ magnons, the XXX Bethe equations are
\begin{equation}
\label{eq:further-xxx-bethe-equations}
    \left(
        \frac{\lambda_j+\frac{\ii}{2}}
             {\lambda_j-\frac{\ii}{2}}
    \right)^L
    =
    \prod_{\substack{\ell=1\\ \ell\neq j}}^{M}
    \frac{\lambda_j-\lambda_\ell+\ii}
         {\lambda_j-\lambda_\ell-\ii},
    \qquad
    j=1,\ldots,M .
\end{equation}
The corresponding two-body scattering amplitude is
\begin{equation}
\label{eq:further-xxx-two-body-S}
    S_{\rm XXX}(\lambda_j-\lambda_\ell)
    =
    \frac{\lambda_j-\lambda_\ell+\ii}
         {\lambda_j-\lambda_\ell-\ii}.
\end{equation}
Thus the many-body quantization condition still has the universal Bethe form
\begin{equation}
\label{eq:further-general-bethe-quantization}
    \ee^{\ii k_jL}
    =
    \prod_{\ell\neq j}S(k_j,k_\ell),
\end{equation}
but the scattering phase has become rational rather than trigonometric.

Using \cref{eq:further-xxx-momentum-rapidity}, one finds
\begin{equation}
\label{eq:further-xxx-cos-rapidity}
    \cos k
    =
    \frac{\lambda^2-\frac14}{\lambda^2+\frac14}.
\end{equation}
Therefore the energy in the normalization of \cref{eq:further-xxx-limit} is
\begin{equation}
\label{eq:further-xxx-energy}
    E
    =
    4J\sum_{j=1}^{M}(\cos k_j-1)
    =
    -2J\sum_{j=1}^{M}
    \frac{1}{\lambda_j^2+\frac14}.
\end{equation}
The total lattice momentum is
\begin{equation}
\label{eq:further-xxx-total-momentum}
    \ee^{\ii P}
    =
    \prod_{j=1}^{M}
    \frac{\lambda_j+\frac{\ii}{2}}
         {\lambda_j-\frac{\ii}{2}}.
\end{equation}

\subsection{Enhanced SU(2) symmetry}
\label{subsec:further-xxx-su2}

For generic $\Delta$, the XXZ chain conserves only the $U(1)$ charge
$S^z_{\rm tot}$.  At $\Delta=1$, the symmetry is enhanced to full spin
rotation symmetry:
\begin{equation}
\label{eq:further-xxx-su2-commutator}
    [H_{\rm XXX}^{(+)},S^\alpha_{\rm tot}]=0,
    \qquad
    S^\alpha_{\rm tot}=\frac12\sum_{j=1}^{L}\sigma_j^\alpha,
    \qquad
    \alpha=x,y,z .
\end{equation}
This enhancement has an important consequence.  Bethe states constructed above
the all-up pseudovacuum are usually highest-weight states of $SU(2)$ multiplets.
The remaining members of the multiplet are obtained by acting with
$S^-_{\rm tot}$.  This is why completeness questions in the XXX chain are
slightly more delicate than merely counting finite rapidity solutions of
\cref{eq:further-xxx-bethe-equations}.

\begin{remark}[Infinite rapidities]
In the XXX chain, acting with the global spin-lowering operator can be viewed as
adding a Bethe root at infinity.  This is a useful algebraic way to understand
how $SU(2)$ descendants fit into the Bethe-ansatz framework.
\end{remark}

\section{Hubbard model and nested Bethe ansatz}
\label{sec:further-hubbard-model}

The one-dimensional Hubbard model is the canonical example of an interacting
fermion model that is exactly solvable but not reducible to a single-species
magnon Bethe ansatz.  Its local degrees of freedom are
\begin{equation}
\label{eq:further-hubbard-local-space}
    \ket{0},
    \qquad
    \ket{\uparrow},
    \qquad
    \ket{\downarrow},
    \qquad
    \ket{\uparrow\downarrow},
\end{equation}
so each site carries both charge and spin.  The Hamiltonian on a periodic chain
is
\begin{equation}
\label{eq:further-hubbard-hamiltonian}
    H_{\rm Hub}
    =
    -t\sum_{j=1}^{L}\sum_{\sigma=\uparrow,\downarrow}
    \left(
        c_{j\sigma}^\dagger c_{j+1,\sigma}
        +c_{j+1,\sigma}^\dagger c_{j\sigma}
    \right)
    +U\sum_{j=1}^{L}n_{j\uparrow}n_{j\downarrow} .
\end{equation}
Here
\begin{equation}
    n_{j\sigma}=c_{j\sigma}^\dagger c_{j\sigma},
    \qquad
    c_{L+1,\sigma}=c_{1\sigma}.
\end{equation}
The hopping amplitude is $t$, and $U$ is the on-site interaction.

\subsection{Why a nested ansatz is needed}
\label{subsec:further-why-nested}

In the XXZ chain, the Bethe particles are down spins moving in an all-up
reference state.  There is only one species of quasiparticle, so every collision
is described by a scalar two-body phase.  In the Hubbard model, a particle
carries charge momentum and spin.  When two fermions scatter, their spin labels
can be rearranged.  The two-body scattering matrix therefore acts on an internal
spin space.  Diagonalizing the charge problem leaves a second spin problem to be
solved.  This second step is the nested Bethe ansatz.

The exact solution is parametrized by two sets of variables:
\begin{equation}
\label{eq:further-hubbard-variables}
    \{k_j\}_{j=1}^{N},
    \qquad
    \{\lambda_\alpha\}_{\alpha=1}^{M},
\end{equation}
where $N$ is the total number of electrons and $M$ is the number of down-spin
electrons.  The $k_j$ are charge momenta, while the $\lambda_\alpha$ are spin
rapidities.

\begin{principle}[Nested Bethe ansatz]
When the two-body scattering matrix carries an internal index, Bethe ansatz
becomes a two-level problem.  First one writes a coordinate Bethe ansatz for
charge motion.  Then one diagonalizes the internal spin scattering problem by a
second Bethe ansatz.  The final answer contains coupled Bethe equations for
charge and spin rapidities.
\end{principle}

\subsection{Lieb--Wu equations}
\label{subsec:further-lieb-wu-equations}

Introduce the dimensionless interaction parameter
\begin{equation}
\label{eq:further-hubbard-u-parameter}
    u=\frac{U}{4t}.
\end{equation}
For the Hamiltonian \cref{eq:further-hubbard-hamiltonian}, the Lieb--Wu Bethe
equations are
\begin{align}
\label{eq:further-lieb-wu-charge}
    \ee^{\ii k_j L}
    &=
    \prod_{\alpha=1}^{M}
    \frac{\sin k_j-\lambda_\alpha+\ii u}
         {\sin k_j-\lambda_\alpha-\ii u},
    \qquad
    j=1,\ldots,N,\\
\label{eq:further-lieb-wu-spin}
    \prod_{j=1}^{N}
    \frac{\lambda_\alpha-\sin k_j+\ii u}
         {\lambda_\alpha-\sin k_j-\ii u}
    &=
    \prod_{\substack{\beta=1\\ \beta\neq\alpha}}^{M}
    \frac{\lambda_\alpha-\lambda_\beta+2\ii u}
         {\lambda_\alpha-\lambda_\beta-2\ii u},
    \qquad
    \alpha=1,\ldots,M .
\end{align}
The energy and total momentum are
\begin{equation}
\label{eq:further-hubbard-energy-momentum}
    E=-2t\sum_{j=1}^{N}\cos k_j,
    \qquad
    P=\sum_{j=1}^{N}k_j
    \quad (\mathrm{mod}\ 2\pi),
\end{equation}
for the unshifted convention in \cref{eq:further-hubbard-hamiltonian}.  If one
uses the particle-hole symmetric Hamiltonian
\begin{equation}
\label{eq:further-shifted-hubbard}
    H'_{\rm Hub}
    =
    -t\sum_{j,\sigma}
    \left(
        c_{j\sigma}^\dagger c_{j+1,\sigma}
        +c_{j+1,\sigma}^\dagger c_{j\sigma}
    \right)
    +U\sum_j
    \left(n_{j\uparrow}-\frac12\right)
    \left(n_{j\downarrow}-\frac12\right),
\end{equation}
then the same Bethe roots give the shifted energy
\begin{equation}
\label{eq:further-shifted-hubbard-energy}
    E'=E-\frac{U}{2}N+\frac{U}{4}L .
\end{equation}

\begin{remark}[Why $\sin k$ appears]
The Hubbard equations are not obtained by simply replacing the XXZ rapidity by
$k$.  The charge dispersion is that of a tight-binding lattice fermion,
$-2t\cos k$, while the spin scattering depends on $\sin k$.  This mixture is a
characteristic feature of the Hubbard nested solution.
\end{remark}

\subsection{Half filling and the Mott gap}
\label{subsec:further-hubbard-half-filling}

At half filling, $N=L$.  For repulsive interaction $U>0$, the exact solution
shows that the charge sector is gapped while the spin sector remains gapless.
This is the one-dimensional Mott-insulating mechanism in its cleanest exact
form.  The Hilbert space still allows charge fluctuations, but the low-energy
sector separates into
\begin{equation}
\label{eq:further-hubbard-spin-charge-structure}
    \text{gapped charge excitations}
    \qquad\text{and}\qquad
    \text{gapless spin excitations}.
\end{equation}
For large positive $U/t$, virtual hopping produces the effective antiferromagnetic
Heisenberg chain
\begin{equation}
\label{eq:further-hubbard-large-U-heisenberg}
    H_{\rm eff}
    =
    J_{\rm ex}\sum_{j=1}^{L}
    \left(
        \bm{S}_j\cdot\bm{S}_{j+1}-\frac14
    \right),
    \qquad
    J_{\rm ex}=\frac{4t^2}{U},
\end{equation}
inside the subspace with one electron per site.  Thus the XXX chain appears as
the spin-sector low-energy limit of the repulsive half-filled Hubbard model.

\begin{example}[From Hubbard to Heisenberg]
Consider two neighboring sites with one electron on each site.  A hopping event
creates a virtual doubly occupied site and costs energy $U$.  A second hopping
event returns to the singly occupied subspace.  Second-order perturbation theory
therefore gives an exchange scale $t^2/U$.  Fermionic statistics select the
antiferromagnetic exchange, producing $J_{\rm ex}=4t^2/U$.
\end{example}

\subsection{Integrability structure}
\label{subsec:further-hubbard-integrability}

The Hubbard model is integrable, but its integrability is more subtle than that
of the six-vertex model.  The relevant $R$-matrix is not the elementary
six-vertex $R$-matrix; it is a higher structure compatible with two coupled
spin species.  Equivalently, one may view the model as having two intertwined
XX-type hopping problems connected by the on-site interaction.  The important
lesson for these notes is the following:
\boxedpar{%
The Hubbard model is solvable not because it is quadratic, but because its
many-body scattering is still factorized after accounting for the nested charge
and spin structure.}

\section{Long-range integrable models}
\label{sec:further-long-range-integrable-models}

The nearest-neighbor XXX, XXZ, and XYZ chains teach that local Hamiltonians can
be integrable when they arise from logarithmic derivatives of commuting transfer
matrices.  There is another important class of integrable models whose defining
feature is not nearest-neighbor locality, but a very special long-range
interaction.  The two standard examples are the Haldane--Shastry spin chain and
the Calogero--Sutherland particle model.

\subsection{Haldane--Shastry chain}
\label{subsec:further-haldane-shastry}

Place $L$ spin-$1/2$ degrees of freedom at equally spaced points on the unit
circle,
\begin{equation}
\label{eq:further-hs-roots-of-unity}
    z_j=\ee^{2\pi\ii j/L},
    \qquad
    j=1,\ldots,L .
\end{equation}
The Haldane--Shastry Hamiltonian is
\begin{equation}
\label{eq:further-haldane-shastry-hamiltonian}
    H_{\rm HS}
    =
    J\left(\frac{\pi}{L}\right)^2
    \sum_{1\leq j<\ell\leq L}
    \frac{\bm{S}_j\cdot\bm{S}_\ell-\frac14}
         {\sin^2\left(\frac{\pi(j-\ell)}{L}\right)} .
\end{equation}
The interaction decays as the inverse square of the chord distance.  Although
this Hamiltonian is long-ranged, it is more exactly solvable than a generic
nearest-neighbor interacting spin chain.

At half filling of down spins, an exact ground-state wave function can be written
in terms of the down-spin coordinates $z_1,\ldots,z_{L/2}$ as a lattice Jastrow
state,
\begin{equation}
\label{eq:further-hs-ground-state}
    \Psi_{\rm HS}(z_1,\ldots,z_{L/2})
    =
    \prod_{a<b}(z_a-z_b)^2
    \prod_{a=1}^{L/2} z_a .
\end{equation}
This is the spin-chain analogue of a Laughlin-type wave function.  The square
Jastrow factor enforces strong correlations between spin flips.

The elementary excitations are spinons.  They carry spin $1/2$ rather than spin
$1$, and they obey fractional exclusion statistics.  This is qualitatively
different from the magnons of the ferromagnetic Bethe ansatz, which are spin-flip
particles above a polarized vacuum.

\begin{remark}[Yangian symmetry]
The hidden symmetry of the Haldane--Shastry chain is a Yangian extension of
$SU(2)$.  This enlarged symmetry is responsible for large spectral degeneracies
and for the unusually simple spinon structure.  It is the long-range analogue of
the algebraic constraints imposed by the Yang--Baxter equation in nearest-neighbor
chains.
\end{remark}

\subsection{Calogero--Sutherland model}
\label{subsec:further-calogero-sutherland}

The Calogero--Sutherland model is a continuum many-body system on a circle with
inverse-square interactions.  A common normalization is
\begin{equation}
\label{eq:further-calogero-sutherland-hamiltonian}
    H_{\rm CS}
    =
    -\sum_{j=1}^{N}\frac{\partial^2}{\partial x_j^2}
    +
    g(g-1)\left(\frac{\pi}{L}\right)^2
    \sum_{1\leq j<\ell\leq N}
    \frac{1}
         {\sin^2\left(\frac{\pi(x_j-x_\ell)}{L}\right)} .
\end{equation}
The exact ground state has the Jastrow form
\begin{equation}
\label{eq:further-cs-ground-state}
    \Psi_0(x_1,\ldots,x_N)
    =
    \prod_{j<\ell}
    \left|
        \sin\frac{\pi(x_j-x_\ell)}{L}
    \right|^{g}.
\end{equation}
The excited states are obtained by multiplying this ground state by symmetric
polynomials.  More precisely, Jack polynomials provide the natural eigenbasis.
Thus the Calogero--Sutherland model is solvable by special-function and
representation-theoretic methods rather than by a nearest-neighbor coordinate
Bethe ansatz.

The coupling $g$ controls the statistical interaction between particles.  At
$g=0$ or $g=1$, the inverse-square potential vanishes, corresponding to free
bosonic or fermionic endpoints depending on the imposed symmetry.  For general
$g$, the model realizes fractional exclusion statistics in a continuum setting.

\subsection{Freezing trick and relation to spin chains}
\label{subsec:further-freezing-trick}

The Haldane--Shastry chain and the Calogero--Sutherland model are not unrelated.
In the strong-coupling limit of spin Calogero--Sutherland-type models, the
particle coordinates freeze near equally spaced classical equilibrium positions
on the circle.  The remaining low-energy dynamics is carried by the internal
spin variables.  This limiting procedure produces inverse-square spin-chain
Hamiltonians of Haldane--Shastry type.

Schematically,
\begin{equation}
\label{eq:further-freezing-trick-schematic}
    \text{spin Calogero--Sutherland particles}
    \xrightarrow[\text{coordinates freeze}]{\text{large coupling}}
    \text{Haldane--Shastry spin chain}.
\end{equation}
This relation is called the freezing trick.  It is another example of the
mother-model principle: a continuum integrable particle model contains a lattice
integrable spin chain as a controlled limit.

\section{Comparison of integrable structures}
\label{sec:further-comparison-integrable-structures}

The examples in this chapter illustrate several distinct meanings of exact
solvability.  They are summarized in \cref{tab:further-integrability-comparison}.

\begin{table}[htbp]
\centering
\small
\begin{tabularx}{\linewidth}{@{}L{0.22\linewidth}L{0.25\linewidth}Y@{}}
\toprule
Model & Solvability mechanism & Main spectral data \\
\midrule
XXX chain & Rational $R$-matrix; Bethe ansatz; enhanced $SU(2)$ symmetry & Rapidities satisfying rational Bethe equations \\
XXZ chain & Trigonometric six-vertex $R$-matrix; coordinate/algebraic Bethe ansatz & Rapidities satisfying trigonometric Bethe equations \\
XYZ chain & Elliptic eight-vertex $R$-matrix; Baxter $TQ$ relation & Elliptic spectral parameter and functional equations \\
Hubbard chain & Nested Bethe ansatz; factorized scattering with charge and spin variables & Charge momenta $k_j$ and spin rapidities $\lambda_\alpha$ satisfying Lieb--Wu equations \\
Haldane--Shastry chain & Inverse-square interaction; Yangian symmetry; spinon basis & Fractional-statistics spinons and exact Jastrow ground state \\
Calogero--Sutherland model & Inverse-square continuum interaction; Jack-polynomial eigenfunctions & Pseudomomenta and Jack-polynomial labels \\
\bottomrule
\end{tabularx}
\caption{Representative integrable models beyond the basic XXZ coordinate
Bethe ansatz.  The word ``integrable'' refers to different but related
structures: rational, trigonometric, or elliptic $R$-matrices; nested Bethe
ansatz; Yangian symmetry; or inverse-square special-function solvability.}
\label{tab:further-integrability-comparison}
\end{table}

The following hierarchy is useful:
\begin{equation}
\label{eq:further-integrability-hierarchy}
\begin{gathered}
    \text{free fermion solvability}
    \subset
    \text{Bethe/Yang--Baxter integrability}
    \subset
    \text{broader exact solvability},\\
    \text{but none of these inclusions should be interpreted as equality.}
\end{gathered}
\end{equation}
Free-fermion models are solvable because the Hamiltonian becomes quadratic.
Bethe-ansatz models are solvable because many-body scattering factorizes.
Long-range inverse-square models are solvable because the interaction is tuned
to special symmetry and polynomial structures.  Commuting-projector models,
which appear in the next part of the notes, are solvable for yet another reason:
the Hamiltonian decomposes into mutually compatible local constraints.

\section{What should be remembered}
\label{sec:further-what-to-remember}

The main points of this chapter are:
\begin{enumerate}[label=(\roman*)]
    \item The XXX chain is the rational limit of the XXZ chain.  Its Bethe
    equations are obtained by replacing trigonometric scattering data by rational
    scattering data.
    \item The XXX point has enhanced $SU(2)$ symmetry.  Bethe roots at infinity
    and highest-weight multiplets are natural consequences of this symmetry.
    \item The Hubbard model requires a nested Bethe ansatz because the particles
    carry both charge and spin.  Its solution is encoded in the coupled
    Lieb--Wu equations.
    \item At half filling and repulsive interaction, the Hubbard model has a
    gapped charge sector and a gapless spin sector.  In the large-$U$ limit, the
    low-energy spin dynamics reduces to an antiferromagnetic Heisenberg chain.
    \item Long-range inverse-square models such as the Haldane--Shastry and
    Calogero--Sutherland models are integrable for reasons different from the
    nearest-neighbor transfer-matrix construction.  Their exact wave functions
    and spectra are governed by Jastrow factors, Yangian symmetry, and special
    polynomials.
\end{enumerate}

\boxedpar{%
The common lesson is that exact solvability is a structural property.  Quadratic
fermions, Yang--Baxter algebras, nested scattering, inverse-square interactions,
and commuting projectors are different mechanisms, but each replaces the generic
many-body problem by a constrained algebraic or analytic problem.}


% -----------------------------------------------------------------------------
\part{Two-Dimensional Quantum Exact Solvability}
% -----------------------------------------------------------------------------

\chapter{Commuting-Projector Models and the Toric Code}
\label{chap:commuting-projector-toric-code}

The previous part studied interacting one-dimensional quantum chains whose exact
solvability comes from factorized scattering and the Yang--Baxter equation.  In
this part we move to two spatial dimensions.  The first solvable structure is
very different from Bethe-ansatz integrability: the Hamiltonian is a sum of
local terms that commute with one another.  Such models are exactly solvable
because all local conserved projectors can be diagonalized simultaneously.

The toric code is the canonical example.  It is simultaneously
\begin{enumerate}[label=(\roman*)]
    \item a stabilizer Hamiltonian;
    \item a commuting-projector model;
    \item a lattice \(\ZZ_2\) gauge theory at zero gauge coupling;
    \item the simplest fixed-point model of intrinsic topological order.
\end{enumerate}
It is therefore an ideal first model for two-dimensional quantum exact
solvability.  Its solution does not require Fourier transforms, Bogoliubov
rotations, or Bethe rapidities.  Instead, the spectrum follows from a finite set
of mutually commuting \(\ZZ_2\)-valued operators, while its physical content is
encoded in nonlocal string operators.

\section{Commuting-projector solvability}
\label{sec:commuting-projector-principle}

Let \(\calH\) be a many-body Hilbert space and suppose that a Hamiltonian can be
written as
\begin{equation}
\label{eq:commuting-projector-general-form}
    H=\sum_{\alpha} E_{\alpha} P_{\alpha},
    \qquad
    P_{\alpha}^2=P_{\alpha},
    \qquad
    [P_{\alpha},P_{\beta}]=0 .
\end{equation}
Then all \(P_{\alpha}\) can be diagonalized simultaneously.  The many-body
spectrum is reduced to the problem of finding which binary eigenvalue patterns
\(p_{\alpha}\in\{0,1\}\) are compatible with the algebraic constraints among the
projectors.

Equivalently, if each local term is written in terms of an involution
\(Q_{\alpha}^2=\id\),
\begin{equation}
\label{eq:involution-projector-relation}
    P_{\alpha}=\frac{1-Q_{\alpha}}{2},
    \qquad
    Q_{\alpha}=\id-2P_{\alpha},
\end{equation}
then exact solvability is the statement that the commuting stabilizers
\(Q_{\alpha}\) admit simultaneous eigenvalues \(q_{\alpha}=\pm1\).

\begin{definition}[Commuting-projector Hamiltonian]
A local Hamiltonian is called a commuting-projector Hamiltonian if it can be
written as a sum of local projectors satisfying \cref{eq:commuting-projector-general-form}.
If the ground state minimizes every local projector separately, the Hamiltonian
is called frustration-free.
\end{definition}

\begin{remark}[Exact solvability versus integrability]
A commuting-projector model has an extensive set of strictly local conserved
operators.  This is stronger than ordinary quantum integrability in one
dimension, where conserved charges are usually quasilocal or generated by a
transfer matrix.  However, the strength of the solvability is also more
specialized: a generic perturbation that destroys commutativity usually cannot be
solved by the same method.  The toric code should therefore be viewed as an
exactly solvable fixed point, not as a representative of a continuously soluble
family in the Bethe-ansatz sense.
\end{remark}

The simplest nontrivial commuting-projector models are stabilizer Hamiltonians.
There one has \(N\) qubits and an Abelian subgroup \(\mathcal{S}\) of the
Pauli group, generated by mutually commuting Pauli strings \(S_a\).  The common
\(+1\) eigenspace
\begin{equation}
\label{eq:stabilizer-code-space}
    \calH_{\rm code}
    =
    \{\ket{\psi}\in(\CC^2)^{\otimes N}:S_a\ket{\psi}=\ket{\psi}
    \text{ for all }a\}
\end{equation}
is the ground space of a positive-semidefinite Hamiltonian
\begin{equation}
\label{eq:stabilizer-positive-hamiltonian}
    H_{\rm stab}=\sum_a \frac{1-S_a}{2} .
\end{equation}
If \(r\) of the stabilizer generators are independent, then
\begin{equation}
\label{eq:stabilizer-dimension-counting}
    \dim\calH_{\rm code}=2^{N-r} .
\end{equation}
This counting formula will give the toric-code ground-state degeneracy on a
closed surface.

\section{Lattice, Hilbert space, and Hamiltonian}
\label{sec:toric-code-hamiltonian}

Consider an oriented square lattice embedded on a closed orientable surface
\(\Sigma_g\) of genus \(g\).  The set of vertices, edges, and plaquettes is
denoted by
\begin{equation}
\label{eq:cellulation-sets}
    V,
    \qquad
    E,
    \qquad
    P,
\end{equation}
with cardinalities \(N_v=\abs{V}\), \(N_e=\abs{E}\), and \(N_p=\abs{P}\).
A spin-\(1/2\) degree of freedom lives on every edge \(\ell\in E\).  Thus
\begin{equation}
\label{eq:toric-code-hilbert-space}
    \calH_{\rm TC}
    =\bigotimes_{\ell\in E}\CC^2_{\ell},
    \qquad
    \dim\calH_{\rm TC}=2^{N_e} .
\end{equation}
We write Pauli operators on edge \(\ell\) as \(\sigma^x_{\ell}\),
\(\sigma^y_{\ell}\), and \(\sigma^z_{\ell}\).

For each vertex \(s\in V\), define the star operator
\begin{equation}
\label{eq:toric-code-star-operator}
    A_s=
    \prod_{\ell\ni s}\sigma^x_{\ell},
\end{equation}
where the product is over the edges incident on \(s\).  For each plaquette
\(p\in P\), define the plaquette operator
\begin{equation}
\label{eq:toric-code-plaquette-operator}
    B_p=
    \prod_{\ell\in\partial p}\sigma^z_{\ell},
\end{equation}
where the product is over the boundary edges of \(p\).  On a square lattice each
star and plaquette contains four edges, but the algebraic construction only
requires that the lattice be a cellulation of a closed oriented surface.

The toric-code Hamiltonian is
\begin{equation}
\label{eq:toric-code-hamiltonian}
    H_{\rm TC}
    =
    -J_e\sum_{s\in V}A_s
    -J_m\sum_{p\in P}B_p,
    \qquad
    J_e,J_m>0 .
\end{equation}
Equivalently, up to an additive constant,
\begin{equation}
\label{eq:toric-code-projector-form}
    H_{\rm TC}
    =
    2J_e\sum_{s\in V}\frac{1-A_s}{2}
    +2J_m\sum_{p\in P}\frac{1-B_p}{2}
    -(J_eN_v+J_mN_p).
\end{equation}
The terms \((1-A_s)/2\) and \((1-B_p)/2\) are projectors onto violated star and
plaquette constraints.

\begin{remark}[Naming convention]
The notation \(J_e\) and \(J_m\) anticipates the excitation content.  A violated
star, \(A_s=-1\), is called an electric charge or \(e\) particle.  A violated
plaquette, \(B_p=-1\), is called a magnetic flux or \(m\) particle.  The model
is self-dual under exchange of direct and dual lattices together with
\(A_s\leftrightarrow B_p\), \(J_e\leftrightarrow J_m\), and
\(\sigma^x\leftrightarrow\sigma^z\).
\end{remark}

\section{Commutativity and simultaneous diagonalization}
\label{sec:toric-code-commutativity}

The toric code is exactly solvable because every star and every plaquette term
commutes.

\begin{proposition}[Toric-code stabilizer algebra]
The operators \(A_s\) and \(B_p\) satisfy
\begin{equation}
\label{eq:toric-code-stabilizer-algebra}
    A_s^2=B_p^2=\id,
    \qquad
    [A_s,A_{s'}]=0,
    \qquad
    [B_p,B_{p'}]=0,
    \qquad
    [A_s,B_p]=0 .
\end{equation}
\end{proposition}

\begin{proof}
The first identity follows from \((\sigma^x)^2=(\sigma^z)^2=\id\).  Two star
operators contain only \(\sigma^x\) operators, hence commute.  Two plaquette
operators contain only \(\sigma^z\) operators, hence commute.

It remains to check \([A_s,B_p]=0\).  The operators \(A_s\) and \(B_p\) can only
fail to commute on edges shared by the star \(s\) and the plaquette \(p\).  On a
shared edge,
\begin{equation}
\label{eq:pauli-anticommute-local}
    \sigma^x_{\ell}\sigma^z_{\ell}
    =
    -\sigma^z_{\ell}\sigma^x_{\ell} .
\end{equation}
On a square lattice embedded without boundary, a star and a plaquette share
either zero or two edges.  Therefore moving \(A_s\) through \(B_p\) produces the
factor \((-1)^n\), where \(n\in\{0,2\}\).  Since \((-1)^n=+1\), the two
operators commute.
\end{proof}

Because of \cref{eq:toric-code-stabilizer-algebra}, the Hamiltonian can be
diagonalized by simultaneous eigenstates
\begin{equation}
\label{eq:toric-code-simultaneous-eigenvalues}
    A_s\ket{\psi}=a_s\ket{\psi},
    \qquad
    B_p\ket{\psi}=b_p\ket{\psi},
    \qquad
    a_s,b_p\in\{+1,-1\} .
\end{equation}
The energy of such a state is
\begin{equation}
\label{eq:toric-code-energy-eigenvalue}
    E[\{a_s\},\{b_p\}]
    =
    -J_e\sum_{s\in V}a_s
    -J_m\sum_{p\in P}b_p .
\end{equation}
Thus the full spectrum is determined once we know the allowed stabilizer
eigenvalue patterns and their degeneracies.

On a closed surface there are two global relations:
\begin{equation}
\label{eq:toric-code-global-relations}
    \prod_{s\in V} A_s=\id,
    \qquad
    \prod_{p\in P} B_p=\id .
\end{equation}
The first relation holds because every edge touches exactly two vertices and
therefore appears twice in \(\prod_s A_s\).  The second holds because every edge
belongs to exactly two plaquette boundaries and therefore appears twice in
\(\prod_p B_p\).

It follows that the number of star violations and the number of plaquette
violations must both be even:
\begin{equation}
\label{eq:toric-code-even-parity-constraints}
    \prod_{s\in V}a_s=1,
    \qquad
    \prod_{p\in P}b_p=1 .
\end{equation}
In particular, on a closed surface a single isolated \(e\) or \(m\) excitation
cannot be created by a finite local operator.  Excitations are created in pairs,
or else one must use open boundaries or punctures.

\section{Ground states as flat \texorpdfstring{\(\ZZ_2\)}{Z2} gauge fields}
\label{sec:toric-code-ground-states}

The ground-state conditions are
\begin{equation}
\label{eq:toric-code-ground-state-conditions}
    A_s\ket{\Omega}=\ket{\Omega},
    \qquad
    B_p\ket{\Omega}=\ket{\Omega}
\end{equation}
for all vertices \(s\) and plaquettes \(p\).  These equations have a useful gauge
theory interpretation.  The \(\sigma^z\)-basis assigns a \(\ZZ_2\) variable to
each edge.  The plaquette constraint \(B_p=+1\) says that the \(\ZZ_2\) flux
through every plaquette is trivial.  The star operator \(A_s\) flips all edge
variables adjacent to a vertex; it acts as a local \(\ZZ_2\) gauge
transformation.  Thus the toric-code ground states are gauge-invariant
superpositions of flat \(\ZZ_2\) connections.

A concrete construction is obtained as follows.  Let \(
\ket{0_z}\) be the
state with \(\sigma^z_{\ell}=+1\) for all edges.  It already satisfies
\(B_p\ket{0_z}=\ket{0_z}\) for every plaquette.  The projector onto the
\(+1\)-eigenspace of all star operators is
\begin{equation}
\label{eq:toric-code-star-projector}
    \Pi_A
    =
    \prod_{s\in V}\frac{1+A_s}{2} .
\end{equation}
Therefore
\begin{equation}
\label{eq:toric-code-plane-ground-state}
    \ket{\Omega_0}
    \propto
    \Pi_A\ket{0_z}
\end{equation}
is a ground state in the topologically trivial sector.  In the \(\sigma^z\)
basis, this state is an equal-amplitude superposition of closed-loop
configurations.  The word ``closed'' means that every vertex is touched by an
even number of flipped edges; this is precisely the no-charge condition
\(A_s=+1\).

On a torus there are additional sectors distinguished by noncontractible
\(\ZZ_2\) fluxes.  Local star operations cannot change these global fluxes, so
one obtains four independent ground states rather than one.  We derive this
counting in \cref{sec:toric-code-topological-degeneracy}.

\section{String operators and anyon excitations}
\label{sec:toric-code-anyons}

The elementary excitations are defects of the stabilizer constraints.

\begin{definition}[Electric and magnetic excitations]
A vertex \(s\) with \(A_s=-1\) is an electric charge, denoted \(e\).  A plaquette
\(p\) with \(B_p=-1\) is a magnetic flux, denoted \(m\).  The composite excitation
\(\epsilon=e\times m\) is obtained by placing an \(e\) and an \(m\) at the same
position in the continuum topological description.
\end{definition}

The energy cost of one violated star is \(2J_e\), because its contribution to
\cref{eq:toric-code-hamiltonian} changes from \(-J_e\) to \(+J_e\).  Similarly,
one violated plaquette costs \(2J_m\).  On a closed surface the minimal locally
created excitations are pairs, with energies
\begin{equation}
\label{eq:toric-code-pair-energy-costs}
    \Delta E_{e\text{-pair}}=4J_e,
    \qquad
    \Delta E_{m\text{-pair}}=4J_m .
\end{equation}

\subsection{Electric strings}

Let \(\gamma\) be a path on the direct lattice, i.e. a sequence of edges.  Define the electric string operator
\begin{equation}
\label{eq:toric-code-electric-string}
    W_e(\gamma)
    =
    \prod_{\ell\in\gamma}\sigma^z_{\ell} .
\end{equation}
Since \(W_e(\gamma)\) is made of \(\sigma^z\) operators, it commutes with all
plaquette operators \(B_p\).  It anticommutes with a star operator \(A_s\) if and
only if \(s\) is an endpoint of \(\gamma\).  Therefore an open electric string
creates or moves a pair of \(e\) excitations at its endpoints:
\begin{equation}
\label{eq:e-string-endpoint-action}
    A_s W_e(\gamma)
    =
    \begin{cases}
        -W_e(\gamma)A_s, & s\in\partial\gamma,\\
        \phantom{-}W_e(\gamma)A_s, & s\notin\partial\gamma .
    \end{cases}
\end{equation}
For a closed path, \(\partial\gamma=\varnothing\), the string operator commutes
with the Hamiltonian.

\subsection{Magnetic strings}

Let \(\gamma^\ast\) be a path on the dual lattice.  It crosses a set of direct
lattice edges.  Define
\begin{equation}
\label{eq:toric-code-magnetic-string}
    W_m(\gamma^\ast)
    =
    \prod_{\ell\perp\gamma^\ast}\sigma^x_{\ell} .
\end{equation}
This operator commutes with all star operators \(A_s\), but it anticommutes with
the two plaquette operators at the endpoints of \(\gamma^\ast\).  Thus an open
magnetic string creates or moves a pair of \(m\) fluxes.

\subsection{Fusion rules}

Because \(e\) and \(m\) excitations are \(\ZZ_2\)-valued defects, bringing two
identical excitations together annihilates them.  The fusion rules are
\begin{equation}
\label{eq:toric-code-fusion-rules}
    e\times e=1,
    \qquad
    m\times m=1,
    \qquad
    e\times m=\epsilon,
    \qquad
    \epsilon\times\epsilon=1 .
\end{equation}
The set of anyon types is therefore
\begin{equation}
\label{eq:toric-code-anyon-types}
    \{1,e,m,\epsilon\}\cong \ZZ_2\times\ZZ_2
\end{equation}
under fusion.

\section{Mutual statistics}
\label{sec:toric-code-mutual-statistics}

The nontrivial topological content of the toric code is not the energy of the
excitations, which is completely local, but their braiding statistics.  Let
\(\gamma\) be a direct-lattice path and \(\gamma^\ast\) a dual-lattice path.  The
corresponding string operators obey
\begin{equation}
\label{eq:toric-code-string-crossing-algebra}
    W_e(\gamma)W_m(\gamma^\ast)
    =
    (-1)^{I(\gamma,\gamma^\ast)}
    W_m(\gamma^\ast)W_e(\gamma),
\end{equation}
where \(I(\gamma,\gamma^\ast)\) is the number of intersections modulo \(2\).
The reason is local: at each crossing the two strings act on the same edge by
\(\sigma^z\) and \(\sigma^x\), which anticommute.

If an \(e\) particle is adiabatically moved once around an \(m\) particle, the
worldline process is represented algebraically by a closed electric string that
intersects a magnetic string once.  Hence the wave function acquires the phase
\begin{equation}
\label{eq:toric-code-mutual-minus-one}
    \theta_{em}=-1 .
\end{equation}
Thus \(e\) and \(m\) are mutual semions.

The individual self-statistics are bosonic,
\begin{equation}
\label{eq:toric-code-self-statistics}
    \theta_e=+1,
    \qquad
    \theta_m=+1,
\end{equation}
because exchanging two identical \(e\)'s or two identical \(m\)'s can be
represented without producing a nontrivial string crossing.  Their composite
\(\epsilon=e\times m\) has fermionic self-statistics,
\begin{equation}
\label{eq:toric-code-epsilon-fermion}
    \theta_{\epsilon}=-1 .
\end{equation}
This is one of the simplest examples where a purely spin Hamiltonian supports an
emergent fermionic quasiparticle.

\begin{remark}[Why this is topological]
The phase in \cref{eq:toric-code-string-crossing-algebra} depends only on the
intersection number modulo \(2\), not on the metric length or detailed shape of
the paths.  This is the operational meaning of topological statistics in the
toric-code fixed-point model.
\end{remark}

\section{Topological degeneracy on closed surfaces}
\label{sec:toric-code-topological-degeneracy}

We now count the ground-state degeneracy on a closed orientable surface
\(\Sigma_g\).  The torus corresponds to \(g=1\).

The toric code has \(N_e\) physical qubits.  The stabilizer group is generated by
\(N_v\) star operators and \(N_p\) plaquette operators, but the two relations in
\cref{eq:toric-code-global-relations} remove two independent generators.  Hence
the number of independent stabilizers is
\begin{equation}
\label{eq:toric-code-independent-stabilizers}
    r=N_v+N_p-2 .
\end{equation}
By the stabilizer counting formula \cref{eq:stabilizer-dimension-counting},
\begin{equation}
\label{eq:toric-code-gsd-counting-first}
    \dim\calH_{\rm GS}
    =
    2^{N_e-r}
    =
    2^{N_e-N_v-N_p+2} .
\end{equation}
Using the Euler characteristic
\begin{equation}
\label{eq:euler-characteristic-genus}
    N_v-N_e+N_p=2-2g,
\end{equation}
we obtain
\begin{equation}
\label{eq:toric-code-gsd-genus-g}
    \boxed{
    \dim\calH_{\rm GS}(\Sigma_g)=2^{2g} .
    }
\end{equation}
On a torus, \(g=1\), this gives
\begin{equation}
\label{eq:toric-code-gsd-torus}
    \dim\calH_{\rm GS}(T^2)=4 .
\end{equation}

\begin{theorem}[Toric-code topological degeneracy]
On a closed orientable surface of genus \(g\), the toric-code Hamiltonian has
\(2^{2g}\) exactly degenerate ground states.  In particular, on the torus the
ground-state degeneracy is four.
\end{theorem}

\begin{remark}[Why the degeneracy is not symmetry breaking]
The degeneracy in \cref{eq:toric-code-gsd-genus-g} is not caused by a local order
parameter.  It is tied to noncontractible cycles of the spatial manifold.  On a
sphere, \(g=0\), the ground state is unique.  On a torus, the four ground states
cannot be distinguished by any strictly local operator in the thermodynamic
limit; they are distinguished by nonlocal Wilson loop operators.
\end{remark}

\section{Logical Wilson loops on the torus}
\label{sec:toric-code-logical-loops}

Let \(\alpha\) and \(\beta\) be the two noncontractible cycles of a torus.  Choose
direct-lattice cycles \(\gamma_{\alpha}\), \(\gamma_{\beta}\) and dual-lattice
cycles \(\gamma^\ast_{\alpha}\), \(\gamma^\ast_{\beta}\).  The corresponding
closed string operators are
\begin{equation}
\label{eq:toric-code-logical-loops}
    \overline{Z}_{\alpha}=W_e(\gamma_{\alpha}),
    \qquad
    \overline{Z}_{\beta}=W_e(\gamma_{\beta}),
    \qquad
    \overline{X}_{\alpha}=W_m(\gamma^\ast_{\alpha}),
    \qquad
    \overline{X}_{\beta}=W_m(\gamma^\ast_{\beta}) .
\end{equation}
All four operators commute with \(H_{\rm TC}\), because their strings are closed.
However, they do not all commute with one another.  Their algebra is determined
by intersection numbers:
\begin{equation}
\label{eq:toric-code-logical-loop-algebra}
    \overline{Z}_{\alpha}\overline{X}_{\beta}
    =
    -\overline{X}_{\beta}\overline{Z}_{\alpha},
    \qquad
    \overline{Z}_{\beta}\overline{X}_{\alpha}
    =
    -\overline{X}_{\alpha}\overline{Z}_{\beta},
\end{equation}
while pairs with zero intersection commute.  Thus the ground-state space on the
torus carries two encoded qubits.  One may choose \(\overline{Z}_{\alpha}\) and
\(\overline{Z}_{\beta}\) as a commuting set of labels, while the corresponding
\(\overline{X}\)-loops act as logical bit flips.

This is the stabilizer-code interpretation of the toric code.  Local errors
create anyon pairs, while noncontractible error strings implement logical
operators.  The code distance is set by the shortest noncontractible cycle of
the lattice.

\section{Local indistinguishability and robustness}
\label{sec:toric-code-local-indistinguishability}

Let \(\ket{\Omega_a}\) and \(\ket{\Omega_b}\) be two different ground states on a
large torus, distinguished by Wilson-loop eigenvalues.  A local operator
\(O_{\rm loc}\) supported on a contractible region cannot measure a
noncontractible flux.  Within the exactly solvable fixed-point model, this gives
the selection rule
\begin{equation}
\label{eq:toric-code-local-indistinguishability}
    \bra{\Omega_a}O_{\rm loc}\ket{\Omega_b}
    =
    C(O_{\rm loc})\delta_{ab} .
\end{equation}
This holds for operators whose support is small compared with the system size, up to
finite-size corrections when perturbations are added.

Physically, this is why the toric-code degeneracy is stable against weak local
perturbations in the thermodynamic limit.  A perturbation can split the ground
states only through virtual processes that create an anyon pair, move one anyon
around a noncontractible cycle, and annihilate the pair.  The amplitude of such a
process is exponentially small in the linear system size in the gapped phase.

\begin{principle}[Exact fixed point of topological order]
The toric code is not merely a solvable Hamiltonian.  It is a fixed-point
representative of a stable gapped phase.  Its universal data are the anyon types
\(\{1,e,m,\epsilon\}\), their fusion rules, their mutual statistics, and the
genus-dependent ground-state degeneracy \(2^{2g}\).
\end{principle}

\section{Summary and relation to later chapters}
\label{sec:toric-code-summary}

The exact solution of the toric code is based on four elementary facts:
\begin{enumerate}[label=(\roman*)]
    \item the Hamiltonian is a sum of commuting stabilizers;
    \item star and plaquette violations define \(e\) and \(m\) excitations;
    \item open strings create anyon pairs, while closed noncontractible strings
    act within the ground-state space;
    \item the ground-state degeneracy on \(\Sigma_g\) is \(2^{2g}\).
\end{enumerate}
This gives the simplest example of intrinsic topological order in these notes.
The following chapters generalize the same structural ideas in two directions.
Quantum double models replace the group \(\ZZ_2\) by a finite group \(G\), leading
in general to non-Abelian anyons.  The Kitaev honeycomb model keeps spin-\(1/2\)
degrees of freedom in two dimensions but solves them by mapping to Majorana
fermions moving in a static \(\ZZ_2\) gauge field.  Thus the toric code is the
commuting-projector endpoint of a broader family of exactly solvable
spin-liquid models.


\chapter{Kitaev Quantum Double Models}
\label{chap:kitaev-quantum-double-models}

The toric code is the quantum double model for the finite group
\(G=\ZZ_2\).  The purpose of this chapter is to explain the general finite-group construction.  The resulting models, usually denoted by \(D(G)\), are
commuting-projector Hamiltonians whose local Hilbert spaces are group algebras
rather than qubits.  They are exactly solvable for every finite group \(G\), and
for non-Abelian \(G\) they give fixed-point lattice models with non-Abelian
anyon excitations.

There are two guiding ideas.
\begin{enumerate}[label=(\roman*)]
    \item The toric-code star operator is a local \(\ZZ_2\) gauge
    transformation.  For a general finite group \(G\), the vertex term is the
    projector onto local gauge-invariant states.
    \item The toric-code plaquette operator enforces zero \(\ZZ_2\) flux.  For a
    general finite group \(G\), the plaquette term projects onto states with
    trivial group holonomy around each plaquette.
\end{enumerate}
The exact solvability follows from the same mechanism as in the toric code:
local gauge projectors and flatness projectors commute.

\section{From \texorpdfstring{\(\ZZ_2\)}{Z2} to finite groups}
\label{sec:quantum-double-from-z2-to-g}

Let \(G\) be a finite group with identity element \(e\).  The local Hilbert
space on each oriented edge \(\ell\) is the group algebra
\begin{equation}
\label{eq:edge-hilbert-group-algebra}
    \calH_{\ell}=\CC[G]
    =\operatorname{span}\{\ket{g}:g\in G\},
    \qquad
    \braket{g}{g'}=\delta_{g,g'} .
\end{equation}
The total Hilbert space on a lattice \(\Gamma=(V,E,P)\), with vertices \(V\),
edges \(E\), and plaquettes \(P\), is
\begin{equation}
\label{eq:quantum-double-total-hilbert-space}
    \calH_G=\bigotimes_{\ell\in E}\CC[G]_{\ell},
    \qquad
    \dim \calH_G=\abs{G}^{\abs{E}} .
\end{equation}
A basis state is a group-valued lattice gauge field
\begin{equation}
\label{eq:group-lattice-configuration}
    \ket{\{g_{\ell}\}}
    =\bigotimes_{\ell\in E}\ket{g_{\ell}}_{\ell} .
\end{equation}
Each edge is assigned an arbitrary reference orientation.  If a path traverses
an edge \(\ell\) in the reference direction, it contributes \(g_{\ell}\); if it
traverses \(\ell\) in the opposite direction, it contributes \(g_{\ell}^{-1}\).
This is the lattice analogue of parallel transport by a gauge connection.

For every \(h\in G\), define left and right multiplication operators on
\(\CC[G]\) by
\begin{equation}
\label{eq:left-right-regular-actions}
    L^h\ket{g}=\ket{hg},
    \qquad
    R^h\ket{g}=\ket{g h^{-1}} .
\end{equation}
The inverse in the definition of \(R^h\) is convenient because both \(L^h\) and
\(R^h\) define honest group actions:
\begin{equation}
\label{eq:left-right-group-actions}
    L^hL^k=L^{hk},
    \qquad
    R^hR^k=R^{hk},
    \qquad
    [L^h,R^k]=0 .
\end{equation}

\begin{remark}[Orientation convention]
A local gauge transformation at a vertex \(v\) should multiply an outgoing edge
variable from the left and an incoming edge variable from the right.  With the
convention \cref{eq:left-right-regular-actions}, this means
\begin{equation}
\label{eq:gauge-transform-edge-variable}
    g_{\ell}\longmapsto
    \begin{cases}
        h g_{\ell}, & \ell\text{ starts at }v,\\
        g_{\ell}h^{-1}, & \ell\text{ ends at }v.
    \end{cases}
\end{equation}
This is the lattice version of \(U_{ij}\mapsto h_iU_{ij}h_j^{-1}\).
\end{remark}

\section{Vertex gauge transformations}
\label{sec:quantum-double-vertex-operators}

For a vertex \(v\in V\) and a group element \(h\in G\), define
\begin{equation}
\label{eq:vertex-gauge-transformation}
    A_v^h
    =
    \prod_{\ell:\,s(\ell)=v} L_{\ell}^{h}
    \prod_{\ell:\,t(\ell)=v} R_{\ell}^{h},
\end{equation}
where \(s(\ell)\) and \(t(\ell)\) denote the source and target of the oriented
edge \(\ell\).  Thus \(A_v^h\) implements a gauge transformation by \(h\) at the
single vertex \(v\).  The vertex projector is the group average
\begin{equation}
\label{eq:quantum-double-vertex-projector}
    A_v=\frac{1}{\abs{G}}\sum_{h\in G}A_v^h .
\end{equation}

\begin{proposition}[Vertex projectors]
The operators \(A_v\) are mutually commuting projectors:
\begin{equation}
\label{eq:vertex-projector-algebra}
    A_v^2=A_v,
    \qquad
    [A_v,A_w]=0 .
\end{equation}
\end{proposition}

\begin{proof}
For a fixed vertex \(v\), the map \(h\mapsto A_v^h\) is a representation of
\(G\), so
\begin{equation}
\label{eq:vertex-projector-proof-1}
    A_v^2
    =\frac{1}{\abs{G}^2}\sum_{h,k\in G}A_v^{hk}
    =\frac{1}{\abs{G}}\sum_{q\in G}A_v^q
    =A_v .
\end{equation}
For \(v\neq w\), the only possible overlap is an edge connecting \(v\) and
\(w\).  On such an edge one operator acts by left multiplication and the other
by right multiplication.  These commute by \cref{eq:left-right-group-actions}.
Therefore \([A_v^h,A_w^k]=0\) for all \(h,k\), and hence
\([A_v,A_w]=0\).
\end{proof}

The condition
\begin{equation}
\label{eq:gauge-invariant-state-condition}
    A_v\ket{\psi}=\ket{\psi}
\end{equation}
is the lattice Gauss-law constraint.  Thus the vertex terms impose local gauge
invariance.

\section{Plaquette holonomy and flatness projectors}
\label{sec:quantum-double-plaquette-operators}

Choose an orientation and a base vertex \(v\in\partial p\) for each plaquette
\(p\).  Let \(\partial p=(\ell_1,\ldots,\ell_n)\) be the ordered boundary path
starting and ending at \(v\).  Define the plaquette holonomy
\begin{equation}
\label{eq:plaquette-holonomy}
    H_{p,v}(\{g_{\ell}\})
    =
    g_{\ell_1}^{\epsilon_1}
    g_{\ell_2}^{\epsilon_2}
    \cdots
    g_{\ell_n}^{\epsilon_n},
    \qquad
    \epsilon_j=\begin{cases}
        +1, & \ell_j\text{ is traversed with the reference orientation},\\
        -1, & \ell_j\text{ is traversed against it}.
    \end{cases}
\end{equation}
For \(k\in G\), define a plaquette flux projector by
\begin{equation}
\label{eq:plaquette-flux-projector-k}
    B_p^k\ket{\{g_{\ell}\}}
    =
    \delta_{H_{p,v}(\{g_{\ell}\}),k}\ket{\{g_{\ell}\}} .
\end{equation}
The Hamiltonian uses the trivial-flux projector
\begin{equation}
\label{eq:quantum-double-plaquette-projector}
    B_p:=B_p^e .
\end{equation}

\begin{remark}[Basepoint dependence]
For non-Abelian \(G\), the holonomy around a plaquette depends on the basepoint
by conjugation.  If the starting point is changed, \(H_{p,v}\) is replaced by
\(xH_{p,v}x^{-1}\) for some boundary transport \(x\).  Therefore the condition
\(H_{p,v}=e\) is basepoint-independent, while the condition \(H_{p,v}=k\) is not
basepoint-independent unless it is stated up to conjugacy.
\end{remark}

The plaquette projectors are diagonal in the group basis and obey
\begin{equation}
\label{eq:plaquette-projector-algebra}
    (B_p^k)^2=B_p^k,
    \qquad
    B_p^kB_p^{k'}=\delta_{k,k'}B_p^k,
    \qquad
    \sum_{k\in G}B_p^k=\id .
\end{equation}
The physical plaquette term \(B_p=B_p^e\) enforces flatness of the lattice gauge
field at \(p\).

\section{The quantum double Hamiltonian}
\label{sec:quantum-double-hamiltonian}

The Kitaev quantum double Hamiltonian is
\begin{equation}
\label{eq:quantum-double-hamiltonian}
    H_{D(G)}
    =
    -J_e\sum_{v\in V}A_v
    -J_m\sum_{p\in P}B_p,
    \qquad
    J_e,J_m>0 .
\end{equation}
Equivalently, up to an additive constant,
\begin{equation}
\label{eq:quantum-double-projector-form}
    H_{D(G)}
    =
    J_e\sum_{v\in V}(\id-A_v)
    +J_m\sum_{p\in P}(\id-B_p)
    -(J_e\abs{V}+J_m\abs{P}) .
\end{equation}
The exact ground states are those satisfying
\begin{equation}
\label{eq:quantum-double-ground-state-conditions}
    A_v\ket{\psi}=\ket{\psi},
    \qquad
    B_p\ket{\psi}=\ket{\psi},
    \qquad
    \forall v\in V,\ p\in P .
\end{equation}
Thus the ground space consists of gauge-invariant wavefunctions supported on
flat \(G\)-connections.

\begin{theorem}[Commuting-projector structure]
The terms of \(H_{D(G)}\) commute:
\begin{equation}
\label{eq:quantum-double-commuting-projectors}
    [A_v,A_w]=0,
    \qquad
    [B_p,B_q]=0,
    \qquad
    [A_v,B_p]=0 .
\end{equation}
Therefore \(H_{D(G)}\) is an exactly solvable commuting-projector Hamiltonian.
\end{theorem}

\begin{proof}
The first equality was proved in \cref{sec:quantum-double-vertex-operators}.  The
second equality follows because all \(B_p\) are diagonal multiplication
operators in the group basis.

It remains to check \([A_v,B_p]=0\).  If \(v\notin\partial p\), the operators act
on disjoint local data or act in a way that does not change the plaquette
holonomy, so they commute.  If \(v\in\partial p\), then a gauge transformation by
\(h\) at \(v\) conjugates the plaquette holonomy based at \(v\):
\begin{equation}
\label{eq:holonomy-conjugation-under-gauge-transform}
    H_{p,v}\longmapsto hH_{p,v}h^{-1} .
\end{equation}
The property \(H_{p,v}=e\) is invariant under conjugation.  Hence
\(A_v^hB_p^e=B_p^eA_v^h\) for every \(h\), and averaging over \(h\) gives
\([A_v,B_p]=0\).
\end{proof}

\begin{remark}[Why the identity flux is special]
The operator \(B_p^k\) with fixed noncentral \(k\) does not commute with a gauge
transformation at the basepoint.  Instead,
\begin{equation}
\label{eq:flux-projector-conjugation}
    A_v^h B_p^k A_v^{h^{-1}}=B_p^{hkh^{-1}} .
\end{equation}
The gauge-invariant flux observable is therefore the conjugacy-class projector
\begin{equation}
\label{eq:conjugacy-class-flux-projector}
    B_p^C=\sum_{k\in C}B_p^k,
\end{equation}
where \(C\subset G\) is a conjugacy class.  The Hamiltonian chooses the trivial
class \(C=\{e\}\).
\end{remark}

\section{Recovery of the toric code}
\label{sec:quantum-double-recovers-toric-code}

Set \(G=\ZZ_2=\{0,1\}\) with addition modulo two.  The group algebra
\(\CC[\ZZ_2]\) is a qubit Hilbert space with computational basis
\(\ket{0},\ket{1}\).  Left multiplication by the nontrivial group element flips
the basis state:
\begin{equation}
\label{eq:z2-left-multiplication-sigma-x}
    L^1\ket{g}=\ket{g+1}=\sigma^x\ket{g} .
\end{equation}
The vertex projector becomes
\begin{equation}
\label{eq:z2-vertex-projector-toric-code}
    A_v
    =\frac{1}{2}\left(\id+\prod_{\ell\ni v}\sigma^x_{\ell}\right).
\end{equation}
The plaquette flatness condition says that the sum of edge variables around a
plaquette is zero modulo two.  In the basis where
\(\sigma^z\ket{g}=(-1)^g\ket{g}\), this is
\begin{equation}
\label{eq:z2-plaquette-projector-toric-code}
    B_p
    =\frac{1}{2}\left(\id+\prod_{\ell\in\partial p}\sigma^z_{\ell}\right).
\end{equation}
Therefore \cref{eq:quantum-double-hamiltonian} becomes the toric-code
Hamiltonian up to an additive constant and an overall normalization of the
couplings.  In this precise sense,
\begin{equation}
\label{eq:toric-code-as-dz2}
    \text{toric code}=D(\ZZ_2) .
\end{equation}

\section{Ground states as flat gauge fields}
\label{sec:quantum-double-ground-states-flat-connections}

Let the lattice be embedded on a closed orientable surface \(\Sigma_g\).  The
constraint \(B_p=1\) imposes trivial holonomy around every contractible
plaquette.  Therefore a ground-state configuration is a flat lattice
\(G\)-connection.  The constraint \(A_v=1\) then identifies configurations
related by local gauge transformations.

Equivalently, the ground-state data are representations of the fundamental
group into \(G\), modulo global conjugation:
\begin{equation}
\label{eq:flat-g-connections-moduli}
    \operatorname{Flat}_G(\Sigma_g)
    \simeq
    \operatorname{Hom}(\pi_1(\Sigma_g),G)/G .
\end{equation}
Here \(G\) acts by simultaneous conjugation.  The ground-state Hilbert space is
the space of gauge-invariant complex functions on flat connections:
\begin{equation}
\label{eq:quantum-double-ground-space-functions}
    \calH_{\rm GS}(\Sigma_g)
    \simeq
    \operatorname{Fun}\!\bigl(\operatorname{Hom}(\pi_1(\Sigma_g),G)\bigr)^G .
\end{equation}
Thus
\begin{equation}
\label{eq:quantum-double-ground-degeneracy-orbits}
    \dim \calH_{\rm GS}(\Sigma_g)
    =
    \#\bigl(\operatorname{Hom}(\pi_1(\Sigma_g),G)/G\bigr),
\end{equation}
where the right-hand side denotes the number of conjugation orbits of flat
\(G\)-connections.

For a torus \(T^2\),
\begin{equation}
\label{eq:torus-fundamental-group-commuting-pairs}
    \pi_1(T^2)=\langle a,b:ab=ba\rangle .
\end{equation}
Therefore a flat connection is a commuting pair \((x,y)\in G\times G\), modulo
simultaneous conjugation:
\begin{equation}
\label{eq:torus-ground-state-commuting-pairs}
    \dim\calH_{\rm GS}(T^2)
    =
    \#\{(x,y)\in G\times G:xy=yx\}/G .
\end{equation}
For Abelian \(G\), conjugation is trivial, and this reduces to
\begin{equation}
\label{eq:abelian-quantum-double-ground-degeneracy}
    \dim\calH_{\rm GS}(\Sigma_g)=\abs{G}^{2g} .
\end{equation}
Taking \(G=\ZZ_2\) gives the toric-code result \(2^{2g}\).

\begin{remark}[Ground-state counting is not stabilizer counting]
For \(G=\ZZ_2\), the model is a Pauli stabilizer code and the degeneracy can be
counted by independent stabilizers.  For a general finite non-Abelian group,
there is no Pauli stabilizer structure.  The correct replacement is the counting
of gauge-invariant flat connections, or equivalently the representation theory
of the quantum double algebra.
\end{remark}

\section{Electric, magnetic, and dyonic excitations}
\label{sec:quantum-double-excitations}

Excitations are local violations of the commuting projectors.
\begin{itemize}[leftmargin=2.2em]
    \item A violation of \(A_v\) is an electric charge.  It transforms
    nontrivially under local gauge transformations.
    \item A violation of \(B_p\) is a magnetic flux.  Its gauge-invariant label
    is a conjugacy class of plaquette holonomy.
    \item In a non-Abelian group, charge and flux cannot in general be separated
    into independent labels.  The elementary superselection sectors are dyons.
\end{itemize}
The precise label of a dyon is a pair
\begin{equation}
\label{eq:quantum-double-anyon-label-pair}
    (C,\rho),
\end{equation}
where \(C\subset G\) is a conjugacy class and \(\rho\) is an irreducible
representation of the centralizer of a representative element of \(C\).

To define this label, choose an element \(c\in C\).  Its centralizer is
\begin{equation}
\label{eq:centralizer-definition}
    Z_c
    =
    \{h\in G:hc=ch\} .
\end{equation}
If \(c' = xcx^{-1}\) is another representative of the same conjugacy class, then
\(Z_{c'}=xZ_cx^{-1}\), so the set of irreducible representations is naturally
identified up to conjugation.  Hence the anyon label is well-defined as
\begin{equation}
\label{eq:anyon-label-set-dg}
    \mathcal{A}_{D(G)}
    =
    \{(C,\rho): C\in\operatorname{Conj}(G),\
    \rho\in\operatorname{Irr}(Z_c),\ c\in C\} .
\end{equation}

\begin{definition}[Quantum-double anyon]
An anyon type of the quantum double model \(D(G)\) is a pair \((C,\rho)\), where
\(C\) is a magnetic flux conjugacy class and \(\rho\) is an electric charge
representation of the centralizer of that flux.
\end{definition}

This formula has several important special cases.
\begin{enumerate}[label=(\roman*)]
    \item The vacuum is \((\{e\},\mathbf{1})\), where \(\mathbf{1}\) is the
    trivial representation of \(G\).
    \item Pure electric charges have trivial flux:
    \begin{equation}
    \label{eq:pure-electric-labels}
        (\{e\},\rho),
        \qquad
        \rho\in\operatorname{Irr}(G).
    \end{equation}
    \item Pure magnetic fluxes have nontrivial conjugacy class but trivial
    centralizer representation:
    \begin{equation}
    \label{eq:pure-magnetic-labels}
        (C,\mathbf{1}_{Z_c}) .
    \end{equation}
    \item Genuine dyons have both nontrivial flux and nontrivial charge
    representation.
\end{enumerate}

\section{The local quantum double algebra}
\label{sec:local-quantum-double-algebra}

The origin of the label \((C,\rho)\) is the finite-dimensional Hopf algebra known
as the Drinfeld double \(D(G)\).  At a puncture or excitation site, there are two
kinds of local operations: measuring flux and applying gauge transformations.
They obey a crossed relation.

Let \(P^g\) denote the projector onto flux \(g\), and let \(T^h\) denote a gauge
transformation by \(h\).  Then
\begin{equation}
\label{eq:local-double-crossed-relation}
    T^h P^g = P^{hgh^{-1}}T^h .
\end{equation}
Thus flux labels are not invariant under gauge transformations; only conjugacy
classes are invariant.  The algebra generated by \(P^g\) and \(T^h\), with
\begin{equation}
\label{eq:local-double-algebra-relations}
    P^gP^{g'}=\delta_{g,g'}P^g,
    \qquad
    T^hT^{h'}=T^{hh'},
    \qquad
    \sum_{g\in G}P^g=\id,
\end{equation}
is the quantum double algebra \(D(G)\) in its standard lattice form.

\begin{proposition}[Irreducible sectors of \(D(G)\)]
The irreducible representations of the local quantum double algebra are labelled
by pairs \((C,\rho)\), where \(C\) is a conjugacy class of \(G\), and \(\rho\) is
an irreducible representation of the centralizer \(Z_c\) of any representative
\(c\in C\).
\end{proposition}

\begin{proof}[Sketch]
First diagonalize the commuting flux projectors \(P^g\).  The gauge action
\(T^h\) moves a flux \(g\) to the conjugate flux \(hgh^{-1}\).  Therefore an
irreducible representation cannot be supported on a single noncentral group
element; it is supported on an entire conjugacy class \(C\).  After choosing a
representative \(c\in C\), the transformations that preserve this representative
are precisely the centralizer elements \(Z_c\).  The remaining internal charge
quantum numbers must therefore form an irreducible representation \(\rho\) of
\(Z_c\).  This gives the pair \((C,\rho)\).
\end{proof}

The quantum dimension of the anyon \((C,\rho)\) is
\begin{equation}
\label{eq:quantum-double-quantum-dimension}
    d_{(C,\rho)}=\abs{C}\,\dim\rho .
\end{equation}
The total quantum dimension is
\begin{equation}
\label{eq:quantum-double-total-quantum-dimension}
    \mathcal{D}
    =
    \sqrt{\sum_{(C,\rho)}d_{(C,\rho)}^2}
    =\abs{G} .
\end{equation}
Indeed,
\begin{equation}
\label{eq:quantum-double-total-dimension-check}
    \sum_{(C,\rho)}d_{(C,\rho)}^2
    =
    \sum_C \abs{C}^2\sum_{\rho\in\operatorname{Irr}(Z_c)}(\dim\rho)^2
    =
    \sum_C \abs{C}^2\abs{Z_c}
    =
    \sum_C \abs{C}\abs{G}
    =
    \abs{G}^2 .
\end{equation}

The topological spin of \((C,\rho)\) is
\begin{equation}
\label{eq:quantum-double-topological-spin}
    \theta_{(C,\rho)}
    =
    \frac{\chi_{\rho}(c)}{\dim\rho},
    \qquad c\in C,
\end{equation}
where \(\chi_{\rho}\) is the character of \(\rho\).  Since \(c\) lies in the
center of its own centralizer \(Z_c\), Schur's lemma ensures that this ratio is a
phase for a unitary irreducible representation.

\section{Abelian case and the toric-code anyons}
\label{sec:abelian-quantum-double-anyons}

If \(G\) is Abelian, every conjugacy class contains one element and every
centralizer is the whole group.  Therefore
\begin{equation}
\label{eq:abelian-anyon-labels}
    \mathcal{A}_{D(G)}
    =
    G\times \widehat{G},
\end{equation}
where \(\widehat{G}=\operatorname{Irr}(G)\) is the character group.  An anyon is
labelled by a magnetic flux \(g\in G\) and an electric character
\(\chi\in\widehat{G}\).  All anyons have quantum dimension one:
\begin{equation}
\label{eq:abelian-quantum-double-dimensions}
    d_{(g,\chi)}=1,
    \qquad
    \mathcal{D}=\abs{G} .
\end{equation}
The braiding phase between \((g,\chi)\) and \((h,\psi)\) is determined by the
Aharonov--Bohm pairing between charge and flux:
\begin{equation}
\label{eq:abelian-double-mutual-braiding}
    M_{(g,\chi),(h,\psi)}
    =
    \chi(h)\,\psi(g) .
\end{equation}

For \(G=\ZZ_2\), there are two group elements \(0,1\) and two characters
\(\mathbf{1},\chi\), with \(\chi(1)=-1\).  The four anyons are
\begin{equation}
\label{eq:toric-code-anyon-labels-from-dz2}
    1=(0,\mathbf{1}),
    \qquad
    e=(0,\chi),
    \qquad
    m=(1,\mathbf{1}),
    \qquad
    \epsilon=(1,\chi).
\end{equation}
This reproduces the toric-code fusion rules
\begin{equation}
\label{eq:toric-code-fusion-from-dz2}
    e^2=m^2=\epsilon^2=1,
    \qquad
    e\times m=\epsilon .
\end{equation}
The mutual phase between \(e\) and \(m\) is \(-1\), because
\begin{equation}
\label{eq:toric-code-mutual-phase-from-character}
    M_{e,m}=\chi(1)=-1 .
\end{equation}

\section{Ribbon operators}
\label{sec:quantum-double-ribbon-operators}

In the toric code, open string operators create pairs of anyons at their
endpoints.  In the non-Abelian quantum double model, the correct generalization
is a ribbon operator.  A ribbon is a strip made from a direct-lattice path and a
nearby dual-lattice path.  It carries both electric and magnetic information.

A ribbon operator has two conceptual components.
\begin{enumerate}[label=(\roman*)]
    \item A magnetic component changes group variables on crossed edges and
    creates plaquette fluxes at the ribbon endpoints.
    \item An electric component parallel-transports representation indices along
    the direct path and creates gauge-charge violations at the endpoints.
\end{enumerate}
Closed ribbon operators commute with the Hamiltonian and measure topological
charge.  Open ribbon operators commute with all projectors away from their
endpoints and create a pair of excitations whose total topological charge is
trivial.

For pure electric probes, one obtains Wilson-loop operators.  If \(\gamma\) is a
closed oriented loop and \(R\) is a representation of \(G\), define
\begin{equation}
\label{eq:quantum-double-wilson-loop}
    W_R(\gamma)
    =
    \chi_R\!\left(\prod_{\ell\in\gamma}g_{\ell}^{\epsilon_{\ell}}\right),
\end{equation}
where \(\chi_R\) is the character of \(R\).  The trace makes
\(W_R(\gamma)\) invariant under conjugation of the loop holonomy.  On a
contractible loop in the ground state, flatness forces \(W_R(\gamma)=\dim R\).
On noncontractible loops, Wilson operators distinguish different topological
sectors.

\begin{remark}[Why ribbons replace strings]
For Abelian \(G\), electric and magnetic string operators can be separated
cleanly, and their algebra is essentially a phase.  For non-Abelian \(G\), moving
flux changes the charge frame by conjugation.  A ribbon keeps track of both the
flux transport and the charge indices.  This is why the excitation label is not
``flux plus representation of \(G\)'', but rather ``flux conjugacy class plus
representation of the corresponding centralizer.''
\end{remark}

\section{Example: \texorpdfstring{\(D(S_3)\)}{D(S3)}}
\label{sec:quantum-double-s3-example}

The smallest non-Abelian example is \(G=S_3\), the permutation group of three
objects.  Its conjugacy classes are
\begin{equation}
\label{eq:s3-conjugacy-classes}
    C_e=\{e\},
    \qquad
    C_t=\{(12),(23),(13)\},
    \qquad
    C_c=\{(123),(132)\} .
\end{equation}
The corresponding centralizers have sizes
\begin{equation}
\label{eq:s3-centralizer-sizes}
    \abs{Z_e}=6,
    \qquad
    \abs{Z_{(12)}}=2,
    \qquad
    \abs{Z_{(123)}}=3 .
\end{equation}
The irreducible representations are therefore distributed as follows:
\begin{equation}
\label{eq:s3-anyon-count}
    \#\mathcal{A}_{D(S_3)}
    =
    \#\operatorname{Irr}(S_3)
    +\#\operatorname{Irr}(\ZZ_2)
    +\#\operatorname{Irr}(\ZZ_3)
    =3+2+3=8 .
\end{equation}
Thus \(D(S_3)\) has eight anyon types.  Their quantum dimensions are
\begin{equation}
\label{eq:s3-quantum-dimensions-list}
    1,1,2,
    \qquad
    3,3,
    \qquad
    2,2,2,
\end{equation}
where the first group comes from pure charges over \(C_e\), the second from
transposition fluxes, and the third from three-cycle fluxes.  The dimension
check gives
\begin{equation}
\label{eq:s3-total-quantum-dimension-check}
    1^2+1^2+2^2+3^2+3^2+2^2+2^2+2^2=36=\abs{S_3}^2 .
\end{equation}
This example illustrates the genuinely non-Abelian feature: some anyons have
quantum dimension larger than one.

\section{Relation to topological quantum field theory}
\label{sec:quantum-double-tqft-viewpoint}

The model \(D(G)\) is the Hamiltonian lattice realization of finite-group
Dijkgraaf--Witten gauge theory with trivial cocycle.  Its input data are purely
group-theoretic, but its low-energy physics is topological.  The ground states
only depend on the topology of the spatial surface, and quasiparticles are
classified by the representation category of the Drinfeld double,
\begin{equation}
\label{eq:quantum-double-category}
    \text{anyon category of }D(G)
    =
    \operatorname{Rep}(D(G)) .
\end{equation}
The toric code corresponds to \(\operatorname{Rep}(D(\ZZ_2))\).

The commuting-projector Hamiltonian should be viewed as a fixed-point
Hamiltonian in the same way that the toric code is a fixed-point Hamiltonian of
\(\ZZ_2\) topological order.  Perturbing the Hamiltonian generally destroys
exact solvability, but small local perturbations cannot immediately remove the
topological phase as long as the bulk gap remains open.

\section{Summary and comparison with the toric code}
\label{sec:quantum-double-summary}

The quantum double construction replaces every toric-code \(\ZZ_2\) ingredient
by its finite-group analogue.

{\small
\begin{longtable}{L{0.27\textwidth}L{0.30\textwidth}L{0.32\textwidth}}
\caption{Toric code versus finite-group quantum double models.}
\label{tab:toric-code-vs-quantum-double}\\
\toprule
Structure & Toric code & Quantum double \(D(G)\) \\
\midrule
\endfirsthead
\toprule
Structure & Toric code & Quantum double \(D(G)\) \\
\midrule
\endhead
Local edge Hilbert space & Qubit \(\CC[\ZZ_2]\) & Group algebra \(\CC[G]\) \\
Vertex term & \(\ZZ_2\) star constraint & Gauge-invariance projector \(A_v=\abs{G}^{-1}\sum_h A_v^h\) \\
Plaquette term & \(\ZZ_2\) zero-flux constraint & Flatness projector \(B_p^e\) \\
Ground states & Flat \(\ZZ_2\) gauge fields modulo gauge transformations & Flat \(G\)-connections modulo gauge transformations \\
Anyon labels & \(1,e,m,\epsilon\) & Pairs \((C,\rho)\) with \(C\) a conjugacy class and \(\rho\in\operatorname{Irr}(Z_c)\) \\
Abelian or non-Abelian & Abelian & Abelian iff \(G\) is Abelian; generally non-Abelian \\
Total quantum dimension & \(2\) & \(\abs{G}\) \\
\bottomrule
\end{longtable}}

\boxedpar{%
The toric code is not an isolated construction.  It is the \(G=\ZZ_2\) member of
the finite-group quantum double family.  The general replacement is
\[
\ZZ_2\text{ gauge theory}
\quad\longrightarrow\quad
G\text{ lattice gauge theory},
\]
with electric charges, magnetic fluxes, and dyons unified by the representation
theory of \(D(G)\).
}

The next chapter studies a different kind of two-dimensional exact solvability:
the Kitaev honeycomb model.  There the Hamiltonian is not a commuting-projector
Hamiltonian.  Instead, spin interactions fractionalize into free Majorana
fermions moving in a static \(\ZZ_2\) gauge field.


\chapter{The Kitaev Honeycomb Model}
\label{chap:kitaev-honeycomb-model}

The toric code and the quantum double models are exactly solvable because their
Hamiltonians are sums of commuting projectors.  The Kitaev honeycomb model is
solvable by a different, but closely related, mechanism.  It is not a
commuting-projector Hamiltonian.  Instead, its spin degrees of freedom can be
represented by Majorana fermions coupled to a static \(\ZZ_2\) gauge field.  In
each fixed gauge-flux sector, the Hamiltonian reduces to a free Majorana hopping
problem.

This chapter therefore serves as the bridge between two ideas developed earlier
in these notes:
\begin{enumerate}[label=(\roman*)]
    \item free-fermion/BdG solvability, as in the transverse-field XY chain and
    the Kitaev chain;
    \item gauge-sector solvability, as in the toric code and finite-group
    quantum double models.
\end{enumerate}
The special feature of the honeycomb model is that these two structures coexist
in a two-dimensional spin system.  The resulting phase diagram contains gapped
Abelian topological phases, a gapless Majorana spin liquid, and, after a
small time-reversal-breaking perturbation, a gapped non-Abelian Ising-anyon
phase.

\section{Honeycomb lattice and bond-dependent exchange}
\label{sec:kitaev-honeycomb-hamiltonian}

Consider a honeycomb lattice.  Each vertex carries a spin-\(1/2\) degree of
freedom, and the nearest-neighbour bonds are colored by three labels
\(\alpha=x,y,z\).  Every site touches exactly one bond of each type.  The Kitaev
honeycomb Hamiltonian is
\begin{equation}
\label{eq:kitaev-honeycomb-spin-hamiltonian}
    H_{\rm KHC}
    =
    -J_x\sum_{\langle ij\rangle_x}\sigma_i^x\sigma_j^x
    -J_y\sum_{\langle ij\rangle_y}\sigma_i^y\sigma_j^y
    -J_z\sum_{\langle ij\rangle_z}\sigma_i^z\sigma_j^z .
\end{equation}
Here \(\langle ij\rangle_\alpha\) denotes a nearest-neighbour bond of type
\(\alpha\).  Unless otherwise stated, we take
\begin{equation}
\label{eq:kitaev-couplings-positive}
    J_x,J_y,J_z\geq0 .
\end{equation}
Changing signs of the couplings can often be absorbed by sublattice-dependent
spin rotations, but keeping the positive-coupling convention makes the phase
diagram transparent.

\begin{remark}[Why the model is not trivially solvable]
Each local term in \cref{eq:kitaev-honeycomb-spin-hamiltonian} is an Ising
exchange, but different bonds use different spin components.  Therefore terms
sharing one site usually do not commute.  For example, if an \(x\)-bond and a
\(y\)-bond meet at site \(i\), then the local Pauli matrices
\(\sigma_i^x\) and \(\sigma_i^y\) anticommute.  Thus the model is not a
commuting-projector model.  Its exact solvability is more subtle: noncommuting
spin interactions become quadratic Majorana hopping terms in a background of
static \(\ZZ_2\) link variables.
\end{remark}

The honeycomb lattice is bipartite.  We denote the two sublattices by
\(A\) and \(B\).  Choose every bond orientation from \(A\) to \(B\).  A unit cell
contains one \(A\)-site and one \(B\)-site.  Let \(\bm a_1\) and \(\bm a_2\) be
primitive lattice vectors chosen so that, in the vortex-free gauge used below,
the three bonds from an \(A\)-site at cell \(\bm r\) connect to
\begin{equation}
\label{eq:honeycomb-neighbour-convention}
    B_{\bm r},
    \qquad
    B_{\bm r+\bm a_1},
    \qquad
    B_{\bm r+\bm a_2},
\end{equation}
with couplings \(J_z\), \(J_x\), and \(J_y\), respectively.  This choice fixes
the later Bloch function
\begin{equation}
\label{eq:kitaev-bloch-function-def-preview}
    f(\bm k)=J_z+J_x\ee^{\ii \bm k\cdot \bm a_1}
             +J_y\ee^{\ii \bm k\cdot \bm a_2} .
\end{equation}
Other conventions only multiply \(f(\bm k)\) by phase factors or relabel the
momenta; the spectrum is unchanged.

\section{Plaquette fluxes in the spin language}
\label{sec:kitaev-spin-fluxes}

Before introducing Majorana fermions, one can already identify conserved
plaquette operators directly in the spin model.  Let \(p\) be a hexagonal
plaquette and label its six vertices cyclically by \(1,2,3,4,5,6\).  With the
usual bond coloring around a hexagon, define
\begin{equation}
\label{eq:kitaev-plaquette-spin-operator}
    W_p
    =
    \sigma_1^x\sigma_2^y\sigma_3^z
    \sigma_4^x\sigma_5^y\sigma_6^z,
\end{equation}
up to a cyclic relabeling determined by the local bond convention.  The precise
sequence of \(x,y,z,x,y,z\) may be shifted, but the operator is always the
product of the spin component not associated with the outgoing external bond at
each vertex of the plaquette.

These operators obey
\begin{equation}
\label{eq:kitaev-flux-algebra-spin}
    W_p^2=\id,
    \qquad
    [W_p,W_{p'}]=0,
    \qquad
    [W_p,H_{\rm KHC}]=0 .
\end{equation}
Thus every plaquette carries a conserved \(\ZZ_2\) flux eigenvalue
\begin{equation}
\label{eq:kitaev-flux-eigenvalue}
    W_p=w_p,
    \qquad
    w_p=\pm1 .
\end{equation}
A plaquette with \(w_p=-1\) is called a vortex or \(\ZZ_2\) flux excitation.

\begin{remark}[Fluxes versus commuting projectors]
The fluxes \(W_p\) commute with the Hamiltonian, but the Hamiltonian itself is
not a sum of the \(W_p\).  They are conserved labels of superselection sectors,
not the complete set of local energy projectors.  After fixing the flux sector,
one still has to diagonalize itinerant Majorana fermions.
\end{remark}

On a closed surface, the plaquette fluxes are not all independent.  For example,
on a torus,
\begin{equation}
\label{eq:kitaev-product-flux-constraint}
    \prod_p W_p=\id,
\end{equation}
so vortices are created in pairs on a closed lattice.  There are also nonlocal
Wilson-loop sectors associated with the two noncontractible cycles of the torus.
These global sectors will later become important for the finite-size ground
state degeneracy.

\section{Four-Majorana representation}
\label{sec:kitaev-four-majorana-representation}

The exact solution begins by enlarging the Hilbert space.  At every site \(i\),
introduce four Hermitian Majorana operators
\begin{equation}
\label{eq:kitaev-four-majoranas}
    b_i^x,
    \qquad
    b_i^y,
    \qquad
    b_i^z,
    \qquad
    c_i,
\end{equation}
satisfying
\begin{equation}
\label{eq:kitaev-majorana-algebra}
    \{b_i^\alpha,b_j^\beta\}=2\delta_{ij}\delta_{\alpha\beta},
    \qquad
    \{c_i,c_j\}=2\delta_{ij},
    \qquad
    \{b_i^\alpha,c_j\}=0 .
\end{equation}
The spin operators are represented as
\begin{equation}
\label{eq:kitaev-spin-majorana-representation}
    \sigma_i^\alpha=\ii b_i^\alpha c_i,
    \qquad
    \alpha=x,y,z .
\end{equation}

This representation is redundant.  Four Majoranas span a four-dimensional local
Hilbert space, whereas one spin-\(1/2\) has dimension two.  The physical
subspace is selected by the local gauge constraint
\begin{equation}
\label{eq:kitaev-local-gauge-constraint}
    D_i\ket{\psi}_{\rm phys}=\ket{\psi}_{\rm phys},
    \qquad
    D_i=b_i^x b_i^y b_i^z c_i .
\end{equation}
The operator \(D_i\) satisfies \(D_i^2=\id\).  It generates a local
\(\ZZ_2\) gauge transformation:
\begin{equation}
\label{eq:kitaev-gauge-flip-majoranas}
    D_i b_i^\alpha D_i=-b_i^\alpha,
    \qquad
    D_i c_i D_i=-c_i,
\end{equation}
while leaving Majoranas on other sites invariant.  Since each spin operator
\(\sigma_i^\alpha=\ii b_i^\alpha c_i\) contains two Majoranas at site \(i\), it
is gauge invariant.

The projector onto the physical spin Hilbert space is
\begin{equation}
\label{eq:kitaev-physical-projector}
    P_{\rm phys}
    =
    \prod_i \frac{1+D_i}{2} .
\end{equation}
On finite lattices, especially with periodic boundary conditions, this projection
must be imposed carefully because it constrains the allowed fermion parity in
each gauge sector.  For deriving the bulk spectrum and the phase diagram, it is
usually sufficient to diagonalize in the enlarged Majorana space and then project
physical states at the end.

\begin{proposition}[Pauli algebra on the physical subspace]
The operators \(\sigma_i^\alpha=\ii b_i^\alpha c_i\) satisfy the spin-\(1/2\)
Pauli algebra on the subspace \(D_i=+1\):
\begin{equation}
\label{eq:kitaev-pauli-algebra-physical}
    \sigma_i^\alpha\sigma_i^\beta
    =
    \delta_{\alpha\beta}\id
    +\ii\sum_{\gamma}\epsilon_{\alpha\beta\gamma}\sigma_i^\gamma .
\end{equation}
\end{proposition}

\begin{proof}
For \(\alpha=\beta\), one has
\((\ii b_i^\alpha c_i)^2=\id\).  For distinct \(\alpha,\beta,\gamma\),
\begin{equation}
\label{eq:kitaev-pauli-proof-product}
    \sigma_i^\alpha\sigma_i^\beta
    =
    (\ii b_i^\alpha c_i)(\ii b_i^\beta c_i)
    =
    b_i^\alpha b_i^\beta .
\end{equation}
Using \(D_i=b_i^x b_i^y b_i^z c_i=+1\), one obtains
\begin{equation}
\label{eq:kitaev-pauli-proof-constraint-use}
    b_i^\alpha b_i^\beta
    =
    \ii\epsilon_{\alpha\beta\gamma}\,\ii b_i^\gamma c_i
    =
    \ii\epsilon_{\alpha\beta\gamma}\sigma_i^\gamma,
\end{equation}
with the sign fixed by the ordering \((x,y,z)\).  This gives
\cref{eq:kitaev-pauli-algebra-physical}.
\end{proof}

\section{Static \texorpdfstring{\(\ZZ_2\)}{Z2} gauge field}
\label{sec:kitaev-static-z2-gauge-field}

Substituting \cref{eq:kitaev-spin-majorana-representation} into the spin
Hamiltonian gives
\begin{align}
\label{eq:kitaev-majorana-hamiltonian-derivation}
    -J_\alpha\sigma_i^\alpha\sigma_j^\alpha
    &=-J_\alpha(\ii b_i^\alpha c_i)(\ii b_j^\alpha c_j)  \notag\\
    &=\ii J_\alpha u_{ij} c_i c_j,
    \qquad
    \langle ij\rangle_\alpha,
\end{align}
where the link operator is
\begin{equation}
\label{eq:kitaev-link-variable}
    u_{ij}=\ii b_i^\alpha b_j^\alpha,
    \qquad
    \langle ij\rangle_\alpha .
\end{equation}
The orientation convention is important: throughout this chapter \(i\in A\) and
\(j\in B\), so \(u_{ji}=-u_{ij}\).  Since \(u_{ij}^2=\id\), its eigenvalues are
\begin{equation}
\label{eq:kitaev-link-eigenvalues}
    u_{ij}=\pm1 .
\end{equation}
The Hamiltonian becomes
\begin{equation}
\label{eq:kitaev-majorana-hamiltonian-static-gauge}
    H_{\rm KHC}
    =
    \ii\sum_{\langle ij\rangle_\alpha}J_\alpha u_{ij}c_i c_j .
\end{equation}

\begin{proposition}[Static link variables]
The link operators satisfy
\begin{equation}
\label{eq:kitaev-u-commutation}
    [u_{ij},u_{kl}]=0,
    \qquad
    [u_{ij},H_{\rm KHC}]=0 .
\end{equation}
Therefore each \(u_{ij}\) is a conserved \(\ZZ_2\) variable in the enlarged
Majorana Hilbert space.
\end{proposition}

\begin{proof}
Two link variables either act on disjoint Majoranas or share one site but use
Majoranas of different bond types.  In both cases the total number of Majorana
interchanges needed to swap the two bilinears is even, so they commute.

To prove \([u_{ij},H_{\rm KHC}]=0\), note first that \(u_{ij}\) trivially
commutes with the corresponding bond term
\(\ii J_\alpha u_{ij}c_i c_j\).  For any other bond term, either the two bonds
are disjoint, or they share one site but involve different \(b\)-Majoranas and a
single \(c\)-Majorana at that site.  Again the relevant operator products differ
by an even number of Majorana interchanges.  Hence \(u_{ij}\) commutes with
every term in \cref{eq:kitaev-majorana-hamiltonian-static-gauge}.
\end{proof}

The conserved \(u_{ij}\) are not themselves gauge-invariant observables.  Under
the local gauge transformation generated by \(D_i\), all link variables adjacent
to site \(i\) change sign:
\begin{equation}
\label{eq:kitaev-link-gauge-transformation}
    u_{ij}\longmapsto -u_{ij},
    \qquad
    u_{ki}\longmapsto -u_{ki} .
\end{equation}
Gauge-invariant information is contained in products around closed loops.  For a
plaquette \(p\), define
\begin{equation}
\label{eq:kitaev-majorana-flux}
    W_p
    =
    \prod_{\langle ij\rangle\in\partial p} u_{ij},
\end{equation}
where the product follows the oriented boundary, with the sign convention fixed
by the chosen bond orientations.  This operator is the Majorana representation
of the spin flux in \cref{eq:kitaev-plaquette-spin-operator}.  The eigenvalue
\(W_p=+1\) is called the vortex-free or zero-flux sector, and \(W_p=-1\) is a
\(\ZZ_2\) vortex.

\section{Reduction to free Majorana fermions}
\label{sec:kitaev-free-majorana-reduction}

Because all \(u_{ij}\) commute with the Hamiltonian, the Hilbert space splits
into static gauge sectors.  In a fixed configuration
\(u=\{u_{ij}=\pm1\}\), the Hamiltonian is quadratic:
\begin{equation}
\label{eq:kitaev-fixed-sector-hamiltonian}
    H[u]
    =
    \ii\sum_{\langle ij\rangle_\alpha}J_\alpha u_{ij}c_i c_j .
\end{equation}
Equivalently, it can be written in the canonical Majorana quadratic form
\begin{equation}
\label{eq:kitaev-majorana-quadratic-form}
    H[u]
    =
    \frac{\ii}{4}\sum_{i,j}A_{ij}[u]c_i c_j,
    \qquad
    A_{ij}[u]=-A_{ji}[u]\in\RR .
\end{equation}
For an \(\alpha\)-bond oriented from \(i\in A\) to \(j\in B\), the matrix entries
are
\begin{equation}
\label{eq:kitaev-a-matrix-entry}
    A_{ij}[u]=2J_\alpha u_{ij},
    \qquad
    A_{ji}[u]=-2J_\alpha u_{ij} .
\end{equation}
Diagonalizing the real antisymmetric matrix \(A[u]\) gives pairs of eigenvalues
\(\pm\epsilon_n\) of \(\ii A[u]\), with \(\epsilon_n\geq0\).  The Hamiltonian in
that sector has the form
\begin{equation}
\label{eq:kitaev-diagonal-majorana-spectrum}
    H[u]
    =
    \sum_{n}\epsilon_n\left(f_n^\dagger f_n-\frac{1}{2}\right),
\end{equation}
up to the physical projection discussed in \cref{eq:kitaev-physical-projector}.
Thus the original interacting spin model is solved by the following sequence:
\begin{align}
\label{eq:kitaev-solution-sequence}
    \text{spin model}
    &\longrightarrow
    \text{Majoranas plus static }\ZZ_2\text{ gauge field}
    \notag\\
    &\longrightarrow
    \text{free Majorana band problem in each flux sector}
    \notag\\
    &\longrightarrow
    \text{physical projection} .
\end{align}

\begin{principle}[Majorana--gauge solvability]
A spin model is Majorana--gauge solvable if, after an enlarged Majorana
representation and a local gauge constraint, its Hamiltonian becomes quadratic in
itinerant Majoranas with coefficients given by conserved \(\ZZ_2\) link
variables.  The Kitaev honeycomb model is the canonical two-dimensional example
of this mechanism.
\end{principle}

\section{Vortex-free sector and Bloch Hamiltonian}
\label{sec:kitaev-vortex-free-bloch}

For the honeycomb lattice with positive couplings, the ground state lies in the
vortex-free sector.  This is a consequence of a flux theorem for reflection-
symmetric planar lattices.  We will not prove the theorem here; instead, we use
it to obtain the bulk phase diagram.

Choose the translationally invariant gauge
\begin{equation}
\label{eq:kitaev-vortex-free-gauge}
    u_{ij}=+1
    \qquad
    \text{for every bond oriented from }A\text{ to }B .
\end{equation}
Using the unit-cell convention in \cref{eq:honeycomb-neighbour-convention}, the
Hamiltonian becomes
\begin{equation}
\label{eq:kitaev-real-space-vortex-free}
    H_0
    =
    \ii\sum_{\bm r}
    \left(
    J_z c_{A,\bm r}c_{B,\bm r}
    +J_x c_{A,\bm r}c_{B,\bm r+\bm a_1}
    +J_y c_{A,\bm r}c_{B,\bm r+\bm a_2}
    \right).
\end{equation}
After Fourier transformation, the only gauge-invariant spectral data are encoded
in
\begin{equation}
\label{eq:kitaev-bloch-function}
    f(\bm k)
    =
    J_z+J_x\ee^{\ii \bm k\cdot\bm a_1}
       +J_y\ee^{\ii \bm k\cdot\bm a_2} .
\end{equation}
The Majorana quasiparticle dispersion is
\begin{equation}
\label{eq:kitaev-majorana-dispersion}
    \epsilon(\bm k)=2\abs{f(\bm k)} .
\end{equation}

\begin{remark}[Factor of two convention]
The factor of two in \cref{eq:kitaev-majorana-dispersion} follows from the
Hamiltonian normalization in \cref{eq:kitaev-honeycomb-spin-hamiltonian}.  Some
references absorb this factor into the definition of \(J_\alpha\) or into the
canonical Majorana matrix \(A\).  The physically invariant statement is the
location of zeros of \(f(\bm k)\) and the relative gap scale.
\end{remark}

The gap closes if and only if there exists a momentum \(\bm k\) such that
\begin{equation}
\label{eq:kitaev-gap-closing-condition-f}
    J_z+J_x\ee^{\ii \bm k\cdot\bm a_1}
       +J_y\ee^{\ii \bm k\cdot\bm a_2}=0 .
\end{equation}
Geometrically, this means that three complex numbers of lengths
\(J_x,J_y,J_z\) can form a triangle.  Therefore the gapless condition is
\begin{equation}
\label{eq:kitaev-triangle-inequalities}
    J_x\leq J_y+J_z,
    \qquad
    J_y\leq J_z+J_x,
    \qquad
    J_z\leq J_x+J_y .
\end{equation}
Inside this region, the spectrum has two Dirac cones.  Outside it, one coupling
dominates and the spectrum is gapped.

\begin{proposition}[Bulk phase criterion]
For positive couplings, the vortex-free Majorana spectrum is gapless precisely
inside the triangle region \cref{eq:kitaev-triangle-inequalities}.  It is gapped
in the three regions
\begin{equation}
\label{eq:kitaev-a-phase-regions}
    J_x>J_y+J_z,
    \qquad
    J_y>J_z+J_x,
    \qquad
    J_z>J_x+J_y .
\end{equation}
\end{proposition}

\begin{proof}
The only possible gap closing is \(f(\bm k)=0\).  Since the phases
\(\ee^{\ii \bm k\cdot\bm a_1}\) and
\(\ee^{\ii \bm k\cdot\bm a_2}\) can vary over the unit circle subject to the
Brillouin-zone periodicity, the equation \(f(\bm k)=0\) is equivalent to closing
a triangle with side lengths \(J_x,J_y,J_z\).  Such a triangle exists if and
only if the three triangle inequalities hold.  If one inequality is violated,
\(\abs{f(\bm k)}\) has a strictly positive lower bound.
\end{proof}

At the isotropic point
\begin{equation}
\label{eq:kitaev-isotropic-point}
    J_x=J_y=J_z,
\end{equation}
the system lies in the gapless phase.  The low-energy theory consists of two
cones of massless Majorana fermions coupled to a gapped \(\ZZ_2\) gauge field.
Since the gauge field is static in the exactly solvable model, the low-energy
fermions are free within each flux sector, but flux excitations remain massive.

\section{Gapped Abelian phases and the toric-code limit}
\label{sec:kitaev-abelian-phases-toric-code-limit}

The three gapped regions in \cref{eq:kitaev-a-phase-regions} are usually called
\(A_x\), \(A_y\), and \(A_z\).  In the \(A_z\) phase, for example,
\begin{equation}
\label{eq:kitaev-az-limit-condition}
    J_z\gg J_x,J_y .
\end{equation}
The unperturbed Hamiltonian consists of decoupled \(z\)-dimers.  Treating the
\(x\)- and \(y\)-bond interactions perturbatively, the first nontrivial
plaquette-dependent term appears at fourth order and has the form
\begin{equation}
\label{eq:kitaev-effective-toric-code}
    H_{\rm eff}^{(4)}
    =
    -J_{\rm eff}\sum_p W_p+\text{constant},
    \qquad
    J_{\rm eff}\sim \frac{J_x^2J_y^2}{J_z^3} .
\end{equation}
More precisely, the numerical prefactor depends on the perturbative convention,
but the structure is fixed: the low-energy Hamiltonian is a \(\ZZ_2\)
commuting-projector theory generated by plaquette fluxes.

This shows that the gapped \(A\) phases are Abelian topological phases in the
same universality class as the toric code.  The elementary excitations can be
identified as
\begin{enumerate}[label=(\roman*)]
    \item \(\ZZ_2\) fluxes, corresponding to plaquettes with \(W_p=-1\);
    \item fermionic matter excitations from the Majorana band;
    \item bound states of fluxes and Majorana fermions.
\end{enumerate}
Their fusion and braiding reduce, in the deep Abelian limit, to the toric-code
anyon structure.

\begin{remark}[Why the toric code reappears]
The toric code describes a deconfined \(\ZZ_2\) gauge theory with gapped charges
and fluxes.  The honeycomb model in a strongly anisotropic gapped phase has the
same long-distance topological order, but its microscopic degrees of freedom are
spins on sites rather than qubits on edges.  The emergence of the toric code is
therefore an infrared statement, not an exact microscopic equality at generic
couplings.
\end{remark}

\section{Time-reversal breaking and the non-Abelian phase}
\label{sec:kitaev-non-abelian-phase}

The gapless phase is protected by the combination of lattice symmetries and
time-reversal symmetry.  A small magnetic field
\begin{equation}
\label{eq:kitaev-magnetic-field-perturbation}
    V_h=-\sum_i\left(h_x\sigma_i^x+h_y\sigma_i^y+h_z\sigma_i^z\right)
\end{equation}
breaks time reversal.  The magnetic field itself destroys exact solvability, but
its leading nonvanishing effect within perturbation theory is a solvable
three-spin interaction of the schematic form
\begin{equation}
\label{eq:kitaev-effective-three-spin-term}
    H_\kappa
    =
    -\kappa\sum_{\langle i j k\rangle}
    \sigma_i^\alpha\sigma_j^\beta\sigma_k^\gamma,
    \qquad
    \kappa\propto \frac{h_xh_yh_z}{J^2},
\end{equation}
where the ordered triples \(\langle i j k\rangle\) are length-two paths on the
honeycomb lattice and \(\alpha,\beta,\gamma\) are fixed by the adjacent bond
types.  In the Majorana representation this term adds next-nearest-neighbour
Majorana hopping.  It opens a mass gap at the Dirac cones.

The resulting gapped phase has a nonzero Majorana Chern number,
\begin{equation}
\label{eq:kitaev-majorana-chern-number}
    \nu=\pm1,
\end{equation}
with the sign determined by the sign of \(\kappa\).  The vortices of the
\(\ZZ_2\) gauge field bind Majorana zero modes.  Consequently, vortices become
non-Abelian anyons of Ising type.

The Ising-anyon fusion rules are
\begin{equation}
\label{eq:ising-anyon-fusion-rules}
    \sigma\times\sigma=1+\psi,
    \qquad
    \sigma\times\psi=\sigma,
    \qquad
    \psi\times\psi=1 .
\end{equation}
Here \(\sigma\) denotes a vortex carrying a Majorana zero mode, and \(\psi\) is
the fermionic Majorana excitation.  The quantum dimensions are
\begin{equation}
\label{eq:ising-anyon-quantum-dimensions}
    d_1=1,
    \qquad
    d_\psi=1,
    \qquad
    d_\sigma=\sqrt{2} .
\end{equation}
The noninteger quantum dimension of \(\sigma\) is the hallmark of non-Abelian
statistics.

\begin{remark}[Relation to chiral topological superconductors]
After fixing a gauge sector, the itinerant Majoranas form a two-dimensional
class-D superconductor.  The time-reversal-breaking mass converts the gapless
Majorana semimetal into a chiral Majorana topological superconductor.  Coupling
this band topology to the dynamical \(\ZZ_2\) fluxes produces the non-Abelian
Ising topological order.
\end{remark}

\section{Fluxes, vortices, and physical projection}
\label{sec:kitaev-fluxes-physical-projection}

The Majorana solution contains two levels of description.  First, one fixes a
configuration of link variables \(u_{ij}\) and diagonalizes the free Majorana
Hamiltonian \(H[u]\).  Second, one projects the resulting states onto the
physical spin Hilbert space using \(P_{\rm phys}\).  This projection has two
important consequences.

\begin{enumerate}[label=(\roman*)]
    \item Gauge-equivalent link configurations describe the same physical flux
    sector.  A local sign flip generated by \(D_i\) changes adjacent \(u_{ij}\)
    and \(c_i\), but leaves spin observables invariant.
    \item On a finite torus, the physical projection fixes a relation between
    the gauge-field sector and the fermion parity of the matter Majoranas.  Thus
    not every free-Majorana eigenstate in the enlarged Hilbert space represents
    a physical spin state.
\end{enumerate}

For bulk thermodynamic properties, the projection usually does not change the
local dispersion or the phase boundaries.  For exact finite-size spectra and
for topological degeneracy on a torus, however, it is essential.

A useful way to summarize the structure is
\begin{equation}
\label{eq:kitaev-sector-label-summary}
    \text{physical sector}
    =
    \frac{\text{link variables }u_{ij}=\pm1}{\text{local }\ZZ_2\text{ gauge transformations}}
    +
    \text{projected Majorana matter states} .
\end{equation}
Gauge-invariant fluxes \(W_p\) and noncontractible Wilson loops label the gauge
part, while the free Majorana occupation numbers label the matter part.

\section{Spin correlations and fractionalization}
\label{sec:kitaev-spin-correlations-fractionalization}

A single spin operator is not diagonal in the flux sector.  From
\(\sigma_i^\alpha=\ii b_i^\alpha c_i\), acting with \(\sigma_i^\alpha\) changes
link variables adjacent to site \(i\).  Equivalently, in the spin language it
creates or moves a pair of fluxes on the two plaquettes adjacent to the
\(\alpha\)-bond touching site \(i\).  Therefore ordinary spin correlations are
strongly constrained.

In the exactly solvable zero-field model, the equal-time two-spin correlator
vanishes beyond nearest neighbours:
\begin{equation}
\label{eq:kitaev-spin-correlation-range}
    \langle \sigma_i^\alpha\sigma_j^\beta\rangle=0
    \qquad
    \text{unless } i,j \text{ are nearest neighbours on an }\alpha\text{-bond}
    \text{ and }\alpha=\beta .
\end{equation}
This short-ranged spin correlation is not a sign of a trivial phase.  It reflects
spin fractionalization: a physical spin flip fractionalizes into an itinerant
Majorana excitation and a pair of \(\ZZ_2\) fluxes.

By contrast, gauge-invariant Majorana bilinears and flux operators can encode
long-distance information about the Majorana band structure.  In the gapless
phase, the itinerant Majoranas have algebraic correlations, but local spin
operators couple to them only together with flux excitations.

\section{Comparison with the toric code and quantum double models}
\label{sec:kitaev-comparison-toric-code-quantum-double}

The following table summarizes the relation between the solvability mechanisms
of the previous two chapters and the honeycomb model.

\begin{center}
\begin{adjustbox}{max width=\textwidth}
\begin{tabular}{@{}L{0.24\textwidth}L{0.28\textwidth}L{0.32\textwidth}@{}}
\toprule
Model class & Exact-solvability mechanism & Excitations \\
\midrule
Toric code &
Commuting stabilizers \(A_s,B_p\) &
Static Abelian anyons \(1,e,m,\epsilon\) \\
\addlinespace
Quantum double \(D(G)\) &
Commuting gauge and flatness projectors \(A_v,B_p\) &
Electric, magnetic, and dyonic anyons labeled by \((C,\rho)\) \\
\addlinespace
Kitaev honeycomb model &
Static \(\ZZ_2\) gauge sectors plus free Majorana matter &
Fluxes, itinerant Majorana fermions, and in the chiral phase non-Abelian Ising anyons \\
\bottomrule
\end{tabular}
\end{adjustbox}
\end{center}

The toric code and quantum double models are fixed-point commuting-projector
Hamiltonians.  The Kitaev honeycomb model is instead a microscopic spin model
with dispersive fractionalized matter.  Its gapped Abelian regions flow to
toric-code topological order at long distances, while its time-reversal-broken
gapless region flows to Ising topological order.

\section{Summary}
\label{sec:kitaev-honeycomb-summary}

The Kitaev honeycomb model is exactly solvable because the spin Hamiltonian can
be rewritten as free Majorana fermions moving in a static \(\ZZ_2\) gauge field.
The essential steps are:
\begin{enumerate}[label=(\roman*)]
    \item enlarge each spin into four Majoranas,
    \begin{equation}
    \label{eq:kitaev-summary-majorana-rep}
        \sigma_i^\alpha=\ii b_i^\alpha c_i,
        \qquad
        D_i=b_i^xb_i^yb_i^zc_i=+1;
    \end{equation}
    \item define conserved link variables,
    \begin{equation}
    \label{eq:kitaev-summary-link-variable}
        u_{ij}=\ii b_i^\alpha b_j^\alpha,
        \qquad
        u_{ij}=\pm1;
    \end{equation}
    \item reduce the Hamiltonian in every gauge sector to
    \begin{equation}
    \label{eq:kitaev-summary-free-majorana}
        H[u]=\ii\sum_{\langle ij\rangle_\alpha}J_\alpha u_{ij}c_i c_j;
    \end{equation}
    \item identify gauge-invariant fluxes
    \begin{equation}
    \label{eq:kitaev-summary-flux}
        W_p=\prod_{\langle ij\rangle\in\partial p}u_{ij};
    \end{equation}
    \item diagonalize the vortex-free Majorana band with dispersion
    \begin{equation}
    \label{eq:kitaev-summary-dispersion}
        \epsilon(\bm k)=2\abs{J_z+J_x\ee^{\ii\bm k\cdot\bm a_1}
        +J_y\ee^{\ii\bm k\cdot\bm a_2}} .
    \end{equation}
\end{enumerate}
The bulk phase diagram follows from the triangle inequalities
\cref{eq:kitaev-triangle-inequalities}.  One obtains gapped Abelian phases when
one coupling dominates, a gapless Majorana spin liquid inside the triangle, and
a non-Abelian Ising-anyon phase when the gapless phase is gapped by a
small time-reversal-breaking perturbation.



\chapter{Rokhsar--Kivelson Points and Quantum Dimer Models}
\label{chap:rk-qdm}

The toric code and the quantum double models are exactly solvable because their
Hamiltonians are sums of mutually commuting projectors.  The Kitaev honeycomb
model is solvable because the spin Hamiltonian reduces to free Majorana
fermions in every static gauge sector.  Rokhsar--Kivelson (RK) quantum dimer
models provide a third mechanism in two dimensions: the full Hamiltonian need
not be diagonalizable, but a special point in parameter space has an exact
frustration-free ground state whose amplitudes are controlled by a classical
statistical model.

The central idea is
\boxedpar{%
classical constrained configurations
\(\mathcal{C}\)
\(\Longrightarrow\)
quantum Hilbert space \(\operatorname{span}\{\ket{C}:C\in\mathcal{C}\}\)
\(\Longrightarrow\)
local resonance Hamiltonian
\(\Longrightarrow\)
exact equal-amplitude or weighted ground state at the RK point.}
This chapter explains this mechanism for quantum dimer models.  The model is
not usually exactly solvable in the strong spectral sense.  Its exact content is
more specific: the ground-state wavefunction, equal-time ground-state
correlators, and topological-sector structure can be obtained from a classical
dimer ensemble.

\section{Classical dimer coverings and the quantum Hilbert space}
\label{sec:qdm-hilbert-space}

Let \(\Lambda\) be a two-dimensional lattice with vertex set \(V\) and edge set
\(E\).  A \emph{dimer covering}, or perfect matching, is a subset
\(D\subset E\) such that every vertex is incident on exactly one occupied edge.
Equivalently,
\begin{equation}
\label{eq:perfect-matching-constraint}
    \sum_{\ell\ni i} n_{\ell}=1,
    \qquad
    n_{\ell}\in\{0,1\},
    \qquad
    i\in V .
\end{equation}
The classical configuration space is therefore
\begin{equation}
\label{eq:dimer-configuration-space}
    \mathcal{C}_{\rm dim}
    =
    \{D\subset E:\text{\(D\) satisfies \cref{eq:perfect-matching-constraint}}\}.
\end{equation}
The quantum dimer Hilbert space is the vector space spanned by orthonormal basis
states labelled by dimer coverings:
\begin{equation}
\label{eq:qdm-hilbert-space}
    \calH_{\rm QDM}
    =
    \operatorname{span}\{\ket{D}:D\in\mathcal{C}_{\rm dim}\},
    \qquad
    \braket{D}{D'}=\delta_{D,D'} .
\end{equation}
This Hilbert space is already constrained.  It is not the full tensor product of
independent qubits on edges.  The local condition \cref{eq:perfect-matching-constraint}
plays a role analogous to a Gauss-law constraint: every site must be touched by
exactly one dimer.

\begin{remark}[Dimer constraint as an emergent gauge constraint]
On a bipartite lattice one can orient every edge from the \(A\) sublattice to
the \(B\) sublattice and define a background charge.  The dimer constraint can
then be rewritten as a lattice divergence equation.  This is why height fields,
Coulomb phases, and gauge-theory language naturally appear in quantum dimer
models.  The constraint is not optional; it defines the Hilbert space.
\end{remark}

The elementary local move in the square-lattice quantum dimer model is a
plaquette flip.  A plaquette \(p\) is called \emph{flippable} if it contains two
parallel dimers on opposite edges.  There are two such local states, denoted
abstractly by
\begin{equation}
\label{eq:local-dimer-plaquette-states}
    \ket{h_p},
    \qquad
    \ket{v_p},
\end{equation}
where \(h\) and \(v\) stand for the two parallel dimer orientations on the
plaquette.  The plaquette flip exchanges these two states and leaves all other
edges unchanged.

\section{The square-lattice quantum dimer Hamiltonian}
\label{sec:square-lattice-qdm-hamiltonian}

The standard square-lattice Rokhsar--Kivelson quantum dimer Hamiltonian is
\begin{equation}
\label{eq:qdm-hamiltonian-square}
    H_{\rm QDM}
    =
    \sum_p H_p,
\end{equation}
with local plaquette term
\begin{equation}
\label{eq:qdm-local-hamiltonian}
    H_p
    =
    -t
    \left(
        \ket{h_p}\bra{v_p}
        +
        \ket{v_p}\bra{h_p}
    \right)
    +
    v
    \left(
        \ket{h_p}\bra{h_p}
        +
        \ket{v_p}\bra{v_p}
    \right).
\end{equation}
The first term is kinetic: it resonates two dimer coverings that differ by one
plaquette flip.  The second term is potential: it counts flippable plaquettes.
The parameters \(t\) and \(v\) are usually taken real, and we will assume
\(t>0\).

In the two-dimensional subspace
\(\operatorname{span}\{\ket{h_p},\ket{v_p}\}\), the local operator is the matrix
\begin{equation}
\label{eq:qdm-local-matrix}
    H_p
    \widehat{=}
    \begin{pmatrix}
        v & -t \\
        -t & v
    \end{pmatrix}.
\end{equation}
Thus the symmetric and antisymmetric local combinations have energies
\begin{equation}
\label{eq:qdm-local-eigenvalues}
    E_{+}=v-t,
    \qquad
    E_{-}=v+t,
\end{equation}
where
\begin{equation}
\label{eq:qdm-local-symmetric-antisymmetric}
    \ket{+_p}=\frac{1}{\sqrt{2}}(\ket{h_p}+\ket{v_p}),
    \qquad
    \ket{-_p}=\frac{1}{\sqrt{2}}(\ket{h_p}-\ket{v_p}).
\end{equation}
The RK point is the special point
\begin{equation}
\label{eq:rk-point-condition}
    v=t .
\end{equation}
At this point, the local Hamiltonian becomes a positive semidefinite projector
up to a scalar factor:
\begin{equation}
\label{eq:qdm-local-rk-projector}
    H_p^{\rm RK}
    =
    t
    \left(
        \ket{h_p}-\ket{v_p}
    \right)
    \left(
        \bra{h_p}-\bra{v_p}
    \right).
\end{equation}
Therefore
\begin{equation}
\label{eq:qdm-rk-positive-semidefinite}
    H_{\rm QDM}^{\rm RK}\geq0 .
\end{equation}
This is the frustration-free structure hidden in the model.

\begin{remark}[Not a commuting-projector model]
Although each \(H_p^{\rm RK}\) is a local positive-semidefinite projector, the
different plaquette terms do not generally commute.  Thus the RK point is not a
toric-code-type commuting-projector point.  Exact solvability here means exact
ground-state solvability, not simultaneous diagonalization of the full spectrum.
\end{remark}

\section{The equal-amplitude RK ground state}
\label{sec:equal-amplitude-rk-ground-state}

The plaquette flip dynamics decomposes \(\mathcal{C}_{\rm dim}\) into connected
components.  On a torus these components are labelled, in part, by winding
numbers.  Let \(\mathcal{S}\subset\mathcal{C}_{\rm dim}\) be one connected
sector under local plaquette flips.  Define
\begin{equation}
\label{eq:rk-equal-amplitude-state}
    \ket{\Psi_{\mathcal{S}}}
    =
    \frac{1}{\sqrt{N_{\mathcal{S}}}}
    \sum_{D\in\mathcal{S}}\ket{D},
    \qquad
    N_{\mathcal{S}}=\abs{\mathcal{S}} .
\end{equation}
This is the equal-amplitude superposition of all dimer coverings in the sector
\(\mathcal{S}\).

\begin{proposition}[RK ground state]
At \(v=t\), the state \(\ket{\Psi_{\mathcal{S}}}\) is annihilated by every local
RK plaquette term:
\begin{equation}
\label{eq:rk-state-annihilated}
    H_p^{\rm RK}\ket{\Psi_{\mathcal{S}}}=0
    \qquad
    \text{for all }p .
\end{equation}
Consequently,
\begin{equation}
\label{eq:rk-state-ground-energy}
    H_{\rm QDM}^{\rm RK}\ket{\Psi_{\mathcal{S}}}=0 .
\end{equation}
\end{proposition}

\begin{proof}
Fix a plaquette \(p\).  The only basis states on which
\(H_p^{\rm RK}\) acts nontrivially are pairs of dimer coverings
\((D_h,D_v)\) that differ by flipping \(p\).  In the equal-amplitude state,
\(\ket{D_h}\) and \(\ket{D_v}\) occur with the same coefficient.  The local
operator \cref{eq:qdm-local-rk-projector} projects onto the antisymmetric
combination \(\ket{D_h}-\ket{D_v}\).  Therefore the contribution from each
flippable pair cancels.  Nonflippable configurations are annihilated because
neither \(\ket{h_p}\) nor \(\ket{v_p}\) is present locally.  Hence every
\(H_p^{\rm RK}\) annihilates the state.
\end{proof}

The proof also explains the logic of the RK construction.  The Hamiltonian is
built so that every allowed local move imposes an equality of amplitudes between
configurations connected by that move.  If the sector \(\mathcal{S}\) is
connected under the local moves, those local equality conditions force a unique
equal-amplitude state within that sector.

\begin{theorem}[Sector-wise RK ground space]
Assume a connected component \(\mathcal{S}\) of the dimer configuration graph is
connected by the plaquette flips appearing in the Hamiltonian.  At \(v=t>0\),
the zero-energy ground state in that component is the equal-amplitude state
\(\ket{\Psi_{\mathcal{S}}}\).  The full RK ground space is spanned by one such
state for each dynamically disconnected sector, together with possible isolated
frozen configurations that contain no flippable plaquettes.
\end{theorem}

\begin{remark}[Frozen states]
Some lattices or boundary conditions admit staggered dimer configurations with
no flippable plaquettes.  These states are also annihilated by
\(H_{\rm QDM}^{\rm RK}\), because every local term acts only on flippable
plaquettes.  They are exact zero-energy states, but they are not liquid-like RK
superpositions.  When discussing universal long-distance physics, one usually
focuses on the liquid sectors rather than these frozen sectors.
\end{remark}

\section{Topological sectors and winding numbers}
\label{sec:qdm-topological-sectors}

On a simply connected domain with fixed boundary conditions, the local plaquette
flips typically connect all dimer coverings compatible with the boundary.  On a
torus, however, local flips cannot change global winding numbers.

For a bipartite square lattice, choose a reference dimer configuration \(D_0\).
Given another dimer covering \(D\), the superposition of \(D\) and \(D_0\) forms
oriented transition loops.  The net number of transition-loop crossings through
two noncontractible cuts defines two winding numbers
\begin{equation}
\label{eq:qdm-winding-numbers}
    (W_x,W_y)\in\ZZ\times\ZZ .
\end{equation}
A local plaquette flip changes the dimer configuration only inside a
contractible region.  Therefore it cannot change \((W_x,W_y)\).

\begin{proposition}[Conservation of winding sectors]
For the square-lattice quantum dimer model on a torus,
\begin{equation}
\label{eq:qdm-winding-conservation}
    [H_{\rm QDM},W_x]=[H_{\rm QDM},W_y]=0,
\end{equation}
where \(W_x\) and \(W_y\) denote the winding-sector labels regarded as operators
that are diagonal in the dimer basis.
\end{proposition}

\begin{proof}
Each Hamiltonian term flips dimers around a single plaquette.  Such a move adds
or removes an elementary contractible loop in the transition-graph
representation.  A contractible modification has zero intersection number with
any noncontractible cut.  Hence the winding numbers are unchanged by every
local term.
\end{proof}

The existence of winding sectors is similar in spirit to the topological sectors
of the toric code, but the physical meaning is different.  In the toric code the
fourfold degeneracy on a torus is robust throughout a topologically ordered
phase.  In the square-lattice RK dimer model, the number of winding sectors can
grow with system size, and the RK point is critical rather than a stable gapped
\(\ZZ_2\) topological phase.

\section{Relation to classical dimer partition functions}
\label{sec:rk-classical-dimer-correspondence}

The RK wavefunction converts quantum equal-time expectation values into
classical statistical averages.  Let \(\widehat{O}\) be an operator diagonal in
the dimer basis:
\begin{equation}
\label{eq:diagonal-dimer-operator}
    \widehat{O}\ket{D}=O(D)\ket{D} .
\end{equation}
Then in the equal-amplitude RK state,
\begin{equation}
\label{eq:rk-diagonal-observable-classical-average}
    \bra{\Psi_{\mathcal{S}}}\widehat{O}\ket{\Psi_{\mathcal{S}}}
    =
    \frac{1}{N_{\mathcal{S}}}
    \sum_{D\in\mathcal{S}}O(D).
\end{equation}
Thus diagonal ground-state correlations are exactly the correlations of a
classical dimer model with uniform Boltzmann weight in the same sector.

More generally, one may introduce positive classical weights \(w(D)>0\) and the
wavefunction
\begin{equation}
\label{eq:weighted-rk-wavefunction}
    \ket{\Psi_w}
    =
    \frac{1}{\sqrt{Z_w}}
    \sum_{D\in\mathcal{C}_{\rm dim}}
    \sqrt{w(D)}\ket{D},
    \qquad
    Z_w=
    \sum_{D\in\mathcal{C}_{\rm dim}}w(D).
\end{equation}
Then diagonal observables satisfy
\begin{equation}
\label{eq:weighted-rk-classical-average}
    \bra{\Psi_w}\widehat{O}\ket{\Psi_w}
    =
    \frac{1}{Z_w}
    \sum_{D\in\mathcal{C}_{\rm dim}}w(D)O(D).
\end{equation}
This is the broader classical-to-quantum RK construction.

\begin{principle}[RK classical--quantum correspondence]
If a quantum ground state has amplitudes equal to the square roots of classical
Boltzmann weights, then the norm of the quantum state is the classical partition
function, and all equal-time diagonal quantum correlators are classical thermal
correlators.
\end{principle}

For the uniform dimer RK state,
\begin{equation}
\label{eq:rk-norm-classical-partition-function}
    \braket{\Psi}{\Psi}
    =
    \sum_{D\in\mathcal{C}_{\rm dim}}1
    =
    Z_{\rm dimer} .
\end{equation}
This relation is the reason why classical dimer techniques, especially
Kasteleyn orientations and Pfaffians, enter the quantum problem.

\section{Kasteleyn orientation and Pfaffian solvability}
\label{sec:kasteleyn-pfaffian-dimers}

Classical close-packed dimers on a planar graph are exactly solvable by the
Kasteleyn--Pfaffian method.  Assign an antisymmetric matrix \(K\) to the graph,
with entries
\begin{equation}
\label{eq:kasteleyn-matrix-definition}
    K_{ij}
    =
    \begin{cases}
        \phantom{-}\kappa_{ij}, & \text{if the oriented edge is }i\to j,\\
        -\kappa_{ij}, & \text{if the oriented edge is }j\to i,\\
        0, & \text{if }i\text{ and }j\text{ are not adjacent},
    \end{cases}
\end{equation}
where \(\kappa_{ij}>0\) is the dimer fugacity on edge \((ij)\).  A Kasteleyn
orientation is an orientation for which each face has an odd number of clockwise
oriented edges, with a standard modification for boundary and topology.

For a planar graph with a Kasteleyn orientation, the classical dimer partition
function is
\begin{equation}
\label{eq:dimer-partition-pfaffian}
    Z_{\rm dimer}=\abs{\Pf K} .
\end{equation}
Since
\begin{equation}
\label{eq:pfaffian-determinant-relation}
    (\Pf K)^2=\det K,
\end{equation}
one can compute the free energy from a free-fermion determinant.  Moreover,
dimer occupation correlations are obtained from entries of \(K^{-1}\).  For
example, for two edges \(e=(ij)\) and \(e'=(kl)\), connected dimer correlations
can be expressed by finite determinants built from matrix elements of
\(K^{-1}\).

\begin{remark}[Why Pfaffians matter for the quantum RK point]
The RK Hamiltonian itself is not a free-fermion Hamiltonian.  However, its
exact ground-state norm and diagonal equal-time correlators reduce to a
classical dimer problem.  If that classical dimer problem is Pfaffian-solvable,
then those quantum ground-state quantities are exactly computable by Pfaffian
methods.  This is a ground-state correspondence, not a full spectral
fermionization of the quantum Hamiltonian.
\end{remark}

On a torus, one needs four Pfaffians corresponding to different boundary signs
around the two noncontractible cycles.  Schematically,
\begin{equation}
\label{eq:torus-dimer-four-pfaffians}
    Z_{\rm dimer}^{\rm torus}
    =
    \frac{1}{2}
    \left|
        -\Pf K_{00}
        +\Pf K_{10}
        +\Pf K_{01}
        +\Pf K_{11}
    \right|,
\end{equation}
where the subscripts indicate twists along the two cycles.  The precise sign
convention depends on the chosen Kasteleyn orientation, but the need for four
spin structures is universal.

\section{Height representation and the square-lattice RK point}
\label{sec:qdm-height-representation}

On bipartite planar lattices, close-packed dimer configurations can be mapped to
integer-valued height configurations on the dual lattice.  The construction is
local: crossing an edge changes the height by an amount that depends on whether
the edge is occupied by a dimer and on the sublattice convention.  The exact
microscopic increments are conventional, but the long-distance structure is not.
The coarse-grained height field \(h(\bm r)\) has an effective classical action
of Gaussian form
\begin{equation}
\label{eq:classical-dimer-height-action}
    S_{\rm cl}[h]
    =
    \frac{K}{2}
    \int \dd^2 r\,
    (\nabla h)^2
    +\text{irrelevant or locking terms} .
\end{equation}
For the square-lattice uniform dimer model, this height theory is critical.
Consequently, the square-lattice RK ground state has algebraic equal-time dimer
correlations.

At the RK point, the quantum dynamics adds a kinetic term to the height field.
The long-wavelength quantum theory is the quantum Lifshitz theory
\begin{equation}
\label{eq:quantum-lifshitz-action}
    S_{\rm QL}[h]
    =
    \frac{1}{2}
    \int \dd\tau\,\dd^2 r
    \left[
        (\partial_{\tau}h)^2
        +
        \kappa^2(\nabla^2 h)^2
    \right] .
\end{equation}
The dispersion relation is therefore
\begin{equation}
\label{eq:quantum-lifshitz-dispersion}
    \omega(\bm q)\sim \kappa \abs{\bm q}^2,
\end{equation}
so the dynamic critical exponent is
\begin{equation}
\label{eq:rk-dynamical-exponent}
    z=2 .
\end{equation}

\begin{remark}[Why this is not a relativistic critical point]
In most quantum critical points discussed earlier in one-dimensional free
fermion chains, the low-energy theory has \(z=1\) and spacetime scaling is
relativistic.  The square-lattice RK point is different: its effective action
contains \((\nabla^2 h)^2\), not \((\nabla h)^2\), in the quantum potential
energy.  The resulting \(\omega\sim q^2\) dispersion is a hallmark of RK-type
criticality.
\end{remark}

The square-lattice phase diagram near the RK point is often summarized as
follows.  For \(v/t>1\), flippable plaquettes cost energy and staggered frozen
states dominate.  For sufficiently smaller \(v/t\), crystalline valence-bond
solid phases appear.  The RK point \(v=t\) is a fine-tuned critical point
separating the staggered regime from resonating dimer physics.

\section{Triangular lattice and gapped \texorpdfstring{\(\ZZ_2\)}{Z2} dimer liquid}
\label{sec:triangular-rk-liquid}

The physics of the RK point depends strongly on the lattice.  On the triangular
lattice, the elementary move again flips two parallel dimers around a rhombus,
but the lattice is nonbipartite.  The absence of a bipartite height field changes
the long-distance behavior.  At the triangular-lattice RK point, the equal-weight
dimer liquid is gapped and realizes \(\ZZ_2\) topological order rather than a
critical height phase.

The distinction may be summarized as
\begin{center}
\begin{tabular}{L{0.28\linewidth}L{0.30\linewidth}L{0.30\linewidth}}
\toprule
Lattice type & Classical dimer ensemble & RK quantum phase \\
\midrule
Bipartite square lattice & Height/Coulomb phase with algebraic correlations & Critical RK point with \(z=2\) \\
Nonbipartite triangular lattice & Short-ranged dimer correlations & Gapped \(\ZZ_2\) topological dimer liquid \\
\bottomrule
\end{tabular}
\end{center}

On a torus, a gapped \(\ZZ_2\) dimer liquid has four topological sectors,
analogous to the four ground states of the toric code.  A useful way to label
them is by the parity of the number of dimers crossing two noncontractible cuts:
\begin{equation}
\label{eq:z2-dimer-winding-parities}
    (P_x,P_y)\in\ZZ_2\times\ZZ_2 .
\end{equation}
Local plaquette flips preserve these parities.  In the thermodynamic limit, the
four corresponding RK ground states become locally indistinguishable if the
phase is gapped and topologically ordered.

\begin{remark}[Relation to the toric code]
The triangular-lattice RK dimer liquid and the toric code are not identical
microscopic Hamiltonians.  However, they lie in the same broad topological
universality class when the dimer liquid is gapped and has \(\ZZ_2\) topological
order.  The toric code is the fixed-point commuting-projector representative;
the quantum dimer model is a constrained Hilbert-space realization with
resonance dynamics.
\end{remark}

\section{General RK construction from stochastic dynamics}
\label{sec:general-rk-stochastic-construction}

The dimer RK point is a special case of a general construction.  Let
\(\mathcal{C}\) be a classical configuration space with positive weights
\(w(C)>0\).  Suppose there is a set of local reversible moves
\(C\leftrightarrow C'\).  Define local operators of the form
\begin{equation}
\label{eq:general-rk-local-projector}
    H_{C,C'}
    =
    \gamma_{C,C'}
    \left(
        \sqrt{\frac{w(C')}{w(C)}}\ket{C}
        -
        \sqrt{\frac{w(C)}{w(C')}}\ket{C'}
    \right)
    \left(
        \sqrt{\frac{w(C')}{w(C)}}\bra{C}
        -
        \sqrt{\frac{w(C)}{w(C')}}\bra{C'}
    \right),
\end{equation}
with \(\gamma_{C,C'}>0\).  Then the state \(\ket{\Psi_w}\) in
\cref{eq:weighted-rk-wavefunction} is annihilated by every such term.  Hence
\begin{equation}
\label{eq:general-rk-hamiltonian}
    H_{\rm RK}[w]
    =
    \sum_{(C,C')}H_{C,C'}
\end{equation}
is frustration-free and has \(\ket{\Psi_w}\) as an exact ground state in every
connected component of the move graph.

This construction is closely related to classical Markov dynamics satisfying
detailed balance.  A Markov generator \(M\) with stationary distribution
\begin{equation}
\label{eq:rk-stationary-distribution}
    \pi(C)=\frac{w(C)}{Z_w}
\end{equation}
can often be transformed by a similarity transformation into a Hermitian RK
Hamiltonian:
\begin{equation}
\label{eq:markov-to-rk-similarity}
    H_{\rm RK}
    =
    -\Pi^{-1/2}M\Pi^{1/2},
    \qquad
    \Pi_{C,C'}=\pi(C)\delta_{C,C'} .
\end{equation}
Consequently, relaxation modes of the classical stochastic process map to
excitation modes of the quantum RK Hamiltonian.

\begin{principle}[Stochastic-matrix form]
RK Hamiltonians often have a stochastic-matrix form: after a diagonal similarity
transformation, the quantum Hamiltonian becomes a classical Markov generator.
The exact ground state is the square root of the stationary distribution, while
low-energy quantum modes mirror slow classical relaxation modes.
\end{principle}

\section{Excitations: resonons, visons, and monomers}
\label{sec:qdm-excitations}

Although the RK point gives an exact ground state, its excitations are generally
more difficult than those of a free-fermion or commuting-projector model.
Nevertheless, several important excitation types can be identified.

\paragraph{Resonons.}
On the square-lattice RK point, low-energy collective modes of the height field
are sometimes called resonons.  Their energy vanishes as
\begin{equation}
\label{eq:resonon-dispersion}
    E(\bm q)\sim \abs{\bm q}^2,
\end{equation}
consistent with the quantum Lifshitz exponent \(z=2\).

\paragraph{Visons.}
In a gapped \(\ZZ_2\) dimer liquid, a vison is a \(\ZZ_2\) vortex excitation.
Operationally, a vison can be detected by a string operator that measures the
parity of dimers crossing a path on the dual lattice.  Its mutual statistics
with a monomer excitation is analogous to the mutual statistics of \(m\) and
\(e\) particles in the toric code.

\paragraph{Monomers.}
A monomer is a violation of the close-packing constraint: a site is left
unmatched or has the wrong dimer occupancy.  In a strict quantum dimer Hilbert
space, monomers are excluded.  They can be introduced by enlarging the Hilbert
space.  Their confinement or deconfinement diagnoses whether the dimer phase is
crystalline, critical, or topologically ordered.

\section{Comparison with earlier exact models}
\label{sec:qdm-comparison-with-earlier-models}

Quantum dimer models clarify an important distinction in these notes.  Exact
solvability can refer to the full spectrum, to a sector-wise reduction, or only
to the ground-state wavefunction and its equal-time correlations.  The RK point
belongs to the last category.

\begin{center}
\small
\begin{tabular}{L{0.23\linewidth}L{0.33\linewidth}L{0.31\linewidth}}
\toprule
Model class & Solvable structure & What is exactly controlled \\
\midrule
TFIM/XY/Kitaev chain & Jordan--Wigner plus BdG diagonalization & Full spectrum and Gaussian correlators \\
XXZ/XYZ chains & Yang--Baxter algebra and Bethe ansatz & Spectrum through Bethe equations \\
Toric code & Commuting stabilizers/projectors & Full spectrum and anyon algebra at the fixed point \\
Kitaev honeycomb & Static \(\ZZ_2\) gauge fields plus free Majoranas & Free-fermion spectrum in each flux sector \\
RK quantum dimer model & Frustration-free stochastic-matrix form & Exact ground state and classical equal-time correlators \\
\bottomrule
\end{tabular}
\end{center}

Thus the RK model is exactly solvable in a weaker but conceptually important
sense.  It shows that a nontrivial many-body wavefunction can be known exactly
even when the full many-body spectrum is not.

\section{Summary}
\label{sec:qdm-summary}

The essential points of this chapter are:
\begin{enumerate}[label=(\roman*)]
    \item A quantum dimer model is defined on the constrained Hilbert space of
    perfect matchings.
    \item The square-lattice QDM has kinetic plaquette flips and a potential
    term that counts flippable plaquettes.
    \item At the RK point \(v=t\), the Hamiltonian becomes a sum of positive
    semidefinite local terms.
    \item In every connected sector, the equal-amplitude superposition of dimer
    coverings is an exact zero-energy ground state.
    \item The norm and diagonal equal-time correlators of the RK state are
    those of a classical dimer model.
    \item Classical dimer Pfaffian methods therefore compute quantum
    ground-state correlations at the RK point.
    \item Bipartite square-lattice RK physics is critical and described by a
    height/quantum-Lifshitz theory with \(z=2\).
    \item Nonbipartite triangular-lattice RK physics can realize a gapped
    \(\ZZ_2\) topological dimer liquid.
\end{enumerate}

The next chapter moves from RK ground-state solvability to more general exact
ground-state constructions, including valence-bond-solid states, stabilizer
cluster states, and parent Hamiltonians of tensor-network states.


\chapter[Exact Ground-State Models]{Exact Ground-State Models: AKLT, Cluster, and Parent Hamiltonians}
\label{chap:exact-ground-state-models}

The preceding chapters presented several different notions of exact solvability
in two-dimensional quantum many-body systems.  The toric code and quantum double
models are commuting-projector models, so their full spectra reduce to local
constraint eigenvalues.  The Kitaev honeycomb model is not a commuting-projector
model, but it reduces to free Majorana fermions in each static gauge sector.
Rokhsar--Kivelson quantum dimer models provide exact ground states at special
frustration-free points, while the full spectrum is generally not solved.

This chapter develops a broader exact-ground-state viewpoint.  The central
objects are Hamiltonians whose ground states are known exactly because the local
terms are projectors annihilating a prescribed many-body state.  The paradigms
are
\begin{enumerate}[label=(\roman*)]
    \item the AKLT chain, whose valence-bond-solid ground state is an exact
    matrix-product state;
    \item the cluster model, whose ground state is a stabilizer state and a
    simple example of symmetry-protected topological order;
    \item the general parent-Hamiltonian construction for MPS and PEPS.
\end{enumerate}
The common mechanism is
\boxedpar{%
exactly known wavefunction
\(\Longrightarrow\)
local reduced-density-matrix support
\(\Longrightarrow\)
positive local projector
\(\Longrightarrow\)
frustration-free Hamiltonian.}
This is a different sense of solvability from Bethe ansatz or free-fermion
diagonalization.  It is exact and nonperturbative, but it usually controls the
ground-state subspace rather than every excited level.

\section{Frustration-free exact ground-state solvability}
\label{sec:frustration-free-ground-state-solvability}

Let a local Hamiltonian be written as
\begin{equation}
\label{eq:ff-general-hamiltonian}
    H=\sum_X h_X,
    \qquad
    h_X\geq 0,
\end{equation}
where each term acts on a finite region \(X\) of the lattice.  A state
\(\ket{\Psi}\) is a frustration-free ground state if
\begin{equation}
\label{eq:ff-ground-condition}
    h_X\ket{\Psi}=0
    \qquad
    \text{for every local term }X.
\end{equation}
Then
\begin{equation}
\label{eq:ff-zero-energy-ground-state}
    H\ket{\Psi}=0,
    \qquad
    H\geq0,
\end{equation}
so \(\ket{\Psi}\) is a ground state with energy zero.  If the kernel of
\(H\) is known exactly, the ground-state problem is solved.

\begin{definition}[Frustration-free Hamiltonian]
A positive local Hamiltonian \(H=\sum_X h_X\), with \(h_X\geq0\), is called
frustration-free if
\begin{equation}
\label{eq:ff-kernel-intersection}
    \ker H
    =
    \bigcap_X \ker h_X .
\end{equation}
Equivalently, a global ground state minimizes every local term separately.
\end{definition}

This definition should be compared with commuting-projector solvability.  If all
local projectors commute, then their common eigenspaces give the full spectrum.
If the terms are frustration-free but noncommuting, the ground space can still be
known exactly, but the excited spectrum is generally nontrivial.

\begin{remark}[Three strengths of solvability]
It is useful to distinguish three levels.
\begin{enumerate}[label=(\alph*)]
    \item \emph{Full-spectrum commuting-projector solvability}: all local terms
    commute, as in the toric code.
    \item \emph{Free-particle or Bethe-ansatz spectral solvability}: the full
    spectrum is generated by quasiparticles or rapidities, as in the XY and XXZ
    chains.
    \item \emph{Exact ground-state solvability}: the ground-state manifold is
    exactly known, but the full spectrum is not generally diagonalized.  AKLT,
    cluster, RK, and generic MPS/PEPS parent Hamiltonians mostly belong to this
    class.
\end{enumerate}
\end{remark}

The exact-ground-state construction is especially powerful because it produces
nontrivial phases with simple wavefunctions.  In one dimension, it naturally
leads to matrix-product states.  In two dimensions, it leads to projected
entangled-pair states, string-net states, dimer liquid states, and other tensor-
network fixed points.

\section{The AKLT chain}
\label{sec:aklt-chain}

The Affleck--Kennedy--Lieb--Tasaki (AKLT) chain is the canonical exactly soluble
spin-1 model.  It realizes the Haldane phase by an explicit valence-bond-solid
(VBS) wavefunction.  Its solvability is not based on free fermions or Bethe
rapidities.  Instead, it follows from a projector decomposition of two adjacent
spin-1 degrees of freedom.

\subsection{Spin-1 degrees of freedom and two-site projectors}
\label{subsec:aklt-two-site-projectors}

Let \(\bm S_j=(S_j^x,S_j^y,S_j^z)\) be spin-1 operators satisfying
\begin{equation}
\label{eq:spin-one-algebra}
    [S_j^\alpha,S_\ell^\beta]
    =
    \ii\delta_{j\ell}\epsilon^{\alpha\beta\gamma}S_j^\gamma,
    \qquad
    \bm S_j^2=2\id .
\end{equation}
For two neighboring spin-1 sites, the tensor product decomposes as
\begin{equation}
\label{eq:spin-one-two-site-decomposition}
    1\otimes1 = 0\oplus1\oplus2 .
\end{equation}
Let \(P_{j,j+1}^{(S)}\) denote the projector onto total spin \(S\) in the pair
\((j,j+1)\).  Since
\begin{equation}
\label{eq:spin-dot-eigenvalues-spin-one}
    \bm S_j\cdot\bm S_{j+1}
    =
    \frac12
    \left[
        (\bm S_j+\bm S_{j+1})^2
        -\bm S_j^2
        -\bm S_{j+1}^2
    \right],
\end{equation}
its eigenvalues in the \(S=0,1,2\) sectors are
\begin{equation}
\label{eq:spin-one-dot-eigenvalues}
    x_0=-2,
    \qquad
    x_1=-1,
    \qquad
    x_2=1 .
\end{equation}
Therefore the projector onto the total-spin-2 sector is the polynomial
\begin{equation}
\label{eq:spin-two-projector-polynomial}
    P_{j,j+1}^{(2)}
    =
    \frac{1}{6}
    \left(
        \bm S_j\cdot\bm S_{j+1}+2
    \right)
    \left(
        \bm S_j\cdot\bm S_{j+1}+1
    \right).
\end{equation}
Equivalently,
\begin{equation}
\label{eq:spin-two-projector-expanded}
    P_{j,j+1}^{(2)}
    =
    \frac16(\bm S_j\cdot\bm S_{j+1})^2
    +\frac12\bm S_j\cdot\bm S_{j+1}
    +\frac13 .
\end{equation}

The AKLT Hamiltonian on a periodic chain is
\begin{equation}
\label{eq:aklt-hamiltonian-projector}
    H_{\rm AKLT}
    =
    \sum_{j=1}^{L} P_{j,j+1}^{(2)},
    \qquad
    L+1\equiv1 .
\end{equation}
Up to an overall factor and additive constant, this is the more familiar
bilinear-biquadratic form
\begin{equation}
\label{eq:aklt-hamiltonian-bilinear-biquadratic}
    H_{\rm AKLT}'
    =
    \sum_{j=1}^{L}
    \left[
        \bm S_j\cdot\bm S_{j+1}
        +\frac13(\bm S_j\cdot\bm S_{j+1})^2
        +\frac23
    \right]
    =
    2H_{\rm AKLT} .
\end{equation}
The projector form \cref{eq:aklt-hamiltonian-projector} is the most transparent
for exact solvability.

\subsection{Valence-bond-solid construction}
\label{subsec:aklt-vbs-construction}

The VBS construction replaces every physical spin-1 by two virtual spin-
\(1/2\) degrees of freedom.  At site \(j\), the two virtual spins are projected
onto their symmetric triplet subspace:
\begin{equation}
\label{eq:spin-one-as-symmetric-two-spin-half}
    \calH^{S=1}_j
    =
    \operatorname{Sym}
    \left(
        \CC^2_{j,L}\otimes\CC^2_{j,R}
    \right).
\end{equation}
Neighboring virtual spins form singlets
\begin{equation}
\label{eq:virtual-singlet}
    \ket{\omega}_{j,R;j+1,L}
    =
    \frac{1}{\sqrt2}
    \left(
        \ket{\uparrow}_{j,R}\ket{\downarrow}_{j+1,L}
        -
        \ket{\downarrow}_{j,R}\ket{\uparrow}_{j+1,L}
    \right).
\end{equation}
The AKLT state is obtained by placing singlets on all virtual bonds and applying
the triplet projection at every site:
\begin{equation}
\label{eq:aklt-vbs-state}
    \ket{\Psi_{\rm AKLT}}
    =
    \prod_j P^{\rm sym}_j
    \prod_j
    \ket{\omega}_{j,R;j+1,L} .
\end{equation}
For an open chain, two unpaired virtual spin-\(1/2\) degrees of freedom remain
at the ends.  They produce a four-dimensional edge multiplet before boundary
terms are added.  This is the simplest microscopic origin of the spin-
\(1/2\) edge states in the Haldane phase.

\begin{proposition}[AKLT ground state]
The VBS state \(\ket{\Psi_{\rm AKLT}}\) is annihilated by every local projector
\(P_{j,j+1}^{(2)}\).  Hence it is a zero-energy ground state of
\(H_{\rm AKLT}\).
\end{proposition}

\begin{proof}
Consider two neighboring physical spin-1 sites.  In the virtual construction,
the right virtual spin of site \(j\) and the left virtual spin of site \(j+1\)
form a singlet.  Thus the total spin carried by the four virtual spin-
\(1/2\) variables on the two sites cannot reach total spin 2.  After projection
to the physical triplet subspace at each site, the two physical spin-1's
therefore have no component in the total-spin-2 sector.  Consequently
\(P_{j,j+1}^{(2)}\ket{\Psi_{\rm AKLT}}=0\) for every bond.  Since the Hamiltonian
is a sum of positive projectors, the state is a ground state.
\end{proof}

\subsection{AKLT state as a matrix-product state}
\label{subsec:aklt-mps}

The VBS state can be written as an MPS with bond dimension \(D=2\).  In the
basis \(\ket{+1},\ket{0},\ket{-1}\) of \(S^z\), one convenient convention is
\begin{equation}
\label{eq:aklt-mps-matrices}
    A^{+1}=\sqrt{\frac23}\,\sigma^+,
    \qquad
    A^{0}=-\frac{1}{\sqrt3}\,\sigma^z,
    \qquad
    A^{-1}=-\sqrt{\frac23}\,\sigma^- .
\end{equation}
The periodic-chain wavefunction is
\begin{equation}
\label{eq:aklt-mps-periodic}
    \ket{\Psi_{\rm AKLT}}
    =
    \sum_{m_1,\ldots,m_L}
    \Tr\left(A^{m_1}A^{m_2}\cdots A^{m_L}\right)
    \ket{m_1m_2\cdots m_L} .
\end{equation}
For an open chain, the trace is replaced by boundary spinors in the virtual
space.  The different boundary spinors describe the edge spin-
\(1/2\) degrees of freedom.

The MPS transfer channel is
\begin{equation}
\label{eq:aklt-transfer-channel}
    \mathbb{E}(X)
    =
    \sum_{m=-1}^{1} A^m X (A^m)^\dagger .
\end{equation}
With the convention \cref{eq:aklt-mps-matrices}, the identity is the leading
fixed point and the three Pauli matrices form subleading eigenoperators with
eigenvalue \(-1/3\).  Thus connected spin correlations decay as
\begin{equation}
\label{eq:aklt-correlation-length}
    \langle S_i^\alpha S_j^\alpha\rangle_c
    \propto
    \left(-\frac13\right)^{\abs{i-j}},
    \qquad
    \xi_{\rm AKLT}=\frac{1}{\log3} .
\end{equation}
This is an exact statement obtained directly from the finite-dimensional MPS
transfer matrix.

\begin{remark}[Hidden order]
Although ordinary spin correlations are short-ranged, the AKLT state has
nonzero string order.  A standard string correlator is
\begin{equation}
\label{eq:aklt-string-order}
    \mathcal{O}_{\rm string}^z
    =
    -\lim_{\abs{i-j}\to\infty}
    \left\langle
        S_i^z
        \exp\left(\ii\pi\sum_{k=i+1}^{j-1}S_k^z\right)
        S_j^z
    \right\rangle .
\end{equation}
For the AKLT state this quantity is finite.  It detects the hidden antiferromag-
netic order associated with the Haldane symmetry-protected topological phase.
\end{remark}

\section{The cluster model}
\label{sec:cluster-model}

The cluster model is the simplest stabilizer example of an exact ground-state
model with symmetry-protected topological structure.  Unlike the AKLT chain,
its local terms commute, so it is also full-spectrum soluble as a stabilizer
Hamiltonian.  It is nevertheless conceptually useful in this chapter because its
exact ground state is a graph state and therefore a tensor-network fixed point.

\subsection{One-dimensional cluster Hamiltonian}
\label{subsec:one-dimensional-cluster-hamiltonian}

Consider a chain of qubits with Pauli operators \(X_j,Y_j,Z_j\).  On a periodic
chain, define the stabilizers
\begin{equation}
\label{eq:cluster-stabilizers-periodic}
    K_j
    =
    Z_{j-1}X_jZ_{j+1},
    \qquad
    j=1,\ldots,L,
\end{equation}
with site labels understood modulo \(L\).  The one-dimensional cluster
Hamiltonian is
\begin{equation}
\label{eq:cluster-hamiltonian-periodic}
    H_{\rm cl}
    =
    -\sum_{j=1}^{L}K_j .
\end{equation}
The stabilizers commute:
\begin{equation}
\label{eq:cluster-stabilizer-commutation}
    [K_i,K_j]=0,
    \qquad
    K_j^2=\id .
\end{equation}
Indeed, if \(K_i\) and \(K_j\) overlap nontrivially, they either share only
commuting Pauli matrices or have two anticommuting overlaps, whose signs cancel.
Therefore the full spectrum is labelled by eigenvalues \(k_j=\pm1\).

The cluster ground state \(\ket{C}\) is the unique common \(+1\) eigenstate on a
periodic chain:
\begin{equation}
\label{eq:cluster-ground-state-stabilizer}
    K_j\ket{C}=\ket{C},
    \qquad
    j=1,\ldots,L .
\end{equation}
Equivalently, it is obtained by applying controlled-\(Z\) gates to the product
state \(\ket{+}^{\otimes L}\):
\begin{equation}
\label{eq:cluster-state-cz-construction}
    \ket{C}
    =
    \prod_{\langle j,j+1\rangle} CZ_{j,j+1}\,
    \ket{+}^{\otimes L},
    \qquad
    X_j\ket{+}_j=\ket{+}_j .
\end{equation}
This formula follows because conjugation by \(CZ_{j,j+1}\) maps
\begin{equation}
\label{eq:cz-pauli-conjugation}
    X_j
    \longmapsto
    Z_{j-1}X_jZ_{j+1}
\end{equation}
when all nearest-neighbor controlled-\(Z\) gates are applied.

\begin{remark}[Open-chain edge degeneracy]
For an open chain, one may define bulk stabilizers
\(K_j=Z_{j-1}X_jZ_{j+1}\) for \(2\leq j\leq L-1\).  Then the two boundary qubits
are not fully fixed by the bulk terms, and protected edge degrees of freedom
appear.  This parallels the AKLT edge spin-
\(1/2\)'s, although the cluster model is built from qubits and commuting
stabilizers rather than spin-1 projectors.
\end{remark}

\subsection{Symmetry-protected topological interpretation}
\label{subsec:cluster-spt}

The one-dimensional cluster state is a fixed point of a nontrivial
\(\ZZ_2\times\ZZ_2\) symmetry-protected topological phase.  A convenient choice
of symmetry generators is
\begin{equation}
\label{eq:cluster-z2xz2-symmetry}
    G_{\rm even}=\prod_{j\,\text{even}}X_j,
    \qquad
    G_{\rm odd}=\prod_{j\,\text{odd}}X_j .
\end{equation}
On a closed chain, these commute with the Hamiltonian.  On an open chain, the
symmetry acts projectively on the edge degrees of freedom.  This projective edge
action is the one-dimensional SPT invariant.

The cluster state also has a string order parameter.  For example, products of
stabilizers imply that
\begin{equation}
\label{eq:cluster-string-order}
    Z_i
    \left(\prod_{k=i+1}^{j-1}X_k\right)
    Z_j
\end{equation}
has long-range expectation value in the cluster ground state for suitable parity
choices of \(i\) and \(j\).  The precise endpoint operators depend on whether
one chooses the even or odd symmetry sector, but the qualitative structure is
robust: nonlocal string order replaces ordinary local order.

\subsection{Graph-state generalization}
\label{subsec:graph-state-generalization}

The cluster construction extends immediately to any graph \(\Gamma=(V,E)\).
Place a qubit on every vertex and define
\begin{equation}
\label{eq:graph-state-stabilizers}
    K_v
    =
    X_v\prod_{u:(u,v)\in E}Z_u .
\end{equation}
The graph-state Hamiltonian is
\begin{equation}
\label{eq:graph-state-hamiltonian}
    H_\Gamma=-\sum_{v\in V}K_v .
\end{equation}
All \(K_v\)'s commute, and the graph state is the common \(+1\) eigenstate.
For a two-dimensional square-lattice graph, this gives the two-dimensional
cluster state, an important stabilizer resource state for measurement-based
quantum computation.

\begin{remark}[Cluster model versus toric code]
Both the cluster model and the toric code are stabilizer Hamiltonians, but their
orders are different.  The cluster state is short-range entangled and becomes a
trivial product state after a finite-depth circuit if the protecting symmetry is
ignored.  The toric code is long-range entangled and has intrinsic topological
order.  Thus stabilizer solvability by itself does not determine the phase;
one must also examine entanglement, symmetry, and nonlocal operators.
\end{remark}

\section{Parent Hamiltonians of matrix-product states}
\label{sec:mps-parent-hamiltonians}

The AKLT chain is not an isolated miracle.  It is an example of a general MPS
parent-Hamiltonian construction.  The input is an exact tensor-network state;
the output is a local frustration-free Hamiltonian for which that tensor network
is an exact ground state.

\subsection{Matrix-product states}
\label{subsec:mps-definition-parent}

Let each site have physical dimension \(d\), with basis
\(\ket{s}\), \(s=1,\ldots,d\).  A translation-invariant MPS with bond dimension
\(D\) on a periodic chain is
\begin{equation}
\label{eq:mps-periodic-definition}
    \ket{\Psi_A}
    =
    \sum_{s_1,\ldots,s_L=1}^{d}
    \Tr\left(A^{s_1}A^{s_2}\cdots A^{s_L}\right)
    \ket{s_1s_2\cdots s_L},
\end{equation}
where each \(A^s\) is a \(D\times D\) matrix.  The transfer channel is
\begin{equation}
\label{eq:mps-transfer-channel-general}
    \mathbb{E}(X)
    =
    \sum_{s=1}^{d}A^sX(A^s)^\dagger .
\end{equation}
The state is called injective, possibly after blocking a finite number of sites,
if the blocked tensors generate the full matrix algebra on the virtual space.
More explicitly, for some finite length \(\ell\), the map
\begin{equation}
\label{eq:mps-injectivity-map}
    \Gamma_\ell:
    \operatorname{Mat}_{D\times D}(\CC)
    \to
    (\CC^d)^{\otimes \ell},
    \qquad
    X
    \mapsto
    \sum_{s_1,\ldots,s_\ell}
    \Tr\left(XA^{s_1}\cdots A^{s_\ell}\right)
    \ket{s_1\cdots s_\ell}
\end{equation}
is injective.

Injectivity is the tensor-network analogue of having a unique locally generated
ground state.  Noninjective MPS describe symmetry-breaking or topologically
ordered fixed-point structures and require a slightly refined parent
construction.

\subsection{Local support and the parent projector}
\label{subsec:mps-parent-projector}

Choose a block length \(\ell\) at least as large as the injectivity length.  Let
\(\mathcal{S}_\ell\) be the support subspace of the \(\ell\)-site reduced density
matrix of \(\ket{\Psi_A}\).  Equivalently,
\begin{equation}
\label{eq:mps-local-support-subspace}
    \mathcal{S}_\ell
    =
    \operatorname{Im}\Gamma_\ell .
\end{equation}
Let \(\Pi_\ell\) be the orthogonal projector onto \(\mathcal{S}_\ell\).  The local
parent-Hamiltonian term is
\begin{equation}
\label{eq:mps-parent-local-term}
    h_{j,j+\ell-1}
    =
    \id-\Pi_{j,j+\ell-1},
\end{equation}
where \(\Pi_{j,j+\ell-1}\) is the translated projector onto the corresponding
support on sites \(j,\ldots,j+\ell-1\).  The parent Hamiltonian is
\begin{equation}
\label{eq:mps-parent-hamiltonian}
    H_{\rm par}
    =
    \sum_j h_{j,j+\ell-1} .
\end{equation}
Because every local reduced density matrix of \(\ket{\Psi_A}\) is supported in
\(\mathcal{S}_\ell\), one has
\begin{equation}
\label{eq:mps-parent-annihilation}
    h_{j,j+\ell-1}\ket{\Psi_A}=0
    \qquad
    \text{for all }j .
\end{equation}
Thus \(\ket{\Psi_A}\) is an exact frustration-free ground state.

\begin{theorem}[Parent Hamiltonian of an injective MPS]
For an injective translation-invariant MPS, the parent Hamiltonian
\cref{eq:mps-parent-hamiltonian} has the MPS as its unique ground state on a
periodic chain, for sufficiently large system size.  On an open chain, the
ground-state space is controlled by the virtual boundary degrees of freedom.
\end{theorem}

\begin{remark}[AKLT as an MPS parent Hamiltonian]
For the AKLT MPS, the two-site physical Hilbert space decomposes into total spin
\(S=0,1,2\).  The two-site support of the VBS state contains only \(S=0\) and
\(S=1\), while the \(S=2\) subspace is absent.  Therefore the parent local term
is precisely the spin-2 projector \(P_{j,j+1}^{(2)}\).  The AKLT Hamiltonian is
an MPS parent Hamiltonian.
\end{remark}

\subsection{Spectral gap and correlation length}
\label{subsec:mps-gap-correlation}

For injective MPS in one dimension, the parent Hamiltonian is generically
expected to be gapped, and the standard injective parent construction indeed
produces a finite spectral gap in the thermodynamic limit.  The finite bond
dimension also implies exponential decay of connected correlations.  Both
features are controlled by the transfer channel \(\mathbb{E}\): if the leading
eigenvalue is normalized to one and the subleading eigenvalue has magnitude
\(\abs{\lambda_2}<1\), then typical connected correlators decay as
\begin{equation}
\label{eq:mps-correlation-length-general}
    C(r)
    \sim
    \lambda_2^r,
    \qquad
    \xi=-\frac{1}{\log\abs{\lambda_2}} .
\end{equation}
The parent Hamiltonian and the transfer channel are therefore two sides of the
same exact-ground-state structure: the parent Hamiltonian fixes the state as a
ground state, while the transfer channel computes its equal-time observables.

\section{Parent Hamiltonians of PEPS}
\label{sec:peps-parent-hamiltonians}

Projected entangled-pair states are the natural two-dimensional generalization
of MPS.  A PEPS assigns a tensor to each vertex of a lattice.  Its physical
index labels the local Hilbert space, while its virtual indices are contracted
along edges.  The AKLT construction itself has a two-dimensional analogue on
suitable lattices: place virtual spin-
\(1/2\)'s on incident bonds, form virtual singlets on edges, and project the
virtual degrees of freedom at each site into a physical spin representation.

The PEPS parent-Hamiltonian construction parallels the MPS case.  For a finite
region \(R\), let \(\mathcal{S}_R\) be the support of the reduced density matrix
of the PEPS on \(R\).  Define a local term
\begin{equation}
\label{eq:peps-parent-local-term}
    h_R=\id-\Pi_R,
\end{equation}
where \(\Pi_R\) projects onto \(\mathcal{S}_R\).  Then
\begin{equation}
\label{eq:peps-parent-hamiltonian}
    H_{\rm PEPS}=
    \sum_R h_R
\end{equation}
is frustration-free and has the PEPS as a ground state.

In two dimensions, however, the physics is more subtle than in injective
one-dimensional MPS.  PEPS parent Hamiltonians can describe
\begin{enumerate}[label=(\roman*)]
    \item short-range-entangled trivial phases;
    \item symmetry-protected topological phases;
    \item intrinsic topological order, when the tensors have virtual gauge
    symmetries;
    \item critical states, where correlations may decay algebraically;
    \item gapless or gapped parent Hamiltonians, depending on the tensor.
\end{enumerate}
Thus the PEPS parent construction gives an exact ground state, but it does not
by itself guarantee a simple spectrum or even a finite gap.

\begin{remark}[Virtual symmetries and topological order]
In PEPS, intrinsic topological order is often encoded by symmetries acting on
the virtual indices.  The toric code, string-net fixed points, and many quantum
double states admit PEPS descriptions with such virtual gauge structures.  In
this language, topological degeneracy arises from nontrivial virtual boundary
conditions rather than from ordinary local order parameters.
\end{remark}

\section{Comparison of AKLT, cluster, RK, and toric-code solvability}
\label{sec:exact-ground-state-comparison}

The models in this part use related but distinct exact mechanisms.  The
following table summarizes the differences.

\begin{center}
\begin{adjustbox}{max width=\textwidth}
\begin{tabular}{L{0.21\textwidth}L{0.25\textwidth}L{0.24\textwidth}L{0.22\textwidth}}
\toprule
Model or construction & Exact mechanism & What is exactly controlled & Full spectrum? \\
\midrule
Toric code & Commuting stabilizer projectors & Ground space, anyon sectors, all stabilizer energies & Yes, in stabilizer variables \\
Cluster model & Commuting graph-state stabilizers & Unique or edge-degenerate stabilizer ground state, SPT edge structure & Yes, in stabilizer variables \\
AKLT chain & Noncommuting projector parent Hamiltonian & VBS/MPS ground state, edge spin-
\(1/2\)'s, correlations & Not by the parent construction alone \\
RK dimer model & Frustration-free projector at RK point & Equal-amplitude or weighted dimer ground state, classical correlators & Generally no \\
Injective MPS parent Hamiltonian & Local support projector & Unique exact MPS ground state and exponential correlations & Generally no \\
PEPS parent Hamiltonian & Local support projector with possible virtual symmetry & Exact tensor-network ground space, possibly topological sectors & Generally no \\
\bottomrule
\end{tabular}
\end{adjustbox}
\end{center}

The important lesson is that projector methods are broader than commuting-
projector methods.  Commutativity solves the spectrum.  Frustration-free parent
Hamiltonians solve the ground state.  Tensor-network transfer matrices then
compute equal-time observables, much as the row-to-row transfer matrix computed
classical correlations in earlier chapters.

\section{A unified parent-Hamiltonian principle}
\label{sec:unified-parent-hamiltonian-principle}

The parent-Hamiltonian construction can be stated abstractly.

\begin{principle}[Parent-Hamiltonian principle]
Given a many-body state \(\ket{\Psi}\), choose local regions \(R\) and let
\(\Pi_R\) project onto the support of the reduced density matrix
\(\rho_R=\Tr_{\overline R}\ket{\Psi}\bra{\Psi}\).  Then
\begin{equation}
\label{eq:abstract-parent-hamiltonian}
    H_{\rm par}[\Psi]
    =
    \sum_R(\id-\Pi_R)
\end{equation}
is positive semidefinite and satisfies
\begin{equation}
\label{eq:abstract-parent-annihilation}
    H_{\rm par}[\Psi]\ket{\Psi}=0 .
\end{equation}
If the intersection of the local supports is exactly the desired ground-state
space, then \(H_{\rm par}[\Psi]\) is an exact frustration-free Hamiltonian for
that state or phase.
\end{principle}

This principle explains why the AKLT chain, graph states, PEPS string-net fixed
points, and many RK-like constructions fit naturally into the same lecture-note
architecture.  They are not solved by the same algebra, but they share the same
logic:
\boxedpar{%
local tensor or local constraint
\(\Longrightarrow\)
finite-dimensional support condition
\(\Longrightarrow\)
projector Hamiltonian
\(\Longrightarrow\)
exact ground state and controlled correlations.}

\section{Common pitfalls}
\label{sec:aklt-cluster-parent-pitfalls}

Several points are worth keeping separate.

\begin{enumerate}[label=(\roman*)]
    \item \emph{Frustration-free does not mean commuting.}  The AKLT local
    projectors do not commute on overlapping bonds, although the exact ground
    state is annihilated by all of them.
    \item \emph{Exact wavefunction does not imply exact spectrum.}  Parent
    Hamiltonians usually determine the ground-state subspace, not all excited
    eigenstates.
    \item \emph{Short-range entangled does not mean trivial under symmetry.}
    AKLT and cluster states are short-range entangled, but they represent
    nontrivial SPT phases when the relevant symmetries are preserved.
    \item \emph{Tensor-network exactness is algebraic, not variational.}  In an
    exact parent-Hamiltonian construction, the tensor network is not merely a
    good approximation; it is annihilated exactly by local projectors.
    \item \emph{Two-dimensional parent Hamiltonians are more delicate than
    one-dimensional ones.}  PEPS may have topological sectors, virtual gauge
    symmetries, critical correlations, or gapless parent Hamiltonians.
\end{enumerate}

\section{Connection to the rest of the notes}
\label{sec:exact-ground-state-models-connections}

The AKLT, cluster, and parent-Hamiltonian constructions complete the set of
solvability mechanisms emphasized in these notes.  They should be read in
parallel with the previous chapters as follows:
\begin{equation}
\label{eq:solvability-mechanism-chain-final}
\begin{array}{ccl}
\text{XY/Kitaev/SSH} &:& \text{quadratic fermions and BdG diagonalization},\\[2pt]
\text{2D Ising} &:& \text{transfer matrices and free-fermion structure},\\[2pt]
\text{XXZ/six-vertex} &:& \text{Yang--Baxter algebra and Bethe ansatz},\\[2pt]
\text{toric code/quantum double} &:& \text{commuting projectors and topological constraints},\\[2pt]
\text{Kitaev honeycomb} &:& \text{Majorana fermions in static gauge sectors},\\[2pt]
\text{RK/AKLT/cluster/parent models} &:& \text{exact frustration-free ground states}.
\end{array}
\end{equation}
Together these examples show that exact solvability is not a single technique.
It is a family of algebraic reductions, each exposing a different kind of
structure in many-body quantum systems.


% -----------------------------------------------------------------------------
% -----------------------------------------------------------------------------
% Concluding synthesis chapter for
% Solvable Quantum Many-Body Models in Condensed Matter
% This file is intended to be included by \input{} from main.tex.
% -----------------------------------------------------------------------------

\part{Concluding Synthesis}

\chapter{Mechanisms of Exact Solvability in Quantum Many-Body Models}
\label{ch:mechanisms-exact-solvability}

The previous chapters developed several families of exactly solvable models: one-dimensional free-fermion spin chains, transfer-matrix formulations of two-dimensional classical systems, Yang--Baxter integrable quantum chains, commuting-projector topological phases, Majorana--gauge solvable models, Rokhsar--Kivelson points, and parent-Hamiltonian constructions.  The purpose of this final chapter is to put those examples into one conceptual map.

The central lesson is simple but important:
\boxedpar{%
Exact solvability is not a single technique.  It is the manifestation of hidden structure.}
Different models are solvable for different reasons.  Sometimes the hidden structure is a nonlocal fermionization.  Sometimes it is a commuting family of transfer matrices.  Sometimes it is factorized scattering.  Sometimes it is a local stabilizer algebra, a static gauge field, a frustration-free projector structure, or a tensor-network parent construction.  These mechanisms are not equivalent, but they are comparable.  Together they form a small set of organizing principles for a large part of exactly solvable many-body physics.

\section{Several Meanings of Exact Solvability}
\label{sec:several-meanings-solvability}

The phrase \emph{exactly solvable} is often used broadly.  In a many-body context it is useful to distinguish several levels of exactness.

\begin{definition}[Exact solvability]
A many-body model is exactly solvable, in a specified sense, if some nonperturbative part of its many-body structure can be determined in closed or systematically controllable form.  The controlled object may be the full spectrum, a transfer-matrix spectrum, the partition function, ground states, correlation functions, quasiparticle scattering, topological sectors, or a complete set of commuting projectors.
\end{definition}

This definition is deliberately plural.  For example, the transverse-field Ising chain and the AKLT chain are both exactly solvable, but not in the same sense.  The transverse-field Ising chain can be mapped to a quadratic fermion Hamiltonian and spectrally diagonalized.  The AKLT chain has an exactly known valence-bond-solid ground state and a frustration-free parent Hamiltonian, but its full excitation spectrum is not solved by the same elementary diagonalization.

\begin{longtable}{L{0.22\textwidth}L{0.43\textwidth}L{0.25\textwidth}}
\caption{Several useful meanings of exact solvability.}\label{tab:exact-solvability-meanings}\\
\toprule
Type of solvability & Controlled structure & Representative examples \\
\midrule
\endfirsthead
\toprule
Type of solvability & Controlled structure & Representative examples \\
\midrule
\endhead
Full quadratic spectral solution & The Hamiltonian is reduced to independent fermionic or bosonic modes by an exact transformation. & TFIM, XY chain, Kitaev chain \\
Transfer-matrix solution & The partition function or quantum spectrum is encoded in the eigenvalues of a transfer matrix. & 2D Ising model, six-vertex model, eight-vertex model \\
Yang--Baxter integrability & There is a commuting family of transfer matrices and factorized many-body scattering. & XXX, XXZ, XYZ, Hubbard chain \\
Commuting-projector solution & The Hamiltonian is a sum of mutually commuting local terms, so the spectrum is organized by simultaneous eigenvalues. & Toric code, quantum double models, cluster stabilizers \\
Majorana--gauge solution & The spin problem decomposes into static gauge sectors, and each sector contains a quadratic Majorana Hamiltonian. & Kitaev honeycomb model \\
Frustration-free ground-state solution & The ground state is annihilated by every local positive term. & RK point, AKLT chain, MPS/PEPS parent Hamiltonians \\
\bottomrule
\end{longtable}

The distinction matters because different exact solutions answer different physical questions.  A free-fermion solution gives precise quasiparticle dispersions and correlation functions.  A Bethe-ansatz solution gives interacting quasiparticles and thermodynamic Bethe equations.  A commuting-projector solution gives topological degeneracy and anyonic superselection sectors with almost no conventional quasiparticle dispersion.  A parent-Hamiltonian construction gives exact ground states and entanglement structure, but usually not a closed expression for all excited states.

\begin{principle}[Solvability must specify what is solved]
A statement that a model is exactly solvable should be accompanied by the object that is exactly controlled: the spectrum, the ground state, the partition function, the transfer-matrix eigenvalues, the scattering matrix, the correlation functions, or the topological sector decomposition.
\end{principle}

\section{The Free-Fermion Mechanism}
\label{sec:free-fermion-mechanism}

The first major solvability mechanism in these notes is nonlocal fermionization followed by quadratic diagonalization.  In one dimension, the Jordan--Wigner transformation converts spin flips into fermionic operators by attaching a parity string:
\begin{equation}
\label{eq:conclusion-jw-convention}
    c_j
    =
    \left(\prod_{\ell=1}^{j-1}\sigma_\ell^z\right)\sigma_j^+,
    \qquad
    c_j^\dagger
    =
    \left(\prod_{\ell=1}^{j-1}\sigma_\ell^z\right)\sigma_j^- .
\end{equation}
The string is not an optional decoration.  It is the operator that converts spin commutation relations on different sites into fermionic anticommutation relations.

After fermionization, many spin chains in these notes take the general quadratic form
\begin{equation}
\label{eq:conclusion-quadratic-bdg}
    H
    =
    \sum_{i,j}
    \left(
        c_i^\dagger A_{ij}c_j
        +\frac12 c_i^\dagger B_{ij}c_j^\dagger
        +\frac12 c_i B_{ji}^*c_j
    \right)
    +E_{\rm const},
\end{equation}
where \(A=A^\dagger\) and \(B=-B^{\mathsf T}\).  Equivalently, in Nambu notation,
\begin{equation}
\label{eq:conclusion-bdg-form}
    H
    =
    \frac12
    \begin{pmatrix}
        c^\dagger & c
    \end{pmatrix}
    \begin{pmatrix}
        A & B\\
        -B^* & -A^*
    \end{pmatrix}
    \begin{pmatrix}
        c\\ c^\dagger
    \end{pmatrix}
    +E_0 .
\end{equation}
The remaining steps are linear algebra: Fourier transform, block diagonalization in momentum space, and Bogoliubov rotation.

\subsection{Jordan--Wigner as nonlocal linearization}

The Jordan--Wigner transformation solves an apparent paradox.  The transverse-field Ising and XY chains are interacting spin systems, but their interactions are special: nearest-neighbor spin terms become quadratic fermion hopping and pairing terms.  For example, spin operators on adjacent sites combine with the Jordan--Wigner strings so that the nonlocal strings cancel locally.  This is why the method is powerful but also dimension-specific.  In two spatial dimensions there is no simple one-dimensional ordering for which all local spin interactions remain local quadratic fermion terms.

The essential pattern is
\begin{equation}
\label{eq:conclusion-spin-to-fermion-pattern}
    \text{local spin interaction}
    \quad\xrightarrow{\rm JW}\quad
    \text{quadratic fermion hopping/pairing},
\end{equation}
provided the spin interaction belongs to the appropriate Ising/XY family.  Once the Hamiltonian is quadratic, its eigenstates are fermionic Gaussian states.

\subsection{Bogoliubov diagonalization and criticality}

For translation-invariant chains, the quadratic Hamiltonian decomposes into independent momentum blocks.  A typical BdG block has the form
\begin{equation}
\label{eq:conclusion-bdg-d-vector}
    \mathcal{H}(k)
    =
    d_x(k)\tau^x+d_y(k)\tau^y+d_z(k)\tau^z,
\end{equation}
with quasiparticle energy
\begin{equation}
\label{eq:conclusion-free-spectrum}
    \epsilon(k)=\sqrt{d_x(k)^2+d_y(k)^2+d_z(k)^2}.
\end{equation}
A quantum critical point occurs when the gap closes:
\begin{equation}
\label{eq:conclusion-gap-closing}
    \min_k \epsilon(k)=0.
\end{equation}
This is the common mechanism behind the critical transverse-field Ising chain, the gapless XX chain, the anisotropic XY chain, and the topological transition of the Kitaev chain.  The universal low-energy theory is obtained by expanding \(\epsilon(k)\) near the gapless momentum.  Linear gap closing gives relativistic Majorana or Dirac modes with dynamical exponent \(z=1\); quadratic band touching gives \(z=2\) behavior.

\subsection{Boundary sectors and parity}

The free-fermion solution is exact only after the boundary and parity sectors are handled correctly.  On a ring, the Jordan--Wigner string crossing the boundary produces a fermion-parity dependence:
\begin{equation}
\label{eq:conclusion-parity-boundary}
    c_{L+1}=\pm c_1,
\end{equation}
where the sign is fixed by the total fermion parity sector.  Thus the finite-size spectrum is not obtained by a single Fourier transform with an arbitrary momentum grid.  One must solve the even and odd parity sectors separately and then impose the physical spin boundary condition.

This caveat is not a minor technicality.  It is the finite-size remnant of the fact that fermionization is nonlocal.  In the thermodynamic limit, the difference between the sectors often disappears in bulk observables, but it remains essential for exact finite-size spectra, degeneracies, and edge-mode physics.

\subsection{Classical--quantum correspondence}

The same free-fermion mechanism appears in the two-dimensional classical Ising model.  The row-to-row transfer matrix \(T\) can be interpreted as an imaginary-time evolution operator,
\begin{equation}
\label{eq:conclusion-transfer-hamiltonian}
    T\sim \ee^{-\delta\tau H_{\rm eff}} .
\end{equation}
In the strongly anisotropic limit, the two-dimensional classical Ising model reduces to the one-dimensional transverse-field Ising chain with coupling dictionary of the form
\begin{equation}
\label{eq:conclusion-ising-tfim-dictionary}
    K_x=\delta\tau J,
    \qquad
    K_y=-\frac12\log\tanh(\delta\tau h).
\end{equation}
The Onsager critical curve
\begin{equation}
\label{eq:conclusion-onsager-critical}
    \sinh(2K_x)\sinh(2K_y)=1
\end{equation}
then becomes the quantum critical point \(h=J\) in the Hamiltonian limit.

This illustrates a recurring principle:
\boxedpar{%
A \(d\)-dimensional quantum model at zero temperature is often related to a \((d+1)\)-dimensional classical statistical model through imaginary time.}
The relation is not merely formal.  It identifies quantum criticality with anisotropic classical criticality and explains why transfer matrices are natural bridges between statistical mechanics and quantum many-body theory.

\section{Transfer Matrices, Yang--Baxter Structure, and Bethe Ansatz}
\label{sec:transfer-yang-baxter-bethe-synthesis}

The second major solvability mechanism is integrability by commuting transfer matrices.  This mechanism is more general than free fermions.  The XXZ chain is interacting after Jordan--Wigner transformation, yet it remains exactly solvable because its many-body scattering factorizes into two-body scattering processes.

\subsection{Transfer matrices as generating functions}

The basic construction starts from an \(R\)-matrix acting on the tensor product of two two-dimensional spaces.  The monodromy matrix is
\begin{equation}
\label{eq:conclusion-monodromy}
    M_a(u)
    =
    R_{aL}(u)R_{a,L-1}(u)\cdots R_{a1}(u),
\end{equation}
where \(a\) is an auxiliary space and \(1,\ldots,L\) are quantum spaces.  The transfer matrix is the auxiliary trace
\begin{equation}
\label{eq:conclusion-transfer-matrix}
    \tau(u)=\Tr_a M_a(u).
\end{equation}
The Yang--Baxter equation implies
\begin{equation}
\label{eq:conclusion-commuting-transfer}
    [\tau(u),\tau(v)]=0
    \qquad
    \text{for all }u,v.
\end{equation}
Therefore \(\tau(u)\) is a generating function for commuting conserved quantities:
\begin{equation}
\label{eq:conclusion-log-transfer-charges}
    Q_n
    \propto
    \left.
    \frac{\dd^n}{\dd u^n}\log\tau(u)
    \right|_{u=u_0} .
\end{equation}
For a regular \(R\)-matrix satisfying \(R(0)\propto P\), the first logarithmic derivative gives a local quantum Hamiltonian.

\subsection{Coordinate Bethe ansatz}

In coordinate Bethe ansatz, one works directly in the basis of magnon positions.  In a sector with \(M\) down spins, the wavefunction is written as a superposition of plane waves in each ordering region:
\begin{equation}
\label{eq:conclusion-cba-wavefunction}
    \psi(x_1,\ldots,x_M)
    =
    \sum_{P\in S_M}
    A(P)
    \exp\left(\ii\sum_{a=1}^M k_{P_a}x_a\right),
    \qquad
    x_1<\cdots<x_M .
\end{equation}
The amplitudes \(A(P)\) are related by two-body scattering.  For the XXZ normalization used in these notes,
\begin{equation}
\label{eq:conclusion-xxz-hamiltonian}
    H_{\rm XXZ}^{(+)}
    =
    J\sum_{j=1}^L
    \left(
        \sigma_j^x\sigma_{j+1}^x
        +\sigma_j^y\sigma_{j+1}^y
        +\Delta(\sigma_j^z\sigma_{j+1}^z-1)
    \right),
\end{equation}
with one-magnon dispersion
\begin{equation}
\label{eq:conclusion-one-magnon-dispersion}
    \varepsilon(k)=4J(\cos k-\Delta).
\end{equation}
Periodic boundary conditions give Bethe equations
\begin{equation}
\label{eq:conclusion-bethe-equations-momentum}
    \ee^{\ii k_jL}
    =
    \prod_{\ell\neq j}S(k_j,k_\ell),
    \qquad j=1,\ldots,M.
\end{equation}
The conceptual content of \cref{eq:conclusion-bethe-equations-momentum} is that taking one magnon around the ring is equivalent to scattering it through all other magnons.

\subsection{Algebraic Bethe ansatz}

Algebraic Bethe ansatz repackages the same physics in operator language.  Write
\begin{equation}
\label{eq:conclusion-aba-monodromy-entries}
    M_a(u)
    =
    \begin{pmatrix}
        A(u) & B(u)\\
        C(u) & D(u)
    \end{pmatrix}_a .
\end{equation}
A reference state \(\ket{0}\) is chosen so that
\begin{equation}
\label{eq:conclusion-pseudovacuum}
    C(u)\ket{0}=0,
    \qquad
    A(u)\ket{0}=\alpha(u)\ket{0},
    \qquad
    D(u)\ket{0}=\delta(u)\ket{0}.
\end{equation}
Bethe vectors are then created by the operator \(B(u)\):
\begin{equation}
\label{eq:conclusion-bethe-vector}
    \ket{\lambda_1,\ldots,\lambda_M}
    =
    B(\lambda_1)\cdots B(\lambda_M)\ket{0}.
\end{equation}
The RTT algebra determines the action of \(\tau(u)=A(u)+D(u)\) on this state.  The unwanted terms vanish precisely when the rapidities satisfy the Bethe equations.  In the hyperbolic XXZ convention used earlier, these equations take the representative form
\begin{equation}
\label{eq:conclusion-aba-bethe-equations}
    \left[
    \frac{\sinh(\lambda_j+\eta)}{\sinh\lambda_j}
    \right]^L
    =
    \prod_{\ell\neq j}
    \frac{\sinh(\lambda_j-\lambda_\ell+\eta)}
         {\sinh(\lambda_j-\lambda_\ell-\eta)} .
\end{equation}

\subsection{One structure, several models}

The transfer-matrix viewpoint explains why classical vertex models and quantum spin chains appear in pairs.

\begin{longtable}{L{0.23\textwidth}L{0.23\textwidth}L{0.34\textwidth}}
\caption{Transfer-matrix and Bethe-ansatz families.}\label{tab:integrability-families}\\
\toprule
Classical/statistical object & Quantum Hamiltonian limit & Integrability structure \\
\midrule
\endfirsthead
\toprule
Classical/statistical object & Quantum Hamiltonian limit & Integrability structure \\
\midrule
\endhead
Six-vertex model & XXZ chain & Trigonometric \(R\)-matrix; \(U(1)\) conservation \\
Rational six-vertex limit & XXX chain & Rational \(R(u)=u\id+P\); enhanced \(SU(2)\) symmetry \\
Eight-vertex model & XYZ chain & Elliptic \(R\)-matrix; Baxter \(TQ\) equation \\
Hubbard vertex structure & Hubbard chain & Nested Bethe ansatz; spin and charge rapidities \\
Long-range inverse-square structure & Haldane--Shastry / Calogero--Sutherland & Yangian symmetry; fractional statistics \\
\bottomrule
\end{longtable}

The most compact statement is:
\boxedpar{%
Bethe-ansatz integrability is the exact solvability of interacting many-body motion by factorized scattering.}
Coordinate Bethe ansatz makes the factorized scattering visible in wavefunctions.  Algebraic Bethe ansatz makes the same structure visible in the \(R\)-matrix algebra.

\section{Topological Exact Solvability in Two Dimensions}
\label{sec:topological-exact-solvability-synthesis}

Two-dimensional quantum models introduce a different class of exact solvability.  Instead of diagonalizing a one-dimensional quadratic Hamiltonian or solving one-dimensional factorized scattering, one often exploits local constraints, gauge structure, and topological sectors.

\subsection{Commuting projectors and the toric code}

The toric code is the cleanest example.  Spins live on links of a two-dimensional lattice, and the Hamiltonian is
\begin{equation}
\label{eq:conclusion-toric-code}
    H_{\rm TC}
    =
    -J_e\sum_s A_s
    -J_m\sum_p B_p,
    \qquad
    A_s=\prod_{\ell\ni s}\sigma_\ell^x,
    \qquad
    B_p=\prod_{\ell\in\partial p}\sigma_\ell^z .
\end{equation}
The exact solution follows from
\begin{equation}
\label{eq:conclusion-toric-commuting}
    [A_s,A_{s'}]=[B_p,B_{p'}]=[A_s,B_p]=0.
\end{equation}
The ground-state conditions are
\begin{equation}
\label{eq:conclusion-toric-ground-conditions}
    A_s=+1,
    \qquad
    B_p=+1.
\end{equation}
On a torus, two global constraints reduce the number of independent stabilizers, leaving a four-dimensional ground-state space:
\begin{equation}
\label{eq:conclusion-toric-degeneracy}
    \dim\calH_{\rm GS}(\mathbb{T}^2)=4.
\end{equation}
This degeneracy is not caused by spontaneous symmetry breaking.  It is topological: local operators cannot distinguish the four ground states in the thermodynamic limit.

The elementary excitations are \(e\) charges, \(m\) fluxes, and their composite \(\epsilon=e\times m\).  Their fusion rules are
\begin{equation}
\label{eq:conclusion-toric-fusion}
    e\times e=m\times m=\epsilon\times\epsilon=1,
    \qquad
    e\times m=\epsilon.
\end{equation}
The mutual statistics of \(e\) and \(m\) is semionic: braiding one around the other gives a phase \(-1\).

\subsection{Quantum doubles}

Kitaev's quantum double models generalize the toric code from \(G=\ZZ_2\) to an arbitrary finite group \(G\).  Each oriented link carries a group algebra Hilbert space
\begin{equation}
\label{eq:conclusion-qdouble-link-hilbert}
    \calH_\ell=\CC[G].
\end{equation}
Vertex terms impose gauge invariance, while plaquette terms impose flatness.  The Hamiltonian has the same schematic form:
\begin{equation}
\label{eq:conclusion-quantum-double-hamiltonian}
    H_{D(G)}
    =
    -J_e\sum_v A_v
    -J_m\sum_p B_p .
\end{equation}
The exact solution again follows from commuting projectors.  On a closed surface, ground states are gauge-invariant flat \(G\)-connections.  On a torus they are labeled, roughly, by commuting pairs of group elements modulo simultaneous conjugation.

The anyons of the untwisted quantum double \(D(G)\) are labeled by
\begin{equation}
\label{eq:conclusion-quantum-double-labels}
    (C,\rho),
\end{equation}
where \(C\) is a conjugacy class of \(G\), and \(\rho\) is an irreducible representation of the centralizer of a representative element of \(C\).  Their quantum dimensions are
\begin{equation}
\label{eq:conclusion-qdouble-dimensions}
    d_{(C,\rho)}=|C|\dim\rho,
    \qquad
    \mathcal{D}=|G|.
\end{equation}
The toric code is the Abelian special case where all conjugacy classes and irreducible representations are one-dimensional.

\subsection{Kitaev honeycomb model}

The Kitaev honeycomb model is not a commuting-projector Hamiltonian.  Its solvability is subtler.  It uses spin fractionalization into four Majorana fermions,
\begin{equation}
\label{eq:conclusion-khc-majorana-representation}
    \sigma_i^\alpha=\ii b_i^\alpha c_i,
    \qquad
    D_i=b_i^x b_i^y b_i^z c_i=+1,
\end{equation}
with link variables
\begin{equation}
\label{eq:conclusion-khc-link-variable}
    u_{ij}=\ii b_i^\alpha b_j^\alpha
    \qquad
    \text{on an }\alpha\text{-type bond}.
\end{equation}
These link variables commute with the Hamiltonian, so they define static \(\ZZ_2\) gauge sectors.  In each fixed sector, the matter Majorana Hamiltonian is quadratic:
\begin{equation}
\label{eq:conclusion-khc-free-majorana}
    H[u]
    =
    \ii\sum_{\langle ij\rangle_\alpha}
    J_\alpha u_{ij}c_i c_j .
\end{equation}
Thus the honeycomb model combines two earlier mechanisms:
\begin{equation}
\label{eq:conclusion-khc-two-mechanisms}
    \text{static gauge-field decomposition}
    +
    \text{free Majorana diagonalization}.
\end{equation}
Its gapped Abelian phases reduce, in the strong anisotropy limit, to toric-code topological order.  With time-reversal-breaking perturbations, the gapless phase can be gapped into a chiral non-Abelian Ising topological phase.

\begin{principle}[Two-dimensional exact solvability]
In two dimensions, exact solvability often comes from decomposing the Hilbert space into gauge or topological sectors before diagonalizing the remaining degrees of freedom.
\end{principle}

\section{Frustration-Free and Parent-Hamiltonian Solvability}
\label{sec:frustration-free-parent-synthesis}

Another form of exact solvability controls the ground state rather than the entire spectrum.  The general structure is frustration-free:
\begin{equation}
\label{eq:conclusion-frustration-free}
    H=\sum_\alpha h_\alpha,
    \qquad
    h_\alpha\geq0,
    \qquad
    h_\alpha\ket{\Psi_0}=0
    \quad
    \text{for every }\alpha.
\end{equation}
It follows immediately that \(\ket{\Psi_0}\) is an exact ground state with energy zero.

\subsection{Rokhsar--Kivelson points}

For a quantum dimer model, the Hilbert space consists of close-packed dimer coverings.  On a flippable plaquette \(p\), the RK Hamiltonian at \(v=t\) is
\begin{equation}
\label{eq:conclusion-rk-projector}
    H_p^{\rm RK}
    =
    t
    (\ket{h_p}-\ket{v_p})(\bra{h_p}-\bra{v_p}).
\end{equation}
The equal-amplitude state in a fixed topological or winding sector \(\mathcal{S}\),
\begin{equation}
\label{eq:conclusion-rk-ground-state}
    \ket{\Psi_{\mathcal{S}}}
    =
    \frac{1}{\sqrt{N_{\mathcal{S}}}}
    \sum_{D\in\mathcal{S}}\ket{D},
\end{equation}

is annihilated by every local RK projector.  Equal-time quantum correlations in this state reduce to correlations of the corresponding classical dimer ensemble.  On the square lattice this connects the RK point to height models and quantum Lifshitz criticality.  On the triangular lattice the RK point realizes a gapped \(\ZZ_2\) dimer liquid.

\subsection{AKLT chain}

The spin-1 AKLT chain is another frustration-free model.  Its Hamiltonian is a sum of projectors onto the total spin-2 subspace of neighboring spin-1's:
\begin{equation}
\label{eq:conclusion-aklt-hamiltonian}
    H_{\rm AKLT}
    =
    \sum_j P_{j,j+1}^{(2)}.
\end{equation}
Using the spin-1 algebra, this projector can be written as
\begin{equation}
\label{eq:conclusion-aklt-projector-polynomial}
    P_{j,j+1}^{(2)}
    =
    \frac16(\bm S_j\cdot\bm S_{j+1})^2
    +\frac12\bm S_j\cdot\bm S_{j+1}
    +\frac13 .
\end{equation}
The exact ground state is the valence-bond-solid state: each spin-1 is decomposed into two virtual spin-\(1/2\)'s, neighboring virtual spins form singlets, and the two virtual spins on each site are projected back to the symmetric spin-1 subspace.  This gives an exact matrix-product state and explains the spin-\(1/2\) edge modes of an open chain.

\subsection{Cluster states and stabilizer parent Hamiltonians}

The one-dimensional cluster Hamiltonian is
\begin{equation}
\label{eq:conclusion-cluster-hamiltonian}
    H_{\rm cl}
    =
    -\sum_j K_j,
    \qquad
    K_j=Z_{j-1}X_jZ_{j+1}.
\end{equation}
The stabilizers commute and satisfy \(K_j^2=1\).  Hence the ground state is exactly characterized by
\begin{equation}
\label{eq:conclusion-cluster-stabilizer-condition}
    K_j\ket{\Psi_{\rm cl}}=\ket{\Psi_{\rm cl}}
    \qquad
    \text{for all }j.
\end{equation}
This is both a commuting-projector model and a parent Hamiltonian for a graph state.  In one dimension, with suitable symmetry, it is also a canonical example of symmetry-protected topological order.

\subsection{MPS and PEPS parent Hamiltonians}

The parent-Hamiltonian idea generalizes these examples.  Suppose \(\ket{\Psi}\) is an MPS or PEPS.  For a region \(R\), let \(\Pi_R\) be the projector onto the support of the reduced density matrix \(\rho_R\).  A local parent Hamiltonian is
\begin{equation}
\label{eq:conclusion-parent-hamiltonian}
    H_{\rm parent}
    =
    \sum_R (\id-\Pi_R).
\end{equation}
By construction,
\begin{equation}
\label{eq:conclusion-parent-annihilation}
    (\id-\Pi_R)\ket{\Psi}=0
    \qquad
    \text{for every }R.
\end{equation}
The state \(\ket{\Psi}\) is therefore an exact ground state.

The important distinction is:
\boxedpar{%
A parent Hamiltonian gives exact ground-state solvability.  It does not automatically give full spectral solvability.}
This separates AKLT, cluster, RK, and PEPS parent models from free-fermion and Bethe-ansatz systems.

\section{A Map of Mother Models and Their Descendants}
\label{sec:mother-model-map-conclusion}

The organizing philosophy of these notes is that many named models are best understood as descendants of a smaller number of mother structures.  The following table summarizes the resulting map.

\begin{longtable}{L{0.24\textwidth}L{0.22\textwidth}L{0.25\textwidth}L{0.19\textwidth}}
\caption{Mother structures and representative descendants.}\label{tab:mother-model-map-concluding}\\
\toprule
Mother structure & Representative model & Descendants or limits & Main method \\
\midrule
\endfirsthead
\toprule
Mother structure & Representative model & Descendants or limits & Main method \\
\midrule
\endhead
Quadratic spin-chain fermions & Transverse-field XY chain & TFIM, XX chain, Kitaev chain & Jordan--Wigner + Bogoliubov \\
BdG two-band structure & Kitaev / SSH chains & Majorana chain, chiral two-band insulator & Nambu spinor and winding number \\
Classical transfer matrix & 2D classical Ising model & 1D TFIM Hamiltonian limit & Row-to-row transfer matrix \\
Free Majorana/Pfaffian structure & 2D Ising and dimer models & Ising fermions, Kasteleyn dimers & Pfaffians and covariance matrices \\
Trigonometric \(R\)-matrix & Six-vertex / XXZ & XXX rational limit, six-vertex phases & Bethe ansatz \\
Elliptic \(R\)-matrix & Eight-vertex / XYZ & XXZ degeneration & Baxter transfer matrix \\
Nested integrability & Hubbard chain & Spin-charge rapidities & Nested Bethe ansatz \\
Commuting projectors & Toric code & Stabilizer codes, \(D(\ZZ_2)\) & Projector algebra \\
Finite-group gauge theory & Quantum double \(D(G)\) & Abelian and non-Abelian anyons & Flat connections and representation theory \\
Static gauge Majoranas & Kitaev honeycomb model & Toric-code limit, Ising topological order & Majorana + \(\ZZ_2\) gauge field \\
RK projector construction & Quantum dimer models & Quantum Lifshitz point, \(\ZZ_2\) dimer liquid & Classical-to-quantum mapping \\
Tensor-network parent structure & AKLT, cluster, MPS/PEPS & SPT states, exact tensor-network states & Frustration-free parent Hamiltonian \\
\bottomrule
\end{longtable}

The mother-model viewpoint also clarifies what should be learned from a model.  The TFIM teaches fermionization, parity sectors, and quantum criticality.  The XXZ chain teaches interacting integrability and factorized scattering.  The toric code teaches topological degeneracy and anyon statistics.  The AKLT chain teaches exact tensor-network ground states and symmetry-protected edge degrees of freedom.  The Kitaev honeycomb model teaches fractionalization into matter Majoranas and static gauge fluxes.

\section{Criticality, Topology, and Fractionalization}
\label{sec:criticality-topology-fractionalization}

Exact solvability is valuable because it gives sharply defined mechanisms for phases and excitations.  Three themes recur throughout the notes: criticality, topology, and fractionalization.

\subsection{Criticality}

At a quantum critical point, the excitation gap closes:
\begin{equation}
\label{eq:conclusion-gap-closes}
    \Delta_{\rm gap}\to0.
\end{equation}
Consequently, the correlation length and correlation time diverge:
\begin{equation}
\label{eq:conclusion-critical-divergence}
    \xi\to\infty,
    \qquad
    \tau\sim \xi^z\to\infty.
\end{equation}
Exactly solvable models show this mechanism without approximation.  The TFIM has an Ising critical point.  The XY chain contains several gapless regimes and anisotropic transitions.  The 2D classical Ising model gives the exact Onsager critical point.  The XXZ chain is gapless for \(-1<\Delta\leq1\) in the antiferromagnetic convention and is described at low energy by a Luttinger liquid.  The square-lattice RK point realizes a \(z=2\) quantum Lifshitz critical point.

Thus solvable models are not merely examples with closed-form spectra; they are laboratories for the exact relation between gap closing, long-distance field theory, and universality.

\subsection{Topology}

Topological phases do not fit into the Landau paradigm of local order parameters.  The toric code ground states on a torus cannot be distinguished by any local operator.  Quantum double models organize topological charge by conjugacy classes and centralizer representations.  The Kitaev honeycomb model realizes Abelian topological order in its strong-coupling phases and non-Abelian Ising topological order after time-reversal breaking.  The Kitaev chain and SSH chain show one-dimensional topological phases through boundary zero modes and winding numbers.  AKLT and cluster models illustrate symmetry-protected topology rather than intrinsic topological order.

A useful distinction is:
\begin{equation}
\label{eq:conclusion-topology-distinction}
    \text{intrinsic topological order}
    \neq
    \text{symmetry-protected topological order}.
\end{equation}
The former has long-range entanglement and anyonic excitations in two dimensions.  The latter is short-range entangled but nontrivial in the presence of protecting symmetries.

\subsection{Fractionalization}

Fractionalization occurs when the elementary degrees of freedom of the Hamiltonian are not the elementary excitations of the phase.  Several versions appeared in these notes:
\begin{enumerate}[label=(\roman*)]
    \item In Jordan--Wigner solvable chains, spin flips are represented by nonlocal fermions.
    \item In the Kitaev chain, the fermion can be decomposed into Majorana components, and boundary Majorana zero modes appear in the topological phase.
    \item In the toric code, local spin operators create pairs of anyons rather than isolated local particles.
    \item In the Kitaev honeycomb model, a spin operator fractionalizes into a matter Majorana excitation together with a change of \(\ZZ_2\) flux sector.
    \item In Bethe-ansatz chains, stable quasiparticles scatter elastically with nontrivial two-body phase shifts.
    \item In the Haldane--Shastry chain, spinon excitations behave as fractionalized spin-\(1/2\) objects with fractional statistics.
\end{enumerate}

The recurring lesson is that exact solvability often reveals the correct quasiparticles.  These quasiparticles may be nonlocal relative to the microscopic variables.

\section{What Exact Solvability Does Not Mean}
\label{sec:what-solvability-does-not-mean}

It is equally important to state what exact solvability does \emph{not} imply.

\begin{remark}[Exact ground state versus exact spectrum]
An exact ground state does not imply an exact solution of the full spectrum.  AKLT, RK, MPS, and PEPS parent Hamiltonians are exact ground-state models, but generic excitations need not be solvable by free-particle or Bethe-ansatz methods.
\end{remark}

\begin{remark}[Free fermions can be physically nontrivial]
Free-fermion solvability does not mean trivial physics.  The TFIM contains a nontrivial quantum critical point.  The Kitaev chain contains topological boundary modes.  The Kitaev honeycomb model is quadratic only after decomposition into gauge sectors and supports topological order and non-Abelian anyons.
\end{remark}

\begin{remark}[Integrability is fragile]
Bethe-ansatz integrability depends on an infinite set of commuting charges.  Generic perturbations, such as arbitrary next-nearest-neighbor interactions, usually destroy this structure.  The resulting model may still be physically close to an integrable point at intermediate scales, but it is no longer exactly solvable in the same sense.
\end{remark}

\begin{remark}[Commuting projectors are not conventional order]
Commuting-projector models often have zero correlation length in local operators, but they can have nontrivial topological order.  The absence of a local order parameter does not imply a trivial phase.
\end{remark}

\begin{remark}[Thermodynamic and finite-size solutions require different care]
Finite-size exact solutions require attention to parity sectors, boundary conditions, Bethe quantum numbers, string deviations, gauge constraints, and topological sectors.  These details may disappear from bulk thermodynamics, but they are essential for exact finite-size spectra and degeneracies.
\end{remark}

These cautions are part of the discipline of exact solvability.  A correct solution is not just a formal diagonalization; it must also specify the Hilbert space, boundary conditions, constraints, and sectors.

\section{Outlook: From Solvable Models to Modern Many-Body Physics}
\label{sec:outlook-solvable-models-modern-physics}

Exactly solvable models remain central because they provide calibration points for theory, numerics, and physical intuition.  They are used to test field-theory predictions, benchmark tensor-network algorithms, verify quantum Monte Carlo simulations, construct topological quantum codes, study integrability breaking, and train intuition about phases that survive away from the exactly solvable point.

Several modern directions grow naturally from the models studied here.

\paragraph{Conformal field theory.}
Critical solvable models such as the TFIM, XXZ chain, and 2D classical Ising model flow to conformal field theories in the infrared.  Exact lattice solutions identify the microscopic origin of scaling dimensions, central charges, operator content, and universal finite-size spectra.

\paragraph{Tensor networks.}
AKLT, cluster, MPS, and PEPS parent Hamiltonians show how exact wavefunctions can be constructed from local tensors.  This is the conceptual basis for modern tensor-network variational methods.  Exact parent models provide fixed points and benchmarks for numerical algorithms.

\paragraph{Topological quantum matter.}
Toric code, quantum double models, and Kitaev honeycomb models provide microscopic realizations of anyons, topological degeneracy, and fault-tolerant logical operators.  They are foundational models for topological quantum computation and quantum error correction.

\paragraph{Integrability breaking and thermalization.}
The XXZ and related Bethe-ansatz chains are starting points for studying what happens when integrability is weakly broken.  The contrast between integrable dynamics and generic thermalizing dynamics is sharp only because the integrable limit is exactly controlled.

\paragraph{Numerical and machine-learning benchmarks.}
Exactly solvable models provide ground-truth data for algorithms.  Free-fermion correlation matrices, Bethe-ansatz energies, toric-code sectors, and RK wavefunctions give controlled tests for tensor networks, Monte Carlo methods, neural quantum states, and generative models.

The final lesson can be summarized as follows:
\boxedpar{%
Solvability is hidden structure made explicit.}
The structures encountered in these notes---nonlocal fermionization, BdG linear algebra, transfer matrices, Yang--Baxter equations, factorized scattering, stabilizer algebras, gauge constraints, topological sectors, and parent Hamiltonians---are not merely solution tricks.  They are organizing principles for quantum many-body physics.

\section*{Final Summary}
\addcontentsline{toc}{section}{Final Summary}

The models studied in these notes can be arranged into one compact hierarchy:
\begin{equation}
\label{eq:conclusion-final-hierarchy}
\begin{gathered}
\text{spin algebra}
\longrightarrow
\text{Jordan--Wigner fermions}
\longrightarrow
\text{BdG/free-fermion mother models},
\\[3pt]
\text{transfer matrices}
\longrightarrow
\text{classical--quantum correspondence}
\longrightarrow
\text{Ising and vertex models},
\\[3pt]
\text{Yang--Baxter equation}
\longrightarrow
\text{commuting transfer matrices}
\longrightarrow
\text{Bethe-ansatz integrable chains},
\\[3pt]
\text{local projectors and gauge constraints}
\longrightarrow
\text{topological exact solvability},
\\[3pt]
\text{frustration-free tensor structures}
\longrightarrow
\text{exact ground-state parent Hamiltonians}.
\end{gathered}
\end{equation}
Each line represents a different answer to the question: \emph{why can this many-body problem be solved?}

The answer is never simply that the model is special.  It is special because it possesses a structure strong enough to reorganize the Hilbert space.  Recognizing that structure is the main skill developed by studying exactly solvable quantum many-body models.

% -----------------------------------------------------------------------------
\backmatter
% -----------------------------------------------------------------------------

% =============================================================================
% Essential References
% =============================================================================

\chapter{Essential References}
\label{chap:essential-references}

These notes are intended to be largely self-contained.  The references below are
not exhaustive; they are selected as minimal original sources and standard texts
for the main exactly solvable structures discussed in the notes.  The emphasis is
therefore on representative sources for the solvability mechanisms rather than
on a complete historical bibliography.

\begin{thebibliography}{99}

% -----------------------------------------------------------------------------
% Free fermions, Ising models, and topological one-dimensional chains
% -----------------------------------------------------------------------------

\bibitem{JordanWigner1928}
P.~Jordan and E.~Wigner,
``\"Uber das Paulische \"Aquivalenzverbot,''
\emph{Zeitschrift f\"ur Physik} \textbf{47}, 631--651 (1928).

\bibitem{LiebSchultzMattis1961}
E.~H.~Lieb, T.~Schultz, and D.~Mattis,
``Two soluble models of an antiferromagnetic chain,''
\emph{Annals of Physics} \textbf{16}, 407--466 (1961).

\bibitem{Onsager1944}
L.~Onsager,
``Crystal statistics. I. A two-dimensional model with an order-disorder
transition,''
\emph{Physical Review} \textbf{65}, 117--149 (1944).

\bibitem{SSH1979}
W.~P.~Su, J.~R.~Schrieffer, and A.~J.~Heeger,
``Solitons in polyacetylene,''
\emph{Physical Review Letters} \textbf{42}, 1698--1701 (1979).

\bibitem{Kitaev2001}
A.~Y.~Kitaev,
``Unpaired Majorana fermions in quantum wires,''
\emph{Physics-Uspekhi} \textbf{44}, 131--136 (2001).

% -----------------------------------------------------------------------------
% Transfer matrices, vertex models, and Bethe ansatz
% -----------------------------------------------------------------------------

\bibitem{Bethe1931}
H.~Bethe,
``Zur Theorie der Metalle. I. Eigenwerte und Eigenfunktionen der linearen
Atomkette,''
\emph{Zeitschrift f\"ur Physik} \textbf{71}, 205--226 (1931).

\bibitem{BaxterBook}
R.~J.~Baxter,
\emph{Exactly Solved Models in Statistical Mechanics},
Academic Press, London (1982).

\bibitem{Faddeev1996}
L.~D.~Faddeev,
``How algebraic Bethe ansatz works for integrable models,''
in \emph{Quantum Symmetries / Symetries Quantiques},
edited by A.~Connes, K.~Gawedzki, and J.~Zinn-Justin,
North-Holland, Amsterdam (1998), pp.~149--219;
also available as arXiv:hep-th/9605187.

\bibitem{LiebWu1968}
E.~H.~Lieb and F.~Y.~Wu,
``Absence of Mott transition in an exact solution of the short-range,
one-band model in one dimension,''
\emph{Physical Review Letters} \textbf{20}, 1445--1448 (1968).

\bibitem{Haldane1988}
F.~D.~M.~Haldane,
``Exact Jastrow--Gutzwiller resonating-valence-bond ground state of the
spin-$\frac12$ antiferromagnetic Heisenberg chain with $1/r^2$ exchange,''
\emph{Physical Review Letters} \textbf{60}, 635--638 (1988).

\bibitem{Shastry1988}
B.~S.~Shastry,
``Exact solution of an $S=\frac12$ Heisenberg antiferromagnetic chain with
long-ranged interactions,''
\emph{Physical Review Letters} \textbf{60}, 639--642 (1988).

\bibitem{SutherlandBook}
B.~Sutherland,
\emph{Beautiful Models: 70 Years of Exactly Solved Quantum Many-Body Problems},
World Scientific, Singapore (2004).

% -----------------------------------------------------------------------------
% Two-dimensional topological exact solvability
% -----------------------------------------------------------------------------

\bibitem{Kitaev2003}
A.~Y.~Kitaev,
``Fault-tolerant quantum computation by anyons,''
\emph{Annals of Physics} \textbf{303}, 2--30 (2003).

\bibitem{Kitaev2006}
A.~Y.~Kitaev,
``Anyons in an exactly solved model and beyond,''
\emph{Annals of Physics} \textbf{321}, 2--111 (2006).

% -----------------------------------------------------------------------------
% Dimers, Rokhsar--Kivelson points, and exact ground-state models
% -----------------------------------------------------------------------------

\bibitem{Kasteleyn1961}
P.~W.~Kasteleyn,
``The statistics of dimers on a lattice: I. The number of dimer arrangements on
a quadratic lattice,''
\emph{Physica} \textbf{27}, 1209--1225 (1961).

\bibitem{RokhsarKivelson1988}
D.~S.~Rokhsar and S.~A.~Kivelson,
``Superconductivity and the quantum hard-core dimer gas,''
\emph{Physical Review Letters} \textbf{61}, 2376--2379 (1988).

\bibitem{AKLT1987}
I.~Affleck, T.~Kennedy, E.~H.~Lieb, and H.~Tasaki,
``Rigorous results on valence-bond ground states in antiferromagnets,''
\emph{Physical Review Letters} \textbf{59}, 799--802 (1987).

\bibitem{RaussendorfBriegel2001}
R.~Raussendorf and H.~J.~Briegel,
``A one-way quantum computer,''
\emph{Physical Review Letters} \textbf{86}, 5188--5191 (2001).

\bibitem{FNW1992}
M.~Fannes, B.~Nachtergaele, and R.~F.~Werner,
``Finitely correlated states on quantum spin chains,''
\emph{Communications in Mathematical Physics} \textbf{144}, 443--490 (1992).

% -----------------------------------------------------------------------------
% General background texts
% -----------------------------------------------------------------------------

\bibitem{SachdevBook}
S.~Sachdev,
\emph{Quantum Phase Transitions},
2nd ed., Cambridge University Press, Cambridge (2011).

\bibitem{FradkinBook}
E.~Fradkin,
\emph{Field Theories of Condensed Matter Physics},
2nd ed., Cambridge University Press, Cambridge (2013).

\bibitem{DiFrancescoCFT}
P.~Di Francesco, P.~Mathieu, and D.~Senechal,
\emph{Conformal Field Theory},
Springer, New York (1997).

\bibitem{AltlandSimons}
A.~Altland and B.~D.~Simons,
\emph{Condensed Matter Field Theory},
2nd ed., Cambridge University Press, Cambridge (2010).

\end{thebibliography}

\end{document}
